\documentclass[11pt]{article}

\usepackage[margin=1.1in]{geometry}

\usepackage{amssymb,mathtools,natbib}

% Anchored formal environments for causalsmith papers.
% First mandatory argument = obj_id (machine anchor joining the paper to the
% Lean crosswalk; invisible in the rendered PDF). Optional argument = name.
% Each environment also defines \label{obj:<obj_id>} so prose can use
% \cref{obj:T-1}; cleveref derives the printed kind and number from the target.
\usepackage{amsmath}
\usepackage{mathrsfs}
\usepackage{amsthm}
\usepackage{booktabs}
% Long displays may break across pages; let TeX stretch a little harder before an
% overfull \hbox so wide formulas/tables are less likely to run off the page.
\allowdisplaybreaks
\setlength{\emergencystretch}{3em}
\newtheorem{theorem}{Theorem}
\newtheorem{lemma}{Lemma}
\newtheorem{assumption}{Assumption}
\newtheorem{definition}{Definition}
\newtheorem{proposition}{Proposition}
\newtheorem{remark}{Remark}
\newtheorem{citedresult}{Cited result}
% The amsthm counter is `algorithmbox` (not `algorithm`) so a later `\usepackage{algorithm}` cannot clash.
\newcounter{algorithmbox}
\renewcommand{\thealgorithmbox}{\arabic{algorithmbox}}
\NewDocumentEnvironment{theoremv}{m o s}{\begin{theorem}\label{obj:#1}{\normalfont[\texttt{\detokenize{#1}}]\IfValueT{#2}{\ (#2)}\IfBooleanT{#3}{\verificationfootnotemark}\ }}{\end{theorem}}
\NewDocumentEnvironment{lemmav}{m o s}{\begin{lemma}\label{obj:#1}{\normalfont[\texttt{\detokenize{#1}}]\IfValueT{#2}{\ (#2)}\IfBooleanT{#3}{\verificationfootnotemark}\ }}{\end{lemma}}
\NewDocumentEnvironment{assumptionv}{m o s}{\begin{assumption}\label{obj:#1}{\normalfont[\texttt{\detokenize{#1}}]\IfValueT{#2}{\ (#2)}\IfBooleanT{#3}{\verificationfootnotemark}\ }}{\end{assumption}}
\NewDocumentEnvironment{definitionv}{m o s}{\begin{definition}\label{obj:#1}{\normalfont[\texttt{\detokenize{#1}}]\IfValueT{#2}{\ (#2)}\IfBooleanT{#3}{\verificationfootnotemark}\ }}{\end{definition}}
% A `propositionv` is a proved result tiered as a Proposition (theorem-grade, machine-verified).
\NewDocumentEnvironment{propositionv}{m o s}{\begin{proposition}\label{obj:#1}{\normalfont[\texttt{\detokenize{#1}}]\IfValueT{#2}{\ (#2)}\IfBooleanT{#3}{\verificationfootnotemark}\ }}{\end{proposition}}
% A `remarkv` holds NON-load-bearing interpretive / open-question content (e.g. a causalsmith `oeq:`
% open question). It is neither proved nor assumed; no theorem may depend on it.
\NewDocumentEnvironment{remarkv}{m o s}{\begin{remark}\label{obj:#1}{\normalfont[\texttt{\detokenize{#1}}]\IfValueT{#2}{\ (#2)}\IfBooleanT{#3}{\verificationfootnotemark}\ }}{\end{remark}}
% An `algorithmv` is a DEFINITION-kind object presented as a boxed, numbered-step procedure (an
% estimator / selection rule built in stages). Same anchor contract as `definitionv`; its body is
% ordinary LaTeX whose steps are `\begin{enumerate}` items, so tex2html needs no extra machinery.
% Presented in the CONVENTIONAL algorithm style readers expect from `algorithm`+`algorithmic`:
% a rule above, an "Algorithm N (Title)" caption line, a rule under the caption, the numbered
% steps, and a closing rule. It is NOT a float — the anchor contract requires the object to stay
% where the outline places it, and floats drift. The obj-id tag is suppressed here (it is
% bookkeeping, and a caption line carrying it reads as debug output in a finished paper).
\NewDocumentEnvironment{algorithmv}{m o s}%
  {\par\addvspace{\medskipamount}\noindent\rule{\linewidth}{0.8pt}\par\nobreak\vspace{2pt}%
   \refstepcounter{algorithmbox}\label{obj:#1}%
   \noindent\textbf{Algorithm \thealgorithmbox}\IfValueT{#2}{ \textnormal{(#2)}}%
   \IfBooleanT{#3}{\verificationfootnotemark}\par\nobreak\vspace{2pt}%
   \noindent\rule{\linewidth}{0.4pt}\par\nobreak\vspace{2pt}\normalfont}%
  {\par\nobreak\vspace{2pt}\noindent\rule{\linewidth}{0.8pt}\par\addvspace{\medskipamount}}
% A `citedv` is an IMPORTED external result the paper relies on but does NOT prove
% (a causalsmith `gate_class:"cited"` node). It is machine-checked against the cited
% source at F2.5, never proved here; the fixed provenance note makes that explicit
% so the environment can never be read as a contribution of this paper. The \citep
% attribution is supplied by the surrounding prose (P2), keyed to references.bib.
\NewDocumentEnvironment{citedv}{m o s}{\begin{citedresult}\label{obj:#1}{\normalfont[\texttt{\detokenize{#1}}]\IfValueT{#2}{\ (#2)}\IfBooleanT{#3}{\verificationfootnotemark}\ \textit{(imported result; machine-checked against the cited source, not proved here.)}\ }}{\end{citedresult}}
% \leanref{obj_id}{display text}: an inline first-mention link to a from-note object's
% verified Lean code. In the PDF it renders the display text plainly (the Lean link is a
% web-only affordance); tex2html turns it into a drawer-opening span.
\newcommand{\leanref}[2]{#2}
% For a theorem/lemma whose Lean declaration accepts a source-matched published proposition as an
% input, the starred anchored environment puts the mark in its heading; the generated text remains
% outside the frozen mathematical environment. Keep the combined macro for older paper bundles.
\newcommand{\verificationfootnotemark}{\footnotemark}
\newcommand{\verificationfootnotetext}[1]{\footnotetext{#1}}
\newcommand{\verificationfootnote}[1]{\verificationfootnotemark\verificationfootnotetext{#1}}
% xcolor before hyperref (hyperref would otherwise pull in plain `color` for colorlinks, and a
% later xcolor load clashes). The page-1 identity stamp at the end of this file needs a grey.
\usepackage{xcolor}
\usepackage{graphicx}
\usepackage{eso-pic}
% hyperref follows natbib and the environment macros; cleveref must follow hyperref.
\usepackage[colorlinks=true,linkcolor=blue,citecolor=blue,urlcolor=blue]{hyperref}
\usepackage[capitalise,nameinlink,noabbrev]{cleveref}
% Pin the names explicitly: labels are placed inside the underlying amsthm environments, and these
% declarations make the PDF contract independent of cleveref's theorem-name inference. House style:
% reference names always print with a capital first letter ("Theorem 1", "Definition 2"), in any
% sentence position — hence `capitalise` above and capitalized \crefname forms below.
\crefname{theorem}{Theorem}{Theorems}
\Crefname{theorem}{Theorem}{Theorems}
\crefname{lemma}{Lemma}{Lemmas}
\Crefname{lemma}{Lemma}{Lemmas}
\crefname{assumption}{Assumption}{Assumptions}
\Crefname{assumption}{Assumption}{Assumptions}
\crefname{definition}{Definition}{Definitions}
\Crefname{definition}{Definition}{Definitions}
\crefname{proposition}{Proposition}{Propositions}
\Crefname{proposition}{Proposition}{Propositions}
\crefname{remark}{Remark}{Remarks}
\Crefname{remark}{Remark}{Remarks}
\crefname{citedresult}{Cited result}{Cited results}
\Crefname{citedresult}{Cited result}{Cited results}
\crefname{algorithmbox}{Algorithm}{Algorithms}
\Crefname{algorithmbox}{Algorithm}{Algorithms}
% ---------------------------------------------------------------------------
% Page-1 identity stamp (arXiv style) and the pinned title-block date.
%
% Split of responsibility: THIS file owns the FORMAT (placement, font, colour, punctuation),
% the generated `paper_stamp.tex` owns only the VALUES (number+version, area, revised date,
% DOI, link target). P4 refreshes this template on every emit, so a layout fix reaches every
% bundle from a recompile; and because `paper_stamp.tex` is not one of the P2 assembly
% digest's authored inputs, a DOI that arrives after publication needs only that one small
% file rewritten plus a recompile — never a redraft, never a reassembly.
%
% Defaults first, so a bundle WITHOUT a stamp file compiles exactly as it always did:
% `\cspaperdate` falls back to `\today` and no stamp is drawn.
\providecommand*{\cspaperdate}{\today}  % the \date{} target
\providecommand*{\csstampid}{}          % e.g. CSWP-2026-008v3 (empty ⇒ no stamp at all)
\providecommand*{\csstamparea}{}        % e.g. Stat
\providecommand*{\csstampdate}{}        % e.g. 27 Aug 2026
\providecommand*{\csstampdoi}{}         % e.g. 10.5281/zenodo.123 (empty ⇒ DOI part omitted)
\providecommand*{\csstamplink}{}        % href target (DOI url, else the paper's page; may be empty)
\IfFileExists{paper_stamp.tex}{% GENERATED by CausalSmith P4 — paper identity stamp values. Do not hand-edit.
%
% Field order on the stamp (owned by paper_macros.tex):
%   <number>v<version> · doi:<doi> · [<Area>] <date>   — the DOI is never the last token, so a
%   text extractor cannot run it into whatever follows. Without a DOI that part is omitted.
% Format lives in paper_macros.tex; only the values are here, and this file is NOT one of
% the P2 assembly digest's authored inputs. A DOI arriving after publication therefore needs
% this file rewritten plus a `latexmk` recompile — never a redraft or a reassembly.
\renewcommand*{\csstampid}{CSWP-2026-026v3}
\renewcommand*{\csstamparea}{Stat}
\renewcommand*{\csstampdate}{1 Oct 2026}
\renewcommand*{\csstampdoi}{}
\renewcommand*{\csstamplink}{https://causalsmith.org/papers/stat_weakoverlap_unisolvent_rate_v1}
\renewcommand*{\cspaperdate}{1 October 2026}
}{}
\makeatletter
% Field order is deliberate: the DOI is NEVER the last token on the line. A trailing identifier
% runs into whatever a text extractor emits next, which makes it unsafe to match mechanically
% (the Zenodo stamp matcher reads exactly this line); the date is a safe terminator.
%   <number>v<version> · doi:<doi> · [<Area>] <date>      — with a DOI
%   <number>v<version> · [<Area>] <date>                  — without
\newcommand{\cs@stampsep}{\,\textperiodcentered\,}
\newcommand{\cs@stamptext}{%
  \csstampid
  \ifx\csstampdoi\@empty\else\cs@stampsep doi:\csstampdoi\fi
  \ifx\csstamparea\@empty\else\cs@stampsep[\csstamparea]\fi
  \ifx\csstampdate\@empty\else\quad\csstampdate\fi
}
\ifx\csstampid\@empty\else
  \definecolor{csstampgrey}{gray}{0.45}
  % Background shipout material: it occupies no space, so the text block and the \thanks
  % footnote are byte-identical to an unstamped compile. The starred form applies to the
  % NEXT page shipped only — i.e. page 1. Placed in the left margin (the text block starts
  % at 1.1in), reading bottom-to-top, in a Times-like face at 8pt.
  \AddToShipoutPictureBG*{%
    \AtPageLowerLeft{%
      \put(\LenToUnit{0.42in},\LenToUnit{1.1in}){%
        \rotatebox{90}{%
          \fontfamily{ptm}\fontsize{8}{10}\selectfont
          \ifx\csstamplink\@empty
            \textcolor{csstampgrey}{\cs@stamptext}%
          \else
            \expandafter\href\expandafter{\csstamplink}{\textcolor{csstampgrey}{\cs@stamptext}}%
          \fi
        }%
      }%
    }%
  }
\fi
\makeatother


\title{Minimax response recovery under global overlap tails}

\author{CausalSmith\thanks{Machine-generated by the CausalSmith research pipeline. Each displayed formal statement is either mapped to a Lean declaration or marked as a presentation-level definition, and each displayed proof renders a Lean proof, as recorded in the verification appendix (scope, toolchain, commit, and any statement conditional on a published input). Cited work otherwise supplies framing and positioning.}}

\date{\cspaperdate}

\begin{document}

\maketitle

\begin{abstract}
We characterize minimax recovery of the treated causal response under a global propensity-tail restriction. The model combines consistency and conditional exchangeability with a covariate density bounded below on a cube, intrinsic Hölder response smoothness of known order \(\beta>1\), a bounded treated conditional mean, and uniformly conditional sub-Gaussian residuals. If \(d\) is the covariate dimension, \(n\) is the sample size, and \(\gamma>1\) is the overlap exponent, the pointwise and expected spatial-supremum minimax risks have scale
\[
n^{-\beta/(2\beta+d\gamma/(\gamma-1))}.
\]
An explicit estimator fits local polynomials with equal aggregate weight across fixed template microcells and selects its mesh through treated observation counts. With model constants fixed, its expected spatial-supremum guarantee attains this pure-power rate uniformly over fixed compact overlap intervals with distinct endpoints in \((1,\infty)\), using the known smoothness and observed treatment design. A bounded, two-valued potential-outcome experiment supplies the matching pointwise lower bound at every fixed point of the closed cube. In dimensions at least two, a thin-strip construction exhibits admissible treatment geometry with degenerating local treated variation. An affine transfer also establishes the expected spatial-supremum upper guarantee for fixed-constant Gaussian source envelopes satisfying Dorn's retained moment, variance, smoothness, uniform-design, and global-tail restrictions.
\end{abstract}

\section{Introduction}

This paper establishes the minimax rate for recovering a treated causal response throughout a covariate cube when treatment availability is governed by a global propensity tail. Under consistency and conditional exchangeability, the target is the conditional mean of the treated potential outcome, identified through treated outcome regression. We study laws with a covariate density bounded below, a response in an intrinsic cube Hölder ball of known order, a uniformly bounded treated conditional mean, and uniformly conditional sub-Gaussian residuals. Write \(n\) for sample size, \(\beta>1\) for response smoothness, \(d\ge1\) for covariate dimension, and \(\gamma>1\) for the overlap exponent. Within this fixed-constant class, one explicit estimator attains the expected spatial-supremum rate
\[
n^{-\beta/(2\beta+d\gamma/(\gamma-1))}.
\]
The upper guarantee is uniform over any fixed compact exponent interval with distinct endpoints contained in \((1,\infty)\). For each fixed exponent and each fixed evaluation point, a matching lower bound establishes the same minimax scale for pointwise expected absolute error.

Response recovery is a basic component of causal adjustment in the potential-outcome framework \citep{rubin1974,rosenbaum1983,imbens2015}. Its statistical difficulty depends on where treated observations are available. Our overlap restriction bounds the covariate probability of low-propensity regions by a power of the propensity threshold. This restriction quantifies the prevalence of treatment scarcity while accommodating measurable propensities and heterogeneous local treatment geometry. The target remains the continuous treated response over the entire closed cube, and the spatial loss measures the largest estimation error within each sample realization before taking expectation.

The construction separates polynomial conditioning from treatment-count precision. Each partition cube contains scaled copies of a fixed collection of small microcells around polynomial interpolation nodes. The fit gives every microcell the same aggregate weight, averaging treated observations within each cell. This weighting keeps the empirical Gram matrix close to a fixed positive-definite template whenever the cells are occupied. A count selector chooses the finest candidate mesh at which every microcell contains enough treated observations relative to the approximation scale. The resulting estimator uses the observed sample, sample size, dimension, known smoothness, and response bound, as specified in \cref{obj:def:mesh-selector,obj:def:estimator}.

The statistical explanation rests on an ordered information profile. The density lower bound assigns covariate mass to each template microcell, and the global propensity tail controls the treated mass of unions of disjoint microcells. Together, they imply a polynomial lower bound in rank for the weakest treated-cell masses across partition cubes. This profile distinguishes the least-informed cubes from the progressively better-informed cubes. Combining its rank-dependent stochastic control with the common scale cap supplied by count feasibility yields the pure-power expected spatial-supremum guarantee in \cref{obj:thm:adaptive-minimax}. Local polynomial approximation and deterministic template conditioning complete the upper bound.

The converse uses a bounded causal experiment. For every fixed evaluation point, including points on faces and corners, a pair of laws places a smooth response perturbation near that point and shares a radial propensity satisfying the global tail envelope. The laws admit potential outcomes taking values in a common two-point set. Their response separation and bounded relative entropy give the pointwise lower bound in \cref{obj:lem:bounded-lower-pair}. Thus the minimax characterization applies to the full conditional sub-Gaussian response class, with the converse already supported by a bounded subclass.

The closest theorem-level comparison is \citet[Assumptions 1--3 and Theorem 1(ii), p.~12]{Dorn2026GlobalOverlap}. For a fixed Gaussian uniform-design setup satisfying Dorn's retained restrictions and local anti-concentration condition, Dorn establishes a probability upper bound at the same pure-power scale in an essential-supremum criterion. The source describes a count-based estimator computable without the overlap exponent, while its probability upper guarantee fixes the source family, its common constants, and that exponent before selecting the estimator and threshold constant \citep[Lemmas 11 and 13, pp.~23--24; proof, pp.~27--28]{Dorn2026GlobalOverlap}. That construction provides a direct antecedent for mesh selection. Dorn's local anti-concentration condition additionally links every sufficiently small neighborhood's largest selected propensity to a common positive fraction of its covariate mass, with constants common across the fixed family and its selected propensity representatives. Our expected spatial-supremum guarantee is established over the global-tail class in \cref{obj:def:model}, using one estimator and common constants across the stated compact exponent range.

The geometric and source-model comparisons make this distinction precise. In dimensions at least two, \cref{obj:prop:strict-enlargement} constructs a thin-strip law within the causal class whose displayed propensity violates the selected-version local anti-concentration condition and whose normalized treated-design second-coordinate moment tends to zero. Separately, \cref{obj:prop:dorn-a3-free-corollary} transfers the upper guarantee to fixed-constant Gaussian envelopes satisfying the retained source restrictions from \citet[Section 2, Assumptions 1--2, pp.~3--7]{Dorn2026GlobalOverlap}. Affine transport and fixed-cube smoothness completion supply the common full Hölder radius. The transferred guarantee controls expected spatial-supremum error uniformly over those envelopes and the compact exponent range. Its Gaussian geometric witnesses use constants selected separately for each exponent.

Related regression results characterize other aspects of heterogeneous information. \citet[Assumption M and Theorem 1, pp.~3--4]{Gaiffas2005DegenerateRates} derives pointwise rates under locally symmetric, regularly varying design densities in a one-dimensional Gaussian experiment. \citet[Assumption D, equation (2.1), and Theorems 1--2]{gaiffas2009} and \citet[Sections 1 and 3]{schmidthieber2024} retain spatial variation through location-dependent loss normalization. Our criterion instead evaluates expected unweighted spatial-supremum error over a global overlap envelope. The causal target also differs from the smooth treatment contrast studied under strict overlap by \citet[Theorems 1--2 and Remark 11]{KennedyBalakrishnanRobinsWasserman2024}: here smoothness is imposed directly on the treatment-specific response.

An appendix documents the precise scope of the mathematical results.

The exposition proceeds through related work, the model and risk criteria, and the estimator and recovery guarantee. It then develops the local treatment-geometry comparison and Gaussian transfer, followed by discussion and future work. The appendices present auxiliary identification, upper-bound, testing, geometry, and source-comparison arguments, followed by the complete main-result proofs.


\section{Related work}

The closest theorem-level comparator is \citet[Assumptions 1--3 and Theorem 1(ii), p. 12]{Dorn2026GlobalOverlap}. In a fixed setup with uniform covariates on a cube, Gaussian outcomes, and Hölder response curves, satisfying Dorn's retained restrictions and local anti-concentration condition, Dorn establishes an essential-supremum probability upper bound uniformly over the source family. Writing \leanref{sym:d}{\ensuremath{d}} for the covariate dimension, \leanref{sym:beta}{\ensuremath{\beta}} for response smoothness, \leanref{sym:gamma}{\ensuremath{\gamma}} for the overlap exponent, and \leanref{sym:n}{\ensuremath{n}} for sample size, the pure-power rate in this paper's notation is
\[
n^{-\beta/(2\beta+d\gamma/(\gamma-1))}.
\]
For each probability tolerance, the source result supplies a finite threshold constant within that fixed setup. Dorn describes a count-based grid and bandwidth estimator computable from the known smoothness and observed treatment design without the overlap exponent \citep[Lemmas 11 and 13, pp. 23--24; proof, pp. 27--28]{Dorn2026GlobalOverlap}. The probability upper guarantee fixes the source family, its common constants, and the exponent before selecting the estimator and threshold constant. The construction is a direct antecedent for the present mesh selector; the new uniformity statement concerns one estimator and common risk constants across a prescribed compact exponent interval.

The present comparison concerns the guarantee and its model domain. Under the \leanref{S-1}{fixed model indices and class constants}, \cref{obj:thm:adaptive-minimax} gives an expected spatial-supremum guarantee over the envelopes in \cref{obj:def:model}, uniformly across a compact range of overlap exponents, with response smoothness fixed. These envelopes impose a bounded treated conditional mean and uniformly conditional sub-Gaussian residuals, intrinsic cube Hölder smoothness, a covariate density lower bound, and a global propensity-tail bound. Bounded potential outcomes are used for the testing subclass and the thin-strip example, rather than for the full upper-bound class. Dorn's local anti-concentration condition additionally requires a common positive fraction of the covariate mass in every sufficiently small neighborhood to have propensity comparable to that neighborhood's propensity supremum. Here, the global tail and the density lower bound support the equal-cell construction under the conditions in \cref{obj:ass:density,obj:ass:global-tail}. The distinction in admissible local treatment geometry is made precise in \cref{obj:def:dorn-a3,obj:prop:strict-enlargement}, while \cref{obj:prop:dorn-a3-free-corollary} states the corresponding Gaussian-envelope upper guarantee.

Degenerate-design regression supplies a complementary comparison at a single location. In a one-dimensional Gaussian regression experiment, \citet[Assumption M and Theorem 1, pp. 3--4]{Gaiffas2005DegenerateRates} derives pointwise minimax moment rates for a fixed design density that is locally symmetric and regularly varying around the evaluation point. The regression class is specified through a local polynomial approximation modulus, together with a bound on the function; regular variation of the design and modulus yields a power rate multiplied by a slowly varying factor. In the separate adaptive experiment, \citet[Abstract]{Gaiffas2005AdaptiveDegenerateDesign} constructs a local polynomial procedure adaptive to design degeneracy and Hölder smoothness, and identifies the logarithmic cost of that pointwise adaptation. These results characterize local information through the behavior of a specified design near one point. The present envelopes instead organize treatment-specific recovery through a global distributional restriction on the propensity.

Spatially normalized losses provide another way to express heterogeneous information. In one-dimensional Gaussian regression over Hölder classes, \citet[Assumption D, equation (2.1), and Theorems 1--2]{gaiffas2009} studies the supremum of estimation error divided by a location-dependent rate obtained from a local bias--variance balance. The adaptive upper result assumes a continuous design density that is positive everywhere or has finitely many zeros with specified local power behavior; the lower result uses an interval-mass condition and establishes optimality over intervals of prescribed shrinking size. Likewise, \citet[Sections 1 and 3]{schmidthieber2024} analyzes a weighted uniform criterion in one-dimensional Gaussian regression under a local doubling condition on the design distribution. Least squares over the unit-Lipschitz class attains the local rate for truths in a strictly smaller Lipschitz ball. Both comparisons retain the design's spatial variation in the loss normalization. Our criterion in \cref{obj:def:risks} measures expected unweighted spatial-supremum error, with worst-case difficulty indexed by the global overlap envelope.

The estimator combines familiar polynomial approximation with a particular treatment-design weighting rule. Local polynomial fitting approximates a smooth response within a neighborhood \citep{fan1996}; series methods approximate conditional expectations in finite-dimensional spaces, including polynomial and spline spaces \citep{newey1997}. Partitioning-based series estimation develops this approach through locally supported bases \citep{CattaneoFarrellFeng2020}. The construction in \cref{obj:def:template,obj:def:estimator} fits polynomials on partition cubes, assigns equal aggregate weight to the reference microcells, and selects its mesh through the treated-count feasibility rule in \cref{obj:def:mesh-selector}.

The causal estimand also determines the relevant comparison. \citet[Theorems 1--2 and Remark 11]{KennedyBalakrishnanRobinsWasserman2024} studies pointwise expected absolute error for the conditional average treatment effect under strict overlap, with separate Hölder restrictions on the treatment effect, propensity, and control response. Its attainable rates depend on their relative smoothness and on conditions controlling covariate-distribution estimation and the local fitting matrices. Our target is a treatment-specific response curve, with smoothness imposed directly on that response and information scarcity described by a global overlap tail. For functional estimation, \citet[Theorems 1--3]{MouDingWainwrightBartlett2023} derives local minimax mean-squared-error bounds for linear functionals over convex symmetric regression classes, together with attaining procedures for reproducing kernel Hilbert space classes under the stated noise, spectral, and sample-size conditions. That experiment measures accuracy through a specified functional and its observational coverage. The present minimax criterion evaluates recovery of the response throughout the covariate cube.


\section{Setup and assumptions}

Generic positive constants may change between displays; the model constants introduced below remain fixed across laws and sample sizes. For a vector, \(\|\cdot\|_\infty\) denotes the maximum absolute coordinate, and for a function it denotes the supremum of its absolute value over its stated domain. Multi-indices have nonnegative integer coordinates, with \(|\alpha|\) denoting their sum. Observation indices range from \(1\) to \(n\), and \(a\vee b=\max\{a,b\}\). The notation \(\lesssim\) denotes an inequality up to a positive constant depending on the fixed model parameters, and \(\asymp\) denotes inequalities in both directions. For comparisons uniform over the overlap range specified below, these constants are independent of both the sample size and the overlap exponent.

\subsection{Sample and target}

Consider an \leanref{S-2}{observational causal model} indexed by a probability law \leanref{sym:P}{\ensuremath{P}}. The covariate vector \leanref{sym:X}{\ensuremath{X}} takes values in the covariate cube \leanref{sym:mathcalX}{\ensuremath{\mathcal X}}, where \(\mathcal X=[0,1]^d\) and \(d\ge1\). The binary treatment indicator \leanref{sym:A}{\ensuremath{A}} takes values in \(\{0,1\}\). Write \(Y(0)\) for the control potential outcome, \leanref{sym:Y1}{\ensuremath{Y(1)}} for the treated potential outcome, and \leanref{sym:Y}{\ensuremath{Y}} for the observed outcome. The observed triple \leanref{sym:Z}{\ensuremath{Z}} is \(Z=(X,A,Y)\). When specifying the sampling law, \(P\) denotes the observed-data marginal of the causal law, and \(P_X\) denotes its covariate marginal.

The covariate density \leanref{sym:f}{\ensuremath{f}} is the density of \(X\) with respect to Lebesgue measure on \(\mathcal X\), and the propensity \leanref{sym:e}{\ensuremath{e(x)}} is a measurable version of \(P(A=1\mid X=x)\). Our target is the treated causal response \leanref{sym:mu1}{\ensuremath{\mu_1}}, the continuous Hölder representative of
\[
\mu_1(x)=\mathbb E_P[Y(1)\mid X=x].
\]
We write \(\mu_{1,P}\) when the indexing law needs to be explicit.

The \leanref{S-3}{sample and estimator convention} uses the observed sample \leanref{sym:Dn}{\ensuremath{\mathcal D_n}}, defined by \(\mathcal D_n=(Z_i)_{i=1}^n\), where \(Z_i=(X_i,A_i,Y_i)\). Every estimator is measurable with respect to this sample. Scalar-valued estimators estimate the response at a specified covariate point; function-valued estimators recover the response over \(\mathcal X\) and are required to have measurable spatial-supremum loss. Expectations of their losses refer to the observed-sample distribution.

Fix a smoothness order \(\beta>1\), a Hölder radius \leanref{sym:L}{\ensuremath{L}} with \(L>0\), and a treated-response bound \leanref{sym:B}{\ensuremath{B}} with \(B>0\). Also fix a density lower bound \leanref{sym:c_f}{\ensuremath{c_f}} with \(0<c_f\le1\) and a global-tail constant \leanref{sym:C}{\ensuremath{C}} with \(C\ge1\). The overlap exponent ranges over the compact interval \leanref{sym:Gamma}{\ensuremath{\Gamma}}, defined by
\[
\Gamma=[\gamma_{\min},\gamma_{\max}].
\]
Its lower endpoint \leanref{sym:gamma_min}{\ensuremath{\gamma_{\min}}} and upper endpoint \leanref{sym:gamma_max}{\ensuremath{\gamma_{\max}}} satisfy \(1<\gamma_{\min}\le\gamma_{\max}<\infty\). Throughout, \(\gamma\in\Gamma\), and \(n\ge1\).

\subsection{Intrinsic cube smoothness}

Smoothness is measured directly on the closed covariate cube. For the fixed smoothness order, the derivative order \leanref{sym:m}{\ensuremath{m}} is \(m=\lceil\beta\rceil-1\). For a generic function \leanref{sym:g}{\ensuremath{g}} on \(\mathcal X\), \(D^\alpha g\) denotes its coordinate derivative, interpreted through continuous boundary traces as specified below. The intrinsic cube Hölder ball \leanref{sym:HbetaL}{\ensuremath{\mathcal H^\beta(L)}} combines a uniform derivative bound with a modulus restriction on the highest-order derivatives. We first state these two restrictions.

\begin{assumptionv}{ass:holder-derivative-bound}[Uniform bound on derivatives]
\[
\max_{|\alpha|\le m}\|D^\alpha g\|_\infty\le L.
\]
\end{assumptionv}

The uniform derivative bound controls the function itself and all coordinate derivatives through order \(m\), using a common radius. This standard Hölder-ball condition gives a common scale for response levels and local polynomial coefficients \citep{Dorn2026GlobalOverlap}. The corresponding restriction on variation across covariate points is as follows.

\begin{assumptionv}{ass:holder-modulus}[Hölder modulus of derivatives]
\[
\max_{|\alpha|=m}\sup_{x\ne x'}
\frac{|D^\alpha g(x)-D^\alpha g(x')|}
{\|x-x'\|_\infty^{\beta-m}}
\le L.
\]
\end{assumptionv}

The Hölder modulus condition bounds changes in the highest-order derivatives by the covariate distance raised to \(\beta-m\). Together with the derivative bound, this standard restriction controls local approximation error at smoothness order \(\beta\) \citep{Dorn2026GlobalOverlap}. The next definition fixes the boundary convention used throughout.

\begin{definitionv}{def:holder}[Intrinsic cube Hölder ball]
The intrinsic cube Hölder ball is
\[
\mathcal H^\beta(L)
:=
\left\{
g:\mathcal X\to\mathbb R:
g\in C^m(\mathcal X),\
g\text{ satisfies }\cref{obj:ass:holder-derivative-bound,obj:ass:holder-modulus}
\right\},
\qquad
\mathcal X=[0,1]^d,
\qquad
m=\max\{\lceil\beta\rceil-1,0\}.
\]
Here \(d\) is a nonnegative integer and \(\beta,L\in\mathbb R\). The integer \(m\) is the derivative order, and the top-order modulus has exponent \(\beta-m\). The class \(C^m(\mathcal X)\) consists of functions whose classical coordinate partial derivatives through order \(m\) on the interior of the cube admit unique continuous traces on \(\mathcal X\). Every derivative \(D^\alpha g\), supremum, and difference quotient in the defining restrictions uses these traces. At positive integer \(\beta\), this convention uses Lipschitz derivatives of order \(\beta-1\).
\end{definitionv}

For the maintained domain \(\beta>1\), the derivative order in \cref{obj:def:holder} agrees with the order fixed above. Its trace convention makes smoothness restrictions meaningful at boundary points as well as in the interior.

\subsection{Sampling and observational restrictions}

The statistical experiment combines independent sampling with restrictions on covariate coverage, causal identification, treatment availability, and the treated response. We begin with the sampling law.

\begin{assumptionv}{ass:iid}[Independent and identically distributed observations]
The observed sample \(\mathcal D_n\) consists of \(n\) independent observations with common law \(P\), so that
\[
\mathcal D_n\sim P^{\otimes n}.
\]
\end{assumptionv}

Independent and identically distributed sampling gives the product experiment \(P^{\otimes n}\) used to evaluate estimation risk. It is a standard sampling condition for nonparametric regression \citep{Stone1982}. Covariate coverage is controlled separately.

\begin{assumptionv}{ass:density}[Covariate density lower bound]
The covariate density \(f\) satisfies
\[
c_f\le f(x)
\]
for Lebesgue-almost every \(x\in\mathcal X\).
\end{assumptionv}

The random-design density lower bound ensures that each measurable covariate region has probability at least \(c_f\) times its volume. This standard coverage condition supports response recovery throughout the cube \citep{Stone1982}. To connect that response to observed treated outcomes, we impose consistency and conditional exchangeability.

\begin{assumptionv}{ass:consistency}[Consistency of observed outcomes]
The observed outcome \(Y\), binary treatment indicator \(A\), and potential outcomes \(Y(0)\) and \(Y(1)\) satisfy
\[
Y=AY(1)+(1-A)Y(0)
\]
almost surely.
\end{assumptionv}

Consistency links each observed outcome to the potential outcome corresponding to the realized treatment. It is the standard condition that gives the observed treatment groups their potential-outcome interpretation \citep{KennedyBalakrishnanRobinsWasserman2024}. The next condition makes treatment assignment comparable within covariate strata.

\begin{assumptionv}{ass:exchangeability}[Conditional exchangeability of treatment]
The potential outcomes \(Y(0)\) and \(Y(1)\) are jointly independent of the treatment indicator \(A\) conditional on the covariate vector \(X\):
\[
(Y(0),Y(1))\perp A\mid X.
\]
\end{assumptionv}

Conditional exchangeability is the standard identification condition under which conditioning on \(X\) removes dependence between treatment assignment and potential outcomes \citep{KennedyBalakrishnanRobinsWasserman2024}. Together with consistency, it identifies the treated causal response from the treated conditional mean wherever the propensity is positive. Treatment availability is quantified by a global tail restriction.

\begin{assumptionv}{ass:global-tail}[Global propensity tail bound]
The propensity \(e(X)\), evaluated at the covariate vector \(X\), satisfies
\[
P\{e(X)\le t\}\le Ct^{\gamma-1}
\qquad\text{for every }t\in[0,1].
\]
\end{assumptionv}

The polynomial global propensity-tail bound controls the covariate probability assigned to regions with low treatment probability, a standard overlap restriction in this setting \citep{Dorn2026GlobalOverlap}. Since \(\gamma>1\), its value at \(t=0\) implies positive propensity almost surely under \(P_X\). Consequently, consistency and exchangeability identify \(\mu_1(x)\) with \(\mathbb E_P[Y\mid A=1,X=x]\) for \(P_X\)-almost every \(x\). The smoothness convention selects a continuous representative over the entire cube.

We next control the magnitude and fluctuations of the treated conditional outcome distribution.

\begin{assumptionv}{ass:bounded-potential-outcomes}[Treated response magnitude bound]
The observed outcome \(Y\), binary treatment indicator \(A\), and covariate vector \(X\) satisfy
\[
\operatorname*{ess\,sup}_{x\sim P_X}
\left(
 \left|\mathbb E_P[Y\mid A=1,X=x]\right|
 \vee
 \sup_{\lambda\in\mathbb R\setminus\{0\}}
 \frac{\sqrt{2\log \mathbb E_P\!\left[
   \exp\!\left\{\lambda\bigl(Y-\mathbb E_P[Y\mid A=1,X]\bigr)\right\}
   \mid A=1,X=x\right]}}{|\lambda|}
\right)
\le B.
\]
\end{assumptionv}

The uniformly bounded treated conditional mean and conditional sub-Gaussian residual condition supplies a common scale for response levels and sampling fluctuations \citep{Dorn2026GlobalOverlap}. In particular, the treated conditional mean has absolute value at most \(B\), and the centered treated residual has conditional moment-generating function bounded by \(\exp(\lambda^2B^2/2)\). These bounds hold essentially uniformly over the covariate distribution. The response curve also satisfies the following smoothness restriction.

\begin{assumptionv}{ass:holder-response}[Hölder smoothness of response]
The treated causal response \(\mu_1\) belongs to the intrinsic cube Hölder ball \(\mathcal H^\beta(L)\):
\[
\mu_1\in\mathcal H^\beta(L).
\]
\end{assumptionv}

Isotropic Hölder response smoothness imposes the same order \(\beta\) across covariate directions, with the common derivative and modulus bounds specified in \cref{obj:def:holder} \citep{Dorn2026GlobalOverlap}. It supplies the regularity through which nearby treated observations inform the target response. Combined with identification and the density lower bound, continuity extends the treated-mean magnitude bound to \(|\mu_1(x)|\le B\) throughout \(\mathcal X\).

\subsection{Law class and risk scales}

We collect the observational restrictions into a fixed-constant class. The dimension, smoothness order, and constants \(L,B,c_f,C\) remain common across its members; \(\gamma\) indexes the overlap restriction. The observational causal law class \leanref{sym:Pgamma}{\ensuremath{\mathcal P_\gamma}} is defined as follows.

\begin{definitionv}{def:model}[Observational causal law class $\mathcal P_\gamma$]
The class \(\mathcal P_\gamma\) consists of observational causal laws \(P\) satisfying the following restrictions, with common fixed constants:
\begin{itemize}
\item \textbf{(Covariate density.)} The density \(f\) of the covariate vector \(X\) satisfies \cref{obj:ass:density}.
\item \textbf{(Consistency.)} The observed outcome \(Y\), binary treatment indicator \(A\), and potential outcomes \(Y(0)\) and \(Y(1)\) satisfy \cref{obj:ass:consistency}.
\item \textbf{(Exchangeability.)} The potential outcomes and treatment satisfy \cref{obj:ass:exchangeability}.
\item \textbf{(Overlap.)} The propensity \(e(X)\) satisfies \cref{obj:ass:global-tail}.
\item \textbf{(Response bound.)} The treated conditional outcome distribution satisfies \cref{obj:ass:bounded-potential-outcomes}.
\item \textbf{(Response smoothness.)} The treated causal response \(\mu_1\) satisfies \cref{obj:ass:holder-response}.
\end{itemize}
\end{definitionv}

For each member of \cref{obj:def:model}, the sample is drawn according to \cref{obj:ass:iid}. Fix an evaluation point \leanref{sym:x0}{\ensuremath{x_0}} in the closed cube, including its boundary. We measure recovery through the pointwise minimax risk \leanref{sym:Rpt}{\ensuremath{R_{\mathrm{pt}}(n,\gamma,x_0)}} and the spatial-supremum minimax risk \leanref{sym:Rsup}{\ensuremath{R_\infty(n,\gamma)}}. Both criteria use expected absolute error, with the latter evaluating the largest error across the cube within each sample realization.

\begin{definitionv}{def:risks}[Minimax risks $R_{\mathrm{pt}}$ and $R_\infty$]
The minimax expected absolute error at \(x_0\) and the minimax expected spatial-supremum error are, respectively,
\[
R_{\mathrm{pt}}(n,\gamma,x_0)
=
\inf_T\sup_{P\in\mathcal P_\gamma}
\mathbb E_P\left|T-\mu_{1,P}(x_0)\right|
\]
and
\[
R_\infty(n,\gamma)
=
\inf_T\sup_{P\in\mathcal P_\gamma}
\mathbb E_P\left[
\sup_{x\in\mathcal X}\left|T(x)-\mu_{1,P}(x)\right|
\right].
\]
Here \(\mu_{1,P}\) is the treated response under law \(P\), and each infimum ranges over estimators measurable with respect to the observed sample \(\mathcal D_n\). The first infimum uses scalar-valued estimators, and the second uses function-valued estimators.
\end{definitionv}

The order of supremum and expectation in \cref{obj:def:risks} makes the functional criterion sensitive to the largest spatial error in each realized estimate. Its spatial supremum ranges over all points of \(\mathcal X\), using the continuous response representative fixed above.

To express the subsequent risk comparisons, define the overlap-dependent effective dimension \leanref{sym:Dgamma}{\ensuremath{D_\gamma}}, the oracle mesh scale \leanref{sym:hgamman}{\ensuremath{h_{\gamma,n}}}, and the response-recovery scale \leanref{sym:rgamman}{\ensuremath{r_{\gamma,n}}} by
\[
D_\gamma=\frac{d\gamma}{\gamma-1},
\qquad
h_{\gamma,n}=n^{-1/(2\beta+D_\gamma)},
\qquad
r_{\gamma,n}=h_{\gamma,n}^{\beta}
=n^{-\beta/(2\beta+D_\gamma)}.
\]
These deterministic scales depend on the sample size, smoothness order, dimension, and overlap exponent. They satisfy
\[
h_{\gamma,n}^{\beta}
=
\bigl(nh_{\gamma,n}^{D_\gamma}\bigr)^{-1/2},
\]
which expresses the balance between a smoothness scale and an overlap-adjusted sampling scale. The term ``oracle'' indicates that the mesh scale is indexed by \(\gamma\); it provides the benchmark for comparing risk and selecting a mesh from the observed treatment design.


\section{Main results}

The oracle scales introduced in the setup suggest selecting a resolution from the observed treatment information. We implement this choice with local polynomial fits whose conditioning is controlled by a fixed interpolation template. The construction uses the known smoothness order \(\beta>1\), with polynomial degree \(m=\lceil\beta\rceil-1\), and selects a mesh through treated observation counts.

\subsection{Construction}

The tensor nodes \leanref{sym:v_ell}{\ensuremath{v_\ell}} specify the locations around which the fitting observations are collected. Their placement leaves room for small neighborhoods inside the reference cube.

\begin{definitionv}{synth_1}[Tensor interpolation nodes]
The tensor node \(v_\ell\in\mathbb R^d\), for nonnegative integers \(d,m\) and an index vector \(\ell\in\{0,\ldots,m\}^d\), is defined by
\[
(v_\ell)_i=\frac{\ell_i+1}{m+2},
\qquad i=1,\ldots,d.
\]
\end{definitionv}

We use total-degree monomials indexed by \leanref{sym:Im}{\ensuremath{\mathcal I_m}} to form the polynomial vector \leanref{sym:U}{\ensuremath{U}}. The node count \leanref{sym:J}{\ensuremath{J}} determines the equal weights in the reference Gram matrix. The following construction fixes both the polynomial basis and the neighborhoods used to fit it.

\begin{definitionv}{synth_2}[Polynomial template and microcell width]
Fix a dimension \(d\ge1\) and a polynomial degree \(m\in\mathbb N_0\). Let
\[
\mathcal I_m:=\{\alpha\in\mathbb N_0^d:|\alpha|\le m\},
\qquad p:=|\mathcal I_m|,
\qquad U(z):=(z^\alpha)_{\alpha\in\mathcal I_m},
\quad z\in[0,1]^d,
\]
with the coordinates of \(U\) ordered lexicographically. For
\(\ell\in\{0,\ldots,m\}^d\), define
\[
v_\ell:=\left(\frac{\ell_1+1}{m+2},\ldots,
\frac{\ell_d+1}{m+2}\right),
\qquad J:=(m+1)^d,
\]
and set
\[
G_0:=\frac1J\sum_{\ell\in\{0,\ldots,m\}^d}
U(v_\ell)U(v_\ell)^\top,
\qquad \lambda_0:=\lambda_{\min}(G_0).
\]
The template Lipschitz constant is
\[
L_U:=
\sup_{\substack{z,z'\in[0,1]^d\\z\ne z'}}
\frac{\|U(z)U(z)^\top-U(z')U(z')^\top\|_{\mathrm{op}}}
{\|z-z'\|_\infty},
\]
where \(\|\cdot\|_{\mathrm{op}}\) is the matrix operator norm induced by the Euclidean vector norm. This finite constant gives the bound
\[
\|U(z)U(z)^\top-U(z')U(z')^\top\|_{\mathrm{op}}
\le L_U\|z-z'\|_\infty.
\]
Define the microcell width and reference microcubes by
\[
\eta:=\min\left\{\frac1{2(m+2)},\frac{\lambda_0}{1+L_U}\right\},
\qquad
V_\ell:=v_\ell+[-\eta/2,\eta/2]^d.
\]
In particular, for \(z\in V_\ell\), the bound entering the width rule is
\[
\|U(z)U(z)^\top-U(v_\ell)U(v_\ell)^\top\|_{\mathrm{op}}
\le \frac{L_U\eta}{2}\le\frac{\lambda_0}{2}.
\]
\end{definitionv}

The reference Gram matrix \leanref{sym:G0}{\ensuremath{G_0}} is positive definite because the tensor grid identifies every polynomial in the chosen basis. Its least eigenvalue \leanref{sym:lambda0}{\ensuremath{\lambda_0}} provides a deterministic conditioning benchmark. The matrix Lipschitz constant \leanref{sym:LU}{\ensuremath{L_U}} controls the perturbation caused by moving observations away from the nodes. Accordingly, the width \leanref{sym:eta}{\ensuremath{\eta}} makes the reference microcubes \leanref{sym:V_ell}{\ensuremath{V_\ell}} small enough to preserve that benchmark; the conditioning argument is recorded in \cref{obj:lem:template-conditioning}.

For the fitting formulas, we enumerate the tensor nodes by a single index from \(1\) to \(J\). The template notation under this enumeration is collected below.

\begin{definitionv}{def:template}[Polynomial template $G_0$ and microcubes $V_\ell$]
The equally weighted template Gram matrix \(G_0\) and its least eigenvalue \(\lambda_0\) are
\[
G_0=\frac{1}{J}\sum_{\ell=1}^{J}U(v_\ell)U(v_\ell)^\top,
\qquad
\lambda_0=\lambda_{\min}(G_0),
\]
where \(U\) is the polynomial template vector, \(v_\ell\) are its tensor nodes, and \(J\) is the number of nodes. With \(m\) denoting the polynomial order and \(L_U\) the constant entering the width rule, define the microcube width
\[
\eta=\min\left\{\frac{1}{2(m+2)},\frac{\lambda_0}{1+L_U}\right\}.
\]
The disjoint reference microcubes are
\[
V_\ell=v_\ell+[-\eta/2,\eta/2]^d,
\qquad \ell=1,\ldots,J.
\]
\end{definitionv}

Each partition cube receives a scaled copy of this same template. We search over the finite dyadic mesh collection \leanref{sym:Hn}{\ensuremath{\mathcal H_n}}, with \leanref{sym:Qh}{\ensuremath{\mathcal Q_h}} denoting the partition at mesh \(h\).

\begin{definitionv}{synth_3}[Candidate dyadic meshes]
For a sample size \(n\ge1\), dimension \(d\ge1\), and smoothness parameter \(\beta>1\), define the finite candidate mesh collection
\[
\mathcal H_n:=
\left\{2^{-j}:j\in\mathbb N_0,\ 
0\le j\le
\left\lfloor\frac{\log_2 n}{2\beta+d}\right\rfloor
\right\}.
\]
For \(h\in\mathcal H_n\), let \(\mathcal Q_h\) be the dyadic partition of
\(\mathcal X=[0,1]^d\) into cubes of side length \(h\). Each coordinate interval is half-open, with the right endpoint \(1\) included in the last interval of that coordinate. Write \(a_Q\) for the lower-left corner of \(Q\in\mathcal Q_h\).
\end{definitionv}

The corner \leanref{sym:aQ}{\ensuremath{a_Q}} supplies the translation for placing the reference template in a partition cube. The resulting microcells \leanref{sym:EQell}{\ensuremath{E_{Q\ell}}} have treated observation counts \leanref{sym:NQell}{\ensuremath{N_{Q\ell}}}. We require each count to reach the smoothness-dependent threshold \(h^{-2\beta}\). The feasible mesh set \leanref{sym:Fh}{\ensuremath{\mathcal F_n}} collects the resolutions meeting this requirement, and the selected width \leanref{sym:hhat}{\ensuremath{\widehat h}} is the finest feasible resolution, with the exceptional-sample convention specified below.

\begin{algorithmv}{def:mesh-selector}[Count-feasible mesh $\widehat h$]
Given the candidate meshes \(\mathcal H_n\), the observed sample \(\mathcal D_n=((X_i,A_i,Y_i))_{i=1}^n\), the smoothness parameter \(\beta\), and the reference microcubes \(V_\ell\), proceed as follows.
\begin{enumerate}
\item For each \(h\in\mathcal H_n\), form the cube partition \(\mathcal Q_h\). For each partition cube \(Q\), with corner \(a_Q\), define its microcells and treated observation counts by
\[
E_{Q\ell}=a_Q+hV_\ell,
\qquad
N_{Q\ell}=\sum_{i=1}^n\mathbf 1_{\{X_i\in E_{Q\ell},\,A_i=1\}}.
\]
\item Define the set of feasible meshes by
\[
\mathcal F_n
=
\left\{h\in\mathcal H_n:
\min_{Q\in\mathcal Q_h,\,1\le\ell\le J}N_{Q\ell}
\ge h^{-2\beta}\right\},
\]
where \(J\) is the number of reference microcubes.
\item Output the numerically smallest feasible mesh, with zero assigned when the feasible set is empty:
\[
\widehat h=
\begin{cases}
\min\mathcal F_n,&\mathcal F_n\ne\varnothing,\\
0,&\mathcal F_n=\varnothing.
\end{cases}
\]
\end{enumerate}
\end{algorithmv}

Selection depends on the sampled covariates and treatment indicators. The threshold calibrates the amount of treatment information to the approximation scale \(h^\beta\). At a feasible mesh, each microcell receives the same total fitting weight, so the template geometry governs the empirical matrix even when treated observations are distributed unevenly across microcells.

The empirical equal-cell Gram matrix \leanref{sym:GhatQ}{\ensuremath{\widehat G_Q}} and coefficient estimate \leanref{sym:chatQ}{\ensuremath{\widehat c_Q}} make this weighting explicit.

\begin{definitionv}{synth_4}[Equal-cell local polynomial coefficients]
Let \(\mathcal D_n=((X_i,A_i,Y_i))_{i=1}^n\) be the observed sample. Fix a feasible mesh \(h\), a partition cube \(Q\in\mathcal Q_h\) with lower-left corner \(a_Q\), the polynomial template vector \(U\), and its \(J\) reference microcubes \(V_\ell\). Define
\[
E_{Q\ell}:=a_Q+hV_\ell,
\qquad
N_{Q\ell}:=\sum_{i=1}^n
\mathbf 1_{\{A_i=1,\ X_i\in E_{Q\ell}\}}.
\]
Feasibility means that \(N_{Q\ell}\ge h^{-2\beta}\) for every partition cube and template microcell, so these counts are positive. The equal-cell empirical Gram matrix is
\[
\widehat G_Q:=
\frac1J\sum_\ell\frac1{N_{Q\ell}}
\sum_{\substack{1\le i\le n\\ A_i=1,\ X_i\in E_{Q\ell}}}
U\!\left(\frac{X_i-a_Q}{h}\right)
U\!\left(\frac{X_i-a_Q}{h}\right)^\top.
\]
Using the polynomial template and its microcell width, this matrix is invertible at every feasible mesh. Define the equal-cell local polynomial coefficient estimate by
\[
\widehat c_Q:=
\widehat G_Q^{-1}
\frac1J\sum_\ell\frac1{N_{Q\ell}}
\sum_{\substack{1\le i\le n\\ A_i=1,\ X_i\in E_{Q\ell}}}
U\!\left(\frac{X_i-a_Q}{h}\right)Y_i.
\]
Thus each microcell receives weight \(1/J\), and the treated observations within that microcell receive equal weights.
\end{definitionv}

Invertibility follows from the deterministic template perturbation bound whenever the count requirement is met. The fitted polynomial is evaluated throughout its partition cube, including points outside the fitting microcells. Projection onto the response-bound interval completes the response estimator \leanref{sym:muhat}{\ensuremath{\widehat\mu_n}} and controls its magnitude on every sample realization.

\begin{algorithmv}{def:estimator}[Clipped response estimator $\widehat\mu_n$]
Given the selected mesh \(\widehat h\), the polynomial template vector \(U\), and the response bound \(B\), proceed as follows.
\begin{enumerate}
\item If \(\widehat h=0\), output
\[
\widehat\mu_n(x)=0
\qquad(x\in\mathcal X).
\]
\item Otherwise, set
\[
h=\widehat h.
\]
For every \(Q\in\mathcal Q_h\), compute the empirical equal-cell Gram matrix \(\widehat G_Q\) and the equal-cell polynomial coefficient estimate \(\widehat c_Q\).
\item For \(x\in\mathcal X\), let \(Q\) be the unique half-open partition cube containing \(x\), and let \(a_Q\) be its corner. Output
\[
\widehat\mu_n(x)
=
\operatorname{clip}_{[-B,B]}
\left\{
U\left(\frac{x-a_Q}{h}\right)^\top\widehat c_Q
\right\},
\qquad
\operatorname{clip}_{[-B,B]}(t)=\max\{-B,\min\{B,t\}\}.
\]
\end{enumerate}
\end{algorithmv}

\subsection{Recovery guarantee}

The count-selected estimator attains the overlap-dependent oracle rate uniformly over a compact range of overlap exponents. A matching converse holds within the model's bounded, two-valued potential-outcome subclass.

\begin{theoremv}{thm:adaptive-minimax}[Adaptive minimax response recovery]
Fix the model parameters subject to the following conditions:
\begin{itemize}
\item \textbf{(Dimension and smoothness.)} \(d\in\mathbb N\), \(d>0\), and \(\beta>1\).
\item \textbf{(Response and smoothness bounds.)} \(B>0\) and \(L>0\).
\item \textbf{(Overlap and density constants.)} \(C\ge1\) and \(0<c_f\le1\).
\item \textbf{(Adaptation interval.)} The real endpoints satisfy \(1<\gamma_{\min}<\gamma_{\max}\), and the adaptation interval is
\[
\Gamma=[\gamma_{\min},\gamma_{\max}].
\]
\end{itemize}
Let \(\mathcal X=[0,1]^d\) be the covariate cube, and let \(\mathcal P_\gamma\) be the observational causal model of \cref{obj:def:model}, with fixed parameters \(d,\beta,B,L,C,c_f,\gamma\), including the propensity-tail envelope in \cref{obj:ass:global-tail}. Define the effective dimension \(D_\gamma\), oracle mesh \(h_{\gamma,n}\), and oracle response-recovery scale \(r_{\gamma,n}\), for \(n\ge1\), by
\[
D_\gamma=\frac{d\gamma}{\gamma-1},
\qquad
h_{\gamma,n}=n^{-1/(2\beta+D_\gamma)},
\qquad
r_{\gamma,n}=h_{\gamma,n}^{\beta}.
\]

The count-selected equal-cell estimator \(\widehat\mu_n\) of \cref{obj:def:estimator} takes \(n,d,\beta,B\) and the observed sample \(\mathcal D_n\) as inputs. There exist \(K>0\) and \(n_0\in\mathbb N\) such that, for every integer \(n\ge\max\{n_0,1\}\), every \(\gamma\in\Gamma\), and every \(P\in\mathcal P_\gamma\),
\[
\mathbb E_{P^{\otimes n}}
\left[
\sup_{x\in\mathcal X}
\left|\widehat\mu_n(x)-\mu_{1,P}(x)\right|
\right]
\le K r_{\gamma,n}.
\]
Here \(P^{\otimes n}\) is the sampling law of \(n\) independent observed triples \((X,A,Y)\), and \(\mu_{1,P}\) is the treated causal response under \(P\).

For every real \(\gamma>1\) and every \(x_0\in\mathcal X\), there exist finite constants \(0<c_\gamma\le K_\gamma\), possibly depending on \(x_0\) and the fixed model parameters, such that, for every integer \(n\ge1\),
\[
\begin{aligned}
c_\gamma r_{\gamma,n}
&\le
\inf_T
\sup_{\substack{P\in\mathcal P_\gamma\\
P\text{ admits a causal completion with}\\
Y(0),Y(1)\in\{0,B\}\ \text{a.s.}}}
\mathbb E_{P^{\otimes n}}
\left|T(\mathcal D_n,x_0)-\mu_{1,P}(x_0)\right|
\\
&\le R_{\mathrm{pt}}(n,\gamma,x_0)
\le R_\infty(n,\gamma)
\le K_\gamma r_{\gamma,n}.
\end{aligned}
\]
The infimum ranges over sample-measurable estimators \(T\); the pointwise and expected supremum minimax risks \(R_{\mathrm{pt}}\) and \(R_\infty\) are those of \cref{obj:def:risks}.
\end{theoremv}

The rate reflects a balance between polynomial approximation and the treatment information available at a given resolution. Smoothness supplies the approximation scale \(h^\beta\), while the global overlap tail determines the effective dimension \(D_\gamma\). At the oracle mesh these scales satisfy
\[
n h_{\gamma,n}^{D_\gamma}
=
h_{\gamma,n}^{-2\beta}.
\]
The count threshold implements this balance through the realized treatment design, and the fixed template turns feasible counts into stable polynomial fits.

Adaptation in \cref{obj:thm:adaptive-minimax} concerns the overlap exponent at a known smoothness order. For fixed \(d,\beta,B,L,C,c_f\) and a fixed compact interval \(\Gamma\subset(1,\infty)\) with distinct endpoints, one estimator and common constants \(K,n_0\) provide the upper bound for every \(\gamma\in\Gamma\). Its inputs are the observed sample and \(n,d,\beta,B\). Thus the overlap exponent is accommodated through count selection, while \(\beta\) determines the polynomial degree, candidate resolutions, and feasibility threshold.

The upper bound controls expected spatial-supremum error: the largest error over the closed covariate cube is evaluated within each sample realization and then averaged over samples. This criterion matches \cref{obj:def:risks} and covers boundary points as well as interior points. For each fixed \(\gamma>1\), the converse establishes the same rate at every fixed evaluation point, already among laws admitting potential outcomes in \(\{0,B\}\). Consequently, both pointwise recovery and recovery under expected spatial-supremum loss have minimax scale
\[
r_{\gamma,n}
=
n^{-\beta/(2\beta+D_\gamma)}.
\]
The constants in the fixed-exponent comparison may depend on the evaluation point and fixed model parameters, as specified in \cref{obj:thm:adaptive-minimax}.

\subsection{Statistical explanation}

Two features of the construction account for the guarantee. First, equal-cell weighting separates polynomial conditioning from the distribution of treated observations within a cube. Every microcell contributes total weight \(1/J\), and every normalized fitting location remains near its reference node. The empirical Gram matrix therefore stays close to the same positive-definite template at every feasible resolution. \Cref{obj:lem:template-conditioning} supplies this deterministic step, while \cref{obj:lem:selected-mesh} connects count feasibility to the oracle resolution.

Second, the global overlap restriction controls the spatial distribution of scarce treatment information. To describe this distribution, fix a law \(P\) and a candidate mesh \(h\). Let the population treated microcell mass \leanref{sym:qQell}{\ensuremath{q_{Q\ell}}} and the weakest such mass in a partition cube \leanref{sym:qQ}{\ensuremath{q_Q}} be
\[
q_{Q\ell}
=
P\{X\in E_{Q\ell},\,A=1\},
\qquad
q_Q=\min_{1\le\ell\le J}q_{Q\ell}.
\]
Write \leanref{sym:qorder}{\ensuremath{q_{(k)}}} for these weakest-cell masses arranged in ascending order, with \(k=1,\ldots,|\mathcal Q_h|\). This ordering records how treatment information improves as one moves beyond the least-informed cubes.

The density lower bound in \cref{obj:ass:density} assigns covariate probability to each microcell, and the global tail envelope in \cref{obj:ass:global-tail} limits how much of that probability can carry very small propensities. Their combination yields the ordered treatment-information bound
\[
q_{(k)}\ge \kappa_\Gamma h^{D_\gamma}k^{1/(\gamma-1)},
\qquad
D_\gamma=\frac{d\gamma}{\gamma-1},
\]
uniformly over the compact overlap range; this is \cref{obj:lem:ordered-mass}. On the count-concentration event of \cref{obj:lem:selected-mesh}, the conditional variance proxy for every fitted coefficient in the cube of rank \(k\) is bounded, up to a common constant, by
\[
\min\!\left\{
\widehat h^{2\beta},
n^{-1}\widehat h^{-D_\gamma}k^{-1/(\gamma-1)}
\right\}.
\]
The first term is the feasibility cap; the second decreases with rank according to the ordered-mass profile. Their equality marks the ranks at which the envelope leaves the common cap. Applying the ordered maximal inequality across these ranks and comparing the selected mesh with the oracle mesh gives
\[
\mathbb E\!\left[
\max_{Q\in\mathcal Q_{\widehat h}}
\left\|\widehat c_Q-
\mathbb E[\widehat c_Q\mid(X_i,A_i)_{i=1}^n]\right\|_\infty
\,\middle|\,(X_i,A_i)_{i=1}^n
\right]
\le K\widehat h^\beta
\sqrt{1+\max\{0,\log(h_{\gamma,n}/\widehat h)\}}
\le K'r_{\gamma,n}.
\]
Thus the few least-informed cubes are controlled by feasibility, while the improving ranks pay progressively smaller stochastic scales. Deterministic template conditioning then transfers this coefficient bound to expected spatial-supremum response error at the oracle rate. The appendix proves the concentration event and the maximal inequality in full.

The appendix develops the corresponding probability bounds and the bounded testing experiment. Their roles are complementary: the selected-mesh analysis establishes the estimator's recovery guarantee, and \cref{obj:lem:bounded-lower-pair} supplies the bounded-response comparison underlying the minimax converse.


\section{Local treatment geometry and the Gaussian transfer}

\subsection{Global tail control and local treatment geometry}

The recovery guarantee in \cref{obj:thm:adaptive-minimax} rests on global tail control and the equal-cell construction. To clarify the geometric scope of this guarantee, we compare the model with the local anti-concentration condition \leanref{sym:A3Dorn}{\ensuremath{\mathsf A_3^{\mathrm{Dorn}}}}. This condition relates a neighborhood's largest selected propensity to the covariate probability assigned to propensities of comparable size. Its formulation fixes both the family of laws and the pointwise propensity representatives.

\begin{definitionv}{def:dorn-a3}[Local anti-concentration condition $\mathsf A_3^{\mathrm{Dorn}}$]
The family-level condition
\(\mathsf A_3^{\mathrm{Dorn}}(\mathfrak P,(e_P)_{P\in\mathfrak P})\)
is defined for a family and selected propensity versions satisfying:
\begin{itemize}
\item \textbf{(Family.)} \(\mathfrak P\) is a nonempty family of probability laws with common covariate cube \(\mathcal S\).
\item \textbf{(Versions.)} For each \(P\in\mathfrak P\), \(e_P:\mathcal S\to[0,1]\) is a fixed, pointwise-defined measurable version of \(P(A=1\mid X=\cdot)\).
\item \textbf{(Overlap.)} For each \(P\in\mathfrak P\), \(e_P(X)\in(0,1)\) \(P\)-almost surely.
\end{itemize}
The condition holds when there exist constants \(\rho,\nu>0\) and \(h_0>0\), common to the family, such that, for every \(P\in\mathfrak P\), every \(x_0\in\mathcal S\), and every \(0<h\le h_0\),
\[
P_P\left(
e_P(X)\ge
\rho\sup_{\substack{x\in\mathcal S\\\|x-x_0\|_2\le h}}e_P(x)
\,\middle|\,
\|X-x_0\|_2\le h
\right)>\nu.
\]
Here \(P_P\) denotes probability under \(P\), and the supremum is the pointwise supremum of the selected version. Thus the condition is a property of the family together with its selected propensity versions; its truth can depend on their values on \(P_X\)-null sets. The covariate cube is \(\mathcal S=[-1,1]^d\) in Dorn's source and \(\mathcal S=\mathcal X\) after affine transport. For a single displayed law and propensity, \(\mathsf A_3^{\mathrm{Dorn}}(P)\) abbreviates the condition for the singleton family equipped with that selected version.
\end{definitionv}

The common constants in \cref{obj:def:dorn-a3} require this comparison to persist across laws, locations, and sufficiently small neighborhoods. The pointwise supremum also makes the selected representative part of the condition. By contrast, the global tail restriction in \cref{obj:ass:global-tail} controls the probability of small propensity values under the covariate distribution. A thin-strip construction makes the distinction concrete: it combines a radial propensity component with additional treatment probability concentrated in a narrowing strip.

\begin{definitionv}{def:thin-strip}[Thin-strip law $P_{\mathrm{strip}}$]
The observed-data law \(P_{\mathrm{strip}}\) is defined by the following construction, under these conditions:
\begin{itemize}
\item \textbf{(Dimension.)} \(d\ge2\).
\item \textbf{(Covariates.)} \(X\) is uniform on \(\mathcal X\).
\item \textbf{(Strip exponent.)} Choose \(a\) satisfying
\[
1<a<1+\frac{d}{\gamma-1}.
\]
\end{itemize}
Define the cube center, radial distance, and its continuous distribution function by
\[
x^\circ=(1/2,\ldots,1/2),
\qquad
R(x)=\|x-x^\circ\|_\infty,
\qquad
F_R(t)=P(R(X)\le t).
\]
Define the strip and selected propensity by
\[
S=\left\{x\in\mathcal X:
|x_2-x_2^\circ|\le |x_1-x_1^\circ|^a\right\},
\qquad
e_\star(x)=
\min\left\{1,F_R(R(x))^{1/(\gamma-1)}
+\frac14\mathbf 1_S(x)\right\}.
\]
Given \(X\), draw
\[
A\sim\operatorname{Bernoulli}(e_\star(X)).
\]
Set
\[
p_0=\min\left\{\frac14,\frac{L}{4B}\right\}.
\]
Draw \(Y(0)\) and \(Y(1)\) conditionally independently given \(X\), each with conditional distribution \(B\operatorname{Bernoulli}(p_0)\), and define
\[
Y=AY(1)+(1-A)Y(0).
\]
Then \(P_{\mathrm{strip}}\) is the law of \((X,A,Y)\).
\end{definitionv}

The thin-strip law \leanref{sym:Pstrip}{\ensuremath{P_{\mathrm{strip}}}} of \cref{obj:def:thin-strip} has constant potential-outcome means and bounded outcomes. It therefore isolates the treatment geometry while keeping the response smooth. The strip exponent makes the strip narrow relative to the surrounding neighborhood as its radius decreases. Under the displayed exponent range, \cref{obj:prop:strict-enlargement} establishes the coexistence of global tail control and a vanishing normalized treated-design moment.

\begin{propositionv}{prop:strict-enlargement}[Thin-strip enlargement and design degeneration]
Suppose the parameters satisfy:
\begin{itemize}
\item \textbf{(Dimension and smoothness.)} \(d\in\mathbb N\), \(d\ge2\), and \(\beta>1\).
\item \textbf{(Model constants.)} \(B>0\), \(L>0\), \(C\ge1\), and \(0<c_f\le1\).
\item \textbf{(Overlap and strip exponent.)} \(\gamma>1\) and
\[
1<a<1+\frac{d}{\gamma-1}.
\]
\end{itemize}
For every such \(a\), let \(P_{\mathrm{strip}}\) and its displayed selected propensity \(e_\star\) be the construction in \cref{obj:def:thin-strip}, and let \(\mathcal P_\gamma\) denote the model class with the fixed parameters \((d,\beta,B,L,C,c_f,\gamma)\). Then \(P_{\mathrm{strip}}\in\mathcal P_\gamma\), and
\[
P_{\mathrm{strip}}\{e_\star(X)\le t\}\le t^{\gamma-1},
\qquad t\in[0,1].
\]

On the unit cube \(\mathcal X=[0,1]^d\), the numerical anti-concentration clause fails: there are no \(\rho,\nu,h_0>0\) such that, for every \(x_0\in\mathcal X\) and every \(0<h\le h_0\),
\[
\frac{
P_{\mathrm{strip}}\!\left\{
e_\star(X)\ge
\rho\sup_{\substack{x\in\mathcal X\\ \|x-x_0\|_2\le h}}e_\star(x),
\ \|X-x_0\|_2\le h
\right\}
}{
P_{\mathrm{strip}}\{\|X-x_0\|_2\le h\}
}
>\nu.
\]
This failure applies to the numerical clause allowing propensity values equal to one. Consequently, the full singleton-family predicate
\(\mathsf A_{3}^{\mathrm{Dorn}}(P_{\mathrm{strip}})\), evaluated using \(e_\star\), also fails.

Moreover, writing \(x^\circ=(1/2,\ldots,1/2)\), the normalized treated-design second moment in the second coordinate satisfies
\[
\lim_{h\downarrow0}
\frac{
\displaystyle
\int_{\{A=1,\ \|X-x^\circ\|_\infty\le h\}}
\left(\frac{X_2-1/2}{h}\right)^2
\,dP_{\mathrm{strip}}
}{
P_{\mathrm{strip}}\{A=1,\ \|X-x^\circ\|_\infty\le h\}
}
=0.
\]
In particular, \(\mathcal P_\gamma\) contains a law whose full singleton-family Dorn predicate fails for the displayed selected propensity \(e_\star\).
\end{propositionv}

The global tail inequality controls the prevalence of small propensities, while the displayed moment records the shape of the treated design near the cube center. Its limit shows that, after normalization by neighborhood width, treated observations concentrate increasingly close to the strip's central hyperplane. The local second-coordinate variation consequently degenerates even though the law satisfies the global model restrictions. For dimensions at least two, \cref{obj:prop:strict-enlargement} thus establishes a strict enlargement relative to imposing the selected-version predicate within the causal class. The equal-cell estimator's guarantee in \cref{obj:thm:adaptive-minimax} covers this geometry through its template and count-feasibility conditions.

\subsection{Transfer to fixed-constant Gaussian source envelopes}

The geometric distinction also has a source-model consequence. Under the retained restrictions in \citet[Section 2, Assumptions 1--2, pp.~3--7]{Dorn2026GlobalOverlap}, affine transport maps each fixed-constant Gaussian source envelope to the causal class used by \cref{obj:thm:adaptive-minimax}. A common response bound controls the transported mean and noise scale, while fixed-cube smoothness completion supplies a common full Hölder radius. Consequently, one transformed equal-cell estimator attains the expected spatial-supremum rate uniformly over the source envelope and the prescribed compact exponent range. The exact envelope, transport constants, exponent-specific Gaussian witness, and family-richness qualifications are stated in \cref{obj:prop:dorn-a3-free-corollary} and proved in the source-comparison appendix.



The transfer aligns the source restrictions with the components needed for response recovery. Uniform covariates supply the density lower bound after transport, the propensity tail inequality supplies global overlap control, and the Gaussian and treated-moment restrictions supply common response bounds. Fixed-cube smoothness completion converts the source's top-derivative seminorm restriction, together with a function bound, into the full Hölder radius used by the estimator. The precise completion argument is given in \cref{obj:lem:fixed-cube-holder-completion}.

The upper guarantee controls the expected spatial supremum over the entire unit cube, uniformly over the fixed-constant envelopes and the compact exponent range. This loss criterion is stronger than the displayed essential-supremum probability criterion: it controls the magnitude of the largest spatial error in expectation and yields the stated probability guarantee at every positive tolerance. The same constants apply to each source subfamily satisfying the retained restrictions because each such family lies within the corresponding envelope.

The Gaussian existence assertion has a separate quantifier structure. In dimensions at least two, it supplies, for each overlap exponent, a source tuple satisfying the retained restrictions with positive finite constants chosen for that exponent and with local anti-concentration failure for its selected representative. These witnesses establish the geometric scope of the retained Gaussian restrictions; the envelope-wide upper guarantee applies with the fixed constants specified at the start of \cref{obj:prop:dorn-a3-free-corollary}.

The converse statements likewise identify their comparison classes explicitly. A source family with a common treated curve has zero minimax pointwise miss probability, whereas the full causal class in \cref{obj:def:risks} admits the displayed positive expected absolute-error lower bound at the oracle scale, supported by its two-valued potential-outcome subclass. Family richness therefore determines the scope of a minimax converse, while the transferred upper guarantee remains valid for every qualifying source subfamily. The technical source-family comparison is developed in \cref{obj:lem:arbitrary-family-lower-obstruction,obj:lem:dorn-rate-scope}.


\section{Discussion}

The guarantees in \cref{obj:thm:adaptive-minimax} describe response recovery with known smoothness and adaptation to overlap. The known smoothness parameter \(\beta\) determines the polynomial degree and the approximation order used in constructing the estimator. Treatment counts then select a feasible mesh through \cref{obj:def:mesh-selector}. Thus response regularity supplies the approximation benchmark, while the observed treatment design determines the spatial resolution.

The compact range for the overlap parameter \(\gamma\), contained in \((1,\infty)\), specifies the domain over which the risk constants are uniform. With the remaining model constants fixed, a single estimator attains the corresponding oracle scale throughout that range. The pointwise and expected spatial-supremum guarantees therefore share an overlap-dependent scale, including the spatial-supremum guarantee without an additional logarithmic factor.

For comparison, the classical regular-design spatial-supremum benchmark includes a logarithmic factor, with scale \((\log n/n)^{\beta/(2\beta+d)}\) under its regular coverage conditions \citep{Stone1982}. The present theorem instead takes a worst case over the larger global-tail class at each fixed finite \(\gamma\), giving the pure power \(n^{-\beta/(2\beta+d\gamma/(\gamma-1))}\). The theorem supplies one estimator and common risk constants on every prescribed compact interval of finite overlap exponents.

Equal-cell weighting connects the deterministic geometry in \cref{obj:def:template} to the fitted polynomials in \cref{obj:def:estimator}. Each microcell contributes the same total weight, so the conditioning argument is governed by the template geometry, while differences in treatment counts govern stochastic precision; \cref{obj:lem:template-conditioning,obj:lem:ordered-mass} give the two bounds.

Conditional sub-Gaussian residuals supply exponential control of fluctuations after conditioning on the sampled covariates and treatment indicators. That control supports the upper-bound argument and its Gaussian-envelope transfer in \cref{obj:prop:dorn-a3-free-corollary}. Together, the response smoothness, template geometry, and treatment-count information determine approximation, conditioning, and stochastic error, respectively.

\subsection{Limitations and future work}

Joint adaptation would replace the known smoothness benchmark with a comparison across response regularities as well as treatment-count-feasible resolutions. The following proposal describes an architecture for that extension.
\begin{remarkv}{def:joint-adaptation-handle}[Future work: joint adaptation in $\beta$ and $\gamma$]
A proposed direction for future work is to construct a single response estimator with an adaptive-risk envelope over compact rectangles of the smoothness parameter \(\beta\) and the overlap parameter \(\gamma\). The proposed architecture uses sample splitting to compare equal-cell fits across polynomial degrees and treatment-count-feasible meshes, with the count selector controlling adaptation to overlap and held-out Lepski inequalities controlling adaptation to smoothness. The sample split, degree-and-mesh comparison rule, Lepski threshold, estimator, and envelope functional remain to be specified; estimator existence and the corresponding risk bounds remain open.
\end{remarkv}

The proposed split would allocate separate observations to fitting and comparison, while the degree-and-mesh rule would compare candidates along both parameter axes. Lepski comparisons provide a natural perspective on smoothness selection \citep{lepski1997}; smoothness adaptation under degenerate random design, studied by \citet{Gaiffas2005AdaptiveDegenerateDesign}, supplies a related setting in which to interpret the adaptation question. These comparisons motivate the proposal rather than establish its risk envelope.

The absence of an additional spatial-supremum logarithm in the known-smoothness guarantee makes the cost of the proposed two-axis comparison particularly relevant. The question concerns both attainability of the oracle scale and any unavoidable adaptation penalty.
\begin{remarkv}{oeq:joint-adaptation}[Joint smoothness and overlap adaptation]
For a compact rectangle of smoothness and overlap parameters \((\beta,\gamma)\) contained in \((1,\infty)^2\), can a procedure adapting along both parameter axes produce a single estimator whose pointwise and expected spatial-supremum risks attain the oracle scale for every pair in the rectangle? What logarithmic penalty, if any, is minimax unavoidable?
\end{remarkv}

Resolving \cref{obj:oeq:joint-adaptation} requires specifying the procedure and establishing its risk bounds, together with an appropriate converse for any claimed unavoidable logarithmic price. The existing common-class bounded-response converse and Gaussian-envelope upper transfer address distinct comparison domains: the former supplies the minimax lower bound in \cref{obj:thm:adaptive-minimax}, while the latter supplies the upper guarantee in \cref{obj:prop:dorn-a3-free-corollary}. The transferred upper bound does not furnish a matching converse for every Gaussian source subfamily. As \cref{obj:lem:arbitrary-family-lower-obstruction,obj:lem:dorn-rate-scope} clarify, a lower-bound analysis must account for family richness. These distinctions also govern how a future joint-adaptation result could characterize an unavoidable logarithmic penalty.


\appendix

\section*{Appendices}

\section{Identification and upper-bound proofs}

The upper bound in \cref{obj:thm:adaptive-minimax} combines causal identification with control of approximation error, conditional stochastic error, and exceptional samples. The auxiliary statements below follow that order: identification fixes the population target, cube smoothness admits a controlled global extension, the global overlap tail supplies an ordered information profile, and the template and mesh bounds translate that profile into stable response estimation.

\subsection{Identification of the continuous response}

Under the causal conditions in \cref{obj:def:model}, the treated conditional mean identifies the causal response almost everywhere. The density and smoothness restrictions then determine its continuous representative throughout the covariate cube. The following statement records both levels of identification.

\begin{lemmav}{lem:causal-identification}[Identification of the treated response]
Let \(P\in\mathcal P_\gamma\) satisfy the parameter restrictions and causal model conditions of \cref{obj:def:model}, with \(\gamma>1\). Write \(P_X\) for its covariate marginal and \(\mathcal X=[0,1]^d\) for the covariate cube. Then
\[
P_X\{x:e(x)=0\}=0,
\qquad
\mu_1(x)=\mathbb E_P[Y\mid A=1,X=x]
\quad\text{for }P_X\text{-almost every }x.
\]
Moreover, every function \(g:\mathbb R^d\to\mathbb R\) that is continuous on \(\mathcal X\) and satisfies
\[
g(x)=\mathbb E_P[Y\mid A=1,X=x]
\quad\text{for }P_X\text{-almost every }x
\]
obeys
\[
g(x)=\mu_1(x)
\qquad\text{for every }x\in\mathcal X.
\]
\end{lemmav}

\begin{proof}[Proof of \cref{obj:lem:causal-identification}]
\begin{enumerate}
\item Let \(P_c\) be the joint causal law on
\(\mathbb R^d\times\{0,1\}\times\mathbb R^3\):
\[
P_c:=\mathcal L\bigl(X,A,Y(0),Y(1),Y\bigr).
\]
Its observed marginal is \(P\), and its covariate marginal is \(P_X\). Define the joint conditioning law and conditional probability kernels by
\[
\begin{aligned}
\pi_{X,A}&:=\mathcal L_{P_c}(X,A),\\
K^c_{A\mid X}(x)&:=\mathcal L_{P_c}(A\mid X=x),\\
K_{1\mid X}(x)&:=\mathcal L_{P_c}(Y(1)\mid X=x),\\
K^c_{Y\mid X,A}(x,a)&:=\mathcal L_{P_c}(Y\mid X=x,A=a),\\
K^c_{1\mid X,A}(x,a)&:=\mathcal L_{P_c}(Y(1)\mid X=x,A=a),
\end{aligned}
\qquad x\in\mathbb R^d,\quad a\in\{0,1\}.
\]
Choose measurable versions defined at every conditioning input, with arbitrary probability values on exceptional conditioning sets.

By \cref{obj:ass:consistency},
\[
A=1\ \Longrightarrow\ Y=Y(1)
\qquad P_c\text{-almost surely}.
\]
Conditioning this equality on \((X,A)\) gives
\[
a=1\ \Longrightarrow\
K^c_{Y\mid X,A}(x,a)=K^c_{1\mid X,A}(x,a)
\qquad \pi_{X,A}\text{-almost everywhere}.
\]
By \cref{obj:ass:exchangeability}, \(Y(1)\) and \(A\) are conditionally independent given \(X\). Consequently,
\[
K^c_{1\mid X,A}(x,a)=K_{1\mid X}(x)
\qquad \pi_{X,A}\text{-almost everywhere}.
\]
% lean: CausalSmith.Stat.WeakOverlap.mu1_eq_treatedRegression (hfiber, hci, hexch)

\item Applying \cref{obj:ass:global-tail} at \(t=0\), and using \(\gamma-1>0\), yields
\[
0\le P_X\{e\le0\}\le C\,0^{\gamma-1}=0.
\]
Since \(\{e=0\}\subseteq\{e\le0\}\),
\[
P_X\{e=0\}=0.
\]
To express positivity through the conditioning kernel, define
\[
K^o_{A\mid X}(x):=\mathcal L_P(A\mid X=x),
\qquad x\in\mathbb R^d,
\]
again as a measurable probability kernel defined everywhere. The observed marginal relation and the definition of the propensity give
\[
K^o_{A\mid X}(x)=K^c_{A\mid X}(x),
\qquad
e(x)=K^o_{A\mid X}(x)(\{1\})\ge0
\qquad P_X\text{-almost everywhere}.
\]
Combining these identities with the null-set calculation proves
\[
K^c_{A\mid X}(x)(\{1\})=e(x)>0
\qquad P_X\text{-almost everywhere}.
\]
% lean: CausalSmith.Stat.WeakOverlap.positivePropensity_ae

\item Disintegration of the joint conditioning law gives
\[
\pi_{X,A}(dx,da)=P_X(dx)\,K^c_{A\mid X}(x,da).
\]
We use the following consequence of this identity. For any predicate
\(p:\mathbb R^d\times\{0,1\}\to\{\text{true},\text{false}\}\) holding
\(\pi_{X,A}\)-almost everywhere, disintegration implies
\[
K^c_{A\mid X}(x)\{a:\neg p(x,a)\}=0
\qquad P_X\text{-almost everywhere}.
\]
At any such \(x\) with \(K^c_{A\mid X}(x)(\{1\})>0\), failure of \(p(x,1)\) would imply
\[
0<K^c_{A\mid X}(x)(\{1\})
\le K^c_{A\mid X}(x)\{a:\neg p(x,a)\}=0.
\]
Thus \(p(x,1)\) holds \(P_X\)-almost everywhere. Applying this argument to the conjunction of the two conditional-kernel identities above gives
\[
K^c_{Y\mid X,A}(x,1)=K_{1\mid X}(x)
\qquad P_X\text{-almost everywhere}.
\]

Now define the observed outcome kernel and its treated slice by
\[
\begin{aligned}
K^o_{Y\mid X,A}(x,a)&:=\mathcal L_P(Y\mid X=x,A=a),\\
K_{\mathrm{tr}}(x)&:=K^o_{Y\mid X,A}(x,1),
\end{aligned}
\qquad x\in\mathbb R^d,\quad a\in\{0,1\}.
\]
The observed marginal relation gives
\[
K^o_{Y\mid X,A}(x,a)=K^c_{Y\mid X,A}(x,a)
\qquad \pi_{X,A}\text{-almost everywhere}.
\]
Applying the same positive-slice argument once more, we obtain
\[
K_{\mathrm{tr}}(x)
=K^c_{Y\mid X,A}(x,1)
=K_{1\mid X}(x)
\qquad P_X\text{-almost everywhere}.
\]
% lean: CausalSmith.Stat.WeakOverlap.ae_slice_of_ae_compProd_of_pos

\item Define the treated regression on all of \(\mathbb R^d\) by
\[
r_{\mathrm{tr}}(x):=\int_{\mathbb R}y\,K_{\mathrm{tr}}(x,dy),
\]
using value zero wherever the integrand is not integrable. By
\cref{obj:ass:bounded-potential-outcomes} and the kernel equality just proved, for \(P_X\)-almost every \(x\),
\[
|r_{\mathrm{tr}}(x)|\le B,
\qquad
\int_{\mathbb R}
e^{\lambda(y-r_{\mathrm{tr}}(x))}\,K_{1\mid X}(x,dy)
\le e^{B^2\lambda^2/2}
\quad\text{for every }\lambda\in\mathbb R,
\]
and each exponential integrand in this display is integrable.

The choices \(\lambda=1\) and \(\lambda=-1\) control the absolute first moment. Indeed, for every \(y\in\mathbb R\),
\[
1+|y-r_{\mathrm{tr}}(x)|
\le e^{|y-r_{\mathrm{tr}}(x)|}
\le e^{y-r_{\mathrm{tr}}(x)}
   +e^{-(y-r_{\mathrm{tr}}(x))},
\]
so
\[
|y|
\le |y-r_{\mathrm{tr}}(x)|+|r_{\mathrm{tr}}(x)|
\le e^{y-r_{\mathrm{tr}}(x)}
   +e^{-(y-r_{\mathrm{tr}}(x))}
   +|r_{\mathrm{tr}}(x)|.
\]
The right-hand side is integrable under \(K_{1\mid X}(x)\). Integrating the inequality and using the exponential bounds yields
\[
\begin{aligned}
\int_{\mathbb R}|y|\,K_{1\mid X}(x,dy)
&\le
\int_{\mathbb R}e^{y-r_{\mathrm{tr}}(x)}\,K_{1\mid X}(x,dy)
+\int_{\mathbb R}e^{-(y-r_{\mathrm{tr}}(x))}\,K_{1\mid X}(x,dy)
+|r_{\mathrm{tr}}(x)|\\
&\le 2e^{B^2/2}+B.
\end{aligned}
\]
This bound is uniform on a set of full \(P_X\)-measure. Disintegration therefore gives
\[
\int |Y(1)|\,dP_c
=
\int_{\mathbb R^d}
\left(\int_{\mathbb R}|y|\,K_{1\mid X}(x,dy)\right)P_X(dx)
\le 2e^{B^2/2}+B<\infty.
\]
In particular, \(Y(1)\) is integrable under \(P_c\).
% lean: CausalSmith.Stat.WeakOverlap.mu1_eq_treatedRegression (houtcome, hY₁fiber, hY₁norm_bound, hY₁int)

\item Define
\[
k_1(x):=\int_{\mathbb R}y\,K_{1\mid X}(x,dy),
\qquad x\in\mathbb R^d,
\]
with value zero wherever the integrand is not integrable. The conditional-distribution formula for conditional expectation, together with the defining conditional-mean relation for the treated causal response, gives
\[
\mu_1(X)
=\mathbb E_{P_c}[Y(1)\mid\sigma(X)]
=k_1(X)
\qquad P_c\text{-almost surely}.
\]
The kernel integral \(k_1\) is measurable. Continuity of \(\mu_1\) on \(\mathcal X\), supplied by
\cref{obj:ass:holder-response,obj:def:holder}, and
\(P_X\ll\operatorname{Leb}|_{\mathcal X}\) supply a measurable representative of \(\mu_1\) under \(P_X\). Using these representatives, the covariate pushforward identity transfers the preceding equality to
\[
\mu_1(x)=k_1(x)
\qquad P_X\text{-almost everywhere}.
\]
The equality \(K_{\mathrm{tr}}=K_{1\mid X}\) then gives
\[
\mu_1(x)
=k_1(x)
=r_{\mathrm{tr}}(x)
=\mathbb E_P[Y\mid A=1,X=x]
\qquad P_X\text{-almost everywhere}.
\]
% lean: CausalSmith.Stat.WeakOverlap.mu1_eq_treatedRegression (hce, hcomp, hbase, hident)

\item It remains to establish the asserted equality for every point of the cube. Define the restricted Lebesgue measure by
\[
\nu:=\operatorname{Leb}|_{\mathcal X}.
\]
The parameter restrictions give \(c_f>0\). Absolute continuity of \(P_X\) with respect to \(\nu\) and \cref{obj:ass:density} yield
\[
P_X=f\nu,
\qquad
f=\frac{dP_X}{d\nu}\ge c_f>0
\qquad \nu\text{-almost everywhere}.
\]
Consequently,
\[
\nu\{f=0\}=0,
\qquad
\nu\ll f\nu=P_X.
\]
For completeness, if a Borel set \(E\subseteq\mathbb R^d\) satisfies \(P_X(E)=0\), then
\[
\int_E f\,d\nu=P_X(E)=0.
\]
Since \(f\) is nonnegative, this forces \(f=0\) \(\nu\)-almost everywhere on \(E\); its strict positivity \(\nu\)-almost everywhere gives \(\nu(E)=0\).

Let \(g:\mathbb R^d\to\mathbb R\) satisfy the continuity and almost-everywhere agreement required in the statement. The identifying equality already established implies
\[
g(x)=\mu_1(x)
\qquad P_X\text{-almost everywhere},
\]
and hence, by \(\nu\ll P_X\),
\[
g(x)=\mu_1(x)
\qquad \nu\text{-almost everywhere}.
\]
Both functions are continuous on \(\mathcal X\). Thus
\[
g(x)=\mu_1(x)
\qquad\text{for every }x\in\operatorname{int}(\mathcal X):
\]
a nonzero difference at an interior point would, by continuity, persist on an open ball contained in the cube, whose positive Lebesgue measure contradicts the preceding almost-everywhere equality. Finally,
\[
\overline{\operatorname{int}(\mathcal X)}
=\prod_{i=1}^{d}\overline{(0,1)}
=\prod_{i=1}^{d}[0,1]
=\mathcal X.
\]
Continuity therefore extends the equality to every \(x\in\mathcal X\), as required.
% lean: CausalSmith.Stat.WeakOverlap.mu1_eq_treatedRegression (hreverse, hAEvol, hcube_reg)
\end{enumerate}
\end{proof}

The pointwise and spatial-supremum losses in \cref{obj:def:risks} therefore concern a uniquely identified continuous response, including at boundary points. This separates the identification of the target from the statistical difficulty created by small treated probabilities.

\subsection{Global extension of cube smoothness}

The intrinsic smoothness convention in \cref{obj:def:holder} specifies coordinate derivatives through their continuous boundary traces. For an integer derivative order \(m\ge0\) and a top-derivative Hölder exponent \(s\in(0,1]\), the following statement supplies a global representative with controlled Fréchet derivatives. In the response class, these orders are \(m=\lceil\beta\rceil-1\) and \(s=\beta-m\).

\begin{lemmav}{lem:global-holder-extension}[Global Hölder extension bounds]
Fix an integer \(d\ge0\), an integer \(m\ge0\), and an exponent \(0<s\le1\). There exists a constant \(K_{\mathrm{ext}}>0\), depending only on \(d\), \(m\), and \(s\), with the following property. Let \(L\ge0\) and let \(g:[-1,1]^d\to\mathbb R\) satisfy:
\begin{itemize}
\item \textbf{(Intrinsic smoothness.)} The function \(g\) belongs to \(C^m([-1,1]^d)\) in the intrinsic closed-cube sense of \cref{obj:def:holder}, transported to \([-1,1]^d\). Its coordinate derivatives \(D^\alpha g\), for \(|\alpha|\le m\), exist on the interior and admit continuous traces on the closed cube.
\item \textbf{(Derivative bounds.)} With derivatives evaluated through their continuous traces,
\[
\sup_{z\in[-1,1]^d}|D^\alpha g(z)|\le L
\qquad\text{for every }|\alpha|\le m.
\]
\item \textbf{(Hölder continuity.)} For every \(|\alpha|=m\) and all \(z,z'\in[-1,1]^d\),
\[
|D^\alpha g(z)-D^\alpha g(z')|
\le L\|z-z'\|_\infty^{s}.
\]
\end{itemize}
Then there exists \(g_{\mathrm{ext}}\in C^m(\mathbb R^d)\) such that
\[
g_{\mathrm{ext}}(z)=g(z)
\qquad\text{for every }z\in[-1,1]^d,
\]
\[
\sup_{z\in\mathbb R^d}
\|\mathrm D^{k}g_{\mathrm{ext}}(z)\|_{\mathrm{op}}
\le K_{\mathrm{ext}}L
\qquad\text{for every }0\le k\le m,
\]
and
\[
\|\mathrm D^{m}g_{\mathrm{ext}}(z)
-\mathrm D^{m}g_{\mathrm{ext}}(z')\|_{\mathrm{op}}
\le K_{\mathrm{ext}}L\|z-z'\|_\infty^{s}
\qquad\text{for all }z,z'\in\mathbb R^d.
\]
Here \(\mathrm D^{k}\) denotes the Fréchet derivative of order \(k\), with order zero interpreted as the function itself, and \(\|\cdot\|_{\mathrm{op}}\) denotes the operator norm of a \(k\)-linear form on \(\mathbb R^d\) equipped with the supremum norm \(\|\cdot\|_\infty\).
\end{lemmav}

\begin{proof}[Proof of \cref{obj:lem:global-holder-extension}]
Write
\[
\mathcal C:=[-1,1]^d.
\]
All intrinsic derivatives below are evaluated through their continuous traces, as in \cref{obj:def:holder}.

\begin{enumerate}
\item \emph{Order zero.}
Suppose \(m=0\). Define the coordinatewise retraction and the extension by
\[
\pi_{\mathcal C}:\mathbb R^d\longrightarrow\mathcal C,
\qquad
(\pi_{\mathcal C}(z))_i:=\max\{-1,\min\{1,z_i\}\},
\qquad
g_{\mathrm{ext}}(z):=g(\pi_{\mathcal C}(z)).
\]
Projection onto an interval is nonexpansive. Consequently,
\[
\pi_{\mathcal C}(z)=z\quad(z\in\mathcal C),
\qquad
\|\pi_{\mathcal C}(z)-\pi_{\mathcal C}(z')\|_\infty
\le \|z-z'\|_\infty.
\]
The hypotheses therefore give
\[
|g_{\mathrm{ext}}(z)|\le L,
\qquad
|g_{\mathrm{ext}}(z)-g_{\mathrm{ext}}(z')|
\le L\|\pi_{\mathcal C}(z)-\pi_{\mathcal C}(z')\|_\infty^s
\le L\|z-z'\|_\infty^s.
\]
Since \(s>0\), the last bound implies continuity. Thus \(K_{\mathrm{ext}}=1\) works in this case. The same construction applies when \(d=0\), with the empty coordinate tuple interpreted as the unique point of \(\mathbb R^0\).
% lean: Causalean.Mathlib.Analysis.Calculus.CubeExtension.exists_global_holder_extension_order_zero

\item\label{proof:global-holder-extension:coordinate-bounds} \emph{From coordinate bounds to multilinear bounds.}
Henceforth assume \(m>0\). Set
\[
K_{\mathrm{coord}}:=(d+1)^m>0.
\]
For endpoints \(\boldsymbol\ell,\boldsymbol b\in\mathbb R^d\) satisfying
\(\ell_i<b_i\), define
\[
\mathcal B(\boldsymbol\ell,\boldsymbol b)
:=\prod_{i=1}^d[\ell_i,b_i].
\]
For an intrinsic \(C^m\) response \(g_{\mathrm{in}}\) on this box, its full intrinsic derivative is the multilinear form
\[
\mathrm D_{\mathcal B}^{k}g_{\mathrm{in}}(z)
[\xi_1,\ldots,\xi_k]
:=
\sum_{i_1,\ldots,i_k=1}^d
\bigl(\partial_{i_1}\cdots\partial_{i_k}g_{\mathrm{in}}\bigr)(z)
\prod_{\nu=1}^k(\xi_\nu)_{i_\nu},
\qquad 0\le k\le m.
\]
At order zero this is \(g_{\mathrm{in}}(z)\); at positive order in dimension zero the sum is empty and the form is zero.

To see the relevant norm estimate, define the coordinate basis by
\[
(\mathbf e_i)_j:=\mathbf 1_{\{i=j\}}.
\]
If a \(k\)-linear form \(M_{\mathrm{coord}}:(\mathbb R^d)^k\to\mathbb R\) satisfies
\[
|M_{\mathrm{coord}}[\mathbf e_{i_1},\ldots,\mathbf e_{i_k}]|
\le L_{\mathrm{coord}},
\qquad L_{\mathrm{coord}}\ge0,
\]
multilinearity gives
\[
\begin{aligned}
|M_{\mathrm{coord}}[\xi_1,\ldots,\xi_k]|
&\le
\sum_{i_1,\ldots,i_k=1}^d
L_{\mathrm{coord}}\prod_{\nu=1}^k|(\xi_\nu)_{i_\nu}|\\
&\le d^kL_{\mathrm{coord}}
\prod_{\nu=1}^k\|\xi_\nu\|_\infty
\le K_{\mathrm{coord}}L_{\mathrm{coord}}
\prod_{\nu=1}^k\|\xi_\nu\|_\infty.
\end{aligned}
\]
Here the order-zero sum has one term, including when \(d=0\). Applying this estimate to the derivative and to the difference of two top derivatives shows that coordinate bounds of radius \(L_{\mathrm{in}}\ge0\) imply
\[
\|\mathrm D_{\mathcal B}^{k}g_{\mathrm{in}}(z)\|_{\mathrm{op}}
\le K_{\mathrm{coord}}L_{\mathrm{in}},
\qquad 0\le k\le m,\quad z\in\mathcal B,
\]
and
\[
\|\mathrm D_{\mathcal B}^{m}g_{\mathrm{in}}(z)
-\mathrm D_{\mathcal B}^{m}g_{\mathrm{in}}(z')\|_{\mathrm{op}}
\le K_{\mathrm{coord}}L_{\mathrm{in}}\|z-z'\|_\infty^s,
\qquad z,z'\in\mathcal B.
\]
% lean: Causalean.Mathlib.Analysis.Calculus.CubeExtension.exists_coord_multilinear_norm_constant

\item\label{proof:global-holder-extension:affine-transport} \emph{Affine transport between boxes.}
Consider two nondegenerate boxes
\[
\mathcal B=\mathcal B(\boldsymbol\ell,\boldsymbol b),
\qquad
\mathcal B'=\mathcal B(\boldsymbol\ell',\boldsymbol b').
\]
Define their affine identification and its coordinate slopes by
\[
\lambda_{\mathrm{aff},i}
:=\frac{b_i'-\ell_i'}{b_i-\ell_i}>0,
\qquad
(\Phi_{\mathcal B\to\mathcal B'}(z))_i
:=\ell_i'+\lambda_{\mathrm{aff},i}(z_i-\ell_i).
\]
This map carries \(\mathcal B\) bijectively onto \(\mathcal B'\). Define
\[
K_{\mathrm{aff},\mathrm{jet}}
:=1+\sum_{i=1}^d\lambda_{\mathrm{aff},i},
\qquad
K_{\mathrm{aff},\mathrm{dist}}
:=1+\sum_{i=1}^d
\bigl(\lambda_{\mathrm{aff},i}+\lambda_{\mathrm{aff},i}^{-1}\bigr),
\]
and
\[
K_{\mathrm{aff}}(\mathcal B,\mathcal B')
:=
K_{\mathrm{aff},\mathrm{jet}}^m
\bigl(1+K_{\mathrm{aff},\mathrm{dist}}^s\bigr)>0.
\]
Coordinatewise subtraction yields
\[
\|\Phi_{\mathcal B\to\mathcal B'}(z)
-\Phi_{\mathcal B\to\mathcal B'}(z')\|_\infty
\le K_{\mathrm{aff},\mathrm{dist}}\|z-z'\|_\infty,
\]
hence, because \(s>0\),
\[
\|\Phi_{\mathcal B\to\mathcal B'}(z)
-\Phi_{\mathcal B\to\mathcal B'}(z')\|_\infty^s
\le K_{\mathrm{aff},\mathrm{dist}}^s\|z-z'\|_\infty^s.
\]
For a response \(g_{\mathrm{in}}\) with intrinsic radius \(L_{\mathrm{in}}\) on \(\mathcal B'\), define
\[
g_{\mathrm{aff}}:=g_{\mathrm{in}}\circ\Phi_{\mathcal B\to\mathcal B'}.
\]
The affine chain rule, including continuous boundary traces, gives
\[
D^\alpha g_{\mathrm{aff}}(z)
=
\left(\prod_{i=1}^d\lambda_{\mathrm{aff},i}^{\alpha_i}\right)
(D^\alpha g_{\mathrm{in}})
(\Phi_{\mathcal B\to\mathcal B'}(z)).
\]
Since \(K_{\mathrm{aff},\mathrm{jet}}\ge1\), for \(|\alpha|\le m\),
\[
\prod_{i=1}^d\lambda_{\mathrm{aff},i}^{\alpha_i}
\le K_{\mathrm{aff},\mathrm{jet}}^{|\alpha|}
\le K_{\mathrm{aff},\mathrm{jet}}^m.
\]
It follows that
\[
|D^\alpha g_{\mathrm{aff}}(z)|
\le K_{\mathrm{aff},\mathrm{jet}}^mL_{\mathrm{in}}
\le K_{\mathrm{aff}}(\mathcal B,\mathcal B')L_{\mathrm{in}},
\]
and, for \(|\alpha|=m\),
\[
\begin{aligned}
|D^\alpha g_{\mathrm{aff}}(z)-D^\alpha g_{\mathrm{aff}}(z')|
&\le K_{\mathrm{aff},\mathrm{jet}}^mL_{\mathrm{in}}
\|\Phi_{\mathcal B\to\mathcal B'}(z)
-\Phi_{\mathcal B\to\mathcal B'}(z')\|_\infty^s\\
&\le K_{\mathrm{aff}}(\mathcal B,\mathcal B')L_{\mathrm{in}}
\|z-z'\|_\infty^s.
\end{aligned}
\]
Composition with the affine map also preserves intrinsic \(C^m\) regularity. These formulas establish the claimed radius factor over the whole closed box.
% lean: Causalean.Mathlib.Analysis.Calculus.CubeExtension.exists_rectAffine_holder_constant

\item \emph{Moment-matched reflection across a normalized face.}
Choose real coefficients \(a_{\mathrm{ref},r}\), \(1\le r\le m+1\), satisfying
\[
\sum_{r=1}^{m+1}a_{\mathrm{ref},r}(-r)^k=1,
\qquad 0\le k\le m.
\]
Such coefficients exist: the matrix
\[
(\mathsf V_{\mathrm{ref}})_{r,k+1}:=(-r)^k,
\qquad 1\le r\le m+1,\quad 0\le k\le m,
\]
is a Vandermonde matrix at distinct nodes. Its determinant is nonzero, so the system
\[
\mathsf V_{\mathrm{ref}}^{\mathsf T}
(a_{\mathrm{ref},1},\ldots,a_{\mathrm{ref},m+1})^{\mathsf T}
=(1,\ldots,1)^{\mathsf T}
\]
has a solution depending only on \(m\).
% lean: Causalean.Mathlib.Analysis.Calculus.CubeExtension.exists_reflection_coefficients

\item\label{proof:global-holder-extension:face-collar} \emph{Regularity and bounds on a one-face collar.}
Fix a coordinate \(i\), and define
\[
\mathcal C_i^-:=
\left\{z:
-1-\frac1{m+1}\le z_i\le1,\
-1\le z_j\le1\ \text{for }j\ne i
\right\},
\]
\[
\mathcal E_i^-:=\mathcal C_i^-\cap\{z:z_i\le-1\},
\qquad
\mathcal E_i^{-,\circ}:=
\left\{z:
-1-\frac1{m+1}<z_i<-1,\
-1<z_j<1\ \text{for }j\ne i
\right\}.
\]
For each \(1\le r\le m+1\), define an affine inward sample on all of \(\mathbb R^d\) by
\[
(\Phi_{\mathrm{ref},i,r}(z))_j
:=
\begin{cases}
-1+r(-1-z_i),&j=i,\\
z_j,&j\ne i.
\end{cases}
\]
If \(z\in\mathcal E_i^{-,\circ}\), then
\[
0<r(-1-z_i)<1,
\]
so \(\Phi_{\mathrm{ref},i,r}(z)\in(-1,1)^d\). The weak inequalities likewise show that the closed exterior slab maps into \(\mathcal C\).

For an auxiliary response \(g_{\mathrm{in}}:\mathbb R^d\to\mathbb R\) with intrinsic radius \(L_{\mathrm{in}}\) on \(\mathcal C\), define
\[
(\mathscr R_i g_{\mathrm{in}})(z)
:=
\begin{cases}
\displaystyle\sum_{r=1}^{m+1}
a_{\mathrm{ref},r}g_{\mathrm{in}}(\Phi_{\mathrm{ref},i,r}(z)),
&z_i<-1,\\
g_{\mathrm{in}}(z),&z_i\ge-1.
\end{cases}
\]
This defines the response on all of \(\mathbb R^d\), and it agrees with \(g_{\mathrm{in}}\) on \(\mathcal C\).

On the closed exterior slab, the finite sum is intrinsically \(C^m\), by composition with the affine samples. At \(z_i=-1\), every sample equals \(z\), and the moment equation at \(k=0\) identifies that sum with the second branch. Moreover, for \(|\alpha|\le m\), the exterior derivative trace at this face is
\[
\begin{aligned}
(D^\alpha\mathscr R_i g_{\mathrm{in}})_{\mathrm{ext}}(z)
&=
\sum_{r=1}^{m+1}
a_{\mathrm{ref},r}(-r)^{\alpha_i}D^\alpha g_{\mathrm{in}}(z)\\
&=D^\alpha g_{\mathrm{in}}(z).
\end{aligned}
\]
Indeed, the affine chain rule gives the factors \((-r)^{\alpha_i}\), and continuity of the original derivative traces gives the limit at the face. Coordinate expansion then identifies the full multilinear traces. This argument includes the edges and corners.

The two closed pieces therefore have matching derivatives through order \(m\). They glue to an intrinsic \(C^m\) response on \(\mathcal C_i^-\): at a common face point, the derivative remainders on both pieces use the same linear approximation and are \(o(\|\xi\|_\infty)\); on their union the same remainder estimate holds. Applying this observation successively to the derivative fields, and using continuity of the matching fields, establishes regularity through order \(m\).

For the quantitative estimates, define the linear part of each sample by
\[
(\Lambda_{\mathrm{ref},i,r}\xi)_j
:=
\begin{cases}
-r\xi_i,&j=i,\\
\xi_j,&j\ne i.
\end{cases}
\]
Then
\[
\|\Lambda_{\mathrm{ref},i,r}\|_{\mathrm{op}}\le r,
\qquad
\|\Phi_{\mathrm{ref},i,r}(z)-\Phi_{\mathrm{ref},i,r}(z')\|_\infty
\le r\|z-z'\|_\infty.
\]
On the open exterior slab, the chain rule gives
\[
\mathrm D^k(g_{\mathrm{in}}\circ\Phi_{\mathrm{ref},i,r})(z)
=
\mathrm D^k g_{\mathrm{in}}(\Phi_{\mathrm{ref},i,r}(z))
\circ\Lambda_{\mathrm{ref},i,r}^{\otimes k}.
\]
Using the full intrinsic bounds from \cref{proof:global-holder-extension:coordinate-bounds} at the inward samples yields
\[
\|\mathrm D^k(g_{\mathrm{in}}\circ\Phi_{\mathrm{ref},i,r})(z)\|_{\mathrm{op}}
\le K_{\mathrm{coord}}L_{\mathrm{in}}r^k
\le K_{\mathrm{coord}}r^mL_{\mathrm{in}},
\]
and
\[
\begin{aligned}
&\|\mathrm D^m(g_{\mathrm{in}}\circ\Phi_{\mathrm{ref},i,r})(z)
-\mathrm D^m(g_{\mathrm{in}}\circ\Phi_{\mathrm{ref},i,r})(z')\|_{\mathrm{op}}\\
&\qquad\le
K_{\mathrm{coord}}L_{\mathrm{in}}r^m
\|\Phi_{\mathrm{ref},i,r}(z)-\Phi_{\mathrm{ref},i,r}(z')\|_\infty^s\\
&\qquad\le K_{\mathrm{coord}}r^mr^sL_{\mathrm{in}}\|z-z'\|_\infty^s.
\end{aligned}
\]
Define
\[
K_{\mathrm{ref},\mathrm{ext},D}
:=1+\sum_{r=1}^{m+1}|a_{\mathrm{ref},r}|K_{\mathrm{coord}}r^m,
\qquad
K_{\mathrm{ref},\mathrm{ext},H}
:=1+\sum_{r=1}^{m+1}|a_{\mathrm{ref},r}|K_{\mathrm{coord}}r^mr^s.
\]
Differentiating the finite reflection sum and applying the triangle inequality gives
\[
\|\mathrm D^k\mathscr R_i g_{\mathrm{in}}(z)\|_{\mathrm{op}}
\le K_{\mathrm{ref},\mathrm{ext},D}L_{\mathrm{in}},
\]
and
\[
\|\mathrm D^m\mathscr R_i g_{\mathrm{in}}(z)
-\mathrm D^m\mathscr R_i g_{\mathrm{in}}(z')\|_{\mathrm{op}}
\le K_{\mathrm{ref},\mathrm{ext},H}L_{\mathrm{in}}\|z-z'\|_\infty^s
\]
on the open exterior slab. Its closure is \(\mathcal E_i^-\). Continuity of the within-collar derivatives therefore extends the first bound to that closure, and successive limits in the two points extend the second bound.

On \(\mathcal C\), agreement of the responses identifies their derivatives in the interior. Continuity extends this identity to the closed cube:
\[
\mathrm D_{\mathcal C_i^-}^k\mathscr R_i g_{\mathrm{in}}(z)
=\mathrm D_{\mathcal C}^k g_{\mathrm{in}}(z),
\qquad z\in\mathcal C,\quad k\le m.
\]
Thus the corresponding full derivative and top-modulus bounds on the cube have coefficient \(K_{\mathrm{coord}}\).

For a pair \(z\in\mathcal E_i^-\), \(z'\in\mathcal C\), choose
\[
t_{\mathrm{face}}\in[0,1],
\qquad
(1-t_{\mathrm{face}})z_i+t_{\mathrm{face}}z_i'=-1,
\qquad
z_{\mathrm{face}}:=(1-t_{\mathrm{face}})z+t_{\mathrm{face}}z'.
\]
Existence follows from \(z_i\le-1\le z_i'\); if both coordinates equal \(-1\), take \(t_{\mathrm{face}}=0\). The other coordinates remain in \([-1,1]\), so
\[
z_{\mathrm{face}}\in\mathcal E_i^-\cap\mathcal C,
\qquad
\|z-z_{\mathrm{face}}\|_\infty\le\|z-z'\|_\infty,
\qquad
\|z_{\mathrm{face}}-z'\|_\infty\le\|z-z'\|_\infty.
\]
The two side estimates and the triangle inequality now give
\[
\begin{aligned}
&\|\mathrm D_{\mathcal C_i^-}^m\mathscr R_i g_{\mathrm{in}}(z)
-\mathrm D_{\mathcal C_i^-}^m\mathscr R_i g_{\mathrm{in}}(z')\|_{\mathrm{op}}\\
&\quad\le
K_{\mathrm{ref},\mathrm{ext},H}L_{\mathrm{in}}
\|z-z_{\mathrm{face}}\|_\infty^s
+K_{\mathrm{coord}}L_{\mathrm{in}}
\|z_{\mathrm{face}}-z'\|_\infty^s\\
&\quad\le
(K_{\mathrm{ref},\mathrm{ext},H}+K_{\mathrm{coord}})
L_{\mathrm{in}}\|z-z'\|_\infty^s.
\end{aligned}
\]
Pairs on the same side satisfy the same enlarged bound; reversing the points treats the other crossing order. Evaluating on coordinate basis tuples converts these full-jet estimates to coordinate estimates. Consequently, with
\[
K_{\mathrm{ref},D}
:=K_{\mathrm{ref},\mathrm{ext},D}+K_{\mathrm{coord}},
\qquad
K_{\mathrm{ref},H}
:=K_{\mathrm{ref},\mathrm{ext},H}+K_{\mathrm{coord}},
\qquad
K_{\mathrm{ref}}:=K_{\mathrm{ref},D}+K_{\mathrm{ref},H},
\]
the reflected response has intrinsic \(C^{m,s}\) radius
\(K_{\mathrm{ref}}L_{\mathrm{in}}\) on \(\mathcal C_i^-\). All these constants are positive and depend only on \(d,m,s\).
% lean: Causalean.Mathlib.Analysis.Calculus.CubeExtension.exists_leftFaceReflection_collar_holder_constant

\item\label{proof:global-holder-extension:larger-box} \emph{A fixed larger box.}
For a nondegenerate box
\(\mathcal B=\mathcal B(\boldsymbol\ell,\boldsymbol b)\) and a coordinate \(i\), define its lower enlargement by
\[
\ell_j^-:=
\begin{cases}
\displaystyle\ell_i-\frac{b_i-\ell_i}{2(m+1)},&j=i,\\
\ell_j,&j\ne i,
\end{cases}
\qquad
\mathcal B^-:=\mathcal B(\boldsymbol\ell^-,\boldsymbol b).
\]
Normalize a response on \(\mathcal B\), reflect it, and transport it back:
\[
g_{\mathrm{norm}}
:=g_{\mathrm{in}}\circ\Phi_{\mathcal C\to\mathcal B},
\qquad
g_{\mathrm{lower}}
:=(\mathscr R_i g_{\mathrm{norm}})
\circ\Phi_{\mathcal B^-\to\mathcal C_i^-}.
\]
The chosen collar width gives the affine identity
\[
\Phi_{\mathcal B^-\to\mathcal C_i^-}
=\Phi_{\mathcal B\to\mathcal C}
\quad\text{on }\mathbb R^d.
\]
Therefore
\[
g_{\mathrm{lower}}(z)=g_{\mathrm{in}}(z)
\quad(z\in\mathcal B).
\]
\Cref{proof:global-holder-extension:affine-transport,proof:global-holder-extension:face-collar} give its radius factor
\[
K_{\mathrm{lower}}(\mathcal B,i)
:=
K_{\mathrm{aff}}(\mathcal B^-,\mathcal C_i^-)
K_{\mathrm{ref}}
K_{\mathrm{aff}}(\mathcal C,\mathcal B)>0.
\]

For the upper enlargement, define
\[
b_j^+:=
\begin{cases}
\displaystyle b_i+\frac{b_i-\ell_i}{2(m+1)},&j=i,\\
b_j,&j\ne i,
\end{cases}
\qquad
\mathcal B^+:=\mathcal B(\boldsymbol\ell,\boldsymbol b^+).
\]
Apply the lower construction to the negated box and then negate back. This transport preserves the radius, since
\[
D^\alpha(g_{\mathrm{in}}\circ(-\mathrm{id}))(z)
=(-1)^{|\alpha|}D^\alpha g_{\mathrm{in}}(-z),
\qquad
\|(-z)-(-z')\|_\infty=\|z-z'\|_\infty.
\]
Thus the upper construction agrees on \(\mathcal B\) and has factor
\[
K_{\mathrm{upper}}(\mathcal B,i)
:=K_{\mathrm{lower}}(-\mathcal B,i),
\qquad
-\mathcal B:=\mathcal B(-\boldsymbol b,-\boldsymbol\ell).
\]
Applying the lower construction first and the upper construction to its enlarged box gives agreement on the original box and factor
\[
K_{\mathrm{two}}(\mathcal B,i)
:=K_{\mathrm{upper}}(\mathcal B^-,i)
K_{\mathrm{lower}}(\mathcal B,i)>0.
\]

Iterate this two-face construction through coordinates \(1,\ldots,d\). More precisely, define
\[
\boldsymbol\ell^{(0)}:=(-1,\ldots,-1),
\qquad
\boldsymbol b^{(0)}:=(1,\ldots,1),
\qquad
\mathcal B^{(0)}:=\mathcal C,
\qquad
K_{\mathrm{neigh}}^{(0)}:=1.
\]
At stage \(\nu=1,\ldots,d\), apply the lower enlargement in coordinate \(\nu\), followed by the upper enlargement using the new lower endpoint, and denote the resulting box by \(\mathcal B^{(\nu)}\). Define
\[
K_{\mathrm{neigh}}^{(\nu)}
:=
K_{\mathrm{two}}(\mathcal B^{(\nu-1)},\nu)
K_{\mathrm{neigh}}^{(\nu-1)}.
\]
The lower and upper increments are positive. Previously enlarged coordinates remain unchanged, and every intermediate box contains \(\mathcal C\). Induction consequently gives
\[
\ell_i^{(d)}<-1<1<b_i^{(d)}
\quad(1\le i\le d),
\qquad
K_{\mathrm{neigh}}^{(d)}>0.
\]
Set
\[
\mathcal B_{\mathrm{ext}}
:=\mathcal B^{(d)},
\qquad
\boldsymbol\ell_{\mathrm{ext}}:=\boldsymbol\ell^{(d)},
\qquad
\boldsymbol b_{\mathrm{ext}}:=\boldsymbol b^{(d)},
\qquad
K_{\mathrm{neigh}}:=K_{\mathrm{neigh}}^{(d)}.
\]
For every admissible \(g,L\), the same induction produces a response \(g_{\mathrm{neigh}}\) satisfying
\[
g_{\mathrm{neigh}}=g\quad\text{on }\mathcal C,
\qquad
g_{\mathrm{neigh}}\in C^m(\mathcal B_{\mathrm{ext}}),
\]
\[
|D^\alpha g_{\mathrm{neigh}}(z)|
\le K_{\mathrm{neigh}}L
\quad(|\alpha|\le m,\ z\in\mathcal B_{\mathrm{ext}}),
\]
and
\[
|D^\alpha g_{\mathrm{neigh}}(z)-D^\alpha g_{\mathrm{neigh}}(z')|
\le K_{\mathrm{neigh}}L\|z-z'\|_\infty^s
\quad(|\alpha|=m,\ z,z'\in\mathcal B_{\mathrm{ext}}).
\]
For definiteness, the starting ambient representative is
\[
g^{(0)}(z):=
\begin{cases}
g(z),&z\in\mathcal C,\\
0,&z\notin\mathcal C.
\end{cases}
\]
Subsequent representatives are defined everywhere by the displayed reflection and affine formulas. When \(d=0\), the iteration is empty:
\(\mathcal B_{\mathrm{ext}}=\mathbb R^0\) and \(K_{\mathrm{neigh}}=1\).
The final box and radius factor depend only on \(d,m,s\).
% lean: Causalean.Mathlib.Analysis.Calculus.CubeExtension.exists_fixedCubeNeighborhood_holder_constant

\item\label{proof:global-holder-extension:cutoff} \emph{A fixed smooth cutoff and its zero extension.}
Define
\[
\rho_{\mathrm{box}}:=
\begin{cases}
\displaystyle\min_{1\le i\le d}
\min\{-\ell_{\mathrm{ext},i},b_{\mathrm{ext},i}\},&d>0,\\
2,&d=0,
\end{cases}
\]
and
\[
\rho_{\mathrm{in}}:=\frac{1+\rho_{\mathrm{box}}}{2},
\qquad
\rho_{\mathrm{out}}:=\frac{\rho_{\mathrm{in}}+\rho_{\mathrm{box}}}{2}.
\]
The strict margins imply
\[
1<\rho_{\mathrm{in}}<\rho_{\mathrm{out}}<\rho_{\mathrm{box}}.
\]
Choose a smooth bump \(\chi:\mathbb R^d\to\mathbb R\) with
\[
\chi(z)=1\quad(\|z\|_\infty\le\rho_{\mathrm{in}}),
\qquad
\operatorname{supp}\chi
\subseteq\{z:\|z\|_\infty\le\rho_{\mathrm{out}}\},
\]
where support denotes closed support. Smooth bump existence in a finite-dimensional normed space supplies this choice. In particular,
\[
\chi=1\quad\text{on }\mathcal C,
\qquad
\operatorname{supp}\chi\subset
\mathcal B_{\mathrm{ext}}^\circ,
\]
with
\[
\mathcal B_{\mathrm{ext}}^\circ
:=\prod_{i=1}^d
(\ell_{\mathrm{ext},i},b_{\mathrm{ext},i}).
\]
Compact support and continuity of the derivatives make the constant
\[
K_\chi:=
\max\left\{
1,\ 
\max_{0\le k\le m+1}
\sup_{z\in\mathbb R^d}
\|\mathrm D^k\chi(z)\|_{\mathrm{op}}
\right\}
\]
finite and positive.

For any response \(g_{\mathrm{in}}\) of intrinsic radius \(L_{\mathrm{in}}\ge0\) on \(\mathcal B_{\mathrm{ext}}\), define
\[
(\mathscr G g_{\mathrm{in}})(z)
:=
\begin{cases}
\chi(z)g_{\mathrm{in}}(z),&z\in\mathcal B_{\mathrm{ext}},\\
0,&z\notin\mathcal B_{\mathrm{ext}}.
\end{cases}
\]
At an interior point this is locally a product of \(C^m\) functions. At every point outside the open box, a neighborhood misses the closed support of \(\chi\), so the extension is locally zero. Hence
\[
\mathscr G g_{\mathrm{in}}\in C^m(\mathbb R^d),
\qquad
\mathscr G g_{\mathrm{in}}=g_{\mathrm{in}}\quad\text{on }\mathcal C,
\]
and
\[
\mathrm D^k\mathscr G g_{\mathrm{in}}(z)=0
\quad(z\notin\mathcal B_{\mathrm{ext}}^\circ,\ 0\le k\le m).
\]
% lean: Causalean.Mathlib.Analysis.Calculus.CubeExtension.rectCutoffExtension_contDiff_eqOn_cube

\item\label{proof:global-holder-extension:cutoff-derivatives} \emph{Uniform derivatives of the cutoff extension.}
We make the product estimate explicit. Define
\[
\mathfrak S_k:=\{\text{permutations of }\{1,\ldots,k\}\}.
\]
For an \(r\)-linear form \(M_1\) and a \((k-r)\)-linear form \(M_2\), where \(0\le r\le k\), define
\[
\begin{aligned}
&\mathcal M_{k,r}(M_1,M_2)[\xi_1,\ldots,\xi_k]\\
&\qquad:=
\frac1{k!}\sum_{\sigma\in\mathfrak S_k}
M_1[\xi_{\sigma(1)},\ldots,\xi_{\sigma(r)}]\,
M_2[\xi_{\sigma(r+1)},\ldots,\xi_{\sigma(k)}].
\end{aligned}
\]
At \(k=0\), the permutation set has one element and the expression is ordinary multiplication. Each summand is bounded by
\(\|M_1\|_{\mathrm{op}}\|M_2\|_{\mathrm{op}}
\prod_{\nu=1}^k\|\xi_\nu\|_\infty\), so
\[
\|\mathcal M_{k,r}(M_1,M_2)\|_{\mathrm{op}}
\le\|M_1\|_{\mathrm{op}}\|M_2\|_{\mathrm{op}}.
\]
Repeated differentiation along a line gives the scalar Leibniz formula on repeated arguments. Both sides below are symmetric multilinear forms, and polarization therefore gives the full product identity
\[
\mathrm D^k(\chi g_{\mathrm{in}})(z)
=
\sum_{r=0}^k\binom{k}{r}
\mathcal M_{k,r}
\bigl(\mathrm D^r\chi(z),
\mathrm D_{\mathcal B_{\mathrm{ext}}}^{k-r}g_{\mathrm{in}}(z)\bigr),
\qquad z\in\mathcal B_{\mathrm{ext}}^\circ.
\]
Thus \cref{proof:global-holder-extension:coordinate-bounds} gives
\[
\begin{aligned}
\|\mathrm D^k\mathscr G g_{\mathrm{in}}(z)\|_{\mathrm{op}}
&\le
\sum_{r=0}^k\binom{k}{r}
K_\chi K_{\mathrm{coord}}L_{\mathrm{in}}\\
&=2^kK_\chi K_{\mathrm{coord}}L_{\mathrm{in}}
\le2^mK_\chi K_{\mathrm{coord}}L_{\mathrm{in}}.
\end{aligned}
\]
Outside the open box the derivative is zero. Consequently, defining
\[
K_{\mathrm{cut},D}:=K_\chi K_{\mathrm{coord}}2^m+1>0,
\]
we obtain the global bound
\[
\|\mathrm D^k\mathscr G g_{\mathrm{in}}(z)\|_{\mathrm{op}}
\le K_{\mathrm{cut},D}L_{\mathrm{in}},
\qquad 0\le k\le m,\quad z\in\mathbb R^d.
\]
% lean: Causalean.Mathlib.Analysis.Calculus.CubeExtension.exists_rectCutoffExtension_jet_constant

\item\label{proof:global-holder-extension:global-modulus} \emph{The global top-order modulus.}
For \(z,z'\) in the open box, define
\[
\delta:=\|z-z'\|_\infty.
\]
The box is convex. For \(k<m\), the bound on the next full intrinsic derivative and the mean-value inequality give
\[
\|\mathrm D_{\mathcal B_{\mathrm{ext}}}^{k}g_{\mathrm{in}}(z)
-\mathrm D_{\mathcal B_{\mathrm{ext}}}^{k}g_{\mathrm{in}}(z')\|_{\mathrm{op}}
\le K_{\mathrm{coord}}L_{\mathrm{in}}\delta.
\]
The uniform bounds also give a difference bound
\(2K_{\mathrm{coord}}L_{\mathrm{in}}\). For \(0\le\delta\le1\), the condition \(s\le1\) implies \(\delta\le\delta^s\); for \(\delta>1\), positivity of \(s\) implies \(1\le\delta^s\). Therefore
\[
\|\mathrm D_{\mathcal B_{\mathrm{ext}}}^{k}g_{\mathrm{in}}(z)
-\mathrm D_{\mathcal B_{\mathrm{ext}}}^{k}g_{\mathrm{in}}(z')\|_{\mathrm{op}}
\le2K_{\mathrm{coord}}L_{\mathrm{in}}\delta^s,
\qquad k<m.
\]
For \(k=m\), \cref{proof:global-holder-extension:coordinate-bounds} gives the sharper bound with coefficient
\(K_{\mathrm{coord}}L_{\mathrm{in}}\), and hence the same displayed bound with coefficient \(2K_{\mathrm{coord}}L_{\mathrm{in}}\).

Likewise, the cutoff derivatives through order \(m+1\) give, for \(0\le r\le m\),
\[
\|\mathrm D^r\chi(z)-\mathrm D^r\chi(z')\|_{\mathrm{op}}
\le K_\chi\delta,
\qquad
\|\mathrm D^r\chi(z)-\mathrm D^r\chi(z')\|_{\mathrm{op}}
\le2K_\chi.
\]
The same two distance ranges yield
\[
\|\mathrm D^r\chi(z)-\mathrm D^r\chi(z')\|_{\mathrm{op}}
\le2K_\chi\delta^s.
\]

Bilinearity of the pairings in \cref{proof:global-holder-extension:cutoff-derivatives} splits each product difference into a cutoff difference times a response jet, and a cutoff jet times a response difference. Subtracting the two product identities and applying their contraction bounds gives
\[
\begin{aligned}
&\|\mathrm D^m\mathscr G g_{\mathrm{in}}(z)
-\mathrm D^m\mathscr G g_{\mathrm{in}}(z')\|_{\mathrm{op}}\\
&\quad\le
\sum_{r=0}^m\binom mr
\Bigl(
\|\mathrm D^r\chi(z)-\mathrm D^r\chi(z')\|_{\mathrm{op}}
\|\mathrm D_{\mathcal B_{\mathrm{ext}}}^{m-r}g_{\mathrm{in}}(z)\|_{\mathrm{op}}\\
&\hspace{105pt}
+\|\mathrm D^r\chi(z')\|_{\mathrm{op}}
\|\mathrm D_{\mathcal B_{\mathrm{ext}}}^{m-r}g_{\mathrm{in}}(z)
-\mathrm D_{\mathcal B_{\mathrm{ext}}}^{m-r}g_{\mathrm{in}}(z')\|_{\mathrm{op}}
\Bigr)\\
&\quad\le
\sum_{r=0}^m\binom mr
\left(
2K_\chi\delta^s K_{\mathrm{coord}}L_{\mathrm{in}}
+K_\chi\,2K_{\mathrm{coord}}L_{\mathrm{in}}\delta^s
\right)\\
&\quad=
4K_\chi K_{\mathrm{coord}}2^mL_{\mathrm{in}}\delta^s.
\end{aligned}
\]
Define
\[
K_{\mathrm{cut},H}:=4K_\chi K_{\mathrm{coord}}2^m+1>0.
\]
The open box is dense in the closed box, including in dimension zero. Since \(\mathscr G g_{\mathrm{in}}\) is globally \(C^m\), its top derivative is continuous. Passing to limits in both points, using continuity of the distance power for \(s>0\), gives
\[
\|\mathrm D^m\mathscr G g_{\mathrm{in}}(z)
-\mathrm D^m\mathscr G g_{\mathrm{in}}(z')\|_{\mathrm{op}}
\le K_{\mathrm{cut},H}L_{\mathrm{in}}\|z-z'\|_\infty^s,
\qquad z,z'\in\mathcal B_{\mathrm{ext}}.
\]

Now let \(z\in\mathcal B_{\mathrm{ext}}\) and \(z'\notin\mathcal B_{\mathrm{ext}}\). Some coordinate \(i\) satisfies
\(z_i'<\ell_{\mathrm{ext},i}\) or \(z_i'>b_{\mathrm{ext},i}\). Define
\[
(z_{\mathrm{boundary}})_j:=
\begin{cases}
\ell_{\mathrm{ext},i},&j=i,\ z_i'<\ell_{\mathrm{ext},i},\\
b_{\mathrm{ext},i},&j=i,\ z_i'>b_{\mathrm{ext},i},\\
z_j,&j\ne i.
\end{cases}
\]
Then
\[
z_{\mathrm{boundary}}\in
\mathcal B_{\mathrm{ext}}\setminus\mathcal B_{\mathrm{ext}}^\circ,
\qquad
\|z-z_{\mathrm{boundary}}\|_\infty
\le |z_i-z_i'|
\le\|z-z'\|_\infty.
\]
\Cref{proof:global-holder-extension:cutoff} makes the top derivative zero at both \(z'\) and
\(z_{\mathrm{boundary}}\). Hence
\[
\begin{aligned}
&\|\mathrm D^m\mathscr G g_{\mathrm{in}}(z)
-\mathrm D^m\mathscr G g_{\mathrm{in}}(z')\|_{\mathrm{op}}\\
&\quad=
\|\mathrm D^m\mathscr G g_{\mathrm{in}}(z)
-\mathrm D^m\mathscr G g_{\mathrm{in}}(z_{\mathrm{boundary}})\|_{\mathrm{op}}\\
&\quad\le
K_{\mathrm{cut},H}L_{\mathrm{in}}
\|z-z_{\mathrm{boundary}}\|_\infty^s
\le K_{\mathrm{cut},H}L_{\mathrm{in}}\|z-z'\|_\infty^s.
\end{aligned}
\]
Reversing the points treats the opposite crossing order. If both points lie outside the closed box, both derivatives are zero. Together with the closed-box estimate, these cases establish the displayed modulus for all \(z,z'\in\mathbb R^d\). In dimension zero the outside cases are empty.
% lean: Causalean.Mathlib.Analysis.Calculus.CubeExtension.exists_rectCutoffExtension_modulus_constant

\item \emph{Composition of the two radius factors.}
Define the fixed constants
\[
K_{\mathrm{cut}}:=\max\{K_{\mathrm{cut},D},K_{\mathrm{cut},H}\},
\qquad
K_{\mathrm{ext}}:=K_{\mathrm{cut}}K_{\mathrm{neigh}}>0.
\]
The box, cutoff, and all constants depend only on \(d,m,s\). For the given \(g,L\), choose the neighborhood response from \cref{proof:global-holder-extension:larger-box} and set
\[
L_{\mathrm{in}}:=K_{\mathrm{neigh}}L\ge0,
\qquad
g_{\mathrm{ext}}:=\mathscr G g_{\mathrm{neigh}}.
\]
Its agreement on the cube is
\[
g_{\mathrm{ext}}(z)=g_{\mathrm{neigh}}(z)=g(z),
\qquad z\in\mathcal C,
\]
and \cref{proof:global-holder-extension:cutoff} gives \(g_{\mathrm{ext}}\in C^m(\mathbb R^d)\). \Cref{proof:global-holder-extension:cutoff-derivatives,proof:global-holder-extension:global-modulus} give
\[
\|\mathrm D^k g_{\mathrm{ext}}(z)\|_{\mathrm{op}}
\le K_{\mathrm{cut}}(K_{\mathrm{neigh}}L)
=K_{\mathrm{ext}}L,
\qquad 0\le k\le m,
\]
and
\[
\|\mathrm D^m g_{\mathrm{ext}}(z)-\mathrm D^m g_{\mathrm{ext}}(z')\|_{\mathrm{op}}
\le K_{\mathrm{cut}}(K_{\mathrm{neigh}}L)\|z-z'\|_\infty^s
=K_{\mathrm{ext}}L\|z-z'\|_\infty^s.
\]
Every estimate uses multiplication by nonnegative radii, so the construction and bounds also apply when \(L=0\).
% lean: Causalean.Mathlib.Analysis.Calculus.CubeExtension.exists_global_holder_extension
\end{enumerate}
\end{proof}

The extension preserves every value on the closed cube and converts the coordinate-derivative restrictions into global operator-norm bounds. Its smoothness factor \(K_{\mathrm{ext}}\) depends on dimension and smoothness order, uniformly over functions satisfying the displayed radius bound. This supplies a controlled global representation of intrinsic cube smoothness while preserving the response values relevant to approximation.

\subsection{The ordered treated-mass profile}

For a dyadic mesh \(h\), let \(q_{Q\ell}\) denote the population treated probability of template microcell \(E_{Q\ell}\) in partition cube \(Q\). The quantity \(q_Q\) is the smallest of these microcell probabilities within \(Q\), and \(q_{(k)}\) is the \(k\)th smallest entry among the cube-level quantities \(q_Q\). Ties may be ordered arbitrarily. These masses concern the scaled microcells of \cref{obj:def:template}, rather than the total treated probability of each partition cube.

The global-tail restriction in \cref{obj:ass:global-tail} controls treated probability on every measurable covariate region. Its implication for the ordered weakest-cell masses supplies the information profile used in the stochastic analysis.

\begin{lemmav}{lem:ordered-mass}[Uniform ordered treated masses]
Fix the model parameters and an overlap interval satisfying
\begin{itemize}
\item \textbf{(Parameter domain.)} \(d\in\mathbb N\), \(d\ge1\), \(\beta>1\), \(B>0\), \(L>0\), \(C\ge1\), and \(0<c_f\le1\).
\item \textbf{(Overlap range.)} \(1<\gamma_{\min}<\gamma_{\max}\), with \(\Gamma=[\gamma_{\min},\gamma_{\max}]\).
\end{itemize}
There exists \(\kappa_\Gamma>0\), depending only on these fixed parameters and interval endpoints, such that the following two bounds hold for every \(\gamma\in\Gamma\) and every causal law \(P\in\mathcal P_\gamma\) satisfying \cref{obj:def:model}.

For every measurable \(E\subseteq\mathcal X=[0,1]^d\), writing \(u=P(X\in E)\),
\[
P(A=1,X\in E)
\ge
\frac{\gamma-1}{\gamma}\,
C^{-1/(\gamma-1)}u^{\gamma/(\gamma-1)}.
\]

For each \(j\in\mathbb N\), set
\[
h=2^{-j},\qquad
m=\lceil\beta\rceil-1,\qquad
D_\gamma=\frac{d\gamma}{\gamma-1}.
\]
Let \(\mathcal Q_h\) be the dyadic cube partition of \(\mathcal X\), with \(|\mathcal Q_h|=(2^j)^d\). Using the degree-\(m\) polynomial template microcells \(E_{Q\ell}\) within each \(Q\in\mathcal Q_h\), define their treated masses and each cube's weakest mass by
\[
q_{Q\ell}=P(A=1,X\in E_{Q\ell}),
\qquad
q_Q=\min_{1\le\ell\le J}q_{Q\ell},
\]
where \(J\) is the number of template microcells. Let \(q_{(k)}\) be the \(k\)th entry of the list of these \(q_Q\)'s sorted in nondecreasing order. Then, for every integer \(1\le k\le(2^j)^d\),
\[
q_{(k)}
\ge
\kappa_\Gamma h^{D_\gamma}k^{1/(\gamma-1)}.
\]
\end{lemmav}

\begin{proof}[Proof of \cref{obj:lem:ordered-mass}]
\begin{enumerate}
\item \emph{Uniform coefficient.}
Set
\[
m:=\lceil\beta\rceil-1,
\qquad
a_{\mathrm{tpl}}:=c_f\eta^d,
\]
using the template of \cref{obj:synth_1,obj:synth_2,obj:def:template}. Its polynomial coordinates and node set are
\[
\mathcal I_m=\{\alpha\in\mathbb N_0^d:|\alpha|\le m\},
\qquad
U(z)=(z^\alpha)_{\alpha\in\mathcal I_m},
\qquad
J=(m+1)^d,
\]
\[
\{v_\ell:1\le\ell\le J\}
=
\left\{
\left(\frac{\iota_i+1}{m+2}\right)_{i=1}^d:
\iota\in\{0,\ldots,m\}^d
\right\}.
\]
We first establish positivity of the template width. Define
\[
\mathcal B_m:=
\left\{
b\in\mathbb R^{\mathcal I_m}:
\sum_{\alpha\in\mathcal I_m}b_\alpha^2=1
\right\},
\qquad
\pi_b(z):=\sum_{\alpha\in\mathcal I_m}b_\alpha z^\alpha
\quad(z\in\mathbb R^d).
\]
For \(b\in\mathcal B_m\), at least one coefficient is nonzero, and
\[
\pi_b\not\equiv0,
\qquad
\deg_{z_i}\pi_b\le m
\quad(1\le i\le d).
\]
The \(m+1\) node values in each coordinate are distinct. If \(\pi_b\) vanished at every tensor node, fixing the other coordinates at node values and applying the univariate root bound would make it vanish identically in the first coordinate. Repeating this argument in each coordinate would give \(\pi_b\equiv0\). Thus
\[
\exists\,\ell\in\{1,\ldots,J\}:\ \pi_b(v_\ell)\ne0,
\qquad
b^\top G_0b
=
\frac1J\sum_{\ell=1}^J\pi_b(v_\ell)^2>0.
\]
% lean: CausalSmith.Stat.WeakOverlap.templateGram_unit_pos

The set \(\mathcal B_m\) is nonempty, since it contains the coefficient vector of the constant polynomial \(1\). It is closed and bounded: for every \(b\in\mathcal B_m\),
\[
b_\alpha^2\le
\sum_{\alpha'\in\mathcal I_m}b_{\alpha'}^2=1
\quad(\alpha\in\mathcal I_m).
\]
Consequently it is compact. The continuous quadratic form attains its minimum there, and the Rayleigh characterization gives
\[
\lambda_0
=
\min_{b\in\mathcal B_m}b^\top G_0b>0.
\]
% lean: CausalSmith.Stat.WeakOverlap.templateLambda_pos

Every quotient defining \(L_U\) is nonnegative; the indexing set contains distinct cube points because \(d\ge1\). Hence the width rule supplies
\[
L_U\ge0,
\qquad
0<\eta
=
\min\left\{\frac1{2(m+2)},\frac{\lambda_0}{1+L_U}\right\}
\le\frac1{2(m+2)},
\qquad
a_{\mathrm{tpl}}>0.
\]
% lean: CausalSmith.Stat.WeakOverlap.templateLip_nonneg

For \(\gamma'\in\Gamma\), define
\[
\Phi_\Gamma(\gamma')
:=
\frac{\gamma'-1}{\gamma'}\,
C^{-1/(\gamma'-1)}
a_{\mathrm{tpl}}^{\gamma'/(\gamma'-1)}.
\]
Since \(C\ge1\), \(a_{\mathrm{tpl}}>0\), and
\(\gamma'\ge\gamma_{\min}>1\), this function is continuous and strictly positive on the nonempty compact interval \(\Gamma\). Therefore
\[
\kappa_\Gamma:=\min_{\gamma'\in\Gamma}\Phi_\Gamma(\gamma')>0,
\qquad
\kappa_\Gamma\le
\frac{\gamma-1}{\gamma}\,
C^{-1/(\gamma-1)}
a_{\mathrm{tpl}}^{\gamma/(\gamma-1)}
\quad(\gamma\in\Gamma).
\]
This constant depends only on the fixed parameters and interval endpoints.
% lean: CausalSmith.Stat.WeakOverlap.orderedMass_uniform_coefficient

\item \emph{Setwise treated mass.}
Fix \(\gamma\in\Gamma\) and \(P\in\mathcal P_\gamma\), and write
\[
\nu:=P_X.
\]
Then \(\gamma\ge\gamma_{\min}>1\). Let \(E\subseteq\mathcal X\) be measurable and set
\[
u:=\nu(E),
\qquad 0\le u\le1.
\]
If \(u=0\), the claimed lower bound is zero, while the treated probability is nonnegative. Suppose henceforth that \(u>0\), and define
\[
s_E:=\left(\frac uC\right)^{1/(\gamma-1)}.
\]
The inequalities \(0<u/C\le1\) and \(1/(\gamma-1)>0\) give
\[
0<s_E\le1.
\]
% lean: CausalSmith.Stat.WeakOverlap.orderedMass_threshold_mem_unit

Define the weak level-set mass function on all real thresholds by
\[
F_E(t):=\nu\bigl(E\cap\{x:t\le e(x)\}\bigr),
\qquad t\in\mathbb R.
\]
It is nonnegative and nonincreasing:
\[
0\le F_E(t')\le F_E(t)\le u
\qquad(t\le t').
\]
As a version of the conditional treatment probability, the propensity satisfies
\[
0\le e(x)\le1
\quad\text{for }\nu\text{-almost every }x.
\]
Disintegration of the treatment probability therefore gives
\[
P(A=1,X\in E)=\int_E e(x)\,d\nu(x).
\]
Moreover,
\[
e(x)=\int_0^1\mathbf 1\{t\le e(x)\}\,dt
\quad\text{for }\nu\text{-almost every }x.
\]
Tonelli's theorem yields the weak-level-set representation
\[
P(A=1,X\in E)=\int_0^1F_E(t)\,dt.
\]
The propensity is integrable on \(E\), and the bounded monotone level-set functions are integrable on the threshold intervals used below.
% lean: CausalSmith.Stat.WeakOverlap.orderedMass_treated_real_integral

For comparison, define
\[
F_E^{>}(t):=\nu\bigl(E\cap\{x:t<e(x)\}\bigr),
\qquad t\in\mathbb R.
\]
For \(t\in[0,1]\), \cref{obj:ass:global-tail} and monotonicity of measure give
\[
\nu\bigl(E\cap\{e\le t\}\bigr)
\le \nu(\{e\le t\})
\le Ct^{\gamma-1}.
\]
The two threshold sets cover \(E\), so subadditivity implies
\[
u\le F_E^{>}(t)+\nu\bigl(E\cap\{e\le t\}\bigr).
\]
Since \(\{e>t\}\subseteq\{e\ge t\}\), it follows that
\[
u-Ct^{\gamma-1}\le F_E^{>}(t)\le F_E(t)
\qquad(0\le t\le1).
\]
Integrating over the cutoff interval gives
\[
\int_0^{s_E}(u-Ct^{\gamma-1})\,dt
\le
\int_0^{s_E}F_E^{>}(t)\,dt
\le
\int_0^{s_E}F_E(t)\,dt.
\]
% lean: CausalSmith.Stat.WeakOverlap.orderedMass_layercake_piece_lower

The cutoff balances the tail envelope and set mass:
\[
Cs_E^{\gamma-1}
=
C\left(\frac uC\right)^{(1/(\gamma-1))(\gamma-1)}
=u.
\]
Also \(s_E^\gamma=s_E^{\gamma-1}s_E\). Thus direct integration gives
\[
\begin{aligned}
\int_0^{s_E}(u-Ct^{\gamma-1})\,dt
&=us_E-\frac C\gamma s_E^\gamma\\
&=\frac{\gamma-1}{\gamma}\,us_E\\
&=\frac{\gamma-1}{\gamma}\,
C^{-1/(\gamma-1)}u^{\gamma/(\gamma-1)}.
\end{aligned}
\]
% lean: CausalSmith.Stat.WeakOverlap.orderedMass_integral_at_threshold

Combining these bounds, using \(F_E\ge0\) and \(s_E\le1\), proves
\[
\begin{aligned}
\frac{\gamma-1}{\gamma}\,
C^{-1/(\gamma-1)}u^{\gamma/(\gamma-1)}
&\le \int_0^{s_E}F_E(t)\,dt\\
&\le \int_0^1F_E(t)\,dt\\
&=P(A=1,X\in E).
\end{aligned}
\]
% lean: CausalSmith.Stat.WeakOverlap.orderedMass_setwise_lower

\item \emph{A ranked union and its upper bound.}
Fix \(j\in\mathbb N\) and an integer \(k\) in the stated range, and set
\[
h:=2^{-j}>0,
\qquad
D_\gamma:=\frac{d\gamma}{\gamma-1}.
\]
Use the dyadic partition \(\mathcal Q_h\), whose coordinate intervals are half-open with the global right boundary included. For each cube, write its lower corner and scaled microcells as
\[
a_Q:=\left(\inf_{x\in Q}x_i\right)_{i=1}^d,
\qquad
E_{Q\ell}:=a_Q+hV_\ell
\quad(1\le\ell\le J).
\]
Because the finite collection of microcell masses is nonempty, choose
\[
\ell_{\min}:\mathcal Q_h\longrightarrow\{1,\ldots,J\},
\qquad
q_{Q,\ell_{\min}(Q)}=q_Q
\quad(Q\in\mathcal Q_h).
\]
Define the threshold collection
\[
\mathcal Q_{h,\le k}
:=
\{Q\in\mathcal Q_h:q_Q\le q_{(k)}\}.
\]
The first \(k\) entries of the sorted mass list are all at most \(q_{(k)}\). Sorting preserves the multiplicity of each mass, so
\[
|\mathcal Q_{h,\le k}|\ge k.
\]
Choose a collection and define its selected union by
\[
\mathcal Q_{h,k}\subseteq\mathcal Q_{h,\le k},
\qquad
|\mathcal Q_{h,k}|=k,
\qquad
E_k:=\bigcup_{Q\in\mathcal Q_{h,k}}E_{Q,\ell_{\min}(Q)}.
\]
Subadditivity of the treated probability then gives
\[
\begin{aligned}
P(A=1,X\in E_k)
&\le
\sum_{Q\in\mathcal Q_{h,k}}
q_{Q,\ell_{\min}(Q)}\\
&=
\sum_{Q\in\mathcal Q_{h,k}}q_Q\\
&\le kq_{(k)}.
\end{aligned}
\]
% lean: CausalSmith.Stat.WeakOverlap.orderedMass_ranked_union_upper

\item \emph{The selected union's lower bound.}
The node coordinates and width bound satisfy
\[
\frac1{m+2}\le(v_\ell)_i\le\frac{m+1}{m+2},
\qquad
\frac\eta2\le\frac1{4(m+2)}.
\]
Hence, for every \(z\in V_\ell\) and every coordinate \(i\),
\[
0<(v_\ell)_i-\frac\eta2
\le z_i
\le(v_\ell)_i+\frac\eta2<1.
\]
After scaling, every microcell lies strictly between its cube's coordinate faces:
\[
(a_Q)_i<x_i<(a_Q)_i+h
\quad(x\in E_{Q\ell},\ 1\le i\le d).
\]
In particular, \(E_{Q\ell}\subseteq Q\subseteq\mathcal X\). Distinct dyadic cubes have lower corners differing by at least \(h\) in some coordinate. The strict bounds therefore imply
\[
E_{Q\ell}\cap E_{Q'\ell'}=\varnothing
\quad(Q\ne Q').
\]
% lean: CausalSmith.Stat.WeakOverlap.orderedMass_scaledMicroCell_disjoint_macro

More explicitly, each microcell is the rectangle
\[
E_{Q\ell}
=
\prod_{i=1}^d
\left[
(a_Q)_i+h\left((v_\ell)_i-\frac\eta2\right),
(a_Q)_i+h\left((v_\ell)_i+\frac\eta2\right)
\right].
\]
Each coordinate interval has length \(\eta h\), so
\[
\operatorname{Leb}(E_{Q\ell})=(\eta h)^d.
\]
% lean: CausalSmith.Stat.WeakOverlap.orderedMass_scaledMicroCell_volume

By \cref{obj:ass:density}, integration of the density over each microcell yields
\[
\nu(E_{Q\ell})
=
\int_{E_{Q\ell}}f(x)\,dx
\ge c_f\operatorname{Leb}(E_{Q\ell})
=c_f(\eta h)^d.
\]
The selected microcells are measurable and pairwise disjoint. Thus
\[
\begin{aligned}
\nu(E_k)
&=
\sum_{Q\in\mathcal Q_{h,k}}
\nu(E_{Q,\ell_{\min}(Q)})\\
&\ge k\,c_f(\eta h)^d.
\end{aligned}
\]
% lean: CausalSmith.Stat.WeakOverlap.orderedMass_selectedMicrocells_covariate_lower

The exponent \(\gamma/(\gamma-1)\) and the setwise coefficient are positive. Applying monotonicity of this power and the setwise estimate already established gives
\[
\begin{aligned}
&\frac{\gamma-1}{\gamma}\,
C^{-1/(\gamma-1)}
\bigl[k\,c_f(\eta h)^d\bigr]^{\gamma/(\gamma-1)}\\
&\qquad\le
\frac{\gamma-1}{\gamma}\,
C^{-1/(\gamma-1)}
\nu(E_k)^{\gamma/(\gamma-1)}
\le P(A=1,X\in E_k).
\end{aligned}
\]
% lean: CausalSmith.Stat.WeakOverlap.orderedMass_selectedMicrocells_treated_lower

\item \emph{Power identity and conclusion.}
The factorization and exponent identity are
\[
c_f(\eta h)^d=a_{\mathrm{tpl}}h^d,
\qquad
\frac{\gamma}{\gamma-1}
=
1+\frac1{\gamma-1}.
\]
Using multiplicativity of real powers on positive bases and \(k>0\), we obtain
\[
\begin{aligned}
\frac{\bigl[k\,c_f(\eta h)^d\bigr]^{\gamma/(\gamma-1)}}{k}
&=
a_{\mathrm{tpl}}^{\gamma/(\gamma-1)}
h^{d\gamma/(\gamma-1)}
k^{\gamma/(\gamma-1)-1}\\
&=
a_{\mathrm{tpl}}^{\gamma/(\gamma-1)}
h^{D_\gamma}k^{1/(\gamma-1)}.
\end{aligned}
\]
% lean: CausalSmith.Stat.WeakOverlap.orderedMass_rank_power_identity

Multiplying the uniform coefficient bound by the positive mesh and rank factors therefore gives
\[
\kappa_\Gamma h^{D_\gamma}k^{1/(\gamma-1)}
\le
\frac{
\frac{\gamma-1}{\gamma}\,
C^{-1/(\gamma-1)}
\bigl[k\,c_f(\eta h)^d\bigr]^{\gamma/(\gamma-1)}
}{k}.
\]
The lower and upper bounds on the same selected union give
\[
\frac{\gamma-1}{\gamma}\,
C^{-1/(\gamma-1)}
\bigl[k\,c_f(\eta h)^d\bigr]^{\gamma/(\gamma-1)}
\le P(A=1,X\in E_k)
\le kq_{(k)}.
\]
Dividing by \(k>0\) and combining the two inequalities proves
\[
q_{(k)}
\ge
\kappa_\Gamma h^{D_\gamma}k^{1/(\gamma-1)}.
\]
The choices of \(\gamma\), \(P\), \(j\), and \(k\) were arbitrary in their stated ranges, so the bound holds uniformly as claimed.
% lean: CausalSmith.Stat.WeakOverlap.orderedTreatedMass_lower
\end{enumerate}
\end{proof}

The rank factor quantifies how information improves across the ordered cubes. At a fixed mesh, the weakest treated mass grows at least polynomially with rank, so the corresponding variance envelope decreases with rank. The constant \(\kappa_\Gamma\) is common to the stated compact overlap range; this uniformity carries the mass comparison into the adaptive upper bound.

\subsection{Deterministic conditioning of the template}

Write \(U(x)\) for the vector of monomials of total degree at most \(m\), with \(m=\lceil\beta\rceil-1\) in the statistical application. The tensor nodes \(v_\ell\), reference microcubes \(V_\ell\), and reference Gram matrix \(G_0\) are those of \cref{obj:def:template}. Its least eigenvalue is denoted by \(\lambda_0\). The empirical matrix \(\widehat G_Q\) assigns equal weight to the microcells and averages within each cell over its treated observations. The next statement gives the conditioning bound for every feasible sampled configuration.

\begin{lemmav}{lem:template-conditioning}[Uniform equal-cell Gram conditioning]
Let \(d,m,n\) be nonnegative integers, let \(\beta\in\mathbb R\), and let \(\mathcal D_n=((X_i,A_i,Y_i))_{i=1}^n\) be any sample with \(X_i\in\mathbb R^d\), \(A_i\in\{0,1\}\), and \(Y_i\in\mathbb R\). Write \(U(x)=(x^\alpha)_{|\alpha|\le m}\) for the vector of total-degree-\(m\) monomials. For the polynomial template's tensor nodes \(v_\ell\) and reference microcubes \(V_\ell\), indexed by \(\ell\in\{0,\ldots,m\}^d\), define
\[
J=(m+1)^d,\qquad
G_0=\frac1J\sum_{\ell\in\{0,\ldots,m\}^d}
U(v_\ell)U(v_\ell)^\top,\qquad
\lambda_0=\inf_{\|b\|_2=1}b^\top G_0b.
\]
Here \(\|b\|_2=(\sum_{|\alpha|\le m}b_\alpha^2)^{1/2}\).

Let \(h=2^{-j}\), with \(j\) a nonnegative integer, be a feasible candidate dyadic mesh: \(h\in\mathcal H_n\), the candidate mesh collection, and
\[
N_{Q\ell}\ge h^{-2\beta}
\quad\text{for every }Q\in\mathcal Q_h
\text{ and }\ell\in\{0,\ldots,m\}^d,
\]
where \(\mathcal Q_h\) is the dyadic cube partition at mesh \(h\). For a cube \(Q\) with corner \(a_Q\), the microcell and its treated count are
\[
E_{Q\ell}=a_Q+hV_\ell,\qquad
N_{Q\ell}=\sum_{i=1}^n
\mathbf1\{A_i=1,\ X_i\in E_{Q\ell}\}.
\]
Define the equal-cell empirical Gram matrix by
\[
\widehat G_Q
=\frac1J\sum_{\ell\in\{0,\ldots,m\}^d}
\frac1{N_{Q\ell}}
\sum_{i=1}^n
\mathbf1\{A_i=1,\ X_i\in E_{Q\ell}\}
U\!\left(\frac{X_i-a_Q}{h}\right)
U\!\left(\frac{X_i-a_Q}{h}\right)^\top.
\]
Then \(\lambda_0>0\), and for every \(Q\in\mathcal Q_h\),
\[
\frac{\lambda_0}{2}\|b\|_2^2
\le b^\top\widehat G_Qb
\quad\text{for every }b=(b_\alpha)_{|\alpha|\le m},
\qquad
\det(\widehat G_Q)\ne0,
\qquad
\|\widehat G_Q^{-1}\|_{\mathrm{op}}
\le\frac2{\lambda_0},
\]
where \(\|M\|_{\mathrm{op}}=\sup_{\|b\|_2=1}\|Mb\|_2\).
\end{lemmav}

\begin{proof}[Proof of \cref{obj:lem:template-conditioning}]
\begin{enumerate}
\item Fix \(Q\in\mathcal Q_h\). Since \(h=2^{-j}>0\), feasibility gives
\[
0<h^{-2\beta}\le N_{Q\ell}
\qquad\text{for every }\ell\in\{0,\ldots,m\}^d.
\]
Thus every microcell count is positive, and its reciprocal is well defined. We establish the conditioning bounds using these positive counts.
% lean: CausalSmith.Stat.WeakOverlap.gramHat_minEig_ge

\item Define the coefficient and node index sets, and the polynomial associated with a coefficient vector, by
\[
\mathcal I_m=\{\alpha\in\mathbb N_0^d:|\alpha|\le m\},
\qquad
\mathcal T_m=\{0,\ldots,m\}^d,
\qquad
r_b(z)=\sum_{\alpha\in\mathcal I_m}b_\alpha z^\alpha
\quad(z\in\mathbb R^d).
\]
Expanding the definition of \(G_0\) yields
\[
b^\top G_0b=\frac1J\sum_{\ell\in\mathcal T_m}r_b(v_\ell)^2\ge0.
\]

The nodes in \cref{obj:synth_1} form the Cartesian product of the one-dimensional set
\[
\mathcal S_{\mathrm{node}}
=\left\{\frac{a+1}{m+2}:a=0,\ldots,m\right\},
\qquad
|\mathcal S_{\mathrm{node}}|=m+1.
\]
Each coordinate degree of \(r_b\) is at most \(m\). Consequently, vanishing at every tensor node implies that \(r_b\) is the zero polynomial. Indeed, fixing all but one coordinate at node values gives a univariate polynomial with \(m+1\) distinct roots and degree at most \(m\), so it vanishes identically. Repeating this argument coordinate by coordinate extends the vanishing to all of \(\mathbb R^d\). Hence
\[
\bigl[r_b(v_\ell)=0\text{ for every }\ell\in\mathcal T_m\bigr]
\Longrightarrow r_b\equiv0
\Longrightarrow b_\alpha=0\text{ for every }\alpha\in\mathcal I_m.
\]
For \(d=0\), both index sets have one element and \(r_b\) is constant, so the same implication follows directly. A unit coefficient vector has a nonzero coordinate; therefore at least one summand is strictly positive, and
\[
\|b\|_2=1\quad\Longrightarrow\quad b^\top G_0b>0.
\]

To obtain strict positivity of the infimum, define
\[
\mathcal S_{\mathrm{coef}}
=\left\{b\in\mathbb R^{\mathcal I_m}:
\sum_{\alpha\in\mathcal I_m}b_\alpha^2=1\right\},
\qquad
(b_{\mathrm{const}})_\alpha=
\begin{cases}
1,&\alpha=0,\\
0,&\alpha\ne0.
\end{cases}
\]
The zero multi-index belongs to \(\mathcal I_m\) for every \(d,m\), and
\(\|b_{\mathrm{const}}\|_2=1\). Thus \(\mathcal S_{\mathrm{coef}}\) is nonempty.
It is closed, and each of its coordinates satisfies
\[
|b_\alpha|^2\le\sum_{\alpha'\in\mathcal I_m}b_{\alpha'}^2=1,
\]
so it is bounded and hence compact. The continuous quadratic form attains its minimum there: for some \(b_{\min}\in\mathcal S_{\mathrm{coef}}\),
\[
\lambda_0
=\min_{b\in\mathcal S_{\mathrm{coef}}}b^\top G_0b
=b_{\min}^\top G_0b_{\min}>0.
\]
% lean: CausalSmith.Stat.WeakOverlap.templateLambda_pos

\item Fix an arbitrary \(b\in\mathbb R^{\mathcal I_m}\). Write the supremum norm, including its zero-dimensional convention, as
\[
\|z\|_\infty
=\max\bigl(\{0\}\cup\{|z_a|:1\le a\le d\}\bigr).
\]
The template Lipschitz constant described in \cref{obj:synth_2} is
\[
L_U=
\sup_{\substack{z,z'\in[0,1]^d\\z\ne z'}}
\frac{\|U(z)U(z)^\top-U(z')U(z')^\top\|_{\mathrm{op}}}
{\|z-z'\|_\infty},
\]
with value zero when the set of quotients is empty. For \(d\ge1\),
\cref{obj:synth_2} supplies its finiteness; for \(d=0\), the cube is a singleton and the stated convention applies. Every quotient is nonnegative, so \(L_U\ge0\).

By \cref{obj:def:template}, the width and microcubes satisfy
\[
\eta=\min\left\{\frac1{2(m+2)},\frac{\lambda_0}{1+L_U}\right\}>0,
\qquad
V_\ell=v_\ell+[-\eta/2,\eta/2]^d.
\]
In particular,
\[
\frac{\eta}{2}\le\frac1{4(m+2)}\le\frac1{m+2},
\qquad
\frac1{m+2}\le(v_\ell)_a
\le1-\frac1{m+2}.
\]
The coordinate inequalities are vacuous when \(d=0\). They imply, for every \(z\in V_\ell\),
\[
v_\ell,z\in[0,1]^d,
\qquad
\|z-v_\ell\|_\infty\le\frac{\eta}{2}.
\]

Define the perturbation matrix and the coefficient-dependent error by
\[
M_\ell(z)
=U(v_\ell)U(v_\ell)^\top-U(z)U(z)^\top,
\qquad
\varepsilon_b=L_U\frac{\eta}{2}\|b\|_2^2.
\]
If \(z=v_\ell\), the desired squared-evaluation inequality follows immediately from \(\varepsilon_b\ge0\). If \(z\ne v_\ell\), the definition of \(L_U\) and the preceding distance bound give
\[
\frac{\|M_\ell(z)\|_{\mathrm{op}}}{\|v_\ell-z\|_\infty}
\le L_U,
\qquad
\|M_\ell(z)\|_{\mathrm{op}}
\le L_U\|v_\ell-z\|_\infty
\le L_U\frac{\eta}{2}.
\]
For \(b=0\), the quadratic-form bound below is immediate. For \(b\ne0\), define
\[
b_{\mathrm{unit}}=\frac{b}{\|b\|_2},
\qquad
\|b_{\mathrm{unit}}\|_2=1.
\]
The operator-norm definition and homogeneity give
\[
\|M_\ell(z)b\|_2
=\|b\|_2\|M_\ell(z)b_{\mathrm{unit}}\|_2
\le\|M_\ell(z)\|_{\mathrm{op}}\|b\|_2.
\]
Cauchy--Schwarz therefore yields
\[
\begin{aligned}
r_b(v_\ell)^2-r_b(z)^2
&=b^\top M_\ell(z)b\\
&\le\|b\|_2\|M_\ell(z)b\|_2\\
&\le\|M_\ell(z)\|_{\mathrm{op}}\|b\|_2^2
\le\varepsilon_b.
\end{aligned}
\]
Combining the two cases for \(z\), we obtain
\[
r_b(v_\ell)^2-\varepsilon_b\le r_b(z)^2
\qquad(z\in V_\ell).
\]
% lean: CausalSmith.Stat.WeakOverlap.templatePointwise_square_perturbation

\item Define the normalized observations and the treated index sets by
\[
z_{Qi}=\frac{X_i-a_Q}{h}\quad(1\le i\le n),
\qquad
\mathcal I_{Q\ell}
=\{i\in\{1,\ldots,n\}:A_i=1,\ X_i\in E_{Q\ell}\}.
\]
Then \(|\mathcal I_{Q\ell}|=N_{Q\ell}>0\). Since \(h>0\) and
\(E_{Q\ell}=a_Q+hV_\ell\), every \(i\in\mathcal I_{Q\ell}\) satisfies
\(z_{Qi}\in V_\ell\). Averaging the pointwise inequality within each cell gives
\[
\begin{aligned}
r_b(v_\ell)^2-\varepsilon_b
&=\frac1{N_{Q\ell}}\sum_{i\in\mathcal I_{Q\ell}}
\bigl(r_b(v_\ell)^2-\varepsilon_b\bigr)\\
&\le\frac1{N_{Q\ell}}\sum_{i\in\mathcal I_{Q\ell}}r_b(z_{Qi})^2.
\end{aligned}
\]
Expansion of the empirical Gram matrix, followed by equal averaging across the \(J>0\) cells, now gives
\[
\begin{aligned}
b^\top\widehat G_Qb
&=\frac1J\sum_{\ell\in\mathcal T_m}
\frac1{N_{Q\ell}}\sum_{i\in\mathcal I_{Q\ell}}r_b(z_{Qi})^2\\
&\ge\frac1J\sum_{\ell\in\mathcal T_m}
\bigl(r_b(v_\ell)^2-\varepsilon_b\bigr)\\
&=b^\top G_0b-\varepsilon_b.
\end{aligned}
\]
The last equality uses \(|\mathcal T_m|=J\), so the equally averaged error remains \(\varepsilon_b\).
% lean: CausalSmith.Stat.WeakOverlap.templateGram_quadratic_sub_error_le_gramHat

\item For \(b=0\), the template lower bound is an equality. For \(b\ne0\), the unit vector defined above and the defining infimum give
\[
\lambda_0\le b_{\mathrm{unit}}^\top G_0b_{\mathrm{unit}},
\qquad
\lambda_0\|b\|_2^2
\le\|b\|_2^2b_{\mathrm{unit}}^\top G_0b_{\mathrm{unit}}
=b^\top G_0b.
\]
The width rule also implies
\[
\eta\le\frac{\lambda_0}{1+L_U},
\qquad
L_U\eta\le\eta(1+L_U)\le\lambda_0.
\]
Consequently,
\[
\varepsilon_b
=L_U\frac{\eta}{2}\|b\|_2^2
\le\frac{\lambda_0}{2}\|b\|_2^2.
\]
Combining these inequalities with the equal-cell transfer proves
\[
b^\top\widehat G_Qb
\ge b^\top G_0b-\varepsilon_b
\ge\lambda_0\|b\|_2^2-\frac{\lambda_0}{2}\|b\|_2^2
=\frac{\lambda_0}{2}\|b\|_2^2.
\]
% lean: CausalSmith.Stat.WeakOverlap.gramHat_minEig_ge_of_pos

\item If \(\det(\widehat G_Q)=0\), there is a nonzero coefficient vector \(b_{\ker}\) with
\[
\widehat G_Qb_{\ker}=0,
\qquad
\|b_{\ker}\|_2>0.
\]
Applying the quadratic lower bound gives the contradiction
\[
0<\frac{\lambda_0}{2}\|b_{\ker}\|_2^2
\le b_{\ker}^\top\widehat G_Qb_{\ker}=0.
\]
Thus \(\det(\widehat G_Q)\ne0\).
% lean: CausalSmith.Stat.WeakOverlap.det_ne_zero_of_quadratic_lower

\item For an arbitrary input vector \(b_{\mathrm{in}}\in\mathbb R^{\mathcal I_m}\), define
\[
b_{\mathrm{inv}}=\widehat G_Q^{-1}b_{\mathrm{in}}.
\]
Invertibility gives \(\widehat G_Qb_{\mathrm{inv}}=b_{\mathrm{in}}\). The quadratic lower bound and Cauchy--Schwarz imply
\[
\frac{\lambda_0}{2}\|b_{\mathrm{inv}}\|_2^2
\le b_{\mathrm{inv}}^\top\widehat G_Qb_{\mathrm{inv}}
=b_{\mathrm{inv}}^\top b_{\mathrm{in}}
\le\|b_{\mathrm{inv}}\|_2\|b_{\mathrm{in}}\|_2.
\]
If \(\|b_{\mathrm{inv}}\|_2=0\), the following bound is immediate; otherwise, divide by its positive norm and then by \(\lambda_0/2>0\). In both cases,
\[
\|\widehat G_Q^{-1}b_{\mathrm{in}}\|_2
=\|b_{\mathrm{inv}}\|_2
\le\frac2{\lambda_0}\|b_{\mathrm{in}}\|_2.
\]
The coefficient unit sphere is nonempty, as witnessed by \(b_{\mathrm{const}}\). Taking the supremum over its vectors therefore gives
\[
\|\widehat G_Q^{-1}\|_{\mathrm{op}}
=\sup_{\|b_{\mathrm{in}}\|_2=1}
\|\widehat G_Q^{-1}b_{\mathrm{in}}\|_2
\le\frac2{\lambda_0}.
\]
Since \(Q\) was arbitrary, all the asserted bounds hold throughout \(\mathcal Q_h\).
% lean: CausalSmith.Stat.WeakOverlap.finiteL2OpNorm_inv_le_of_quadratic_lower
\end{enumerate}
\end{proof}

Equal weighting preserves the reference template's conditioning even when the treated counts differ across microcells. The bound is deterministic once feasibility holds, and it controls the amplification of both approximation residuals and outcome fluctuations in the fitted coefficients. Its geometric requirement is supplied by the occupied template microcells and their prescribed widths.

\subsection{Selected-mesh concentration and maximal error}

Let \(\mathcal D_n=((X_i,A_i,Y_i))_{i=1}^n\) be the observed sample. Conditional expectations below condition on the sampled design \((X_i,A_i)_{i=1}^n\), comprising the covariates and treatment indicators. Under this conditioning, the selected mesh of \cref{obj:def:mesh-selector} is fixed. The coefficient vector \(\widehat c_Q\) is the equal-cell polynomial fit used in \cref{obj:def:estimator}; centering it by its conditional expectation isolates the stochastic coefficient error.

The oracle mesh is \(h_{\gamma,n}=n^{-1/(2\beta+D_\gamma)}\), and the oracle response-recovery scale is \(r_{\gamma,n}=h_{\gamma,n}^{\beta}\). The sample event \(\mathcal E_n\) below controls both selection and the comparison of empirical counts with population treated masses. For the accompanying maximal inequality, \(Z_1,\ldots,Z_N\) denotes a generic family of centered real variables under a probability law \(\mu\). The positive quantities \(a\) and \(b\) specify a common scale cap and an ordered decay scale, respectively.

\begin{lemmav}{lem:selected-mesh}[Selected mesh and coefficient envelope]
Fix the dimension \(d\), smoothness parameter \(\beta\), model constants \(B,L,C,c_f\), and adaptation endpoints \(\gamma_{\min},\gamma_{\max}\), satisfying
\begin{itemize}
\item \textbf{(Dimension and smoothness.)} \(d\in\mathbb N\), \(d\ge1\), and \(\beta>1\).
\item \textbf{(Model constants.)} \(B>0\), \(L>0\), \(C\ge1\), and \(0<c_f\le1\).
\item \textbf{(Adaptation interval.)} \(1<\gamma_{\min}<\gamma_{\max}\), with
\[
\Gamma=[\gamma_{\min},\gamma_{\max}].
\]
\end{itemize}
There exist constants \(K,K'>0\) and \(n_0\in\mathbb N\) such that the following holds uniformly for every \(n\ge n_0\), \(\gamma\in\Gamma\), and observed law \(P\in\mathcal P_\gamma\), with the displayed fixed model constants. Let the sample have the iid law \(P^{\otimes n}\) specified in \cref{obj:ass:iid}, and use the mesh selector of \cref{obj:def:mesh-selector}. Write
\[
m=\lceil\beta\rceil-1,
\qquad
h_{\gamma,n}=n^{-1/(2\beta+D_\gamma)},
\]
where \(D_\gamma\) is the effective dimension governing the model's overlap-dependent rate, and let \(r_{\gamma,n}\) denote its oracle response-recovery risk scale.

For each partition cube \(Q\) and template microcell \(E_{Q\ell}\), define the population treated mass, its empirical count, and the weakest treated microcell mass by
\[
q_{Q\ell}=P\{A=1,\ X\in E_{Q\ell}\},
\qquad
N_{Q\ell}=\sum_{i=1}^n
\mathbf1\{A_i=1,\ X_i\in E_{Q\ell}\},
\qquad
q_Q=\min_\ell q_{Q\ell}.
\]
There is a measurable sample event \(\mathcal E_n\) such that
\[
P^{\otimes n}(\mathcal E_n^c)\le n^{-2}.
\]
For every sample in \(\mathcal E_n\), the selected mesh satisfies
\[
0<\widehat h\le K h_{\gamma,n},
\qquad
N_{Q\ell}\ge \frac{nq_{Q\ell}}5
\quad
\text{for every }Q\in\mathcal Q_{\widehat h}
\text{ and every }\ell.
\]

Order the selected cubes by increasing \(q_Q\), with ranks starting at \(1\). Conditional on the sampled design \((X_i,A_i)_{i=1}^n\), evaluate all coefficient estimators at the selected mesh of the given sample. For the cube of rank \(k\), every monomial coefficient indexed by a multiindex \(\alpha\) of total degree at most \(m\), and every \(t\in\mathbb R\),
\[
\begin{aligned}
&\mathbb E\!\left[
\exp\!\left\{
t\left(
\widehat c_{Q,\alpha}
-\mathbb E[\widehat c_{Q,\alpha}\mid(X_i,A_i)_{i=1}^n]
\right)
\right\}
\,\middle|\,(X_i,A_i)_{i=1}^n
\right]
\\
&\hspace{2em}\le
\exp\!\left\{
\frac{Kt^2}{2}
\min\!\left(
\widehat h^{2\beta},
n^{-1}\widehat h^{-D_\gamma}k^{-1/(\gamma-1)}
\right)
\right\}.
\end{aligned}
\]
These conditional exponential moments are integrable for every coefficient and every \(t\), and the maximum coefficient-vector error is conditionally integrable. With \(\|\cdot\|_\infty\) denoting the maximum absolute coordinate over those monomial indices, its conditional expectation satisfies
\[
\begin{aligned}
&\mathbb E\!\left[
\max_{Q\in\mathcal Q_{\widehat h}}
\left\|
\widehat c_Q-
\mathbb E[\widehat c_Q\mid(X_i,A_i)_{i=1}^n]
\right\|_\infty
\,\middle|\,(X_i,A_i)_{i=1}^n
\right]
\\
&\hspace{2em}\le
K\widehat h^\beta
\sqrt{1+\max\{0,\log(h_{\gamma,n}/\widehat h)\}}
\le K'r_{\gamma,n}.
\end{aligned}
\]

Moreover, for every exponent \(\alpha>0\), there exists \(K_\alpha>0\) such that the following maximal inequality holds. Let \(\mu\) be a probability law on a measurable space, let \(N\in\mathbb N\), and let \(Z_1,\ldots,Z_N\) be real random variables satisfying
\begin{itemize}
\item \textbf{(Scales.)} \(a>0\) and \(b>0\).
\item \textbf{(Exponential integrability.)} For every \(1\le k\le N\) and \(t\in\mathbb R\), \(\exp(tZ_k)\) is integrable under \(\mu\).
\item \textbf{(Moment envelope.)} For every \(1\le k\le N\) and \(t\in\mathbb R\),
\[
\mathbb E_\mu[\exp(tZ_k)]
\le
\exp\!\left\{
\frac{t^2}{2}\min(a^2,b^2k^{-\alpha})
\right\}.
\]
\end{itemize}
Then
\[
\mathbb E_\mu\!\left[
\max\bigl(\{0\}\cup\{|Z_k|:1\le k\le N\}\bigr)
\right]
\le
K_\alpha a
\sqrt{1+\max\!\left\{0,\log\!\left((b/a)^{2/\alpha}\right)\right\}}.
\]
\end{lemmav}

\begin{proof}[Proof of \cref{obj:lem:selected-mesh}]
Throughout the proof, write
\[
m=\lceil\beta\rceil-1,\qquad
\mathcal I_m=\{\alpha\in\mathbb N_0^d:|\alpha|\le m\},
\qquad
p_{\mathrm{coef}}=|\mathcal I_m|,
\qquad
J=(m+1)^d,
\]
and use the template width \(\eta\) and least eigenvalue \(\lambda_0\) from
\cref{obj:def:template}. Also write
\[
D_\gamma=\frac{d\gamma}{\gamma-1},
\qquad
r_{\gamma,n}=h_{\gamma,n}^{\beta},
\qquad
\log_+(u)=\max\{0,\log u\}\quad(u>0).
\]

\begin{enumerate}
\item Membership in \(\mathcal P_\gamma\), as specified by
\cref{obj:def:model}, supplies the treated response bound in
\cref{obj:ass:bounded-potential-outcomes}. In particular, with
\[
\overline y_P(x)=\mathbb E_P[Y\mid A=1,X=x],
\]
the treated conditional law satisfies, for \(P_X\)-almost every \(x\),
\[
|\overline y_P(x)|\le B,\qquad
\mathbb E_P\!\left[
 e^{t(Y-\overline y_P(x))}\mid A=1,X=x
\right]\le e^{B^2t^2/2}
\quad(t\in\mathbb R),
\]
and every displayed exponential moment is integrable. These are the
outcome properties used below.
% lean: CausalSmith.Stat.WeakOverlap.selectedMesh_envelope

\item We first construct the uniform count bound. Define the positive mass
coefficient
\[
c_{\mathrm{mass}}(\gamma)
=
\frac{\gamma-1}{\gamma}\,
C^{-1/(\gamma-1)}
(c_f\eta^d)^{\gamma/(\gamma-1)},
\qquad
\kappa_*=\min_{\gamma\in\Gamma}c_{\mathrm{mass}}(\gamma)>0.
\]
The minimum exists and is positive because the displayed function is
continuous and positive on the compact interval \(\Gamma\).
Every template microcell at mesh \(h\) has volume \(\eta^dh^d\).
Consequently, \cref{obj:ass:density} and the setwise bound in
\cref{obj:lem:ordered-mass} give
\[
q_{Q\ell}\ge c_{\mathrm{mass}}(\gamma)h^{D_\gamma}
\ge\kappa_*h^{D_\gamma}.
\]

Choose a fixed enlargement \(A_{\mathrm{mesh}}\ge1\) such that
\[
4\le
\kappa_*A_{\mathrm{mesh}}^{\,2\beta+D_{\gamma_{\max}}}.
\]
Such a choice exists because \(2\beta+D_{\gamma_{\max}}>0\).
Since
\[
D_\gamma=d+\frac d{\gamma-1},
\]
the effective dimension decreases with \(\gamma\), and hence
\[
4\le
\kappa_*A_{\mathrm{mesh}}^{\,2\beta+D_\gamma}
\qquad(\gamma\in\Gamma).
\]

The candidate collection in \cref{obj:synth_3} has largest dyadic index
\[
j_{\max,n}
=
\left\lfloor\frac{\log_2 n}{2\beta+d}\right\rfloor.
\]
The floor inequality yields
\[
2^{-j_{\max,n}}\le 2n^{-1/(2\beta+d)}.
\]
Moreover,
\[
h_{\gamma_{\max},n}\le h_{\gamma,n}
\le h_{\gamma_{\min},n},
\qquad
\frac1{2\beta+d}-\frac1{2\beta+D_{\gamma_{\max}}}>0.
\]
The positive power in the last display tends to infinity when applied to
\(n\). Thus, for all sufficiently large \(n\),
\[
2^{-j_{\max,n}}
\le A_{\mathrm{mesh}}h_{\gamma_{\max},n}
\le A_{\mathrm{mesh}}h_{\gamma,n},
\qquad
A_{\mathrm{mesh}}h_{\gamma,n}
\le A_{\mathrm{mesh}}h_{\gamma_{\min},n}\le1.
\]
Rounding upward within the finite dyadic collection therefore gives a
candidate \(h_{\mathrm{o}}\) satisfying
\[
A_{\mathrm{mesh}}h_{\gamma,n}
\le h_{\mathrm{o}}
\le2A_{\mathrm{mesh}}h_{\gamma,n}.
\]
Indeed, successive candidate widths differ by a factor of two; the
smallest candidate width at least
\(A_{\mathrm{mesh}}h_{\gamma,n}\) satisfies these inequalities, including
when either endpoint is itself a candidate.

For the simultaneous concentration budget, define
\[
M_{\mathrm{count},n}
=
J\sum_{j=0}^{j_{\max,n}}2^{jd},
\qquad
s_n=\log\!\left(2(M_{\mathrm{count},n}+1)n^2\right).
\]
Here \(2^{jd}\) counts the partition cubes and \(J\) counts the microcells
within each cube. Since \(2\beta+d\ge1\),
\[
2^{j_{\max,n}}\le n,\qquad
j_{\max,n}\le2^{j_{\max,n}}\le n,
\]
and therefore
\[
M_{\mathrm{count},n}
\le(j_{\max,n}+1)(2^{j_{\max,n}})^dJ
\le(n+1)n^dJ.
\]
For \(n\ge2\), this implies
\[
2(M_{\mathrm{count},n}+1)n^2
\le8(J+1)n^{d+3},
\qquad
s_n\le(d+3)\log n+\log(8(J+1)).
\]

To make the count threshold dominate this budget, set
\[
U_{\mathrm{budget}}=24(d+3),\qquad
V_{\mathrm{budget}}=24\log(8(J+1)),\qquad
\delta_{\mathrm{count}}
=\frac{2\beta}{2\beta+D_{\gamma_{\min}}}>0,
\qquad
c_{\mathrm{budget}}=(2A_{\mathrm{mesh}})^{-2\beta}>0.
\]
The elementary bound
\[
\log n\le\frac{2}{\delta_{\mathrm{count}}}
n^{\delta_{\mathrm{count}}/2}
\]
and divergence of the positive powers allow one common starting sample
size for which
\[
U_{\mathrm{budget}}\log n
\le\frac{c_{\mathrm{budget}}}{2}n^{\delta_{\mathrm{count}}},
\qquad
V_{\mathrm{budget}}
\le\frac{c_{\mathrm{budget}}}{2}n^{\delta_{\mathrm{count}}}.
\]
Together with the endpoint and candidate bounds, these give
\[
\begin{aligned}
24s_n
&\le U_{\mathrm{budget}}\log n+V_{\mathrm{budget}}\\
&\le(2A_{\mathrm{mesh}}h_{\gamma_{\min},n})^{-2\beta}\\
&\le(2A_{\mathrm{mesh}}h_{\gamma,n})^{-2\beta}
\le h_{\mathrm{o}}^{-2\beta}.
\end{aligned}
\]
Fix \(n_{\mathrm{count}}\ge2\) large enough that all these bounds hold
uniformly for \(\gamma\in\Gamma\).

Finally, the oracle balance identity is
\[
nh_{\gamma,n}^{2\beta+D_\gamma}=1.
\]
Thus every microcell at \(h_{\mathrm{o}}\) satisfies
\[
\begin{aligned}
nq_{Q\ell}h_{\mathrm{o}}^{2\beta}
&\ge n\kappa_*h_{\mathrm{o}}^{2\beta+D_\gamma}\\
&\ge
\kappa_*A_{\mathrm{mesh}}^{2\beta+D_\gamma}
nh_{\gamma,n}^{2\beta+D_\gamma}
\ge4,
\end{aligned}
\]
or equivalently \(nq_{Q\ell}\ge4h_{\mathrm{o}}^{-2\beta}\).
% lean: CausalSmith.Stat.WeakOverlap.selectedCountEvent_uniform

\item Under \cref{obj:ass:iid}, a treated microcell count is a sum of
independent Bernoulli indicators. Each indicator has mean \(q_{Q\ell}\),
its centered absolute value is at most one, and its variance is
\(q_{Q\ell}(1-q_{Q\ell})\le q_{Q\ell}\). The Bernstein bound for these
bounded centered summands gives, for every \(u>0\),
\[
P^{\otimes n}\{|N_{Q\ell}-nq_{Q\ell}|\ge u\}
\le
2\exp\!\left\{-\frac{u^2}{2(2nq_{Q\ell}+u)}\right\}.
\]
For \(u=nq_{Q\ell}/2+4s_n\), the numerical inequality
\[
2s_n\left(2nq_{Q\ell}+u\right)\le u^2
\]
follows by expanding both sides. Hence
\[
P^{\otimes n}\left\{
|N_{Q\ell}-nq_{Q\ell}|
\ge\frac{nq_{Q\ell}}2+4s_n
\right\}\le2e^{-s_n}.
\]
Define
\[
\mathcal E_n^{\mathrm{count}}
=
\bigcap_{h\in\mathcal H_n}
\bigcap_{Q\in\mathcal Q_h}
\bigcap_{\ell=1}^{J}
\left\{
|N_{Q\ell}-nq_{Q\ell}|
<\frac{nq_{Q\ell}}2+4s_n
\right\}.
\]
This finite intersection is measurable. The union bound gives
\[
P^{\otimes n}\!\left((\mathcal E_n^{\mathrm{count}})^c\right)
\le2M_{\mathrm{count},n}e^{-s_n}
=
\frac{M_{\mathrm{count},n}}{M_{\mathrm{count},n}+1}\,n^{-2}
\le n^{-2}.
\]

On this event, the candidate from the preceding step satisfies
\[
N_{Q\ell}
\ge\frac{nq_{Q\ell}}2-4s_n
\ge2h_{\mathrm{o}}^{-2\beta}-4s_n
\ge h_{\mathrm{o}}^{-2\beta}.
\]
It is therefore feasible. The selector in \cref{obj:def:mesh-selector}
then gives
\[
0<\widehat h\le h_{\mathrm{o}}
\le2A_{\mathrm{mesh}}h_{\gamma,n}.
\]
At the selected mesh, feasibility and the budget imply
\[
N_{Q\ell}\ge\widehat h^{-2\beta}
\ge h_{\mathrm{o}}^{-2\beta}\ge24s_n.
\]
Using the upper deviation inequality first gives
\[
nq_{Q\ell}
\ge\frac23(N_{Q\ell}-4s_n)
\ge\frac59N_{Q\ell}
\ge\frac{40}{3}s_n.
\]
The lower deviation inequality now yields
\[
N_{Q\ell}
\ge\frac{nq_{Q\ell}}2-4s_n
\ge\frac{nq_{Q\ell}}5.
\]
These arguments apply to every partition cube, with the boundary
convention of \cref{obj:synth_3}.
% lean: CausalSmith.Stat.WeakOverlap.selectedCountEvent_of_oracle_witness

\item We next obtain the actual conditional coefficient bound. Define
the sampled design and its conditional sample law by
\[
\mathcal G_n=\sigma((X_i,A_i)_{i=1}^n),
\qquad
\nu_\omega
=
\mathcal L\!\left(\mathcal D_n\mid
(X_i,A_i)_{i=1}^n=(X_i(\omega),A_i(\omega))_{i=1}^n\right).
\]
Under the iid product law, disintegration coordinate by coordinate gives
\[
\mathcal L_{\nu_\omega}((Y_i)_{i=1}^n)
=
\bigotimes_{i=1}^n
\mathcal L_P(Y\mid X=X_i(\omega),A=A_i(\omega))
\]
for almost every \(\omega\). To see the product structure, conditional
probabilities of response rectangles factor into the product of the
coordinate conditional probabilities; these rectangles determine the
conditional response law.

For such a design, define the treated residuals on a conditional sample
\(\xi\) by
\[
\varepsilon_i^\omega(\xi)
=
\begin{cases}
Y_i(\xi)-\overline y_P(X_i(\omega)),&A_i(\omega)=1,\\
0,&A_i(\omega)=0.
\end{cases}
\]
They are independent under \(\nu_\omega\), and the treated conditional
bound from the first step, together with the zero residual in the second
case, gives
\[
\mathbb E_{\nu_\omega}\varepsilon_i^\omega=0,\qquad
\mathbb E_{\nu_\omega}e^{t\varepsilon_i^\omega}
\le e^{B^2t^2/2}\quad(t\in\mathbb R),
\]
with integrable exponential moments.

Fix a feasible candidate \(h\), a cube \(Q\in\mathcal Q_h\), and
\(\alpha\in\mathcal I_m\). Define
\[
x_{iQ}=\frac{X_i(\omega)-a_Q}{h},
\qquad
w_{iQ\alpha}
=
\frac1J\sum_{\ell=1}^{J}
\frac{\mathbf1\{A_i(\omega)=1,\ X_i(\omega)\in E_{Q\ell}\}}
     {N_{Q\ell}(\omega)}
[\widehat G_Q(\omega)^{-1}U(x_{iQ})]_\alpha.
\]
The equal-cell formula in \cref{obj:synth_4} is understood with
\(0^{-1}=0\), and with the zero matrix assigned as the inverse of a
singular Gram matrix. On the feasible meshes considered here,
\cref{obj:lem:template-conditioning} supplies the ordinary inverse and
\[
\|\widehat G_Q^{-1}\|_{\mathrm{op}}\le\frac2{\lambda_0},
\qquad \lambda_0>0.
\]

The conditional sample has the same design as \(\omega\) almost surely.
Expanding the coefficient formula therefore gives
\[
\widehat c_{Q,\alpha}(\xi)
=
\sum_{i=1}^n w_{iQ\alpha}
\mathbf1\{A_i(\omega)=1\}\overline y_P(X_i(\omega))
+
\sum_{i=1}^n w_{iQ\alpha}\varepsilon_i^\omega(\xi).
\]
The residuals are integrable and have mean zero, so integration proves
\[
\mathbb E_{\nu_\omega}\widehat c_{Q,\alpha}
=
\sum_{i=1}^n w_{iQ\alpha}
\mathbf1\{A_i(\omega)=1\}\overline y_P(X_i(\omega)).
\]
In particular, defining
\[
Z_{Q,\alpha}^\omega
=
\widehat c_{Q,\alpha}
-\mathbb E_{\nu_\omega}\widehat c_{Q,\alpha},
\qquad
Z_Q^\omega=(Z_{Q,\alpha}^\omega)_{\alpha\in\mathcal I_m},
\]
we have the conditional identity and moment bound
\[
Z_{Q,\alpha}^\omega
=
\sum_{i=1}^n w_{iQ\alpha}\varepsilon_i^\omega,
\qquad
\mathbb E_{\nu_\omega}e^{tZ_{Q,\alpha}^\omega}
=
\prod_{i=1}^n
\mathbb E_{\nu_\omega}e^{tw_{iQ\alpha}\varepsilon_i^\omega}
\le
\exp\!\left\{
\frac{B^2t^2}{2}\sum_{i=1}^n w_{iQ\alpha}^2
\right\}.
\]
Independence and coordinate exponential integrability also show that
every exponential moment in this display is integrable.

The normalized covariates appearing in a microcell lie in
\([0,1]^d\). Since every coordinate of \(U\) has absolute value at most
one there,
\[
\|U(x_{iQ})\|_2\le\sqrt{p_{\mathrm{coef}}}.
\]
Consequently, if
\[
R_{\mathrm{row}}
=
\frac{2\sqrt{p_{\mathrm{coef}}}}{J\lambda_0},
\]
then
\[
\left|
\frac1J[\widehat G_Q^{-1}U(x_{iQ})]_\alpha
\right|\le R_{\mathrm{row}}
\]
for observations counted in a microcell. Applying
\((\sum_{\ell=1}^{J}z_\ell)^2\le J\sum_{\ell=1}^{J}z_\ell^2\)
to the cell contributions to each weight, and then summing over the
observations, gives
\[
\begin{aligned}
\sum_{i=1}^n w_{iQ\alpha}^2
&\le
J\sum_{\ell=1}^{J}
\sum_{\substack{i:A_i=1\\X_i\in E_{Q\ell}}}
\frac{R_{\mathrm{row}}^2}{N_{Q\ell}^2}\\
&=
JR_{\mathrm{row}}^2\sum_{\ell=1}^{J}N_{Q\ell}^{-1}.
\end{aligned}
\]

For the mass comparison, define, at the fixed mesh,
\[
\mathcal S_Q=\{Q'\in\mathcal Q_h:q_{Q'}\le q_Q\},
\qquad
\rho_Q=|\mathcal S_Q|\ge1,
\]
and select weakest microcells by
\[
\ell_*(Q')\in
\operatorname*{arg\,min}_{1\le\ell\le J}q_{Q'\ell},
\qquad
\mathcal U_Q=\bigcup_{Q'\in\mathcal S_Q}E_{Q',\ell_*(Q')}.
\]
The selected microcells in different cubes are disjoint. Their union
has covariate mass at least \(\rho_Qc_f\eta^dh^d\), whereas its treated
mass is at most \(\rho_Qq_Q\). The setwise inequality of
\cref{obj:lem:ordered-mass} thus gives
\[
\rho_Qq_Q
\ge
\frac{\gamma-1}{\gamma}C^{-1/(\gamma-1)}
(\rho_Qc_f\eta^dh^d)^{\gamma/(\gamma-1)}.
\]
Dividing by \(\rho_Q>0\) and collecting powers yields the explicit
mass-rank bound
\[
q_Q\ge
c_{\mathrm{mass}}(\gamma)h^{D_\gamma}
\rho_Q^{1/(\gamma-1)}.
\]
At a feasible mesh satisfying the count lower bound from the preceding
step, each inverse count consequently satisfies
\[
N_{Q\ell}^{-1}
\le
\min\!\left\{
h^{2\beta},
\frac{5}{
n c_{\mathrm{mass}}(\gamma)h^{D_\gamma}
\rho_Q^{1/(\gamma-1)}}
\right\}.
\]
Define the fixed template factor
\[
A_{\mathrm{coef}}
=
B^2J^2R_{\mathrm{row}}^2
=
\frac{4B^2p_{\mathrm{coef}}}{\lambda_0^2}.
\]
The conditional coefficient bound is therefore
\[
\mathbb E_{\nu_\omega}e^{tZ_{Q,\alpha}^\omega}
\le
\exp\!\left[
\frac{A_{\mathrm{coef}}t^2}{2}
\min\!\left\{
h^{2\beta},
\frac{5}{
n c_{\mathrm{mass}}(\gamma)h^{D_\gamma}
\rho_Q^{1/(\gamma-1)}}
\right\}
\right].
\]

Choose a measurable full-measure set
\(\mathcal E_n^{\mathrm{cond}}\) on which the conditional identities,
exponential integrability, and bounds just established hold
simultaneously for all feasible candidate indices and coefficient
coordinates. Define
\[
\mathcal E_n
=
\mathcal E_n^{\mathrm{count}}\cap\mathcal E_n^{\mathrm{cond}},
\qquad
P^{\otimes n}(\mathcal E_n^{\mathrm{cond}})=1.
\]
Intersecting with this full-measure set preserves the failure bound:
\[
P^{\otimes n}(\mathcal E_n^c)\le n^{-2}.
\]
All subsequent conditional assertions hold for every \(\omega\in\mathcal
E_n\), at the selected mesh of that sample.
% lean: CausalSmith.Stat.WeakOverlap.selectedMesh_count_and_raw_mgf

\item We normalize the two variance scales uniformly. Define
\[
F_{\mathrm{mass}}=\max\{1,5/\kappa_*\}.
\]
For \(x\ge0\), \(y>0\), and \(c\ge\kappa_*\), we have
\[
\frac5{cy}\le F_{\mathrm{mass}}y^{-1},
\qquad
\min\{x,5/(cy)\}
\le F_{\mathrm{mass}}\min\{x,y^{-1}\}.
\]
For the second inequality, if \(x\le y^{-1}\), the left side is at most
\(x\le F_{\mathrm{mass}}x\); if \(x>y^{-1}\), it is at most
\(5/(cy)\le F_{\mathrm{mass}}y^{-1}\).
Applying this with the displayed mass coefficient and collecting
inverse powers gives
\[
\begin{aligned}
&A_{\mathrm{coef}}
\min\!\left\{
\widehat h^{2\beta},
\frac5{
n c_{\mathrm{mass}}(\gamma)\widehat h^{D_\gamma}
\rho_Q^{1/(\gamma-1)}}
\right\}\\
&\qquad\le
A_{\mathrm{coef}}F_{\mathrm{mass}}
\min\!\left\{
\widehat h^{2\beta},
n^{-1}\widehat h^{-D_\gamma}\rho_Q^{-1/(\gamma-1)}
\right\}.
\end{aligned}
\]
Set
\[
K_{\mathrm{count}}
=
\max\{2A_{\mathrm{mesh}},A_{\mathrm{coef}}F_{\mathrm{mass}}\}+1>0.
\]
The selected mesh and conditional moments thus obey
\[
0<\widehat h\le K_{\mathrm{count}}h_{\gamma,n},
\qquad
\mathbb E_{\nu_\omega}e^{tZ_{Q,\alpha}^\omega}
\le
\exp\!\left\{
\frac{K_{\mathrm{count}}t^2}{2}
\min\!\left(
\widehat h^{2\beta},
n^{-1}\widehat h^{-D_\gamma}\rho_Q^{-1/(\gamma-1)}
\right)
\right\}.
\]
If the cubes are enumerated in nondecreasing mass order, the first \(k\)
cubes all belong to \(\mathcal S_{Q_{(k)}}\). Hence
\[
\rho_{Q_{(k)}}\ge k,
\qquad
\rho_{Q_{(k)}}^{-1/(\gamma-1)}
\le k^{-1/(\gamma-1)}.
\]
This converts the displayed bound to the rank-\(k\) envelope, including
ties.
% lean: CausalSmith.Stat.WeakOverlap.selectedMesh_count_and_rank_mgf

\item We establish the scalar maximal estimate used to control the
coefficient vectors. Begin with a clipped Gaussian integral. For an
integer \(R_{\mathrm{cut}}\ge1\) and a scale \(a>0\), define
\[
s_{\mathrm{clip}}=\sqrt{1+\log R_{\mathrm{cut}}},
\qquad
r_{\mathrm{clip}}=a s_{\mathrm{clip}},
\qquad
g_{\mathrm{pre}}(u)
=
\min\!\left\{1,\,
2R_{\mathrm{cut}}e^{-u^2/(2a^2)}\right\}.
\]
We have \(s_{\mathrm{clip}}^2\ge1\) and
\(\log(2R_{\mathrm{cut}})\le s_{\mathrm{clip}}^2\).
For \(u\le4r_{\mathrm{clip}}\),
\[
g_{\mathrm{pre}}(u)\le1\le e^{4-u/r_{\mathrm{clip}}}.
\]
For \(u>4r_{\mathrm{clip}}\), put
\[
x_{\mathrm{clip}}=u/r_{\mathrm{clip}}>4.
\]
The nonnegativity of
\[
(s_{\mathrm{clip}}^2-1)(x_{\mathrm{clip}}^2/2-1)
\]
and the elementary quadratic inequality
\(1-x_{\mathrm{clip}}^2/2\le4-x_{\mathrm{clip}}\) give
\[
s_{\mathrm{clip}}^2
-s_{\mathrm{clip}}^2x_{\mathrm{clip}}^2/2
\le4-x_{\mathrm{clip}}.
\]
Therefore
\[
g_{\mathrm{pre}}(u)
\le
\exp\!\left\{
\log(2R_{\mathrm{cut}})
-s_{\mathrm{clip}}^2x_{\mathrm{clip}}^2/2
\right\}
\le e^{4-u/r_{\mathrm{clip}}}.
\]
Integrating this exponential majorant proves
\[
\int_0^\infty g_{\mathrm{pre}}(u)\,du
\le e^4a\sqrt{1+\log R_{\mathrm{cut}}}.
\]
% lean: Causalean.Mathlib.Probability.SubGaussian.clipped_gaussian_integral

\item Fix any exponent \(\alpha_{\mathrm{ord}}>0\). For an integer
\(N\ge1\), retain \(R_{\mathrm{cut}}\ge1\) and define the clipped suffix
envelope
\[
g_{\mathrm{suf}}(u)
=
\min\!\left\{1,\,
\sum_{\substack{1\le k\le N\\k\ge R_{\mathrm{cut}}}}
2\exp\!\left[
-\frac{u^2}{
2a^2(R_{\mathrm{cut}}/k)^{\alpha_{\mathrm{ord}}}}
\right]\right\}.
\]
Partition its indices into dyadic blocks by
\[
L_j=2^jR_{\mathrm{cut}},\qquad
\mathcal B_j
=
\{k:1\le k\le N,\ L_j\le k<2L_j\},
\qquad j\in\mathbb N_0,
\]
and define
\[
g_j(u)
=
\min\!\left\{1,\,
\sum_{k\in\mathcal B_j}
2\exp\!\left[
-\frac{u^2}{
2a^2(R_{\mathrm{cut}}/k)^{\alpha_{\mathrm{ord}}}}
\right]\right\}.
\]
Each block has at most \(L_j\) indices, and within it
\[
a^2(R_{\mathrm{cut}}/k)^{\alpha_{\mathrm{ord}}}
\le a^2(R_{\mathrm{cut}}/L_j)^{\alpha_{\mathrm{ord}}}.
\]
The preceding clipped Gaussian calculation, at scale
\(a(R_{\mathrm{cut}}/L_j)^{\alpha_{\mathrm{ord}}/2}\), consequently gives
\[
\int_0^\infty g_j(u)\,du
\le
e^4a\,2^{-j\alpha_{\mathrm{ord}}/2}
\sqrt{1+\log(2^jR_{\mathrm{cut}})}.
\]
Since \(\log2\le1\) and \(\log R_{\mathrm{cut}}\ge0\),
\[
\sqrt{1+\log(2^jR_{\mathrm{cut}})}
\le(j+1)\sqrt{1+\log R_{\mathrm{cut}}}.
\]
Define
\[
q_{\mathrm{geom}}
=
e^{-\alpha_{\mathrm{ord}}\log2/2}\in(0,1),
\qquad
D_{\mathrm{geom}}(\alpha_{\mathrm{ord}})
=
1+\sum_{j=0}^{\infty}(j+1)q_{\mathrm{geom}}^j<\infty.
\]
Convergence follows from the geometric series and its linearly weighted
series.

The dyadic blocks partition the suffix. For nonnegative summands,
clipping a sum at one is bounded by the sum of the individually clipped
block sums. Only finitely many blocks meet \(\{1,\ldots,N\}\), so
integration gives
\[
\begin{aligned}
\int_0^\infty g_{\mathrm{suf}}(u)\,du
&\le\sum_{j=0}^{N}\int_0^\infty g_j(u)\,du\\
&\le
e^4a\sqrt{1+\log R_{\mathrm{cut}}}
\sum_{j=0}^{N}(j+1)q_{\mathrm{geom}}^j\\
&\le
e^4D_{\mathrm{geom}}(\alpha_{\mathrm{ord}})
a\sqrt{1+\log R_{\mathrm{cut}}}.
\end{aligned}
\]
The finite Gaussian sums and their clipped versions are integrable, so
all these integrations are legitimate.
% lean: Causalean.Mathlib.Probability.SubGaussian.decaying_gaussian_suffix_integral

\item Now let \(\nu\) be a probability law, let \(a,b>0\), and let
\(Z_1,\ldots,Z_N\) have the exponential integrability and moment envelope
with exponent \(\alpha_{\mathrm{ord}}\). Define
\[
W_N=\max\bigl(\{0\}\cup\{|Z_k|:1\le k\le N\}\bigr).
\]
For \(N=0\), this definition gives \(W_N=0\), and the maximal inequality
follows immediately. Suppose \(N\ge1\), and define
\[
v_k=\min\{a^2,b^2k^{-\alpha_{\mathrm{ord}}}\}>0,
\qquad
\tau_{\mathrm{tail}}(u)=\min\!\left\{1,\sum_{k=1}^{N}2e^{-u^2/(2v_k)}\right\}.
\]
Exponential Markov bounds for \(Z_k\) and \(-Z_k\), optimized at
\(u/v_k\), give
\[
\nu\{|Z_k|>u\}\le2e^{-u^2/(2v_k)}
\quad(u\ge0).
\]
The finite union bound, together with the probability bound by one,
therefore yields
\[
\nu\{W_N>u\}\le\tau_{\mathrm{tail}}(u).
\]
Exponential moments at \(1\) and \(-1\) imply integrability of every
\(|Z_k|\), and hence of \(W_N\). The layer-cake identity gives
\[
\mathbb E_\nu W_N
=
\int_0^\infty\nu\{W_N>u\}\,du
\le\int_0^\infty\tau_{\mathrm{tail}}(u)\,du.
\]

To bound the last integral, set
\[
R_{\mathrm{ratio}}=(b/a)^{2/\alpha_{\mathrm{ord}}}>0,
\qquad
R_{\mathrm{cut}}=\lceil R_{\mathrm{ratio}}\rceil\ge1.
\]
Then
\[
R_{\mathrm{ratio}}\le R_{\mathrm{cut}}
\le2\max\{1,R_{\mathrm{ratio}}\},
\qquad
1+\log R_{\mathrm{cut}}
\le2(1+\log_+(R_{\mathrm{ratio}})).
\]
The latter bound follows by taking logarithms of the preceding
inequality, using
\(\log\max\{1,R_{\mathrm{ratio}}\}
=\log_+(R_{\mathrm{ratio}})\) and \(\log2\le1\).
It implies
\[
\sqrt{1+\log R_{\mathrm{cut}}}
\le2\sqrt{1+\log_+(R_{\mathrm{ratio}})}.
\]
Furthermore,
\[
R_{\mathrm{ratio}}^{\alpha_{\mathrm{ord}}}=(b/a)^2,
\qquad
b^2\le a^2R_{\mathrm{cut}}^{\alpha_{\mathrm{ord}}},
\qquad
v_k\le a^2(R_{\mathrm{cut}}/k)^{\alpha_{\mathrm{ord}}}.
\]
Split the sum defining \(\tau_{\mathrm{tail}}\) at \(k<R_{\mathrm{cut}}\) and
\(k\ge R_{\mathrm{cut}}\). The prefix has at most \(R_{\mathrm{cut}}\)
terms, each with variance at most \(a^2\). The suffix has the preceding
decaying variance bound. Since
\[
\min\{1,x+y\}\le\min\{1,x\}+\min\{1,y\}
\qquad(x,y\ge0),
\]
we obtain
\[
\tau_{\mathrm{tail}}(u)\le g_{\mathrm{pre}}(u)+g_{\mathrm{suf}}(u).
\]
The two integral estimates prove
\[
\mathbb E_\nu W_N
\le
K_{\mathrm{sc}}(\alpha_{\mathrm{ord}})
a\sqrt{1+\log_+((b/a)^{2/\alpha_{\mathrm{ord}}})},
\]
where the positive constant may be taken as
\[
K_{\mathrm{sc}}(\alpha_{\mathrm{ord}})
=
2e^4\bigl(1+D_{\mathrm{geom}}(\alpha_{\mathrm{ord}})\bigr).
\]
% lean: Causalean.Mathlib.Probability.SubGaussian.orderedSubGaussianMaximal

\item To obtain a constant uniform over the adaptation interval, define
\[
\alpha_{\min}=\frac1{\gamma_{\max}-1}>0,\qquad
\alpha_{\max}=\frac1{\gamma_{\min}-1},
\qquad
R_{\mathrm{exp}}=\frac{\alpha_{\max}}{\alpha_{\min}}\ge1.
\]
Taking reciprocals of positive denominators shows
\[
\alpha_{\min}\le\frac1{\gamma-1}\le\alpha_{\max}
\qquad(\gamma\in\Gamma).
\]

For any \(\alpha_{\mathrm{ord}}\in[\alpha_{\min},\alpha_{\max}]\) and
\(k\ge1\),
\[
k^{-\alpha_{\mathrm{ord}}}\le k^{-\alpha_{\min}}.
\]
Thus the scalar estimate at \(\alpha_{\min}\) applies to the envelope
with exponent \(\alpha_{\mathrm{ord}}\). Its logarithmic factor satisfies
\[
1+\log_+((b/a)^{2/\alpha_{\min}})
\le
R_{\mathrm{exp}}
\bigl(1+\log_+((b/a)^{2/\alpha_{\mathrm{ord}}})\bigr).
\]
Indeed, if \(\log(b/a)\ge0\), use
\[
\frac2{\alpha_{\min}}
\le R_{\mathrm{exp}}\frac2{\alpha_{\mathrm{ord}}};
\]
if \(\log(b/a)<0\), both positive-part logarithms vanish and the
inequality is \(1\le R_{\mathrm{exp}}\).
Taking square roots establishes a uniform scalar constant
\[
M_\Gamma
=
K_{\mathrm{sc}}(\alpha_{\min})\sqrt{R_{\mathrm{exp}}}>0.
\]

For a finite nonempty family of vectors with coordinates in
\(\mathcal I_m\), the pointwise inequality is
\[
\max_Q\|Z_Q\|_\infty
\le
\sum_{\alpha\in\mathcal I_m}\max_Q|Z_{Q,\alpha}|.
\]
Coordinate exponential integrability implies integrability of these
finite maxima and of the vector maximum. For example,
\[
|Z_{Q,\alpha}|
\le e^{Z_{Q,\alpha}}+e^{-Z_{Q,\alpha}},
\]
and each finite maximum is bounded by the sum of the corresponding
absolute values. Applying the uniform scalar estimate coordinatewise
and integrating the pointwise inequality gives
\[
\mathbb E_\nu\max_Q\|Z_Q\|_\infty
\le
p_{\mathrm{coef}}M_\Gamma a
\sqrt{1+\log_+((b/a)^{2/\alpha_{\mathrm{ord}}})}.
\]
% lean: CausalSmith.Stat.WeakOverlap.orderedSubGaussianMaximal_vector_uniform_of_exp

\item The preceding vector estimate applies to mass-ranked cube
envelopes. Enumerate the cubes in nondecreasing \(q_Q\), breaking ties
by any fixed ordering of the finite cube family. This enumeration is a
bijection, so
\[
\max_Q\|Z_Q\|_\infty
=
\max_{1\le k\le|\mathcal Q_h|}\|Z_{Q_{(k)}}\|_\infty.
\]
As established above, \(\rho_{Q_{(k)}}\ge k\). For any nonnegative
variance factor \(y\),
\[
\min\{x,y\rho_{Q_{(k)}}^{-\alpha_{\mathrm{ord}}}\}
\le
\min\{x,yk^{-\alpha_{\mathrm{ord}}}\}.
\]
Therefore an envelope with variance proxy
\(K_{\mathrm{count}}\min\{x,y\rho_Q^{-\alpha_{\mathrm{ord}}}\}\)
has the ordered envelope required by the vector estimate whenever
\[
a^2=K_{\mathrm{count}}x,\qquad
b^2=K_{\mathrm{count}}y.
\]
The cube family is nonempty at every dyadic mesh, so the vector estimate
applies throughout the selected-mesh argument.
% lean: CausalSmith.Stat.WeakOverlap.selectedMesh_cube_rank_vector_maximal

\item We now fix the final constants. Define
\[
E_\Gamma
=
\max\{1,\,2\beta(\gamma_{\max}-1)+d\gamma_{\max}\},
\]
\[
K=
\max\!\left\{
K_{\mathrm{count}},
p_{\mathrm{coef}}M_\Gamma
\sqrt{K_{\mathrm{count}}}\sqrt{E_\Gamma}
\right\}+1,
\qquad
K'=K\max\{1,K^\beta\},
\qquad
n_0=n_{\mathrm{count}}.
\]
These constants are positive and depend only on the fixed parameters
and interval endpoints.

Fix \(n\ge n_0\), \(\gamma\in\Gamma\), \(P\in\mathcal P_\gamma\), and
\(\omega\in\mathcal E_n\). Conditional on its sampled design, keep
\(\widehat h=\widehat h(\omega)\) fixed and define the scales
\[
a_{\mathrm{sel}}
=
\sqrt{K_{\mathrm{count}}}\widehat h^\beta,
\qquad
b_{\mathrm{sel}}
=
\sqrt{K_{\mathrm{count}}}\,
n^{-1/2}\widehat h^{-D_\gamma/2},
\qquad
\alpha_{\mathrm{sel}}=\frac1{\gamma-1}.
\]
All three are positive, and
\[
a_{\mathrm{sel}}^2
=K_{\mathrm{count}}\widehat h^{2\beta},
\qquad
b_{\mathrm{sel}}^2
=K_{\mathrm{count}}n^{-1}\widehat h^{-D_\gamma},
\qquad
\alpha_{\mathrm{sel}}\in[\alpha_{\min},\alpha_{\max}].
\]
Applying the ranked vector estimate under \(\nu_\omega\) gives
\[
\mathbb E_{\nu_\omega}
\max_{Q\in\mathcal Q_{\widehat h}}\|Z_Q^\omega\|_\infty
\le
p_{\mathrm{coef}}M_\Gamma a_{\mathrm{sel}}
\sqrt{
1+\log_+\bigl((b_{\mathrm{sel}}/a_{\mathrm{sel}})
^{2/\alpha_{\mathrm{sel}}}\bigr)
}.
\]
The same exponential integrability supplies integrability of the
maximum on the left.
% lean: CausalSmith.Stat.WeakOverlap.selectedMesh_envelope

\item The oracle definition and \(2\beta+D_\gamma>0\) imply
\[
n^{-1/2}
=h_{\gamma,n}^{\,\beta+D_\gamma/2}.
\]
After cancelling \(\sqrt{K_{\mathrm{count}}}\), this gives
\[
\frac{b_{\mathrm{sel}}}{a_{\mathrm{sel}}}
=
\left(\frac{h_{\gamma,n}}{\widehat h}\right)^{
\beta+D_\gamma/2}.
\]
Define the resulting logarithmic exponent by
\[
e_{\mathrm{sel}}
=
\left(\beta+\frac{D_\gamma}{2}\right)
\frac2{\alpha_{\mathrm{sel}}}
=
2\beta(\gamma-1)+d\gamma.
\]
Its reciprocal dependence cancels, and
\[
0\le e_{\mathrm{sel}}
\le2\beta(\gamma_{\max}-1)+d\gamma_{\max}
\le E_\Gamma.
\]

For a positive ratio \(u\) and an exponent \(\zeta_{\mathrm{exp}}\ge0\),
\[
1+\log_+(u^{\zeta_{\mathrm{exp}}})
\le\max\{1,\zeta_{\mathrm{exp}}\}(1+\log_+(u)).
\]
If \(\log u\ge0\), this follows from
\(1+\zeta_{\mathrm{exp}}\log u\le\max\{1,\zeta_{\mathrm{exp}}\}(1+\log u)\).
If \(\log u<0\), both positive parts vanish and the assertion is
\(1\le\max\{1,\zeta_{\mathrm{exp}}\}\). Applying these two cases to the oracle ratio proves
\[
\sqrt{
1+\log_+\bigl((b_{\mathrm{sel}}/a_{\mathrm{sel}})
^{2/\alpha_{\mathrm{sel}}}\bigr)
}
\le
\sqrt{E_\Gamma}\,
\sqrt{1+\log_+(h_{\gamma,n}/\widehat h)}.
\]
% lean: CausalSmith.Stat.WeakOverlap.selectedMesh_maximal_oracle_log_ratio

\item Combining the preceding estimates and the definition of \(K\)
yields
\[
\begin{aligned}
\mathbb E_{\nu_\omega}
\max_{Q\in\mathcal Q_{\widehat h}}\|Z_Q^\omega\|_\infty
&\le
p_{\mathrm{coef}}M_\Gamma
\sqrt{K_{\mathrm{count}}}\sqrt{E_\Gamma}\,
\widehat h^\beta
\sqrt{1+\log_+(h_{\gamma,n}/\widehat h)}\\
&\le
K\widehat h^\beta
\sqrt{1+\log_+(h_{\gamma,n}/\widehat h)}.
\end{aligned}
\]
Since \(K\ge K_{\mathrm{count}}\), the mesh bound also becomes
\[
0<\widehat h\le Kh_{\gamma,n}.
\]
The variance proxies are nonnegative, so increasing
\(K_{\mathrm{count}}\) to \(K\) in the conditional exponential bound
preserves that inequality. Together with the rank comparison, this
gives precisely
\[
\mathbb E_{\nu_\omega}e^{tZ_{Q_{(k)},\alpha}^\omega}
\le
\exp\!\left\{
\frac{Kt^2}{2}
\min\!\left(
\widehat h^{2\beta},
n^{-1}\widehat h^{-D_\gamma}k^{-1/(\gamma-1)}
\right)
\right\}.
\]
The count bound and all exponential integrability assertions are
retained on the same event \(\mathcal E_n\).
% lean: CausalSmith.Stat.WeakOverlap.selectedMesh_envelope

\item For the deterministic comparison with the oracle rate, define
\[
s_{\mathrm{mesh}}=\frac{\widehat h}{h_{\gamma,n}},
\qquad
0<s_{\mathrm{mesh}}\le K.
\]
If \(s_{\mathrm{mesh}}\le1\), then
\[
0\le\log(1/s_{\mathrm{mesh}})
\le1/s_{\mathrm{mesh}}-1,
\]
and consequently
\[
\sqrt{1+\log_+(1/s_{\mathrm{mesh}})}
\le\sqrt{1/s_{\mathrm{mesh}}}
\le1/s_{\mathrm{mesh}}.
\]
Since \(\beta>1\),
\[
s_{\mathrm{mesh}}^\beta\le s_{\mathrm{mesh}},
\qquad
s_{\mathrm{mesh}}^\beta
\sqrt{1+\log_+(1/s_{\mathrm{mesh}})}
\le1.
\]
If \(s_{\mathrm{mesh}}>1\), the positive-part logarithm vanishes, and
\[
s_{\mathrm{mesh}}^\beta
\sqrt{1+\log_+(1/s_{\mathrm{mesh}})}
=s_{\mathrm{mesh}}^\beta\le K^\beta.
\]
The two cases give
\[
\begin{aligned}
\widehat h^\beta
\sqrt{1+\log_+(h_{\gamma,n}/\widehat h)}
&=
h_{\gamma,n}^\beta s_{\mathrm{mesh}}^\beta
\sqrt{1+\log_+(1/s_{\mathrm{mesh}})}\\
&\le\max\{1,K^\beta\}r_{\gamma,n}.
\end{aligned}
\]
Multiplying by \(K\) proves the asserted bound with
\(K'=K\max\{1,K^\beta\}\).
% lean: CausalSmith.Stat.WeakOverlap.selectedMesh_oracle_rate_log_bound

\item Finally, for the additional maximal assertion, take
\(\nu=\mu\) and \(\alpha_{\mathrm{ord}}=\alpha>0\) in the scalar
calculation above, and set
\[
K_\alpha=K_{\mathrm{sc}}(\alpha)>0.
\]
Its hypotheses are exactly the stated exponential integrability and
moment envelope. Thus, including the case \(N=0\),
\[
\mathbb E_\mu
\max\bigl(\{0\}\cup\{|Z_k|:1\le k\le N\}\bigr)
\le
K_\alpha a
\sqrt{1+\max\{0,\log((b/a)^{2/\alpha})\}}.
\]
This completes all assertions.
% lean: Causalean.Mathlib.Probability.SubGaussian.orderedSubGaussianMaximal
\end{enumerate}
\end{proof}

The two parts of the coefficient envelope have distinct statistical roles. Count feasibility caps every coefficient's stochastic scale at a constant times \(\widehat h^\beta\). The population-mass comparison and \cref{obj:lem:ordered-mass} additionally supply the rank-dependent scale
\(n^{-1/2}\widehat h^{-D_\gamma/2}k^{-1/(2(\gamma-1))}\).
The ordered maximal inequality combines these two controls. Its logarithm measures the number of ranks at which the common cap is relevant, through the scale ratio \((b/a)^{2/\alpha}\).

For the coefficient application, the decay exponent is \(1/(\gamma-1)\), and, apart from fixed multiplicative constants, the two scales are
\[
a=\widehat h^\beta,
\qquad
b=n^{-1/2}\widehat h^{-D_\gamma/2}.
\]
Their ratio is governed by the selected mesh relative to the oracle mesh. Thus the conditional maximum is controlled at the statistical scale
\(\widehat h^\beta\sqrt{1+\max\{0,\log(h_{\gamma,n}/\widehat h)\}}\).
The final inequality in \cref{obj:lem:selected-mesh} places this expression at the oracle scale even for finer selected meshes. In this application, the decay exponent ranges over
\[
\left[
\frac{1}{\gamma_{\max}-1},
\frac{1}{\gamma_{\min}-1}
\right],
\]
and the lemma supplies common constants \(K,K'\) and a common threshold \(n_0\) over the stated overlap interval. The fixed number of polynomial coordinates is included in those constants.

\subsection{The three components of the upper bound}

The approximation component uses the intrinsic Hölder regularity in \cref{obj:def:holder,obj:ass:holder-response} at degree \(m=\lceil\beta\rceil-1\). The deterministic inverse-Gram bound in \cref{obj:lem:template-conditioning} keeps the fitted conditional mean stable relative to its local polynomial approximation. On \(\mathcal E_n\), the selected-mesh comparison in \cref{obj:lem:selected-mesh} places the resulting mesh-dependent approximation error at the oracle response-recovery scale.

The conditional stochastic component is controlled by the coefficient maximum in \cref{obj:lem:selected-mesh}. On normalized partition cubes, polynomial evaluation transfers this coefficient control to spatial-supremum control, with constants determined by the fixed dimension and degree. Conditioning on the sampled design accommodates the count-based mesh selection while retaining the ordered information profile.

Finally, the clipping rule in \cref{obj:def:estimator} and the response bound in \cref{obj:def:model} bound the loss on exceptional samples by \(2B\). Consequently, the exceptional-sample contribution is at most \(2Bn^{-2}\). These three controls furnish the uniform spatial-supremum upper guarantee in \cref{obj:thm:adaptive-minimax}; the pointwise upper guarantee follows from the same spatial-supremum control.


\section{Bounded testing experiment and geometry proofs}

\subsection{The bounded testing construction}

The minimax converse in \cref{obj:thm:adaptive-minimax} is supported by a pair of causal laws with two-valued potential outcomes. The construction places a smooth response perturbation around the evaluation point and uses a common radial propensity to control the treatment information available there. Its radial distribution accounts for intersection with the covariate cube, so the experiment accommodates evaluation points on faces and corners as well as in the interior. The following definition specifies both the observed laws and their causal completions.

\begin{definitionv}{def:bounded-lower-pair}[Bounded radial testing pair]
The bounded testing construction is the measure-valued family \(P_{j,h}\) on observations \((X,A,Y)\), indexed by \(j\in\{0,1\}\) and \(n\in\mathbb N\), with parameters \(d\in\mathbb N\), \(\beta,B,L,\gamma\in\mathbb R\), \(x_0\in\mathbb R^d\), and constants \(a,c>0\). Set
\[
\mathcal X:=[0,1]^d,\qquad
h:=c h_{\gamma,n},\qquad
\theta_0:=\min\{1/4,L/(4B)\},
\]
where \(h_{\gamma,n}\) is the oracle mesh for the smoothness and overlap parameters. Define the radial distance, cube-truncated radial function, and common propensity by
\[
R(x):=\|x-x_0\|_\infty,\qquad
F_R(t):=\prod_{i=1}^d
\max\{0,\min\{1,x_{0i}+t\}-\max\{0,x_{0i}-t\}\},
\qquad
e_0(x):=F_R(R(x))^{1/(\gamma-1)}.
\]
For \(x_0\in\mathcal X\) and \(t\ge0\), \(F_R(t)\) is the probability of \(R(X)\le t\) under uniform covariates. Define the smooth product bump and treated success parameters by
\[
\psi(z):=\prod_{i=1}^d
\begin{cases}
\exp\{1-(1-z_i^2)^{-1}\},& |z_i|<1,\\
0,& |z_i|\ge1,
\end{cases}
\qquad
p_{0,h}(x):=\theta_0,\qquad
p_{1,h}(x):=\theta_0+ah^\beta\psi((x-x_0)/h).
\]

The totalized construction uses the positive-part weights
\[
w_0(p):=\max\{1-p,0\},\qquad
w_1(p):=\max\{p,0\}.
\]
For each covariate \(x\), take the product of the three two-point measures with weights
\[
\begin{array}{c|cc}
 & \text{value }0 & \text{value }1\\ \hline
A & w_0(e_0(x)) & w_1(e_0(x))\\
Y(0)/B & w_0(\theta_0) & w_1(\theta_0)\\
Y(1)/B & w_0(p_{j,h}(x)) & w_1(p_{j,h}(x))
\end{array}
\]
where the potential-outcome measures place their masses at \(0\) and \(B\), respectively. Integrate this product measure against Lebesgue measure restricted to \(\mathcal X\), and let \(P_{j,h}\) be its image under
\[
(X,A,Y(0),Y(1))
\longmapsto
\bigl(X,A,AY(1)+(1-A)Y(0)\bigr).
\]

The admissible domain consists exactly of the constants \(a,c>0\) for which
\[
e_0(x),p_{0,h}(x),p_{1,h}(x)\in[0,1]
\]
for every \(n\ge1\) and every \(x\in\mathcal X\), with \(h=c h_{\gamma,n}\). On this domain, \(P_{j,h}\) is the observed law obtained from uniform \(X\) and conditionally mutually independent \(A,Y(0),Y(1)\) given \(X\), with
\[
A\mid X\sim\operatorname{Bernoulli}(e_0(X)),\qquad
Y(0)\mid X\sim B\operatorname{Bernoulli}(\theta_0),\qquad
Y(1)\mid X\sim B\operatorname{Bernoulli}(p_{j,h}(X)),
\qquad
Y=AY(1)+(1-A)Y(0).
\]
For fixed admissible \(a,c\), define the sample-size-indexed collection
\[
\mathcal P_{\mathrm{pair}}
:=\{(P_{0,h},P_{1,h}):n\ge1,\ h=c h_{\gamma,n}\}.
\]
\end{definitionv}

The positive-part weights in \cref{obj:def:bounded-lower-pair} specify a common measure-family construction across parameter values. On the admissible domain, each conditional two-point measure has total mass one and the construction yields the displayed Bernoulli probability laws. The sample-size-indexed collection \leanref{sym:Ppair}{\ensuremath{\mathcal P_{\mathrm{pair}}}} uses fixed amplitude and mesh constants; sample size changes the perturbation through the oracle mesh.

Admissibility supplies the probability interpretation, while membership in the model of \cref{obj:def:model} additionally requires the response and overlap restrictions. The next result makes these requirements compatible with small statistical information and separation of the chosen response representatives at the evaluation point.

\begin{lemmav}{lem:bounded-lower-pair}[Bounded testing pair guarantees]
Suppose that:
\begin{itemize}
\item \textbf{(Dimension and exponents.)} \(d\in\mathbb N\), \(d>0\), \(\beta>1\), and \(\gamma>1\).
\item \textbf{(Response constants.)} \(B>0\) and \(L>0\).
\item \textbf{(Model constants.)} \(C\ge1\) and \(0<c_f\le1\).
\item \textbf{(Evaluation point.)} \(x_0\in[0,1]^d\).
\end{itemize}
There exist constants \(a,c,c_0>0\), fixed across sample sizes, such that, for every integer \(n\ge1\), with \(h=ch_{\gamma,n}\), where \(h_{\gamma,n}\) is the oracle mesh, the construction in \cref{obj:def:bounded-lower-pair} satisfies
\[
e_0(x),\ p_{0,h}(x),\ p_{1,h}(x)\in[0,1]
\qquad\text{for every }x\in[0,1]^d.
\]
Both observed laws \(P_{0,h}\) and \(P_{1,h}\) belong to \(\mathcal P_\gamma\), the model class with parameters \(d,\beta,B,L,C,c_f,\gamma\). Under each of their constructed causal completions,
\[
Y(0),Y(1)\in\{0,B\}
\qquad\text{almost surely}.
\]
Moreover, \(P_{0,h}\ll P_{1,h}\), the log likelihood ratio
\(\log(dP_{0,h}/dP_{1,h})\) is integrable under \(P_{0,h}\), and
\[
n\,\operatorname{KL}(P_{0,h},P_{1,h})\le\frac18.
\]
For the chosen Hölder response representatives \(\mu_{1,P_{0,h}}\) and \(\mu_{1,P_{1,h}}\) associated with their model-class memberships,
\[
\left|\mu_{1,P_{1,h}}(x_0)-\mu_{1,P_{0,h}}(x_0)\right|
\ge c_0 r_{\gamma,n},
\qquad
r_{\gamma,n}=h_{\gamma,n}^{\beta},
\]
where \(r_{\gamma,n}\) is the oracle response-recovery risk scale.
\end{lemmav}

\begin{proof}[Proof of \cref{obj:lem:bounded-lower-pair}]
\begin{enumerate}
\item We first obtain a Hölder radius for the scaled bump that is uniform in its scale and center. Set
\[
m:=\max\{\lceil\beta\rceil-1,0\},
\qquad
\sigma_\beta:=\beta-m\in(0,1].
\]
For each derivative order \(\kappa\in\mathbb N\), write
\[
\mathscr J_\kappa(z):=\mathrm d^\kappa\psi(z),
\qquad z\in\mathbb R^d,
\qquad
\mathscr J_0(z)=\psi(z),
\]
where \(\mathrm d^\kappa\) denotes the \(\kappa\)-th Fréchet derivative, with the operator norm induced by the supremum norm.

The coordinate factors defining \(\psi\) are smooth across their endpoints: their derivatives from inside \((-1,1)\) vanish at the endpoints because the exponential factor decays faster than every inverse power of the distance to the endpoint. Consequently,
\[
\psi\in C^\infty(\mathbb R^d),
\qquad
\operatorname{supp}\psi\subseteq[-1,1]^d.
\]
Every \(\mathscr J_\kappa\) is continuous and compactly supported. Thus the constants
\[
M_{\psi,\kappa}:=
\sup_{z\in\mathbb R^d}\|\mathscr J_\kappa(z)\|,
\qquad \kappa\in\mathbb N,
\]
are finite and nonnegative. Define
\[
M_{\psi,*}:=\max\{M_{\psi,m},M_{\psi,m+1}\},
\qquad
M_{\psi,\mathrm{mod}}:=2M_{\psi,*}.
\]
The derivative bound of order \(m+1\), applied along a line segment, gives
\[
\|\mathscr J_m(z)-\mathscr J_m(z')\|
\le M_{\psi,*}\|z-z'\|_\infty.
\]
If \(\|z-z'\|_\infty\le1\), then
\[
\|\mathscr J_m(z)-\mathscr J_m(z')\|
\le M_{\psi,*}\|z-z'\|_\infty^{\sigma_\beta}
\le M_{\psi,\mathrm{mod}}\|z-z'\|_\infty^{\sigma_\beta}.
\]
If \(\|z-z'\|_\infty>1\), then
\[
\|\mathscr J_m(z)-\mathscr J_m(z')\|
\le 2M_{\psi,*}
\le M_{\psi,\mathrm{mod}}\|z-z'\|_\infty^{\sigma_\beta}.
\]
These two cases include \(\sigma_\beta=1\).

For \(0<h\le1\), define the scaled bump by
\[
\psi_h(x):=h^\beta\psi((x-x_0)/h),
\qquad x\in\mathbb R^d.
\]
Translation and dilation give, for \(0\le\kappa\le m\),
\[
\mathrm d^\kappa\psi_h(x)
=h^\beta h^{-\kappa}\mathscr J_\kappa((x-x_0)/h),
\]
and hence
\[
\|\mathrm d^\kappa\psi_h(x)\|
\le h^{\beta-\kappa}M_{\psi,\kappa}
\le M_{\psi,\kappa}.
\]
For the top derivative, the preceding modulus bound yields
\[
\begin{aligned}
\|\mathrm d^m\psi_h(x)-\mathrm d^m\psi_h(x')\|
&\le h^{\beta-m}M_{\psi,\mathrm{mod}}
       \left(\frac{\|x-x'\|_\infty}{h}\right)^{\sigma_\beta}\\
&=M_{\psi,\mathrm{mod}}\|x-x'\|_\infty^{\sigma_\beta},
\end{aligned}
\]
since \(\beta-m=\sigma_\beta\). Set
\[
A_\psi:=
\max\left\{
1,\,
\max\left\{
\sum_{\kappa=0}^{m}M_{\psi,\kappa},
M_{\psi,\mathrm{mod}}
\right\}
\right\}>0.
\]
Coordinate derivatives are evaluations of these multilinear derivatives on coordinate unit vectors, so their magnitudes and differences are bounded by the corresponding operator norms. Restricting the ambient smooth functions to the cube supplies their continuous boundary traces. Therefore, with the convention of \cref{obj:def:holder},
\[
\psi_h|_{\mathcal X}\in\mathcal H^\beta(A_\psi)
\qquad
\text{for every }0<h\le1,
\quad \mathcal X=[0,1]^d.
\]
% lean: Causalean.Mathlib.Analysis.Calculus.CubeExtension.exists_uniform_scaledProductBump_holderBallOn

\item Using the baseline probability from \cref{obj:def:bounded-lower-pair}, choose
\[
\theta_0:=\min\left\{\frac14,\frac{L}{4B}\right\},
\qquad
a:=\min\left\{\frac{\theta_0}{2},\frac{L}{4BA_\psi}\right\}.
\]
Since \(B,L,A_\psi>0\),
\[
0<\theta_0\le\frac14,
\qquad
0<a\le\theta_0.
\]
Define the effective dimension and entropy coefficient by
\[
D_\gamma:=\frac{d\gamma}{\gamma-1},
\qquad
\tau_{\mathrm{KL}}:=2\beta+D_\gamma,
\qquad
K_{\mathrm{KL}}:=\frac{2}{\theta_0}\,2^{D_\gamma}.
\]
Here \(D_\gamma>0\), \(K_{\mathrm{KL}}>0\), and
\(\tau_{\mathrm{KL}}\ge1\). Choose
\[
c:=\frac{1}{8(K_{\mathrm{KL}}a^2+1)},
\qquad
c_0:=Bac^\beta.
\]
Then \(0<c\le1\) and \(c_0>0\). Because \(c\le1\) and
\(\tau_{\mathrm{KL}}\ge1\),
\[
c^{\tau_{\mathrm{KL}}}\le c,
\]
so
\[
K_{\mathrm{KL}}a^2c^{\tau_{\mathrm{KL}}}
\le K_{\mathrm{KL}}a^2c
=\frac{K_{\mathrm{KL}}a^2}{8(K_{\mathrm{KL}}a^2+1)}
\le\frac18.
\]
All these constants are fixed independently of \(n\).
% lean: CausalSmith.Stat.WeakOverlap.lowerPair_choose_kl_scale

\item Fix an integer \(n\ge1\). The oracle mesh and the selected bump scale are
\[
h_{\gamma,n}:=n^{-1/\tau_{\mathrm{KL}}},
\qquad
h:=ch_{\gamma,n}.
\]
Since \(n\ge1\), \(\tau_{\mathrm{KL}}>0\), and \(0<c\le1\),
\[
0<h_{\gamma,n}\le1,
\qquad
0<h\le1.
\]
Use \(R,F_R,e_0,\psi,p_{j,h}\) and \(P_{j,h}\) as defined in
\cref{obj:def:bounded-lower-pair}. Every factor defining \(F_R\) lies in
\([0,1]\), and the exponent \(1/(\gamma-1)\) is positive. Thus
\[
0\le F_R(R(x))\le1,
\qquad
0\le e_0(x)\le1,
\qquad x\in\mathbb R^d.
\]
For each coordinate factor of \(\psi\), the exponent is nonpositive when
the coordinate has absolute value less than one, and the factor is zero
otherwise. Hence
\[
0\le\psi(z)\le1,
\qquad z\in\mathbb R^d.
\]
Together with \(h^\beta\le1\), this gives
\[
0\le ah^\beta\psi((x-x_0)/h)\le a
\]
and therefore
\[
0<\theta_0\le p_{j,h}(x)
\le\theta_0+a
\le2\theta_0
\le\frac12,
\qquad j\in\{0,1\},\quad x\in\mathbb R^d.
\]
In particular, \(a,c\) lie in the admissible domain of
\cref{obj:def:bounded-lower-pair}, simultaneously for every \(n\ge1\).

Write the constructed completion laws as
\[
\widetilde P_{j,h}:=
\mathcal L(X,A,Y(0),Y(1)),
\qquad j\in\{0,1\},
\]
where the joint law is the product-kernel construction in
\cref{obj:def:bounded-lower-pair}. The cube has Lebesgue volume
\[
\operatorname{vol}(\mathcal X)=\prod_{i=1}^{d}(1-0)=1.
\]
The admissible weights sum to one in each two-point factor. The kernels
are measurable, so integrating them over the cube gives probability
laws \(\widetilde P_{j,h}\), whose observed images are \(P_{j,h}\).
% lean: CausalSmith.Stat.WeakOverlap.boundedLowerPairAdmissible_of_small

\item Define the displayed treated-response curves by
\[
\mu_{j,h}(x):=Bp_{j,h}(x),
\qquad x\in\mathbb R^d,\quad j\in\{0,1\}.
\]
The choices of \(\theta_0\) and \(a\) give
\[
|B\theta_0|=B\theta_0\le\frac L4,
\qquad
|Ba|A_\psi=BaA_\psi\le\frac L4,
\]
and consequently
\[
|B\theta_0|+|Ba|A_\psi\le L.
\]
For the baseline curve, all positive-order derivatives vanish and
\[
\|\mu_{0,h}\|_\infty=B\theta_0\le L.
\]
For the perturbed curve,
\[
\mu_{1,h}(x)=B\theta_0+Ba\psi_h(x).
\]
For every multi-index \(\alpha\) with \(|\alpha|\le m\),
\[
D^\alpha\mu_{1,h}(x)
=
\begin{cases}
B\theta_0+Ba\psi_h(x),&|\alpha|=0,\\
BaD^\alpha\psi_h(x),&|\alpha|>0.
\end{cases}
\]
The bounds from the first step imply
\[
\|D^\alpha\mu_{1,h}\|_\infty
\le |B\theta_0|+|Ba|A_\psi
\le L.
\]
For \(|\alpha|=m\), the constant term cancels in the difference, giving
\[
|D^\alpha\mu_{1,h}(x)-D^\alpha\mu_{1,h}(x')|
\le |Ba|A_\psi\|x-x'\|_\infty^{\sigma_\beta}
\le L\|x-x'\|_\infty^{\sigma_\beta}.
\]
Their ambient smoothness supplies the required boundary traces.
It follows from \cref{obj:def:holder} that
\[
\mu_{0,h},\mu_{1,h}\in\mathcal H^\beta(L).
\]

The treated potential-outcome marginal of the conditional product kernel
is \(B\operatorname{Bernoulli}(p_{j,h}(x))\). Thus, under the constructed
completion,
\[
\mathbb E_{\widetilde P_{j,h}}[Y(1)\mid X=x]
=(1-p_{j,h}(x))\,0+p_{j,h}(x)B
=\mu_{j,h}(x)
\]
for Lebesgue-almost every \(x\in\mathcal X\). This establishes the
response smoothness requirement in \cref{obj:ass:holder-response}.
% lean: CausalSmith.Stat.WeakOverlap.pairSuccess_true_holderBall
% lean: CausalSmith.Stat.WeakOverlap.completionBind_holderResponse

\item We next establish the propensity tail bound. For every \(r\ge0\),
the intersection of the radial sublevel set with the cube is
\[
\{x:R(x)\le r\}\cap\mathcal X
=
\prod_{i=1}^{d}
\left[
\max\{0,x_{0i}-r\},
\min\{1,x_{0i}+r\}
\right].
\]
Since \(x_0\in\mathcal X\), these intervals are nonempty. Their product
volume gives
\[
\mathbb P\{R(X)\le r\}=F_R(r)
\]
under uniform \(X\). The displayed formula for \(F_R\) is continuous,
and the nested sublevel sets show that it is nondecreasing on
\([0,\infty)\). Also,
\[
F_R(0)=0,\qquad F_R(1)=1.
\]
The first identity uses \(d>0\); the second follows because every
coordinate interval at radius one is \([0,1]\).

For \(u_{\mathrm{tail}}\in[0,1)\), define
\[
\mathscr R(u_{\mathrm{tail}})
:=
\{r\in[0,\infty):F_R(r)\le u_{\mathrm{tail}}\},
\qquad
r_{\sup}(u_{\mathrm{tail}})
:=\sup\mathscr R(u_{\mathrm{tail}}).
\]
The set contains zero, is closed by continuity, and is bounded above
by one because \(F_R(r)\ge F_R(1)=1\) for \(r\ge1\). Its supremum
therefore belongs to the set. Consequently,
\[
\begin{aligned}
\mathbb P\{F_R(R(X))\le u_{\mathrm{tail}}\}
&\le \mathbb P\{R(X)\le r_{\sup}(u_{\mathrm{tail}})\}\\
&=F_R(r_{\sup}(u_{\mathrm{tail}}))
\le u_{\mathrm{tail}}.
\end{aligned}
\]
For \(u_{\mathrm{tail}}=1\), the same upper bound follows from total
probability one.

For \(t\in[0,1]\), raising \(e_0(x)\le t\) to the positive power
\(\gamma-1\) gives
\[
\{x:e_0(x)\le t\}
\subseteq
\{x:F_R(R(x))\le t^{\gamma-1}\}.
\]
Since \(t^{\gamma-1}\in[0,1]\), the preceding bound yields
\[
\mathbb P\{e_0(X)\le t\}
\le t^{\gamma-1}
\le Ct^{\gamma-1}.
\]
At \(t=0\), this shows
\[
\mathbb P\{e_0(X)=0\}=0.
\]
Together with \(e_0\ge0\), it follows that \(e_0(X)>0\) almost surely.
% lean: CausalSmith.Stat.WeakOverlap.radialCDF_self_sublevel_mass
% lean: CausalSmith.Stat.WeakOverlap.radialPropensity_uniform_tail
% lean: CausalSmith.Stat.WeakOverlap.radialPropensity_uniform_pos_ae

\item We verify the remaining model requirements for the constructed
completions. Their covariate marginal is uniform on the cube, so
\[
f(x)=1\ge c_f
\qquad\text{for Lebesgue-almost every }x\in\mathcal X.
\]
The conditional product kernel has the factorization
\[
\begin{aligned}
&\mathcal L_{\widetilde P_{j,h}}
   \bigl((Y(0),Y(1)),A\mid X=x\bigr)\\
&\qquad=
\left(
B\operatorname{Bernoulli}(\theta_0)
\otimes
B\operatorname{Bernoulli}(p_{j,h}(x))
\right)
\otimes\operatorname{Bernoulli}(e_0(x)).
\end{aligned}
\]
Thus \((Y(0),Y(1))\perp A\mid X\), and
\[
\mathbb P_{\widetilde P_{j,h}}\{A=1\mid X=x\}=e_0(x)
\]
almost everywhere. The observation map enforces
\[
Y=AY(1)+(1-A)Y(0)
\qquad\widetilde P_{j,h}\text{-almost surely}.
\]

The endpoint-supported potential-outcome kernels and this observation
map imply \(0\le Y\le B\) almost surely. To apply this support bound to
the treated conditional distribution, define
\[
\nu_{j,x}:=\mathcal L_{P_{j,h}}(Y\mid X=x,A=1).
\]
Disintegrating the support statement first over \(X\) and then over
\(A\), and using \(e_0(X)>0\) almost surely, gives
\[
\nu_{j,x}([0,B])=1
\qquad\text{for Lebesgue-almost every }x\in\mathcal X.
\]
For such \(x\), put
\[
\overline y_{j,x}:=\int y\,\nu_{j,x}(dy).
\]
Integrating the support bounds gives
\[
0\le\overline y_{j,x}\le B.
\]
The bounded-variable sub-Gaussian moment bound gives, for every
\(\lambda\in\mathbb R\),
\[
\int
\exp\{\lambda(y-\overline y_{j,x})\}\,\nu_{j,x}(dy)
\le
\exp\left\{\frac{B^2\lambda^2}{8}\right\}
\le
\exp\left\{\frac{B^2\lambda^2}{2}\right\}.
\]
The exponential is integrable because \(y\in[0,B]\) almost surely.
These bounds establish
\cref{obj:ass:bounded-potential-outcomes}. The density, consistency,
exchangeability, propensity-tail, and smoothness properties proved
above establish the other requirements of \cref{obj:def:model}.
Therefore, with these exact completion witnesses and response curves,
\[
P_{0,h},P_{1,h}\in\mathcal P_\gamma.
\]
% lean: CausalSmith.Stat.WeakOverlap.completionBind_uniformCube_globalTailModel
% lean: CausalSmith.Stat.WeakOverlap.completionBind_boundedOutcome_of_endpoint_support

\item For completeness, the two potential-outcome factors at every
covariate point assign total mass one to \(\{0,B\}\). Integrating their
product over the cube therefore gives
\[
\widetilde P_{j,h}
\{Y(0)\in\{0,B\},\ Y(1)\in\{0,B\}\}
=
\int_{\mathcal X}1\,dx
=1,
\qquad j\in\{0,1\}.
\]
This is the required endpoint support under each constructed completion.
% lean: CausalSmith.Stat.WeakOverlap.lowerPairCompletion_supported_on_endpoints

\item We now bound the entropy of the conditional observation laws.
For \(x\in\mathbb R^d\), define
\[
p_{\mathrm{Ber},x}:=p_{0,h}(x),
\qquad
q_{\mathrm{Ber},x}:=p_{1,h}(x).
\]
Both parameters lie in \((0,1)\), and
\[
\theta_0\le q_{\mathrm{Ber},x}\le\frac12,
\qquad
q_{\mathrm{Ber},x}(1-q_{\mathrm{Ber},x})
\ge\frac{\theta_0}{2}.
\]
Applying \(\log v\le v-1\), for \(v>0\), to each likelihood-ratio term gives
\[
\begin{aligned}
&\operatorname{KL}\!\left(
\operatorname{Bernoulli}(p_{\mathrm{Ber},x}),
\operatorname{Bernoulli}(q_{\mathrm{Ber},x})
\right)\\
&\quad=
p_{\mathrm{Ber},x}\log\frac{p_{\mathrm{Ber},x}}{q_{\mathrm{Ber},x}}
+(1-p_{\mathrm{Ber},x})
 \log\frac{1-p_{\mathrm{Ber},x}}{1-q_{\mathrm{Ber},x}}\\
&\quad\le
p_{\mathrm{Ber},x}
 \left(\frac{p_{\mathrm{Ber},x}}{q_{\mathrm{Ber},x}}-1\right)
+(1-p_{\mathrm{Ber},x})
 \left(\frac{1-p_{\mathrm{Ber},x}}{1-q_{\mathrm{Ber},x}}-1\right)\\
&\quad=
\frac{(p_{\mathrm{Ber},x}-q_{\mathrm{Ber},x})^2}
     {q_{\mathrm{Ber},x}(1-q_{\mathrm{Ber},x})}\\
&\quad\le
\frac{2}{\theta_0}(p_{\mathrm{Ber},x}-q_{\mathrm{Ber},x})^2.
\end{aligned}
\]
Also,
\[
0\le p_{1,h}(x)-p_{0,h}(x)
=ah^\beta\psi((x-x_0)/h)
\le ah^\beta.
\]

Define the laws of the treatment and binary selected outcome by
\[
\mathsf V_{j,x}:=
(1-e_0(x))\,\delta_0\otimes\operatorname{Bernoulli}(\theta_0)
+
e_0(x)\,\delta_1\otimes\operatorname{Bernoulli}(p_{j,h}(x)),
\]
on \(\{0,1\}^2\). Their scaled and recorded versions are
\[
\mathsf W_{j,x}
:=
\bigl((b_A,b_Y)\mapsto(b_A,Bb_Y)\bigr)_\#\mathsf V_{j,x},
\qquad
\mathsf Q_{j,x}
:=
\bigl((b_A,y)\mapsto(x,b_A,y)\bigr)_\#\mathsf W_{j,x}.
\]
Here the first map has domain \(\{0,1\}^2\), the second has domain
\(\{0,1\}\times\mathbb R\), and \(\#\) denotes image measure.
The product construction shows that \(\mathsf Q_{j,x}\) is the
conditional observation law at covariate \(x\).

The control parameters agree, so the control contribution to entropy
is zero. The treated contribution satisfies the preceding quadratic
bound. Conditioning on the binary treatment therefore gives
\[
\begin{aligned}
\operatorname{KL}(\mathsf V_{0,x},\mathsf V_{1,x})
&=(1-e_0(x))
  \operatorname{KL}\!\left(
  \operatorname{Bernoulli}(\theta_0),
  \operatorname{Bernoulli}(\theta_0)\right)\\
&\quad+
e_0(x)\operatorname{KL}\!\left(
\operatorname{Bernoulli}(p_{0,h}(x)),
\operatorname{Bernoulli}(p_{1,h}(x))\right)\\
&\le \frac{2}{\theta_0}(ah^\beta)^2e_0(x).
\end{aligned}
\]
Applying entropy contraction first to recording the covariate and then
to scaling the binary outcome yields
\[
\operatorname{KL}(\mathsf Q_{0,x},\mathsf Q_{1,x})
\le
\operatorname{KL}(\mathsf W_{0,x},\mathsf W_{1,x})
\le
\operatorname{KL}(\mathsf V_{0,x},\mathsf V_{1,x})
\le
\frac{2}{\theta_0}(ah^\beta)^2e_0(x).
\]
In particular, the conditional entropies are finite.
% lean: CausalSmith.Stat.WeakOverlap.lowerPair_bernoulli_kl_le_quadratic
% lean: CausalSmith.Stat.WeakOverlap.lowerPair_bernoulli_kl_le_baseline_quadratic
% lean: CausalSmith.Stat.WeakOverlap.lowerPair_completionAt_observed_klDiv_le

\item Define the open bump box by
\[
E_h:=\{x\in\mathbb R^d:|x_i-x_{0i}|<h
\text{ for every }i=1,\ldots,d\}.
\]
Outside \(E_h\), at least one coordinate factor of the bump is zero.
Hence
\[
p_{0,h}(x)=p_{1,h}(x),
\qquad
\mathsf Q_{0,x}=\mathsf Q_{1,x},
\qquad x\notin E_h.
\]
For every \(r\ge0\), each truncated coordinate interval has length at
most \(2r\), so
\[
F_R(r)\le(2r)^d.
\]
For \(x\in E_h\), we have \(R(x)\le h\), and thus
\[
e_0(x)
=F_R(R(x))^{1/(\gamma-1)}
\le \bigl((2R(x))^d\bigr)^{1/(\gamma-1)}
\le \bigl((2h)^d\bigr)^{1/(\gamma-1)}.
\]
Moreover,
\[
\operatorname{vol}(E_h\cap\mathcal X)
\le\mathbb P\{R(X)\le h\}
=F_R(h)
\le(2h)^d.
\]

Both observed laws are obtained by integrating \(\mathsf Q_{j,x}\)
against the same uniform covariate law, and each conditional law
records its covariate \(x\) in the first coordinate. The entropy
disintegration identity therefore gives
\[
\operatorname{KL}(P_{0,h},P_{1,h})
=
\int_{\mathcal X}
\operatorname{KL}(\mathsf Q_{0,x},\mathsf Q_{1,x})\,dx.
\]
On \(E_h\), apply the conditional entropy and propensity bounds; on
its complement, the integrand is zero. Consequently,
\[
\begin{aligned}
\operatorname{KL}(P_{0,h},P_{1,h})
&\le
\frac{2}{\theta_0}(ah^\beta)^2
\bigl((2h)^d\bigr)^{1/(\gamma-1)}
\operatorname{vol}(E_h\cap\mathcal X)\\
&\le
\frac{2}{\theta_0}(ah^\beta)^2
\bigl((2h)^d\bigr)^{1/(\gamma-1)}(2h)^d\\
&=
\frac{2}{\theta_0}\,2^{D_\gamma}a^2h^{2\beta+D_\gamma}
=K_{\mathrm{KL}}a^2h^{\tau_{\mathrm{KL}}}.
\end{aligned}
\]
The last equality uses
\[
\frac{d}{\gamma-1}+d=D_\gamma.
\]
The finite entropy bound supplies both
\[
P_{0,h}\ll P_{1,h},
\qquad
\int
\left|\log\frac{dP_{0,h}}{dP_{1,h}}\right|\,dP_{0,h}<\infty.
\]

Finally, the oracle mesh satisfies
\[
h_{\gamma,n}^{\tau_{\mathrm{KL}}}=n^{-1},
\qquad
nh^{\tau_{\mathrm{KL}}}
=n(ch_{\gamma,n})^{\tau_{\mathrm{KL}}}
=c^{\tau_{\mathrm{KL}}}.
\]
The fixed calibration from the second step therefore gives
\[
n\operatorname{KL}(P_{0,h},P_{1,h})
\le K_{\mathrm{KL}}a^2c^{\tau_{\mathrm{KL}}}
\le\frac18.
\]
% lean: CausalSmith.Stat.WeakOverlap.lowerPair_observed_klDiv_le_raw
% lean: CausalSmith.Stat.WeakOverlap.lowerPair_raw_rate_identity
% lean: CausalSmith.Stat.WeakOverlap.lowerPair_observed_guards_and_kl_rate
% lean: CausalSmith.Stat.WeakOverlap.lowerPair_oracle_kl_balance

\item It remains to identify the chosen response representatives at
\(x_0\). The constructed memberships have the continuous Hölder
representatives \(\mu_{j,h}=Bp_{j,h}\). By
\cref{obj:lem:causal-identification}, both these representatives and the
chosen representatives associated with the same observed laws agree
almost everywhere with the treated conditional regression. Its
pointwise uniqueness conclusion gives
\[
\mu_{1,P_{j,h}}(x)=\mu_{j,h}(x)=Bp_{j,h}(x),
\qquad x\in\mathcal X,\quad j\in\{0,1\}.
\]
This includes boundary points of the cube.

Every coordinate factor of \(\psi\) equals one at zero, so
\[
\psi(0)=1,
\qquad
p_{1,h}(x_0)-p_{0,h}(x_0)=ah^\beta.
\]
Writing
\[
r_{\gamma,n}:=h_{\gamma,n}^{\beta},
\]
we conclude that
\[
\begin{aligned}
\left|\mu_{1,P_{1,h}}(x_0)-\mu_{1,P_{0,h}}(x_0)\right|
&=\left|Ba(ch_{\gamma,n})^\beta\right|\\
&=Bac^\beta r_{\gamma,n}
=c_0r_{\gamma,n}.
\end{aligned}
\]
The positivity of \(B,a,c\) justifies the second equality. Since \(n\ge1\)
was arbitrary, the fixed constants \(a,c,c_0\) satisfy all the stated
requirements.
% lean: CausalSmith.Stat.WeakOverlap.responseOf_eq_given_model_response
% lean: CausalSmith.Stat.WeakOverlap.pairSuccess_oracle_separation
% lean: CausalSmith.Stat.WeakOverlap.boundedLowerPair_spec
\end{enumerate}
\end{proof}

\Cref{obj:lem:bounded-lower-pair} supplies the three components of the testing comparison: two laws within the fixed-constant model, a bound on their relative entropy, and response separation at the oracle scale. Here \(\operatorname{KL}\) denotes Kullback--Leibler relative entropy. The absolute-continuity and integrability clauses support the information comparison for independent samples, whose laws are \(P_{0,h}^{\otimes n}\) and \(P_{1,h}^{\otimes n}\). The chosen Hölder representatives specify the pointwise targets, including their values on the cube boundary.

The lower guarantee in \cref{obj:thm:adaptive-minimax} uses this pair through the two-point testing principle \citep[Section~2.2]{tsybakov2009}. The information bound and target separation provide the comparison for expected absolute error within the subclass admitting potential outcomes in \(\{0,B\}\). The risk ordering in \cref{obj:thm:adaptive-minimax} then places the same oracle scale below the pointwise and expected spatial-supremum minimax risks of \cref{obj:def:risks}.

The quantifiers of the construction preserve the pointwise scope of the converse. For each fixed \(\gamma>1\) and each fixed \(x_0\) in the closed cube, the constants are chosen once and apply to every integer \(n\ge1\). Their permitted dependence on the evaluation point and fixed model parameters matches \cref{obj:thm:adaptive-minimax}. Cube truncation in \cref{obj:def:bounded-lower-pair} and the intrinsic smoothness convention in \cref{obj:def:holder} keep boundary evaluation within this same comparison.

\subsection{The thin-strip geometry}

The geometry assertions in \cref{obj:prop:strict-enlargement} concern the strip law of \cref{obj:def:thin-strip}. That construction uses the cube center and a radial-plus-strip propensity; its strip exponent is the parameter specified there, separately from the amplitude constant in the testing pair. Constant potential-outcome means keep response smoothness fixed while the propensity determines the local treated design.

The model-membership and tail clauses of \cref{obj:prop:strict-enlargement} place this law within the global-tail causal class. Its remaining clauses describe two distinct aspects of local geometry. The anti-concentration statement concerns the displayed selected propensity and the pointwise supremum over Euclidean neighborhoods, as specified in \cref{obj:def:dorn-a3}. The numerical clause has a counterexample even when propensity values equal to one are permitted. This establishes the selected-version failure underlying the strict enlargement within the causal class.

The second-moment clause instead concerns treatment-weighted variation in supremum-norm neighborhoods of the cube center. Its denominator is the probability of both treatment and neighborhood membership, so the ratio measures the conditional second moment of the second coordinate after normalization by neighborhood width. The zero limit in \cref{obj:prop:strict-enlargement} records degeneration of that coordinate's local treated variation under the stated strip-exponent range.

These conclusions have complementary roles. The selected-version statement establishes the comparison with the local anti-concentration predicate, while the normalized moment describes the shape of the treated distribution itself. Together with model membership and the global tail inequality, they characterize the treatment geometry covered by \cref{obj:thm:adaptive-minimax}.


\section{Source-model transfer and comparison scope}

The retained source restrictions fix a nonempty family \(\mathfrak P\), common class constants, and selected regression and propensity representatives before stating a uniform recovery guarantee. Under \citet[Section 2, Assumptions 1--2, author PDF pp.~3--7]{Dorn2026GlobalOverlap}, each member has uniform covariates on \([-1,1]^d\) and a selected measurable propensity satisfying
\[
e_P(X)=P(A=1\mid X)\in(0,1)
\qquad P\text{-almost surely}.
\]
Writing \(\mu_P\) for the selected treated regression, the common constants \(q>3\), \(M,\sigma_{\min},C>0\), and finite overlap exponent \(\gamma_0>1\) satisfy
\[
\mathbb E_P\!\left[|Y|^q\mid X,A=1\right]\le M^q,
\qquad
\operatorname{Var}_P(\mu_P(X))\le M,
\qquad
\operatorname{Var}_P(Y\mid X,A)\ge\sigma_{\min}^2,
\]
with conditional inequalities holding almost surely, and
\[
P\{e_P(X)\le t\}\le Ct^{\gamma_0-1}
\qquad(t\in[0,1]).
\]
Conditional outcomes in both treatment arms are Gaussian with their selected regression means and common variance \(\sigma_{\min}^2\). Both regression curves belong to \(\Sigma^{\mathrm D}(\beta,L_0)\), the source smoothness class with common \(\beta,L_0>0\).

The source smoothness convention uses derivative order
\[
m=\max\{k\in\mathbb Z:k<\beta\}=\lceil\beta\rceil-1
\]
and bounds the \((\beta-m)\)-Hölder seminorm of derivatives of order \(m\) by \(L_0\), using Euclidean distance. For the intrinsic cube interpretation used in \cref{obj:prop:dorn-a3-free-corollary}, interior derivatives through order \(m\) have continuous traces on the closed cube. The source upper comparison additionally imposes the family-level local condition in \cref{obj:def:dorn-a3}, with common positive constants and the selected pointwise propensity representatives, as in \citet[Section 2, Assumption 3, author PDF pp.~6--7]{Dorn2026GlobalOverlap}. These restrictions distinguish the retained source envelope from the fixed families to which the published upper comparison applies.

\subsection{Fixed-cube smoothness completion}

For the transfer in \cref{obj:prop:dorn-a3-free-corollary}, the treated-moment restriction supplies the function bound that complements the source seminorm. The full response class in \cref{obj:def:holder} also controls the function and all lower derivatives. The following completion statement connects these two smoothness descriptions for the paper's range \(\beta>1\).

\begin{lemmav}{lem:fixed-cube-holder-completion}[Hölder completion under affine scaling]
Fix an integer \(d\ge 1\) and a real number \(\beta>1\), and set \(m=\lceil\beta\rceil-1\). There exists a constant \(K_{d,\beta}>0\), depending only on \(d\) and \(\beta\), such that the following holds for every \(M_0,L_0\in\mathbb R\) and \(u:\mathbb R^d\to\mathbb R\) satisfying:
\begin{itemize}
\item \textbf{(Nonnegative bounds.)} \(M_0\ge0\) and \(L_0\ge0\).
\item \textbf{(Intrinsic regularity.)} \(u\in C^m([-1,1]^d)\), with interior derivatives through order \(m\) admitting continuous traces on the closed cube.
\item \textbf{(Top derivative modulus.)} For every multi-index \(\alpha\) with \(|\alpha|=m\) and every \(x,y\in[-1,1]^d\), the coordinate derivatives, interpreted through their continuous boundary traces, satisfy
\[
|D^\alpha u(x)-D^\alpha u(y)|
\le L_0\left(\sum_{i=1}^d(x_i-y_i)^2\right)^{(\beta-m)/2}.
\]
\item \textbf{(Response bound.)} \(|u(x)|\le M_0\) for every \(x\in[-1,1]^d\).
\end{itemize}
For the affine transport
\[
T(x)=2x-\mathbf1,\qquad x\in[0,1]^d,
\]
the transported function satisfies
\[
u\circ T\in\mathcal H^\beta\!\left(K_{d,\beta}(M_0+L_0)\right)
\]
on \([0,1]^d\), with the intrinsic Hölder ball defined in \cref{obj:def:holder}.
\end{lemmav}

\begin{proof}[Proof of \cref{obj:lem:fixed-cube-holder-completion}]
\textbf{1. Parameters and derivative notation.}
Set
\[
s:=\beta-m,\qquad
\mathscr C:=[-1,1]^d,\qquad
\mathscr C^\circ:=(-1,1)^d,\qquad
\mathscr U:=[0,1]^d.
\]
Since \(m=\lceil\beta\rceil-1\) and \(\beta>1\),
\[
m\ge1,\qquad
m<\beta\le m+1,\qquad
0<s\le1.
\]
For each integer \(j\ge0\), define the set of ordered coordinate indices and the associated derivative by
\[
\mathscr I_j:=\{1,\ldots,d\}^{\{1,\ldots,j\}},
\qquad
\partial_\iota^j w
:=
\partial_{\iota(1)}\cdots\partial_{\iota(j)}w,
\quad \iota\in\mathscr I_j.
\]
Here \(\mathscr I_0\) contains the empty index and
\(\partial_\iota^0w=w\). On a closed cube, these derivatives use their continuous interior traces. Symmetry of the interior derivatives identifies an ordered derivative with the multi-index derivative having the same coordinate multiplicities; continuity preserves this identity on the boundary.
% lean: CausalSmith.Stat.WeakOverlap.fixedCube_holder_completion

\textbf{2. Completion on the source cube.}
We first derive a constant \(K_0>0\), depending only on \(d,m,s\), with the following property. For arbitrary
\[
u_{\mathrm{aux}}:\mathbb R^d\to\mathbb R,
\qquad
M_{\mathrm{aux}},L_{\mathrm{aux}}\in[0,\infty),
\]
suppose that
\[
u_{\mathrm{aux}}\in C^m(\mathscr C),\qquad
|u_{\mathrm{aux}}(x)|\le M_{\mathrm{aux}}
\quad(x\in\mathscr C),
\]
and
\[
|\partial_\iota^m u_{\mathrm{aux}}(x)
-\partial_\iota^m u_{\mathrm{aux}}(y)|
\le L_{\mathrm{aux}}\|x-y\|_\infty^s
\quad(\iota\in\mathscr I_m,\ x,y\in\mathscr C).
\]
The desired completion estimates are
\[
|\partial_\iota^j u_{\mathrm{aux}}(x)|
\le K_0(M_{\mathrm{aux}}+L_{\mathrm{aux}})
\quad(0\le j\le m,\ \iota\in\mathscr I_j,\ x\in\mathscr C)
\]
and
\[
|\partial_\iota^m u_{\mathrm{aux}}(x)
-\partial_\iota^m u_{\mathrm{aux}}(y)|
\le K_0(M_{\mathrm{aux}}+L_{\mathrm{aux}})
\|x-y\|_\infty^s.
\]

To obtain the interpolation estimate used in this completion, first consider a function whose ordinary coordinate derivatives on \(\mathscr C\) satisfy the displayed top-order modulus. Multilinearity gives, for a direction vector \(v\in\mathbb R^d\),
\[
D^m u_{\mathrm{aux}}(x)[v,\ldots,v]
=
\sum_{\iota\in\mathscr I_m}
\left(\prod_{a=1}^m v_{\iota(a)}\right)
\partial_\iota^m u_{\mathrm{aux}}(x).
\]
Consequently,
\[
\begin{aligned}
&|D^m u_{\mathrm{aux}}(x)[v,\ldots,v]
-D^m u_{\mathrm{aux}}(y)[v,\ldots,v]|\\
&\qquad\le
d^m L_{\mathrm{aux}}\|x-y\|_\infty^s\|v\|_\infty^m\\
&\qquad\le
(d+1)^m L_{\mathrm{aux}}\|x-y\|_\infty^s\|v\|_\infty^m.
\end{aligned}
\]
Define
\[
C_{\mathrm{diag}}:=(d+1)^m,\qquad
\varphi_{x,y}(r):=
u_{\mathrm{aux}}\bigl(x+r(y-x)\bigr),
\quad r\in\mathbb R,
\]
for \(x\in\mathscr C^\circ\) and \(y\in\mathscr C\).
The segment lies in \(\mathscr C\), and its points with \(0\le r<1\) lie in \(\mathscr C^\circ\). The chain rule therefore gives
\[
\varphi_{x,y}^{(k)}(r)
=
D^k u_{\mathrm{aux}}\bigl(x+r(y-x)\bigr)
[y-x,\ldots,y-x],
\qquad 0\le k\le m,\quad 0\le r<1.
\]
Using the diagonal estimate and
\[
\|x+r(y-x)-x-r'(y-x)\|_\infty
=
|r-r'|\|y-x\|_\infty,
\]
we obtain
\[
|\varphi_{x,y}^{(m)}(r)-\varphi_{x,y}^{(m)}(r')|
\le
C_{\mathrm{diag}}L_{\mathrm{aux}}
\|y-x\|_\infty^{m+s}|r-r'|^s,
\qquad r,r'\in[0,1).
\]
Taylor's Lagrange remainder of degree \(m-1\), applied on \([0,1]\), supplies \(\xi\in(0,1)\) such that
\[
\varphi_{x,y}(1)
-\sum_{k=0}^{m-1}\frac{\varphi_{x,y}^{(k)}(0)}{k!}
=
\frac{\varphi_{x,y}^{(m)}(\xi)}{m!}.
\]
After subtracting the degree-\(m\) term,
\[
\varphi_{x,y}(1)
-\sum_{k=0}^{m}\frac{\varphi_{x,y}^{(k)}(0)}{k!}
=
\frac{\varphi_{x,y}^{(m)}(\xi)
-\varphi_{x,y}^{(m)}(0)}{m!}.
\]
Since \(\xi^s\le1\) and \(m!\ge1\), this proves
\[
\left|
u_{\mathrm{aux}}(y)
-\sum_{k=0}^{m}\frac{1}{k!}
D^k u_{\mathrm{aux}}(x)[y-x,\ldots,y-x]
\right|
\le
C_{\mathrm{diag}}L_{\mathrm{aux}}
\|y-x\|_\infty^{m+s}.
\]

We next bound this Taylor polynomial on a fixed tensor grid. Define
\[
\nu_j:=\frac{j}{m+1}\quad(0\le j\le m),\qquad
\mathscr A:=\{0,\ldots,m\}^d,\qquad
z_b:=(\nu_{b_1},\ldots,\nu_{b_d})
\quad(b\in\mathscr A).
\]
Every \(z_b\) belongs to \(\mathscr C\). Compactness supplies a finite number \(D_{\mathrm{cube}}\ge0\) satisfying
\[
\|y-x\|_\infty\le D_{\mathrm{cube}}
\quad(x,y\in\mathscr C).
\]
Set
\[
\kappa_{\mathrm{cube}}:=\max\{1,D_{\mathrm{cube}}\},
\qquad
C_{\mathrm{grid}}
:=
1+C_{\mathrm{diag}}\kappa_{\mathrm{cube}}^{m+s}.
\]
For the Taylor polynomial
\[
p_x(z):=
\sum_{k=0}^{m}\frac{1}{k!}
D^k u_{\mathrm{aux}}(x)[z-x,\ldots,z-x],
\qquad z\in\mathbb R^d,
\]
the response bound and remainder estimate yield
\[
\begin{aligned}
|p_x(z_b)|
&\le |u_{\mathrm{aux}}(z_b)|
   +|u_{\mathrm{aux}}(z_b)-p_x(z_b)|\\
&\le M_{\mathrm{aux}}
   +C_{\mathrm{diag}}L_{\mathrm{aux}}
      \kappa_{\mathrm{cube}}^{m+s}\\
&\le C_{\mathrm{grid}}
      (M_{\mathrm{aux}}+L_{\mathrm{aux}}).
\end{aligned}
\]

Here is the polynomial estimate that converts these grid bounds into derivative bounds. Define the univariate Lagrange polynomials, their coefficients, and the tensor weights by
\[
\Lambda_j(t):=
\prod_{\substack{0\le q\le m\\q\ne j}}
\frac{t-\nu_q}{\nu_j-\nu_q},
\qquad
\omega_{kj}:=[t^k]\Lambda_j(t),
\qquad
W_{ab}:=\prod_{i=1}^d\omega_{a_i b_i},
\]
where \(0\le k,j\le m\) and \(a,b\in\mathscr A\).
For any polynomial \(p_{\mathrm{aux}}\) of degree at most \(m\) in each coordinate, define
\[
c_a(p_{\mathrm{aux}}):=[z^a]p_{\mathrm{aux}}(z),
\qquad
z^a:=\prod_{i=1}^d z_i^{a_i}.
\]
Univariate Lagrange interpolation, applied to each monomial, gives
\[
\sum_{j=0}^m\omega_{kj}\nu_j^q
=
\begin{cases}
1,&k=q,\\
0,&k\ne q,
\end{cases}
\qquad 0\le k,q\le m.
\]
Taking products over the coordinates and then expanding \(p_{\mathrm{aux}}\) proves the coefficient recovery identity
\[
c_a(p_{\mathrm{aux}})
=
\sum_{b\in\mathscr A}W_{ab}p_{\mathrm{aux}}(z_b).
\]

The weights admit a uniform bound. Define the coefficient one-norm of a univariate polynomial by
\[
\|p_{\mathrm{uni}}\|_{\mathrm{coef},1}
:=
\sum_{k\ge0}|[t^k]p_{\mathrm{uni}}(t)|.
\]
The coefficient convolution formula and the triangle inequality imply
\[
\|p_{\mathrm{uni}}q_{\mathrm{uni}}\|_{\mathrm{coef},1}
\le
\|p_{\mathrm{uni}}\|_{\mathrm{coef},1}
\|q_{\mathrm{uni}}\|_{\mathrm{coef},1}.
\]
Since \(0\le\nu_q\le1\), each numerator factor has coefficient one-norm at most \(2\). Splitting the denominator factors below and above \(j\) gives
\[
\left|
\prod_{\substack{0\le q\le m\\q\ne j}}(\nu_j-\nu_q)
\right|
=
\frac{j!(m-j)!}{(m+1)^m}.
\]
Thus
\[
\sum_{k=0}^m|\omega_{kj}|
=
\|\Lambda_j\|_{\mathrm{coef},1}
\le
\frac{2^m(m+1)^m}{j!(m-j)!}.
\]
The binomial identity and the elementary binomial coefficient bounds give
\[
\sum_{j=0}^m\frac{1}{j!(m-j)!}
=\frac{2^m}{m!},
\qquad
\frac{(m+1)^m}{m!}
\le\binom{2m}{m}\le 2^{2m}=4^m.
\]
Consequently,
\[
\sum_{k=0}^m\sum_{j=0}^m|\omega_{kj}|
\le
4^m\frac{(m+1)^m}{m!}
\le 2^{4m}
\le 2^{4m+4}.
\]
Define
\[
C_{\mathrm{coef}}:=(2^{4m+4})^d.
\]
Factoring the finite tensor sums now yields
\[
\sum_{a\in\mathscr A}\sum_{b\in\mathscr A}|W_{ab}|
=
\left(\sum_{k=0}^m\sum_{j=0}^m|\omega_{kj}|\right)^d
\le C_{\mathrm{coef}}.
\]
Hence, whenever \(B_{\mathrm{pol}}\ge0\) and
\(|p_{\mathrm{aux}}(z_b)|\le B_{\mathrm{pol}}\) for every \(b\in\mathscr A\),
\[
|c_a(p_{\mathrm{aux}})|
\le
\sum_{b\in\mathscr A}|W_{ab}|
|p_{\mathrm{aux}}(z_b)|
\le C_{\mathrm{coef}}B_{\mathrm{pol}}.
\]

For each \(a\in\mathscr A\), the monomial \(z\mapsto z^a\) is smooth. Continuity of its iterated differentials on the compact cube, followed by a maximum over the finitely many \(a\) and \(0\le j\le m\), supplies \(M_{\mathrm{mon}}\ge0\) such that
\[
\|D^j(z\mapsto z^a)(x)\|_{\mathrm{op},\infty}
\le M_{\mathrm{mon}}
\quad(a\in\mathscr A,\ 0\le j\le m,\ x\in\mathscr C).
\]
Each coordinate unit vector has supremum norm \(1\), so
\[
|\partial_\iota^j(z\mapsto z^a)(x)|
\le M_{\mathrm{mon}}.
\]
Define
\[
C_{\mathrm{pol}}
:=
(m+1)^d C_{\mathrm{coef}}(M_{\mathrm{mon}}+1).
\]
Differentiating the finite expansion
\(p_{\mathrm{aux}}(z)=\sum_{a\in\mathscr A}c_a(p_{\mathrm{aux}})z^a\)
therefore gives
\[
\begin{aligned}
|\partial_\iota^j p_{\mathrm{aux}}(x)|
&\le
\sum_{a\in\mathscr A}|c_a(p_{\mathrm{aux}})|
|\partial_\iota^j(z\mapsto z^a)(x)|\\
&\le (m+1)^d C_{\mathrm{coef}}B_{\mathrm{pol}}M_{\mathrm{mon}}\\
&\le C_{\mathrm{pol}}B_{\mathrm{pol}},
\end{aligned}
\]
for \(0\le j\le m\), \(\iota\in\mathscr I_j\), and \(x\in\mathscr C\).

The polynomial \(p_x\) has total degree at most \(m\), as is also seen from its coordinate expansion
\[
p_x(z)=
\sum_{k=0}^m\frac{1}{k!}
\sum_{\iota\in\mathscr I_k}
\partial_\iota^k u_{\mathrm{aux}}(x)
\prod_{a=1}^k(z_{\iota(a)}-x_{\iota(a)}).
\]
Its derivatives at the expansion point agree with those of \(u_{\mathrm{aux}}\). Indeed, a homogeneous term of degree \(k<j\) has zero \(j\)-th derivative; a term of degree \(k>j\) retains a vanishing factor at \(x\); and for \(k=j\), symmetry of the differential produces \(j!\) identical contributions, cancelling \(1/j!\). Thus
\[
\partial_\iota^j p_x(x)
=
\partial_\iota^j u_{\mathrm{aux}}(x),
\qquad 0\le j\le m.
\]
Set
\[
K_{\mathrm{int}}:=C_{\mathrm{pol}}C_{\mathrm{grid}}>0.
\]
Combining the grid bound, polynomial derivative bound, and derivative agreement proves
\[
|\partial_\iota^j u_{\mathrm{aux}}(x)|
\le K_{\mathrm{int}}(M_{\mathrm{aux}}+L_{\mathrm{aux}})
\quad(x\in\mathscr C^\circ)
\]
under the ordinary-derivative modulus used above.

To apply this estimate to intrinsic cube derivatives, define the inward contractions
\[
u_{\mathrm{aux},t}(z):=u_{\mathrm{aux}}(tz),
\qquad z\in\mathbb R^d,\quad 0<t<1.
\]
For every \(z\in\mathscr C\), \(tz\in\mathscr C^\circ\). The chain rule on this interior region gives
\[
u_{\mathrm{aux},t}\in C^m(\mathscr C),\qquad
|u_{\mathrm{aux},t}(z)|\le M_{\mathrm{aux}},
\qquad
\partial_\iota^j u_{\mathrm{aux},t}(z)
=t^j\partial_\iota^j u_{\mathrm{aux}}(tz).
\]
The assumed intrinsic top-order modulus implies
\[
\begin{aligned}
&|\partial_\iota^m u_{\mathrm{aux},t}(x)
-\partial_\iota^m u_{\mathrm{aux},t}(y)|\\
&\qquad\le
t^mL_{\mathrm{aux}}\|tx-ty\|_\infty^s\\
&\qquad=
t^m t^s L_{\mathrm{aux}}\|x-y\|_\infty^s
\le L_{\mathrm{aux}}\|x-y\|_\infty^s,
\end{aligned}
\]
because \(t^m\le1\) and \(t^s\le1\). The preceding interpolation estimate therefore applies to \(u_{\mathrm{aux},t}\), yielding
\[
|t^j\partial_\iota^j u_{\mathrm{aux}}(tx)|
\le K_{\mathrm{int}}(M_{\mathrm{aux}}+L_{\mathrm{aux}})
\quad(x\in\mathscr C^\circ,\ 0<t<1).
\]
Continuity of the derivative traces allows \(t\uparrow1\), giving
\[
|\partial_\iota^j u_{\mathrm{aux}}(x)|
\le K_{\mathrm{int}}(M_{\mathrm{aux}}+L_{\mathrm{aux}})
\quad(x\in\mathscr C^\circ).
\]
Since \(\overline{\mathscr C^\circ}=\mathscr C\), the same continuity extends this bound to every \(x\in\mathscr C\).

Finally, define
\[
K_0:=K_{\mathrm{int}}+1>0.
\]
Nonnegativity of the bounds gives
\[
K_{\mathrm{int}}(M_{\mathrm{aux}}+L_{\mathrm{aux}})
\le K_0(M_{\mathrm{aux}}+L_{\mathrm{aux}}),
\qquad
L_{\mathrm{aux}}
\le K_0(M_{\mathrm{aux}}+L_{\mathrm{aux}}).
\]
These inequalities prove both completion estimates stated at the beginning of this step, together with the original \(C^m\) regularity.
% lean: Causalean.Mathlib.Analysis.Calculus.CubeExtension.cube_holder_completion

\textbf{3. Affine transport of a completed radius.}
Define
\[
K_1:=2^{m+1}>0.
\]
Suppose that a function \(u_{\mathrm{aux}}\in C^m(\mathscr C)\) has derivative and top-order modulus bounds of radius \(\rho\ge0\):
\[
|\partial_\iota^j u_{\mathrm{aux}}(x)|\le\rho,
\qquad
|\partial_\iota^m u_{\mathrm{aux}}(x)
-\partial_\iota^m u_{\mathrm{aux}}(y)|
\le\rho\|x-y\|_\infty^s.
\]
The affine map \(T(x)=2x-\mathbf1\) maps \(\mathscr U\) into \(\mathscr C\). Composition preserves the intrinsic \(C^m\) regularity, and the chain rule gives
\[
\partial_\iota^j(u_{\mathrm{aux}}\circ T)(x)
=
2^j\partial_\iota^j u_{\mathrm{aux}}(T(x)),
\qquad x\in\mathscr U,\quad 0\le j\le m.
\]
This identity holds on the interior and extends to the boundary by the continuous traces. Hence
\[
|\partial_\iota^j(u_{\mathrm{aux}}\circ T)(x)|
\le2^j\rho\le K_1\rho.
\]
Moreover,
\[
\|T(x)-T(y)\|_\infty=2\|x-y\|_\infty,
\qquad 2^s\le2,
\]
so
\[
\begin{aligned}
&|\partial_\iota^m(u_{\mathrm{aux}}\circ T)(x)
-\partial_\iota^m(u_{\mathrm{aux}}\circ T)(y)|\\
&\qquad\le
2^m\rho\|T(x)-T(y)\|_\infty^s\\
&\qquad=
2^m2^s\rho\|x-y\|_\infty^s\\
&\qquad\le
K_1\rho\|x-y\|_\infty^s.
\end{aligned}
\]
Thus affine transport multiplies the completed radius by at most \(K_1\).
% lean: Causalean.Mathlib.Analysis.Calculus.CubeExtension.cube_holder_affine_transport

\textbf{4. Conversion of the given modulus.}
Choose
\[
C_{\mathrm{norm}}:=(\sqrt d)^s>0,
\qquad
K_{d,\beta}:=K_1K_0(C_{\mathrm{norm}}+1)>0.
\]
These constants depend only on \(d\) and \(\beta\). For any
\(\iota\in\mathscr I_m\), define its coordinate multiplicities by
\[
\alpha_i:=\#\{a\in\{1,\ldots,m\}:\iota(a)=i\},
\qquad i=1,\ldots,d.
\]
Then
\[
|\alpha|=\sum_{i=1}^d\alpha_i=m,
\qquad
\partial_\iota^m u=D^\alpha u.
\]
The hypothesis therefore bounds every ordered top derivative. For \(x,y\in\mathscr C\),
\[
\sum_{i=1}^d(x_i-y_i)^2
\le d\|x-y\|_\infty^2,
\qquad
\left(\sum_{i=1}^d(x_i-y_i)^2\right)^{1/2}
\le\sqrt d\,\|x-y\|_\infty.
\]
Since \(s>0\), raising the latter inequality to the power \(s\) gives
\[
\left(\sum_{i=1}^d(x_i-y_i)^2\right)^{s/2}
\le C_{\mathrm{norm}}\|x-y\|_\infty^s.
\]
Multiplication by \(L_0\ge0\) now yields
\[
|\partial_\iota^m u(x)-\partial_\iota^m u(y)|
\le
C_{\mathrm{norm}}L_0\|x-y\|_\infty^s.
\]
% lean: CausalSmith.Stat.WeakOverlap.fixedCube_holder_completion

\textbf{5. Completion and transport for \(u\).}
Define the source radius
\[
\rho_0:=K_0(M_0+C_{\mathrm{norm}}L_0)\ge0.
\]
The source-cube completion estimates apply with the response bound \(M_0\) and modulus constant \(C_{\mathrm{norm}}L_0\). They give
\[
u\in C^m(\mathscr C),\qquad
|\partial_\iota^j u(x)|\le\rho_0
\quad(0\le j\le m,\ \iota\in\mathscr I_j,\ x\in\mathscr C),
\]
and
\[
|\partial_\iota^m u(x)-\partial_\iota^m u(y)|
\le\rho_0\|x-y\|_\infty^s.
\]
Applying the affine transport estimate with radius \(\rho_0\) proves
\[
u\circ T\in C^m(\mathscr U),\qquad
|\partial_\iota^j(u\circ T)(x)|\le K_1\rho_0,
\]
and
\[
|\partial_\iota^m(u\circ T)(x)
-\partial_\iota^m(u\circ T)(y)|
\le K_1\rho_0\|x-y\|_\infty^s.
\]
% lean: CausalSmith.Stat.WeakOverlap.fixedCube_holder_completion

\textbf{6. Calibration of the final radius.}
The nonnegative bounds imply
\[
\begin{aligned}
&(C_{\mathrm{norm}}+1)(M_0+L_0)
-(M_0+C_{\mathrm{norm}}L_0)\\
&\qquad=C_{\mathrm{norm}}M_0+L_0\ge0.
\end{aligned}
\]
Multiplying by \(K_1K_0>0\) gives
\[
K_1\rho_0
=
K_1K_0(M_0+C_{\mathrm{norm}}L_0)
\le K_{d,\beta}(M_0+L_0).
\]
Consequently,
\[
\max_{|\alpha|\le m}
\|D^\alpha(u\circ T)\|_\infty
\le K_{d,\beta}(M_0+L_0),
\]
and, for every \(|\alpha|=m\) and \(x,y\in\mathscr U\),
\[
|D^\alpha(u\circ T)(x)-D^\alpha(u\circ T)(y)|
\le
K_{d,\beta}(M_0+L_0)\|x-y\|_\infty^{\beta-m}.
\]
Together with the transported intrinsic \(C^m\) regularity, these are precisely the restrictions in
\cref{obj:ass:holder-derivative-bound,obj:ass:holder-modulus}.
By \cref{obj:def:holder},
\[
u\circ T\in
\mathcal H^\beta\!\left(K_{d,\beta}(M_0+L_0)\right).
\]
All estimates include zero values of either bound, and the argument uses \(0<s\le1\), including \(s=1\).
% lean: CausalSmith.Stat.WeakOverlap.fixedCube_holder_completion
\end{proof}

\begin{propositionv}{prop:dorn-a3-free-corollary}[Transfer of Dorn source families]*
Fix parameters satisfying:
\begin{itemize}
\item \textbf{(Dimension and smoothness.)} \(d\in\mathbb N\), \(d\ge1\), and \(\beta>1\).
\item \textbf{(Source constants.)} \(q>3\) and \(M,\sigma_{\min},C,L_0>0\).
\item \textbf{(Adaptation range.)} \(1<\gamma_{\min}<\gamma_{\max}\), with
\[
\Gamma=[\gamma_{\min},\gamma_{\max}].
\]
\end{itemize}

For each \(\gamma\in\Gamma\), let \(\mathcal D_\gamma\) be the maximal envelope of source-law tuples consisting of an observed law \(P\), treated and control mean functions, and a selected propensity \(e_P\), satisfying the following retained restrictions from \citet[Section 2, Assumptions 1--2, pp.~3--7]{Dorn2026GlobalOverlap}: \(X\) is uniform on \([-1,1]^d\); the selected propensity belongs to \((0,1)\) almost surely; conditional on \(X,A\), the outcome is Gaussian with variance \(\sigma_{\min}^2\); both conditional mean functions belong to \(\Sigma^{\mathrm D}(\beta,L_0)\); and
\[
\mathbb E_P\!\left[|Y|^q\mid X,A=1\right]\le M^q,
\qquad
\operatorname{Var}_P(\mu_P(X))\le M,
\qquad
\operatorname{Var}_P(Y\mid X,A)\ge\sigma_{\min}^2,
\]
\[
P\{e_P(X)\le t\}\le Ct^{\gamma-1}
\qquad (0\le t\le1).
\]
Here \(\mu_P\) is the selected treated mean. The source class \(\Sigma^{\mathrm D}(\beta,L_0)\) consists of functions in \(C^m([-1,1]^d)\), with continuous boundary traces of their derivatives, whose derivatives of order
\[
m=\max\{k\in\mathbb Z:k<\beta\}
\]
have \((\beta-m)\)-Hölder seminorm at most \(L_0\), using the source's convention.

Let \(T\), the affine transport from the unit cube to the source cube, and the common response bound \(B_*\) be
\[
T(x)=2x-\mathbf1,
\qquad
B_*=\max\{M,\sigma_{\min}\}.
\]
Define the common transported Hölder radius and the interpolation constant by
\[
L_*=
\sup_{\substack{u\in\Sigma^{\mathrm D}(\beta,L_0)\\
                 \sup_{x\in[-1,1]^d}|u(x)|\le M}}
\|u\circ T\|_{\mathcal H^\beta},
\]
\[
C_{d,\beta}=
\sup_{\substack{M'>0,\ L_0'>0\\
                 u\in\Sigma^{\mathrm D}(\beta,L_0')\\
                 \sup_{x\in[-1,1]^d}|u(x)|\le M'}}
\frac{\|u\circ T\|_{\mathcal H^\beta}}{M'+L_0'}.
\]
The functional \(\|\cdot\|_{\mathcal H^\beta}\) is the full Hölder-radius functional on \([0,1]^d\), including the function, lower derivatives, and top-derivative modulus. Then
\[
0<C_{d,\beta}<\infty,
\qquad
0<L_*\le C_{d,\beta}(M+L_0)<\infty.
\]

Every member of \(\mathcal D_\gamma\), after replacing \(X\) by \(T^{-1}(X)\), belongs to the global-tail causal class \(\mathcal P_\gamma\) with constants
\[
B=B_*,
\qquad L=L_*,
\qquad C=C,
\qquad c_f=1.
\]
In particular, its transported observed law admits a consistency-and-exchangeability causal completion with these constants.

Let \(\widehat\mu_n\) be the equal-cell estimator of \cref{obj:def:estimator}, with smoothness \(\beta\) and response bound \(B_*\), applied to \(n\) independent transported observations. Write \(r_{\gamma,n}\), the oracle risk scale, as
\[
r_{\gamma,n}=n^{-\beta/(2\beta+d\gamma/(\gamma-1))}
\qquad(n\ge1).
\]
There exist \(K>0\) and \(n_0\in\mathbb N\) such that, for every \(n\ge n_0\), every \(\gamma\in\Gamma\), and every source tuple in \(\mathcal D_\gamma\),
\[
\mathbb E_{P^{\otimes n}}
\left[
\sup_{x\in[0,1]^d}
\left|\widehat\mu_n(x)-\mu_P(T(x))\right|
\right]
\le K r_{\gamma,n}.
\]
Every source family \(\mathfrak P\) satisfying the cited Assumptions 1--2 with the fixed displayed constants is nonempty and contained in \(\mathcal D_\gamma\). The same \(K,n_0\) therefore satisfy
\[
\sup_{P\in\mathfrak P}
\mathbb E_{P^{\otimes n}}
\left[
\sup_{x\in[0,1]^d}
\left|\widehat\mu_n(x)-\mu_P(T(x))\right|
\right]
\le K r_{\gamma,n}
\qquad(n\ge n_0).
\]

For every \(\varepsilon>0\), there exist \(c(\varepsilon)>0\) and \(n_0\in\mathbb N\) such that, simultaneously for all \(n\ge n_0\), \(\gamma\in\Gamma\), and source tuples in \(\mathcal D_\gamma\),
\[
P^{\otimes n}\!\left\{
\left\|\widehat\mu_n\circ T^{-1}-\mu_P\right\|_{L^\infty(P)}
>c(\varepsilon)r_{\gamma,n}
\right\}\le\varepsilon.
\]
Here the norm is the essential supremum under the source covariate law, and the estimator is computed from the transported sample.

If \(d\ge2\), then for each \(\gamma\in\Gamma\) there exist positive finite constants \(C',M',\sigma_{\min}',L_0'\) and a source tuple with uniform covariates, Gaussian conditional outcomes, and a selected pointwise propensity representative satisfying all the retained restrictions with these constants, for which
\[
\neg\,
\mathsf A_{3}^{\mathrm{Dorn}}
\bigl(\{P\},(e_P)\bigr),
\]
where the family-level condition is defined in \cref{obj:def:dorn-a3}. The constants in this existence assertion are selected separately for each \(\gamma\).

For any nonempty family \(\mathfrak P\) of source-law tuples sharing a common treated regression curve \(\mu^\circ\), meaning
\[
\mu_P(x)=\mu^\circ(x)
\qquad(P\in\mathfrak P,\ x\in\mathbb R^d),
\]
the minimax pointwise miss-probability value at any evaluation point \(x_0\), sample size \(n\in\mathbb N\), and positive threshold \(cr\), with \(c,r>0\), equals zero.

Finally, for every \(\gamma>1\) and every \(x_0\in[0,1]^d\), there exists \(c_\gamma>0\) such that
\[
c_\gamma r_{\gamma,n}
\le R_{\mathrm{pt}}(n,\gamma,x_0)
\qquad(n\ge1),
\]
where \(R_{\mathrm{pt}}\) is the pointwise minimax expected absolute-error risk of \cref{obj:def:risks} over the causal class with constants \(B_*,L_*,\max\{C,1\}\), and \(c_f=1\). The displayed lower bound is first established for the two-valued potential-outcome subclass risk; monotonicity under class inclusion transfers it to \(R_{\mathrm{pt}}\). The constant \(c_\gamma\) may depend on the evaluation point and the fixed parameters.
\end{propositionv}
% CAUSALSMITH-CITED-SCOPE-BEGIN prop:dorn-a3-free-corollary
\verificationfootnotetext{\textbf{Formalization scope.} This result uses the cited conclusion \citep{Dorn2026GlobalOverlap} (Section 2 Setting and Notation, Assumptions 1--2; author PDF pages 3--7). The cited source proof is not formalized here: Lean verifies its use and all remaining steps. If that cited conclusion is false, the portions of this result that depend on it are not certified.}
% CAUSALSMITH-CITED-SCOPE-END prop:dorn-a3-free-corollary

The completion constant depends on dimension and smoothness alone. Consequently, a common function bound and a common top-derivative modulus yield a common full Hölder radius after affine transport. Applied to the treated regression curves in \cref{obj:prop:dorn-a3-free-corollary}, this statement supplies the finite radius needed by the equal-cell estimator. The continuous boundary traces make the completion applicable on the entire closed cube, including its faces and corners.

\subsection{Source loss and the fixed-family criterion}

For the source comparisons, the regression target is \(\mu_P(x)=\mathbb E_P[Y\mid X=x,A=1]\), with the selected version understood. Let \(\mathcal D_n\) denote the sample of \(n\) independent observed triples, whose sampling law is \(P^{\otimes n}\). For a regression estimator \(\widehat\mu_n\), the source loss is
\[
\|\widehat\mu_n-\mu_P\|_{L^\infty(P)}
=
\operatorname*{ess\,sup}_{x\sim P_X}
|\widehat\mu_n(\mathcal D_n,x)-\mu_P(x)|.
\]
Thus the essential supremum is taken under the source covariate marginal. For a fixed finite overlap exponent \(\gamma_0>1\), write
\[
r^*_{\mu,n}
=
n^{-\beta/(2\beta+d\gamma_0/(\gamma_0-1))}
\qquad(n\ge1)
\]
for the source comparison scale. The fixed-family probability criterion bounds the probability that the source loss exceeds \(c(\varepsilon)r^*_{\mu,n}\), where \(c(\varepsilon)\) is a positive threshold constant allowed to depend on the fixed setup.

A pointwise minimax miss probability also depends on which regression curves the family contains. The next statement records the common-curve obstruction once; later comparisons invoke it by reference.

\begin{lemmav}{lem:arbitrary-family-lower-obstruction}[Known regression minimax obstruction]
Fix parameters satisfying
\begin{itemize}
\item \textbf{(Dimension.)} \(d\in\mathbb N\) and \(d\ge1\).
\item \textbf{(Smoothness.)} \(\beta>0\) and \(L_0>0\).
\item \textbf{(Overlap.)} \(\gamma\in\mathbb R\) and \(\gamma>1\).
\end{itemize}
Let \(\mathfrak P\) be any nonempty source family equipped with selected treated and control regression functions and propensity representatives. Suppose its selected treated regression \(\mu_P\) is the same known function \(\mu^\circ:\mathbb R^d\to\mathbb R\) throughout the family:
\[
\mu_P(x)=\mu^\circ(x)
\qquad
\text{for every }P\in\mathfrak P
\text{ and every }x\in\mathbb R^d.
\]
For a sample size \(n\in\mathbb N\), let \(\mathcal D_n\) denote the observed sample and \(P^{\otimes n}\) its product sampling law. Then, for every evaluation point \(x_0\in\mathbb R^d\), every \(n\in\mathbb N\), and every \(c,r_n>0\), where \(r_n\) is an arbitrary positive comparison scale,
\[
\inf_{T\ \mathrm{measurable}}
\sup_{P\in\mathfrak P}
P^{\otimes n}
\left\{\left|T(\mathcal D_n)-\mu_P(x_0)\right|>c r_n\right\}
=0.
\]
The infimum ranges over all measurable real-valued sample estimators \(T\), and the minimax value is taken in the extended nonnegative reals.

In particular, consider the singleton family whose observed law has mutually independent coordinates
\[
X\sim\operatorname{Unif}([-1,1]^d),
\qquad
A\sim\operatorname{Bernoulli}(1/2),
\qquad
Y\sim N(0,1),
\]
with selected treated and control regression functions identically zero and selected propensity \(e_P(x)=1/2\). There exist common constants \(q>3\), \(M>0\), and \(C>0\) for which this law satisfies all of the following:
\begin{itemize}
\item \textbf{(Selected propensity.)} The selected propensity is a measurable version of the conditional treatment probability and lies in \((0,1)\) almost surely.
\item \textbf{(Moment and variance bounds.)} The treated conditional absolute \(q\)-moment is integrable and satisfies
\[
\mathbb E_P\!\left[|Y|^q\mid X=x,A=1\right]\le M^q
\quad\text{for }P_X\text{-almost every }x,
\qquad
\mathbb E_P\!\left[
\left(\mu_P(X)-\mathbb E_P[\mu_P(X)]\right)^2
\right]\le M.
\]
\item \textbf{(Source smoothness.)} Both selected regression functions belong to the source top-derivative seminorm class \(\Sigma^{\mathrm D}(\beta,L_0)\) specified in \cref{obj:prop:dorn-a3-free-corollary}.
\item \textbf{(Global overlap.)} The selected propensity satisfies
\[
P_X\{e_P(X)\le t\}\le C t^{\gamma-1}
\qquad\text{for every }t\in[0,1].
\]
\item \textbf{(Conditional Gaussian outcomes.)} For \(P_X\)-almost every \(x\), the conditional outcome distribution in each treatment arm is Gaussian with its selected regression mean and variance one.
\item \textbf{(Common local anti-concentration.)} The singleton satisfies
\[
\mathsf A_{3}^{\mathrm{Dorn}}
\bigl(\mathfrak P,(e_P)_{P\in\mathfrak P}\bigr),
\]
with the local ball-ratio inequalities and common positive constants specified in \cref{obj:def:dorn-a3}.
\end{itemize}
At the cube center \(x_0=0\), for every \(n\in\mathbb N\) and every \(c,r_n>0\), this singleton's displayed minimax miss probability equals zero.
\end{lemmav}

\begin{proof}[Proof of \cref{obj:lem:arbitrary-family-lower-obstruction}]
\begin{enumerate}
\item Fix the family, evaluation point, sample size, and positive constants in the first assertion. Define the constant estimator
\[
T_{\mathrm{known}}(\mathcal D_n):=\mu^\circ(x_0).
\]
It is measurable for every \(n\), including \(n=0\). For every \(P\in\mathfrak P\) and every sample,
\[
\left|T_{\mathrm{known}}(\mathcal D_n)-\mu_P(x_0)\right|
=\left|\mu^\circ(x_0)-\mu^\circ(x_0)\right|=0.
\]
Since \(c r_n>0\), its miss event is empty. Consequently, in the extended nonnegative reals,
\[
0\le
\inf_{T\ \mathrm{measurable}}\sup_{P\in\mathfrak P}
P^{\otimes n}\!\left\{
|T(\mathcal D_n)-\mu_P(x_0)|>c r_n
\right\}
\le
\sup_{P\in\mathfrak P}P^{\otimes n}(\varnothing)
=0.
\]
This proves the first assertion. % lean: CausalSmith.Stat.WeakOverlap.knownRegression_minimax_miss_zero

\item For the explicit witness, define the cube and component measures by
\[
\mathcal S:=[-1,1]^d,\qquad
\lambda_d:=\text{Lebesgue measure on }\mathbb R^d,
\]
\[
\Pi_X:=2^{-d}\lambda_d\!\restriction_{\mathcal S},
\qquad
\Pi_A:=\operatorname{Bernoulli}(1/2),
\qquad
\Pi_Y:=N(0,1),
\qquad
P^\circ:=\Pi_X\otimes(\Pi_A\otimes\Pi_Y).
\]
The cube-volume formula gives
\[
\lambda_d(\mathcal S)=\prod_{i=1}^d 2=2^d,
\qquad
\Pi_X(\mathbb R^d)=2^{-d}2^d=1.
\]
Also,
\[
\Pi_A(\{0,1\})=\frac12+\frac12=1,
\qquad
\Pi_Y(\mathbb R)=1.
\]
Thus \(P^\circ\) is a probability law with the independent coordinates specified in the statement. Its covariate marginal is
\[
P_X^\circ=\Pi_X.
\]
Select the treated regression, control regression, and propensity, respectively, as
\[
\mu_{P^\circ}(x):=0,\qquad
\mu_{0,P^\circ}(x):=0,\qquad
e_{P^\circ}(x):=\frac12
\quad (x\in\mathbb R^d).
\]
% lean: CausalSmith.Stat.WeakOverlap.dornCanonicalLaw_covariate

\item The Gaussian fourth absolute moment is finite. Define
\[
M_{\mathrm{mom}}
:=
\max\left\{
1,\left|\int_{\mathbb R}|y|^4\,\Pi_Y(dy)\right|
\right\}.
\]
Then
\[
1\le M_{\mathrm{mom}},
\qquad
\int_{\mathbb R}|y|^4\,\Pi_Y(dy)
\le
\left|\int_{\mathbb R}|y|^4\,\Pi_Y(dy)\right|
\le M_{\mathrm{mom}}.
\]
Monotonicity of powers in the exponent for a base at least one yields
\[
M_{\mathrm{mom}}
=M_{\mathrm{mom}}^1
\le M_{\mathrm{mom}}^4.
\]
Hence
\[
\int_{\mathbb R}|y|^4\,\Pi_Y(dy)
\le M_{\mathrm{mom}}^4.
\]
Choose the common constants
\[
q:=4,\qquad M:=M_{\mathrm{mom}},\qquad C:=2^{\gamma-1}.
\]
They satisfy \(q>3\), \(M>0\), and \(C>0\). % lean: CausalSmith.Stat.WeakOverlap.arbitraryFamily_lower_obstruction

\item We verify the remaining model restrictions with these constants. The product law gives
\[
\mathcal L_{P^\circ}(X,A)=\Pi_X\otimes\Pi_A.
\]
This is the joint law obtained from \(\Pi_X\) and the constant conditional distribution \(\Pi_A\). Uniqueness of conditional distributions therefore gives
\[
\mathcal L_{P^\circ}(A\mid X=x)=\Pi_A,
\qquad
P^\circ(A=1\mid X=x)=\frac12
\quad\text{for }P_X^\circ\text{-almost every }x.
\]
Thus the selected propensity is a measurable version, with
\[
0<e_{P^\circ}(x)=\frac12<1
\quad (x\in\mathbb R^d).
\]

Similarly, reassociating the product law gives
\[
\mathcal L_{P^\circ}((X,A),Y)
=(\Pi_X\otimes\Pi_A)\otimes\Pi_Y.
\]
Uniqueness of conditional distributions identifies the conditional outcome law with \(\Pi_Y\) for \((\Pi_X\otimes\Pi_A)\)-almost every \((x,a)\). Fubini's theorem gives this identity for \(\Pi_A\)-almost every \(a\), for \(\Pi_X\)-almost every \(x\). Since
\[
\Pi_A(\{0\})=\Pi_A(\{1\})=\frac12>0,
\]
both arms satisfy, simultaneously,
\[
\mathcal L_{P^\circ}(Y\mid X=x,A=a)=\Pi_Y
\quad(a\in\{0,1\}),
\qquad
P_X^\circ\text{-almost everywhere}.
\]
In particular, the treated conditional absolute fourth moment is integrable and
\[
\mathbb E_{P^\circ}\!\left[|Y|^4\mid X=x,A=1\right]
=
\int_{\mathbb R}|y|^4\,\Pi_Y(dy)
\le M^4
\quad\text{for }P_X^\circ\text{-almost every }x.
\]
The regression variance satisfies
\[
\int_{\mathbb R^d}
\left(
\mu_{P^\circ}(x)
-\int_{\mathbb R^d}\mu_{P^\circ}(y)\,P_X^\circ(dy)
\right)^2P_X^\circ(dx)
=0\le M.
\]

Both selected regressions have continuous derivatives of every order, and for every multi-index \(\alpha\) and every \(x,y\in\mathcal S\),
\[
D^\alpha\mu_{P^\circ}(x)
=D^\alpha\mu_{0,P^\circ}(x)=0,
\]
\[
|D^\alpha\mu_{P^\circ}(x)-D^\alpha\mu_{P^\circ}(y)|
=
|D^\alpha\mu_{0,P^\circ}(x)-D^\alpha\mu_{0,P^\circ}(y)|
=0.
\]
Their continuous boundary traces are also zero. Every defining top-derivative difference therefore vanishes, while its prescribed modulus bound is nonnegative. Thus
\[
\mu_{P^\circ},\mu_{0,P^\circ}\in\Sigma^{\mathrm D}(\beta,L_0).
\]

For the global overlap bound, fix \(t\in[0,1]\). If \(t\ge1/2\), monotonicity of the power with positive exponent \(\gamma-1\) gives
\[
(1/2)^{\gamma-1}\le t^{\gamma-1},
\qquad
2^{\gamma-1}(1/2)^{\gamma-1}=1.
\]
The propensity event is then the whole covariate space, so
\[
P_X^\circ\{e_{P^\circ}(X)\le t\}
=1
\le 2^{\gamma-1}t^{\gamma-1}
=Ct^{\gamma-1}.
\]
If \(t<1/2\), that event is empty, and
\[
P_X^\circ\{e_{P^\circ}(X)\le t\}
=0\le Ct^{\gamma-1}.
\]
Finally, the conditional outcome identities already established give Gaussian distributions in both arms with their selected zero means and variance one:
\[
\mathcal L_{P^\circ}(Y\mid X=x,A=1)
=N(\mu_{P^\circ}(x),1),
\qquad
\mathcal L_{P^\circ}(Y\mid X=x,A=0)
=N(\mu_{0,P^\circ}(x),1)
\]
for \(P_X^\circ\)-almost every \(x\). % lean: CausalSmith.Stat.WeakOverlap.arbitraryFamily_lower_obstruction

\item Define the selected singleton family and common local constants by
\[
\mathfrak P^\circ:=\{P^\circ\},
\qquad
\rho:=1,\qquad \nu:=\frac12,\qquad h_0:=1.
\]
The family is nonempty, its selected propensity has the required properties, and
\[
P_X^\circ(\mathcal S)=2^{-d}\lambda_d(\mathcal S)=1.
\]
We verify the local ratio in \cref{obj:def:dorn-a3}. For \(x_0\in\mathcal S\) and \(h>0\), define
\[
\mathcal B_h(x_0)
:=
\left\{x\in\mathbb R^d:
\sum_{i=1}^d(x_i-x_{0,i})^2\le h^2\right\},
\qquad
\mathcal O_h(x_0)
:=
\left\{x\in\mathbb R^d:
\sum_{i=1}^d(x_i-x_{0,i})^2<h^2\right\}.
\]
Since \(x_0\in\mathcal B_h(x_0)\cap\mathcal S\), this intersection is nonempty, and the constant selected propensity has supremum
\[
\sup_{x\in\mathcal B_h(x_0)\cap\mathcal S}e_{P^\circ}(x)
=\frac12.
\]

The denominator in the local ratio is positive, including when \(x_0\) lies on the cube boundary. Indeed, the coordinatewise closure identity
\[
\overline{(-1,1)}=[-1,1]
\]
and the finite-product closure and interior identities give
\[
\mathcal S\subseteq\overline{\operatorname{int}(\mathcal S)}.
\]
The set \(\mathcal O_h(x_0)\) is open and contains \(x_0\), because \(0<h^2\). It therefore meets \(\operatorname{int}(\mathcal S)\). Their intersection is a nonempty open set, so its Lebesgue measure is positive. Since
\[
\mathcal O_h(x_0)\cap\operatorname{int}(\mathcal S)
\subseteq
\mathcal B_h(x_0)\cap\mathcal S,
\]
we obtain
\[
0<
\lambda_d\!\left(
\mathcal O_h(x_0)\cap\operatorname{int}(\mathcal S)
\right)
\le
\lambda_d\!\left(\mathcal B_h(x_0)\cap\mathcal S\right).
\]
Consequently,
\[
0<
P_X^\circ(\mathcal B_h(x_0))
=
2^{-d}\lambda_d\!\left(\mathcal B_h(x_0)\cap\mathcal S\right)
\le1.
\]
Because the propensity is identically \(1/2\) on all of \(\mathbb R^d\),
\[
\left\{x:
e_{P^\circ}(x)\ge
\rho\sup_{y\in\mathcal B_h(x_0)\cap\mathcal S}e_{P^\circ}(y)
\right\}
\cap\mathcal B_h(x_0)
=
\mathcal B_h(x_0).
\]
Thus, for every \(x_0\in\mathcal S\) and \(0<h\le h_0\),
\[
\frac{
P_X^\circ\!\left(
\left\{x:
e_{P^\circ}(x)\ge
\rho\sup_{y\in\mathcal B_h(x_0)\cap\mathcal S}e_{P^\circ}(y)
\right\}
\cap\mathcal B_h(x_0)
\right)}
{P_X^\circ(\mathcal B_h(x_0))}
=1>\frac12=\nu.
\]
This proves
\[
\mathsf A_3^{\mathrm{Dorn}}
\bigl(\mathfrak P^\circ,(e_P)_{P\in\mathfrak P^\circ}\bigr).
\]
% lean: CausalSmith.Stat.WeakOverlap.dornCube_ball2_volume_pos

\item The singleton has the common known treated regression
\[
\mu_P(x)=0
\quad(P\in\mathfrak P^\circ,\ x\in\mathbb R^d).
\]
Applying the constant-estimator argument above with this regression and \(x_0=0\) gives, for every \(n\in\mathbb N\) and every \(c,r_n>0\),
\[
\inf_{T\ \mathrm{measurable}}
\sup_{P\in\mathfrak P^\circ}
P^{\otimes n}
\left\{
|T(\mathcal D_n)-\mu_P(0)|>c r_n
\right\}
=0.
\]
% lean: CausalSmith.Stat.WeakOverlap.arbitraryFamily_lower_obstruction
\end{enumerate}
\end{proof}

The decision underlying \cref{obj:lem:arbitrary-family-lower-obstruction} is the constant sample estimator \(T(\mathcal D_n)=\mu^\circ(x_0)\). Its pointwise error is identically zero throughout the family:
\[
|T(\mathcal D_n)-\mu_P(x_0)|=0
\qquad(P\in\mathfrak P).
\]
The Gaussian singleton shows that the obstruction is compatible with the source moment, smoothness, overlap, and local anti-concentration restrictions.

\subsection{The source upper comparison}

The cited upper comparison concerns an estimator selected for an already fixed family and its common constants. Its probability bound is uniform over members of that family, with the comparison scale determined by the family's fixed overlap exponent. The following proposition retains that quantifier order; its common-regression clause is the application of \cref{obj:lem:arbitrary-family-lower-obstruction} needed to delimit the converse scope.

\begin{propositionv}{lem:dorn-rate-scope}[Fixed-family regression rate scope]*
Fix an integer \(d\ge1\), constants \(\beta,L_0,M,\sigma,C>0\), and \(q>3\). For every fixed finite \(\gamma_0>1\), consider a nonempty source family \(\mathfrak P\) of observed-data laws, equipped with selected treated and control regression curves and selected propensity representatives \((e_P)_{P\in\mathfrak P}\). Suppose that:
\begin{itemize}
\item \textbf{(Design and propensity.)} Under every \(P\in\mathfrak P\), \(X\) is uniform on \([-1,1]^d\), and the selected measurable representative satisfies
\[
e_P(X)=P(A=1\mid X)\in(0,1)
\qquad P\text{-almost surely}.
\]
\item \textbf{(Moment and variance.)} Writing \(\mu_P\) for the selected treated regression curve, every member satisfies
\[
\mathbb E_P\!\left[|Y|^q\mid X,A=1\right]\le M^q,\qquad
\operatorname{Var}_P\!\left(\mu_P(X)\right)\le M,\qquad
\operatorname{Var}_P(Y\mid X,A)\ge\sigma^2,
\]
with the conditional inequalities holding almost surely in each applicable arm.
\item \textbf{(Gaussian regression and smoothness.)} Almost surely with respect to the covariate marginal, the conditional outcome distributions in the treated and control arms are Gaussian with their respective selected regression means and common variance \(\sigma^2\). Both selected regression curves belong to the source seminorm class \(\Sigma^{\mathrm D}(\beta,L_0)\): their derivatives of order
\[
m=\max\{k\in\mathbb Z:k<\beta\}
\]
satisfy Dorn's Hölder-seminorm bound \(L_0\), with exponent \(\beta-m\).
\item \textbf{(Global overlap tail.)} For every \(P\in\mathfrak P\) and \(t\in[0,1]\),
\[
P\{e_P(X)\le t\}\le Ct^{\gamma_0-1}.
\]
\item \textbf{(Family-level local condition.)} On the original cube \([-1,1]^d\), the selected propensity representatives satisfy
\[
\mathsf A_{3}^{\mathrm{Dorn}}
\bigl(\mathfrak P,(e_P)_{P\in\mathfrak P}\bigr)
\]
as defined in \cref{obj:def:dorn-a3}, including its common positive constants \(\rho,\nu,h_0\) and uniform local-ball mass-ratio inequality.
\end{itemize}

Define the fixed-setup comparison scale by
\[
r^*_{\mu,n}
:=n^{-\beta/(2\beta+d\gamma_0/(\gamma_0-1))}.
\]
There exists a regression estimator \(\widehat\mu_n(\mathcal D_n,x)\), where \(\mathcal D_n\) is the sample of \(n\) independent observations, selected after the family, its common constants, and its propensity representatives have been fixed, such that its dependence on the sample is measurable for every \(n\) and \(x\). For every sample size \(n\), member \(P\in\mathfrak P\), and real threshold \(R\), the event
\[
\left\{
\|\widehat\mu_n-\mu_P\|_{L^\infty(P)}
>\max\{Rr^*_{\mu,n},0\}
\right\}
\]
is measurable, where the norm is the essential supremum with respect to the covariate marginal of \(P\). For every \(\varepsilon>0\), there exists a finite constant \(c(\varepsilon)>0\) such that
\[
\limsup_{n\to\infty}\sup_{P\in\mathfrak P}
P^{\otimes n}\!\left\{
\|\widehat\mu_n-\mu_P\|_{L^\infty(P)}
>c(\varepsilon)r^*_{\mu,n}
\right\}
\le\varepsilon,
\]
where \(P^{\otimes n}\) is the product sampling law. The estimator and \(c(\varepsilon)\) may depend on the fixed setup. This upper comparison is the assertion of \citet[Section 2, Assumptions 1--3, Theorem 1(ii)]{Dorn2026GlobalOverlap}.

For every \(\gamma_0>1\), the following two pointwise conclusions also hold. First, let \(\mathfrak P\) be any nonempty source family whose selected treated regression curves equal a common known function \(\mu^\circ\) at every point:
\[
\mu_P(x)=\mu^\circ(x)
\qquad(P\in\mathfrak P,\ x\in\mathbb R^d).
\]
For every evaluation point \(x_0\in\mathbb R^d\), sample size \(n\), and constants \(c,r>0\), its extended nonnegative minimax miss risk is
\[
\inf_{T\ \mathrm{sample\text{-}measurable}}
\sup_{P\in\mathfrak P}
P^{\otimes n}\!\left\{
|T(\mathcal D_n)-\mu_P(x_0)|>cr
\right\}
=0.
\]
Here the infimum ranges over all measurable real-valued functions of the sample.

Second, there exist constants \(q'>3\) and \(M',C'>0\) for which the singleton Gaussian experiment
\[
X\sim\operatorname{Unif}([-1,1]^d),\qquad
A\sim\operatorname{Bernoulli}(1/2),\qquad
Y\sim N(0,1),
\]
with \(X,A,Y\) mutually independent and selected treated mean, control mean, and propensity respectively \(0,0,1/2\), satisfies the displayed design, propensity, moment, variance, Gaussian, smoothness, and global-tail conditions with parameter tuple
\[
(\beta,q',M',1,C',L_0,\gamma_0),
\]
and satisfies the family-level condition of \cref{obj:def:dorn-a3} on the original cube. For this singleton, at \(x_0=0\), every sample size \(n\) and every \(c,r>0\) satisfy
\[
\inf_{T\ \mathrm{sample\text{-}measurable}}
P^{\otimes n}\!\left\{|T(\mathcal D_n)|>cr\right\}
=0.
\]
\end{propositionv}
% CAUSALSMITH-CITED-SCOPE-BEGIN lem:dorn-rate-scope
\verificationfootnotetext{\textbf{Formalization scope.} This result uses the cited conclusion \citep{Dorn2026GlobalOverlap} (Section 2 Notation and Setting, Assumptions 1--3, and Theorem 1 (author PDF pp. 3--7 and 12); Lemma 11 (p. 23); Lemma 13 (p. 24); proof of Theorem 1 (pp. 27--28); Section 5 Conclusion (p. 28)). The cited source proof is not formalized here: Lean verifies its use and all remaining steps. If that cited conclusion is false, the portions of this result that depend on it are not certified.}
% CAUSALSMITH-CITED-SCOPE-END lem:dorn-rate-scope

The source upper assertion in \cref{obj:lem:dorn-rate-scope} is the fixed-setup, uniform-in-\(P\) probability guarantee of \citet[Theorem 1(ii), author PDF p.~12; proof of Theorem 1, pp.~27--28]{Dorn2026GlobalOverlap}. Its displayed scale has no logarithmic factor. The estimator, threshold constant, and overlap exponent remain indexed by the fixed setup, while sample measurability and exceedance-event measurability give the displayed probabilities their ordinary sampling interpretation.

The common-regression clauses give the model-family qualification relevant to the pointwise lower attribution printed in \citet[Theorem 1(i), author PDF p.~12]{Dorn2026GlobalOverlap}. The positive pointwise converse in \cref{obj:prop:dorn-a3-free-corollary} instead concerns the full causal class of \cref{obj:def:risks}, with expected absolute-error loss; its lower bound is supported by the two-valued potential-outcome subclass in \cref{obj:def:bounded-lower-pair,obj:lem:bounded-lower-pair}.

\subsection{The affine transfer}

The transfer component of \cref{obj:prop:dorn-a3-free-corollary} uses the retained source restrictions from \citet[Section 2, Assumptions 1--2, author PDF pp.~3--7]{Dorn2026GlobalOverlap}. The maximal envelopes \(\mathcal D_\gamma\) and the constants \(B_*,L_*\) are defined in that statement. Their role is to place every qualifying source tuple within one response-recovery class after transporting covariates by \(T^{-1}\), where \(T(x)=2x-\mathbf1\).

The ingredients have distinct mathematical functions. Uniform source covariates become uniform covariates on the unit cube, supplying density constant \(c_f=1\). The retained propensity tail controls global overlap at the same exponent. The treated-moment and Gaussian restrictions supply the common response bound \(B_*=\max\{M,\sigma_{\min}\}\), and \cref{obj:lem:fixed-cube-holder-completion} supplies the finite full Hölder radius \(L_*\). The causal-completion clause of \cref{obj:prop:dorn-a3-free-corollary} then identifies the transported observed laws as members of the model in \cref{obj:def:model}.

With this membership established, the equal-cell estimator of \cref{obj:def:estimator} has the expected spatial-supremum guarantee stated in \cref{obj:prop:dorn-a3-free-corollary}, uniformly over the fixed-constant envelopes and the compact overlap range. At \(\gamma=\gamma_0\), its oracle scale \(r_{\gamma,n}\) equals the source scale \(r^*_{\mu,n}\). The accompanying essential-supremum probability bound supplies a comparison in the source loss criterion, with constants common to those envelopes and that exponent range. Every qualifying source subfamily inherits the envelope guarantee through the containment asserted in \cref{obj:prop:dorn-a3-free-corollary}.

\subsection{The Gaussian strictness construction}

The Gaussian existence component of \cref{obj:prop:dorn-a3-free-corollary} concerns the geometric scope of the retained restrictions. For each overlap exponent in the stated range and dimension \(d\ge2\), it supplies a source tuple with uniform covariates, Gaussian conditional outcomes, smooth means, and the required moment and global-tail bounds, together with failure of the selected-version local predicate in \cref{obj:def:dorn-a3}. Its positive finite constants are selected separately for each exponent.

This quantifier structure distinguishes the existence construction from the upper guarantee over a prescribed fixed-constant envelope. The construction establishes that the retained Gaussian restrictions accommodate the selected-propensity geometry described by the corollary. The affine-transfer component supplies recovery guarantees for each envelope with its displayed common constants, while \cref{obj:lem:dorn-rate-scope} supplies the preprint comparator for fixed families satisfying the additional local condition. Together, these components identify the smoothness completion, family scope, and treatment geometry underlying \cref{obj:prop:dorn-a3-free-corollary}.


\section{Verification note}

The theorem and lemma environments carrying declaration links are machine-checked in Lean 4.33.0 (compiler commit \texttt{d8b18978322de05a8f3dba51ef03cf5461676c17}) against mathlib revision \texttt{db584cd6d46c92f209a44c0f1c829460d327499d} and the adjacent Causalean library. The Lean sources are under \texttt{CausalSmith/Stat/STAT\_WeakoverlapUnisolventRate\_Research}; the bundle's \texttt{lean\_snippets.json} and \texttt{formal\_layer.json} record the complete object-to-declaration map. The verification record identifies repository revision \texttt{d07574d7eb0a903c2a049c6ae28702de46193ec0}; this is distinct from the Lean compiler commit and the mathlib revision above. The emitted paper checksum is recorded separately in \texttt{meta.json}.

The checked results include causal identification, ordered treated-cell masses, template conditioning, and mesh selection in \cref{obj:lem:causal-identification,obj:lem:ordered-mass,obj:lem:template-conditioning,obj:lem:selected-mesh}; the bounded testing experiment in \cref{obj:lem:bounded-lower-pair}; and the adaptive minimax conclusion in \cref{obj:thm:adaptive-minimax}. They also include the treatment-geometry and Gaussian-transfer conclusions in \cref{obj:prop:strict-enlargement,obj:prop:dorn-a3-free-corollary}, together with smoothness completion, the family-richness obstruction, and the source comparison scope in \cref{obj:lem:fixed-cube-holder-completion,obj:lem:arbitrary-family-lower-obstruction,obj:lem:dorn-rate-scope}. The synthesized notation environments and explanatory transitions are presentation-level material. The joint-adaptation entries in \cref{obj:def:joint-adaptation-handle,obj:oeq:joint-adaptation} describe future work and assert no proved risk bound.

These checks certify deductions under the stated assumptions. Response recovery uses the sampling, causal, density, global-tail, outcome, and smoothness restrictions defining \cref{obj:def:model}, with the fixed constants and compact overlap range specified in \cref{obj:thm:adaptive-minimax}. The geometry and transfer conclusions retain their stated dimension restrictions, selected propensity versions, and source-family conditions. External-source transcriptions and generated theorem-local dependency footnotes delimit the trust boundary for cited inputs: the checked deductions use the recorded source conclusions under their recorded assumptions, while the source proofs of those conclusions remain external. The consolidated account here complements those local disclosures.


\section{Proofs of the main results}\label{sec:deferred-proofs}

\begin{proof}[Proof of \cref{obj:thm:adaptive-minimax}]
\emph{1. Uniform local approximation.}
We first establish the estimator bound uniformly over the prescribed adaptation interval. Set
\[
m:=\lceil\beta\rceil-1,\qquad
s:=\beta-m,\qquad
\mathcal A_{\mathrm{poly}}
:=\{\alpha\in\mathbb N_0^d:|\alpha|\le m\},
\qquad
p_{\mathrm{coef}}:=|\mathcal A_{\mathrm{poly}}|.
\]
Then \(0<s\le1\) and \(p_{\mathrm{coef}}>0\). By the identification and uniqueness conclusions of \cref{obj:lem:causal-identification}, every Hölder response witnessing membership of \(P\) in \(\mathcal P_\gamma\) agrees with \(\mu_{1,P}\) throughout \(\mathcal X\).

Here is the local approximation bound used below, including at cube boundaries. Define the affine map and transported response by
\[
T_{\mathrm{aff}}(x):=2x-\mathbf1
\quad(x\in\mathbb R^d),
\qquad
g_{\mathrm{aff}}(z)
:=\mu_{1,P}\!\left(\frac{z+\mathbf1}{2}\right)
\quad(z\in[-1,1]^d).
\]
Repeated coordinate differentiation of the affine composition gives, for every multiindex \(\alpha\) with \(|\alpha|\le m\),
\[
D^\alpha g_{\mathrm{aff}}(z)
=
2^{-|\alpha|}D^\alpha\mu_{1,P}\!\left(\frac{z+\mathbf1}{2}\right),
\qquad z\in[-1,1]^d.
\]
The chain rule proves this identity in the interior, and continuity of the derivative traces extends it to the boundary. Thus \cref{obj:ass:holder-derivative-bound} gives
\[
\sup_{z\in[-1,1]^d}|D^\alpha g_{\mathrm{aff}}(z)|
\le 2^{-|\alpha|}L\le L.
\]
For \(|\alpha|=m\) and \(z,z'\in[-1,1]^d\), the identity and \cref{obj:ass:holder-modulus} give
\[
\begin{aligned}
|D^\alpha g_{\mathrm{aff}}(z)-D^\alpha g_{\mathrm{aff}}(z')|
&=2^{-m}\left|
D^\alpha\mu_{1,P}\!\left(\frac{z+\mathbf1}{2}\right)
-D^\alpha\mu_{1,P}\!\left(\frac{z'+\mathbf1}{2}\right)
\right|
\\
&\le 2^{-m}L
\left\|\frac{z-z'}2\right\|_\infty^s
\\
&\le L\|z-z'\|_\infty^s.
\end{aligned}
\]
Affine composition also preserves the continuous derivative traces through order \(m\). Consequently, defining
\[
K_{\mathrm{aff}}:=1,
\]
we obtain intrinsic Hölder radius \(K_{\mathrm{aff}}L\) for \(g_{\mathrm{aff}}\) on \([-1,1]^d\). \Cref{obj:lem:global-holder-extension}, applied with derivative order \(m\), exponent \(s\), and radius \(K_{\mathrm{aff}}L\), supplies a constant \(K_{\mathrm{ext}}>0\) and a function \(g_{\mathrm{ext}}\in C^m(\mathbb R^d)\) satisfying
\[
g_{\mathrm{ext}}(z)=g_{\mathrm{aff}}(z)
\quad(z\in[-1,1]^d),
\]
\[
\max_{0\le k_{\mathrm{der}}\le m}
\sup_{z\in\mathbb R^d}
\|\mathrm D^{k_{\mathrm{der}}}g_{\mathrm{ext}}(z)\|_{\mathrm{op}}
\le K_{\mathrm{ext}}K_{\mathrm{aff}}L,
\]
and
\[
\|\mathrm D^m g_{\mathrm{ext}}(z)
-\mathrm D^m g_{\mathrm{ext}}(z')\|_{\mathrm{op}}
\le
K_{\mathrm{ext}}K_{\mathrm{aff}}L
\|z-z'\|_\infty^s
\quad(z,z'\in\mathbb R^d).
\]
Here \(\mathrm D^{k_{\mathrm{der}}}\) denotes the Fréchet derivative of the indicated order, with order zero interpreted as the function itself.

For a partition cube \(Q\) of mesh \(0<h\le1\), define
\[
\vartheta_{Q,\alpha}
:=
\frac{(2h)^{|\alpha|}}{\alpha!}
D^\alpha g_{\mathrm{ext}}(T_{\mathrm{aff}}(a_Q)),
\qquad \alpha\in\mathcal A_{\mathrm{poly}},
\]
and
\[
\Pi_Q(x)
:=
U\!\left(\frac{x-a_Q}{h}\right)^\top\vartheta_Q,
\qquad x\in Q.
\]
We derive the remainder bound along the expansion segment. Since \(\beta>1\), we have \(m\ge1\). For \(x\in Q\), define
\[
z_{Q,0}:=T_{\mathrm{aff}}(a_Q),
\qquad
v_{Q,x}:=2(x-a_Q),
\qquad
\varphi_{Q,x}(t_{\mathrm{seg}})
:=g_{\mathrm{ext}}(z_{Q,0}+t_{\mathrm{seg}}v_{Q,x})
\quad(t_{\mathrm{seg}}\in\mathbb R).
\]
The chain rule gives
\[
\varphi_{Q,x}^{(k_{\mathrm{der}})}(t_{\mathrm{seg}})
=
\mathrm D^{k_{\mathrm{der}}}g_{\mathrm{ext}}
(z_{Q,0}+t_{\mathrm{seg}}v_{Q,x})
[\underbrace{v_{Q,x},\ldots,v_{Q,x}}_{k_{\mathrm{der}}\text{ arguments}}],
\qquad
0\le k_{\mathrm{der}}\le m,
\quad t_{\mathrm{seg}}\in\mathbb R,
\]
with order zero interpreted as the function value. Expanding the symmetric derivatives in coordinate directions and using the definition of \(\vartheta_Q\) yields
\[
\sum_{k_{\mathrm{der}}=0}^{m}
\frac{\varphi_{Q,x}^{(k_{\mathrm{der}})}(0)}{k_{\mathrm{der}}!}
=
\sum_{\alpha\in\mathcal A_{\mathrm{poly}}}
\frac{D^\alpha g_{\mathrm{ext}}(z_{Q,0})}{\alpha!}
\prod_{i=1}^{d}(v_{Q,x,i})^{\alpha_i}
=\Pi_Q(x),
\qquad
\varphi_{Q,x}(1)=\mu_{1,P}(x).
\]
Taylor's theorem with the order-\(m\) Lagrange remainder gives a point satisfying
\[
\exists\,t_{Q,x}\in(0,1):\qquad
\varphi_{Q,x}(1)
-
\sum_{k_{\mathrm{der}}=0}^{m-1}
\frac{\varphi_{Q,x}^{(k_{\mathrm{der}})}(0)}{k_{\mathrm{der}}!}
=
\frac{\varphi_{Q,x}^{(m)}(t_{Q,x})}{m!}.
\]
Choose such a point and define the degree-\(m\) remainder by
\[
\begin{aligned}
R_{\mathrm{Taylor},Q,x}
&:=\mu_{1,P}(x)-\Pi_Q(x)\\
&=\frac{\varphi_{Q,x}^{(m)}(t_{Q,x})
-\varphi_{Q,x}^{(m)}(0)}{m!}\\
&=\frac1{m!}
\left(
\mathrm D^m g_{\mathrm{ext}}(z_{Q,0}+t_{Q,x}v_{Q,x})
-\mathrm D^m g_{\mathrm{ext}}(z_{Q,0})
\right)
[\underbrace{v_{Q,x},\ldots,v_{Q,x}}_{m\text{ arguments}}].
\end{aligned}
\]
The displayed Hölder modulus for the extended function therefore implies
\[
\begin{aligned}
|R_{\mathrm{Taylor},Q,x}|
&\le
\frac{K_{\mathrm{ext}}K_{\mathrm{aff}}L}{m!}
\|t_{Q,x}v_{Q,x}\|_\infty^s
\|v_{Q,x}\|_\infty^m\\
&=
\frac{K_{\mathrm{ext}}K_{\mathrm{aff}}L}{m!}
 t_{Q,x}^{s}\|v_{Q,x}\|_\infty^{m+s}\\
&\le
\frac{K_{\mathrm{ext}}K_{\mathrm{aff}}L}{m!}
\|v_{Q,x}\|_\infty^\beta,
\end{aligned}
\]
since \(0<t_{Q,x}<1\), \(s>0\), and \(m+s=\beta\). This bound also applies when \(v_{Q,x}=0\). Define
\[
C_{\mathrm{Taylor}}:=\frac1{m!}\ge0.
\]
Because \(\|v_{Q,x}\|_\infty\le2h\), we obtain, uniformly in the expansion centre,
\[
\begin{aligned}
|\mu_{1,P}(x)-\Pi_Q(x)|
&=
\left|
g_{\mathrm{ext}}\!\left(
T_{\mathrm{aff}}(a_Q)+2h\frac{x-a_Q}{h}
\right)
-
U\!\left(\frac{x-a_Q}{h}\right)^\top\vartheta_Q
\right|
\\
&\le
C_{\mathrm{Taylor}}K_{\mathrm{ext}}K_{\mathrm{aff}}L(2h)^\beta
=
C_{\mathrm{bias}}Lh^\beta,
\end{aligned}
\]
where
\[
C_{\mathrm{bias}}
:=
C_{\mathrm{Taylor}}K_{\mathrm{ext}}K_{\mathrm{aff}}2^\beta.
\]
The normalized coordinates belong to \([0,1]^d\), and the affine image of \(x\) belongs to \([-1,1]^d\), which justify the equality and the approximation. This construction also covers integer \(\beta\), for which \(s=1\).
% lean: CausalSmith.Stat.WeakOverlap.exists_uniform_holderBall_local_monoVec_approx

\emph{2. Conditional bias and supremum loss.}
Apply \cref{obj:lem:selected-mesh}. Denote its uniform mesh and coefficient constants by \(K_{\mathrm{mesh}},K_{\mathrm{coef}}>0\), and its sample-size threshold by \(n_{\mathrm{env}}\). Enlarge these constants and the threshold as required by the constructions below, taking \(n_{\mathrm{env}}\ge2\). Define
\[
n_0:=\max\{n_{\mathrm{env}},1\},
\qquad
\mathcal G_n:=\sigma((X_i,A_i)_{i=1}^n).
\]
For \(n\ge n_0\), \(\gamma\in\Gamma\), and \(P\in\mathcal P_\gamma\), choose the simultaneous treated-count event underlying the mesh and coefficient bounds as follows. For all candidate microcells, define
\[
q_{Q\ell}:=P\{A=1,\ X\in E_{Q\ell}\},
\qquad h\in\mathcal H_n,\quad Q\in\mathcal Q_h,\quad 1\le\ell\le J,
\]
\[
M_{\mathrm{count},n}
:=\sum_{h\in\mathcal H_n}J|\mathcal Q_h|,
\qquad
s_{\mathrm{count},n}
:=\log\!\left(2(M_{\mathrm{count},n}+1)n^2\right)>0,
\]
and take
\[
\mathcal E_n
:=
\bigcap_{h\in\mathcal H_n}
\bigcap_{Q\in\mathcal Q_h}
\bigcap_{\ell=1}^J
\left\{
\left|N_{Q\ell}-nq_{Q\ell}\right|
<\frac{nq_{Q\ell}}2+4s_{\mathrm{count},n}
\right\}.
\]
The Bernstein bound for each treated count and the finite union bound give
\[
\begin{aligned}
P^{\otimes n}(\mathcal E_n^c)
&\le 2M_{\mathrm{count},n}e^{-s_{\mathrm{count},n}}
\\
&=\frac{M_{\mathrm{count},n}}{M_{\mathrm{count},n}+1}n^{-2}
\le n^{-2}.
\end{aligned}
\]
We establish the selected-mesh and treated-count bounds on this particular intersection. Let \(\eta\) be the reference microcube width in \cref{obj:def:template}. Each microcell at mesh \(h\) has volume \((\eta h)^d\), so \cref{obj:ass:density} and the measurable-set bound in \cref{obj:lem:ordered-mass} give
\[
\begin{aligned}
P(X\in E_{Q\ell})&\ge c_f(\eta h)^d,\\
q_{Q\ell}
&\ge \frac{\gamma-1}{\gamma}C^{-1/(\gamma-1)}
\bigl(c_f(\eta h)^d\bigr)^{\gamma/(\gamma-1)}
=\kappa_{\mathrm{cell}}(\gamma)h^{D_\gamma},
\end{aligned}
\]
where
\[
\kappa_{\mathrm{cell}}(\gamma)
:=\frac{\gamma-1}{\gamma}C^{-1/(\gamma-1)}
(c_f\eta^d)^{\gamma/(\gamma-1)},
\qquad
\kappa_{\mathrm{cell},\Gamma}
:=\min_{\gamma'\in\Gamma}\kappa_{\mathrm{cell}}(\gamma')>0.
\]
The positivity of the minimum follows from continuity and positivity of the displayed coefficient on the compact interval \(\Gamma\). Choose a fixed enlargement constant satisfying
\[
A_{\mathrm{or}}\ge1,
\qquad
4\le\kappa_{\mathrm{cell},\Gamma}
A_{\mathrm{or}}^{2\beta+D_{\gamma_{\max}}}.
\]
Such a choice exists because \(2\beta+D_{\gamma_{\max}}>0\). Since
\[
D_\gamma=d+\frac d{\gamma-1},
\qquad
D_{\gamma_{\max}}\le D_\gamma\le D_{\gamma_{\min}},
\qquad
h_{\gamma_{\max},n}\le h_{\gamma,n}\le h_{\gamma_{\min},n},
\]
we have
\[
4\le\kappa_{\mathrm{cell},\Gamma}
A_{\mathrm{or}}^{2\beta+D_\gamma}
\quad(\gamma\in\Gamma).
\]

The candidate collection in \cref{obj:synth_3} has largest index
\[
j_{\mathrm{cand},n}
:=\left\lfloor\frac{\log_2 n}{2\beta+d}\right\rfloor,
\qquad
\mathcal H_n=\{2^{-j}:0\le j\le j_{\mathrm{cand},n}\}.
\]
The floor inequality gives
\[
2^{-j_{\mathrm{cand},n}}
\le2n^{-1/(2\beta+d)}.
\]
Because \(d>0\) and \(\gamma_{\max}>1\),
\[
D_{\gamma_{\max}}>d,
\qquad
\frac{2n^{-1/(2\beta+d)}}{A_{\mathrm{or}}h_{\gamma_{\max},n}}
=\frac2{A_{\mathrm{or}}}
 n^{-1/(2\beta+d)+1/(2\beta+D_{\gamma_{\max}})}
\longrightarrow0,
\qquad
A_{\mathrm{or}}h_{\gamma_{\min},n}\longrightarrow0.
\]
After enlarging \(n_{\mathrm{env}}\), these relations ensure, uniformly over \(\gamma\in\Gamma\), that
\[
2^{-j_{\mathrm{cand},n}}
\le A_{\mathrm{or}}h_{\gamma,n}\le1.
\]
Define an oracle candidate by
\[
j_{\mathrm{or}}(n,\gamma)
:=\max\{j\in\{0,\ldots,j_{\mathrm{cand},n}\}:
2^{-j}\ge A_{\mathrm{or}}h_{\gamma,n}\},
\qquad
h_{\mathrm{or},n,\gamma}:=2^{-j_{\mathrm{or}}(n,\gamma)}.
\]
The set defining the maximum contains zero. If its maximum is smaller than \(j_{\mathrm{cand},n}\), the next dyadic width is smaller than \(A_{\mathrm{or}}h_{\gamma,n}\); if the maximum equals \(j_{\mathrm{cand},n}\), the preceding lower bound on the finest width forces equality with the target width. Thus
\[
h_{\mathrm{or},n,\gamma}\in\mathcal H_n,
\qquad
A_{\mathrm{or}}h_{\gamma,n}
\le h_{\mathrm{or},n,\gamma}
\le2A_{\mathrm{or}}h_{\gamma,n}.
\]

We next verify the logarithmic budget. The template and dyadic partition have
\[
J=(m+1)^d,
\qquad
|\mathcal Q_{2^{-j}}|=2^{jd},
\qquad
M_{\mathrm{count},n}
=J\sum_{j=0}^{j_{\mathrm{cand},n}}2^{jd}.
\]
Since \(2\beta+d\ge1\), the definition of \(j_{\mathrm{cand},n}\) gives
\[
j_{\mathrm{cand},n}\le2^{j_{\mathrm{cand},n}}\le n,
\qquad
M_{\mathrm{count},n}\le(n+1)n^dJ.
\]
For \(n\ge2\), using \(n+1\le2n\) and \(n^{d+1}\ge1\), we obtain
\[
2(M_{\mathrm{count},n}+1)n^2
\le8(J+1)n^{d+3},
\qquad
s_{\mathrm{count},n}
\le(d+3)\log n+\log(8(J+1)).
\]
Define
\[
U_{\mathrm{count}}:=24(d+3),
\qquad
V_{\mathrm{count}}:=24\log(8(J+1)),
\qquad
\rho_{\mathrm{budget}}
:=\frac{2\beta}{2\beta+D_{\gamma_{\min}}}>0.
\]
The elementary bound
\[
\log n\le\frac2{\rho_{\mathrm{budget}}}
 n^{\rho_{\mathrm{budget}}/2}
\]
and divergence of positive powers allow a further uniform enlargement of \(n_{\mathrm{env}}\) such that
\[
2U_{\mathrm{count}}\frac2{\rho_{\mathrm{budget}}}
\le(2A_{\mathrm{or}})^{-2\beta}
 n^{\rho_{\mathrm{budget}}/2},
\qquad
2V_{\mathrm{count}}
\le(2A_{\mathrm{or}})^{-2\beta}n^{\rho_{\mathrm{budget}}}.
\]
Consequently,
\[
\begin{aligned}
24s_{\mathrm{count},n}
&\le U_{\mathrm{count}}\log n+V_{\mathrm{count}}\\
&\le(2A_{\mathrm{or}})^{-2\beta}n^{\rho_{\mathrm{budget}}}
=(2A_{\mathrm{or}}h_{\gamma_{\min},n})^{-2\beta}\\
&\le(2A_{\mathrm{or}}h_{\gamma,n})^{-2\beta}
\le h_{\mathrm{or},n,\gamma}^{-2\beta}.
\end{aligned}
\]
The population-mass bound also gives, for every microcell of this oracle candidate,
\[
\begin{aligned}
nq_{Q\ell}
&\ge n\kappa_{\mathrm{cell},\Gamma}
 h_{\mathrm{or},n,\gamma}^{D_\gamma}\\
&=\kappa_{\mathrm{cell},\Gamma}
 n h_{\mathrm{or},n,\gamma}^{2\beta+D_\gamma}
 h_{\mathrm{or},n,\gamma}^{-2\beta}\\
&\ge\kappa_{\mathrm{cell},\Gamma}
 A_{\mathrm{or}}^{2\beta+D_\gamma}
 h_{\mathrm{or},n,\gamma}^{-2\beta}
\ge4h_{\mathrm{or},n,\gamma}^{-2\beta},
\end{aligned}
\]
where \(nh_{\gamma,n}^{2\beta+D_\gamma}=1\).

On \(\mathcal E_n\), the defining deviation inequality therefore yields, at the oracle candidate,
\[
N_{Q\ell}
>\frac{nq_{Q\ell}}2-4s_{\mathrm{count},n}
\ge2h_{\mathrm{or},n,\gamma}^{-2\beta}
 -4s_{\mathrm{count},n}
\ge h_{\mathrm{or},n,\gamma}^{-2\beta}.
\]
Thus this candidate is feasible. The smallest-feasible-mesh rule of \cref{obj:def:mesh-selector} gives
\[
0<\widehat h\le h_{\mathrm{or},n,\gamma}
\le2A_{\mathrm{or}}h_{\gamma,n},
\qquad
24s_{\mathrm{count},n}
\le h_{\mathrm{or},n,\gamma}^{-2\beta}
\le\widehat h^{-2\beta}.
\]
For a microcell of the selected mesh, feasibility and the same deviation inequality imply
\[
N_{Q\ell}\ge\widehat h^{-2\beta}
\ge24s_{\mathrm{count},n},
\qquad
\frac{nq_{Q\ell}}2
<N_{Q\ell}+4s_{\mathrm{count},n}
\le\frac76N_{Q\ell}.
\]
In particular,
\[
N_{Q\ell}\ge\frac{nq_{Q\ell}}5
\quad(Q\in\mathcal Q_{\widehat h},\ 1\le\ell\le J).
\]
Taking \(K_{\mathrm{mesh}}\ge2A_{\mathrm{or}}\), we have established on the displayed count intersection that
\[
0<\widehat h\le K_{\mathrm{mesh}}h_{\gamma,n},
\qquad
N_{Q\ell}\ge\frac{nq_{Q\ell}}5.
\]
% lean: CausalSmith.Stat.WeakOverlap.selectedCountEvent_uniform_design_saturated
% lean: CausalSmith.Stat.WeakOverlap.selectedCount_deterministic_bridge
Every count in the defining finite intersection is measurable with respect to \(\mathcal G_n\), and every population mass is fixed under \(P\). Hence \(\mathcal E_n\in\mathcal G_n\). Conditioning therefore fixes the sampled design, the selected mesh, its treated counts, and membership in \(\mathcal E_n\).

First, the response bound holds everywhere on the cube. Define
\[
g_{\mathrm{clip}}(x)
:=
\operatorname{clip}_{[-B,B]}(\mu_{1,P}(x)),
\qquad x\in\mathcal X.
\]
The conditional-mean bound in \cref{obj:ass:bounded-potential-outcomes} and identification give
\[
g_{\mathrm{clip}}(x)
=
\mu_{1,P}(x)
=
\mathbb E_P[Y\mid A=1,X=x]
\quad\text{for }P_X\text{-almost every }x.
\]
Both functions are continuous on the cube. The uniqueness conclusion of \cref{obj:lem:causal-identification} therefore gives
\[
g_{\mathrm{clip}}(x)=\mu_{1,P}(x),
\qquad
|\mu_{1,P}(x)|\le B
\quad(x\in\mathcal X).
\]

On \(\mathcal E_n\), set
\[
h:=\widehat h,
\qquad
\bar c_Q:=\mathbb E[\widehat c_Q\mid\mathcal G_n],
\qquad Q\in\mathcal Q_h.
\]
At each fixed design, this conditional expectation evaluates the coefficients at that design's selected mesh. Feasibility gives positive treated counts. For each treated observation used in the fit on \(Q\), define
\[
U_{i,Q}:=U\!\left(\frac{X_i-a_Q}{h}\right),
\qquad
R_{\mathrm{res},i,Q}
:=\mu_{1,P}(X_i)-U_{i,Q}^{\top}\vartheta_Q.
\]
The local polynomial approximation gives
\[
|R_{\mathrm{res},i,Q}|\le C_{\mathrm{bias}}Lh^\beta.
\]
Identification and the equal-cell normal equations yield, almost surely,
\[
\bar c_Q-\vartheta_Q
=
\widehat G_Q^{-1}
\frac1J\sum_{\ell=1}^J\frac1{N_{Q\ell}}
\sum_{\substack{i:A_i=1\\X_i\in E_{Q\ell}}}
U_{i,Q}R_{\mathrm{res},i,Q}.
\]
Every monomial coordinate has absolute value at most one on the normalized cube, so
\[
\|U_{i,Q}\|_2\le\sqrt{p_{\mathrm{coef}}}.
\]
Consequently, \cref{obj:lem:template-conditioning} implies
\[
\|\bar c_Q-\vartheta_Q\|_\infty
\le
\|\bar c_Q-\vartheta_Q\|_2
\le
\frac{2\sqrt{p_{\mathrm{coef}}}}{\lambda_0}
C_{\mathrm{bias}}Lh^\beta.
\]
Define the fixed constants
\[
K_{\mathrm{inv}}
:=\frac{2\sqrt{p_{\mathrm{coef}}}}{\lambda_0},
\qquad
K_{\mathrm{bias}}
:=p_{\mathrm{coef}}K_{\mathrm{inv}}C_{\mathrm{bias}}L
+C_{\mathrm{bias}}L.
\]

For the selected design, define the maximum coefficient fluctuation by
\[
S_{\mathrm{coef}}
:=
\max_{Q\in\mathcal Q_h}
\|\widehat c_Q-\bar c_Q\|_\infty.
\]
This maximum uses the selected mesh fixed by the conditioning design. Since
\[
|U(z)^\top(\widehat c_Q-\vartheta_Q)|
\le
p_{\mathrm{coef}}
\bigl(S_{\mathrm{coef}}+K_{\mathrm{inv}}C_{\mathrm{bias}}Lh^\beta\bigr),
\qquad z\in[0,1]^d,
\]
and clipping contracts distance to every value in \([-B,B]\), we obtain
\[
\sup_{x\in\mathcal X}
|\widehat\mu_n(x)-\mu_{1,P}(x)|
\le
p_{\mathrm{coef}}S_{\mathrm{coef}}+K_{\mathrm{bias}}h^\beta.
\]
We derive the conditional coefficient envelope on this same count event. For the selected design, define
\[
q_Q:=\min_{1\le\ell\le J}q_{Q\ell},
\qquad
\rho_\gamma:=\frac1{\gamma-1},
\]
and order the selected cubes so that
\[
\{Q_{(k)}:1\le k\le h^{-d}\}=\mathcal Q_h,
\qquad
q_{(k)}:=q_{Q_{(k)}},
\qquad
q_{(1)}\le\cdots\le q_{(h^{-d})}.
\]
Here \(h^{-d}=|\mathcal Q_h|\ge1\). The ordered-mass conclusion of \cref{obj:lem:ordered-mass} supplies a fixed constant satisfying
\[
\kappa_\Gamma>0,
\qquad
q_{(k)}\ge\kappa_\Gamma h^{D_\gamma}k^{\rho_\gamma}
\quad(1\le k\le h^{-d}).
\]

Define the treated residuals and coefficient weights by
\[
\epsilon_{\mathrm{tr},i}
:=\mathbf1_{\{A_i=1\}}\bigl(Y_i-\mu_{1,P}(X_i)\bigr),
\]
\[
w_{\mathrm{fit},i,Q,\alpha}
:=\frac1J\sum_{\ell=1}^J
\frac{\mathbf1_{\{A_i=1,\ X_i\in E_{Q\ell}\}}}{N_{Q\ell}}
\bigl(\widehat G_Q^{-1}U_{i,Q}\bigr)_\alpha,
\qquad \alpha\in\mathcal A_{\mathrm{poly}}.
\]
These weights are evaluated only at the positive feasible mesh on \(\mathcal E_n\), where every denominator is positive. Conditional on \(\mathcal G_n\), the residuals are independent by iid sampling, and identification together with \cref{obj:ass:bounded-potential-outcomes} gives, almost surely,
\[
\mathbb E[\epsilon_{\mathrm{tr},i}\mid\mathcal G_n]=0,
\qquad
\mathbb E[e^{t\epsilon_{\mathrm{tr},i}}\mid\mathcal G_n]
\le e^{B^2t^2/2}
\quad(t\in\mathbb R).
\]
For an untreated observation the residual is zero. The equal-cell coefficient formula consequently gives
\[
\widehat c_{Q,\alpha}-\bar c_{Q,\alpha}
=\sum_{i=1}^n w_{\mathrm{fit},i,Q,\alpha}
\epsilon_{\mathrm{tr},i}.
\]
The inverse bound in \cref{obj:lem:template-conditioning} and \(\|U_{i,Q}\|_2\le\sqrt{p_{\mathrm{coef}}}\) imply
\[
\left|\frac1J
(\widehat G_Q^{-1}U_{i,Q})_\alpha\right|
\le\frac{K_{\mathrm{inv}}}{J}.
\]
Cauchy--Schwarz over the \(J\) microcells, followed by summation over their treated observations, gives
\[
\begin{aligned}
\sum_{i=1}^n w_{\mathrm{fit},i,Q,\alpha}^2
&\le J\left(\frac{K_{\mathrm{inv}}}{J}\right)^2
\sum_{\ell=1}^J\frac1{N_{Q\ell}}\\
&\le K_{\mathrm{inv}}^2
\min\left(h^{2\beta},\frac5{nq_Q}\right).
\end{aligned}
\]
The last line uses both \(N_{Q\ell}\ge h^{-2\beta}\) and the established bound \(N_{Q\ell}\ge nq_{Q\ell}/5\). Multiplying the conditional exponential moments of the independent weighted residuals yields
\[
\mathbb E\!\left[
 e^{t(\widehat c_{Q,\alpha}-\bar c_{Q,\alpha})}
 \,\middle|\,\mathcal G_n\right]
\le\exp\!\left\{
\frac{B^2t^2}{2}\sum_{i=1}^n
w_{\mathrm{fit},i,Q,\alpha}^2\right\}.
\]
Define
\[
K_{\mathrm{mgf}}
:=B^2K_{\mathrm{inv}}^2\max\{1,5/\kappa_\Gamma\}>0.
\]
The ordered-mass bound and
\(\min\{u,cv\}\le\max\{1,c\}\min\{u,v\}\) for nonnegative \(u,v,c\) now give
\[
\mathbb E\!\left[
 e^{t(\widehat c_{Q_{(k)},\alpha}-\bar c_{Q_{(k)},\alpha})}
 \,\middle|\,\mathcal G_n\right]
\le\exp\!\left\{
\frac{K_{\mathrm{mgf}}t^2}{2}
\min\left(h^{2\beta},n^{-1}h^{-D_\gamma}k^{-\rho_\gamma}\right)
\right\}.
\]
All these exponential moments are finite; their two-sided finiteness also gives conditional integrability of the finite coefficient maxima.
% lean: CausalSmith.Stat.WeakOverlap.selectedMesh_saturated_rank_mgf

To apply the maximal inequality in \cref{obj:lem:selected-mesh} with a constant uniform in \(\gamma\), define
\[
\rho_{\min}:=\frac1{\gamma_{\max}-1},
\qquad
\rho_{\max}:=\frac1{\gamma_{\min}-1},
\qquad
R_{\mathrm{max}}:=\frac{\rho_{\max}}{\rho_{\min}},
\]
\[
a_{\mathrm{max}}:=\sqrt{K_{\mathrm{mgf}}}\,h^\beta,
\qquad
b_{\mathrm{max}}:=\sqrt{K_{\mathrm{mgf}}}\,n^{-1/2}h^{-D_\gamma/2}.
\]
Then
\[
0<\rho_{\min}\le\rho_\gamma\le\rho_{\max},
\qquad R_{\mathrm{max}}\ge1,
\qquad
 a_{\mathrm{max}}^2=K_{\mathrm{mgf}}h^{2\beta},
\qquad
 b_{\mathrm{max}}^2=K_{\mathrm{mgf}}n^{-1}h^{-D_\gamma}.
\]
Since \(k^{-\rho_\gamma}\le k^{-\rho_{\min}}\) for \(k\ge1\), the same moment bound satisfies the maximal inequality at the fixed exponent \(\rho_{\min}\). Let \(K_{\min,\mathrm{max}}>0\) be its constant, and define
\[
M_{\mathrm{max}}:=K_{\min,\mathrm{max}}\sqrt{R_{\mathrm{max}}}.
\]
The logarithmic factors satisfy
\[
\begin{aligned}
1+\max\left\{0,\log\left(
 (b_{\mathrm{max}}/a_{\mathrm{max}})^{2/\rho_{\min}}
\right)\right\}
\le R_{\mathrm{max}}\left[
1+\max\left\{0,\log\left(
 (b_{\mathrm{max}}/a_{\mathrm{max}})^{2/\rho_\gamma}
\right)\right\}\right].
\end{aligned}
\]
Indeed, when \(b_{\mathrm{max}}/a_{\mathrm{max}}\ge1\), expand the logarithms and use \(\rho_\gamma\le\rho_{\max}\); when this ratio is smaller than one, both positive-part logarithms vanish. Applying the scalar maximal inequality to each coefficient and using
\[
S_{\mathrm{coef}}
\le\sum_{\alpha\in\mathcal A_{\mathrm{poly}}}
\max_{Q\in\mathcal Q_h}
|\widehat c_{Q,\alpha}-\bar c_{Q,\alpha}|,
\]
we obtain
\[
\mathbb E[S_{\mathrm{coef}}\mid\mathcal G_n]
\le p_{\mathrm{coef}}M_{\mathrm{max}}a_{\mathrm{max}}
\sqrt{1+\max\left\{0,\log\left(
(b_{\mathrm{max}}/a_{\mathrm{max}})^{2/\rho_\gamma}
\right)\right\}}.
\]

The oracle-mesh identity gives
\[
n^{-1/2}=h_{\gamma,n}^{\beta+D_\gamma/2},
\qquad
\frac{b_{\mathrm{max}}}{a_{\mathrm{max}}}
=\left(\frac{h_{\gamma,n}}h\right)^{\beta+D_\gamma/2}.
\]
Define
\[
\zeta_{\mathrm{log},\gamma}:=2\beta(\gamma-1)+d\gamma,
\qquad
\zeta_{\mathrm{log},\Gamma}
:=\max\{1,2\beta(\gamma_{\max}-1)+d\gamma_{\max}\}.
\]
Then
\[
\log\left((b_{\mathrm{max}}/a_{\mathrm{max}})^{2/\rho_\gamma}\right)
=\zeta_{\mathrm{log},\gamma}\log(h_{\gamma,n}/h),
\qquad
0<\zeta_{\mathrm{log},\gamma}\le\zeta_{\mathrm{log},\Gamma},
\]
and hence
\[
\mathbb E[S_{\mathrm{coef}}\mid\mathcal G_n]
\le p_{\mathrm{coef}}M_{\mathrm{max}}
\sqrt{K_{\mathrm{mgf}}\zeta_{\mathrm{log},\Gamma}}\,
 h^\beta\sqrt{1+\max\{0,\log(h_{\gamma,n}/h)\}}.
\]
To absorb the remaining logarithm, define
\[
z_{\mathrm{mesh}}:=h/h_{\gamma,n},
\qquad
0<z_{\mathrm{mesh}}\le K_{\mathrm{mesh}}.
\]
If \(z_{\mathrm{mesh}}\le1\), then
\[
\log(1/z_{\mathrm{mesh}})\le1/z_{\mathrm{mesh}}-1,
\qquad
\sqrt{1+\log(1/z_{\mathrm{mesh}})}\le1/z_{\mathrm{mesh}},
\qquad
z_{\mathrm{mesh}}^\beta\le z_{\mathrm{mesh}},
\]
so their product is at most one. If \(z_{\mathrm{mesh}}>1\), the positive-part logarithm vanishes and \(z_{\mathrm{mesh}}^\beta\le K_{\mathrm{mesh}}^\beta\). Thus
\[
h^\beta\sqrt{1+\max\{0,\log(h_{\gamma,n}/h)\}}
\le\max\{1,K_{\mathrm{mesh}}^\beta\}r_{\gamma,n}.
\]
Choose the initially enlarged coefficient constant to satisfy
\[
K_{\mathrm{coef}}\ge
p_{\mathrm{coef}}M_{\mathrm{max}}
\sqrt{K_{\mathrm{mgf}}\zeta_{\mathrm{log},\Gamma}}
\max\{1,K_{\mathrm{mesh}}^\beta\}.
\]
We have therefore established, almost surely on the explicitly chosen count event,
\[
\mathbb E[S_{\mathrm{coef}}\mid\mathcal G_n]
\le K_{\mathrm{coef}}r_{\gamma,n}
\quad\text{on }\mathcal E_n.
\]
% lean: CausalSmith.Stat.WeakOverlap.selectedMesh_saturated_conditional_envelope
Also,
\[
h^\beta
\le K_{\mathrm{mesh}}^\beta r_{\gamma,n}
\le \max\{1,K_{\mathrm{mesh}}^\beta\}r_{\gamma,n}.
\]
Thus, defining
\[
K_{\mathrm{cond}}
:=
p_{\mathrm{coef}}K_{\mathrm{coef}}
+
K_{\mathrm{bias}}\max\{1,K_{\mathrm{mesh}}^\beta\}
+1,
\]
we have \(K_{\mathrm{cond}}>0\) and
\[
\mathbb E\!\left[
\sup_{x\in\mathcal X}
|\widehat\mu_n(x)-\mu_{1,P}(x)|
\,\middle|\,\mathcal G_n
\right]
\le K_{\mathrm{cond}}r_{\gamma,n}
\quad\text{almost surely on }\mathcal E_n.
\]
The finite mesh selection and polynomial fits are measurable, and the conditional integrability needed here is supplied by \cref{obj:lem:selected-mesh}.
% lean: CausalSmith.Stat.WeakOverlap.selectedMesh_saturated_conditional_cubeLoss

\emph{3. Uniform expected risk.}
For every sample, including samples with \(\widehat h=0\), the estimator in \cref{obj:def:estimator} takes values in \([-B,B]\). Define, on the entire sample space,
\[
Z_{\mathrm{loss}}
:=
\sup_{x\in\mathcal X}
|\widehat\mu_n(x)-\mu_{1,P}(x)|,
\qquad
F_{\mathrm{loss}}
:=
\mathbb E[Z_{\mathrm{loss}}\mid\mathcal G_n].
\]
The preceding response bound gives
\[
0\le Z_{\mathrm{loss}}\le2B,
\qquad
0\le F_{\mathrm{loss}}\le2B.
\]
The supremum loss is measurable and bounded. The tower property and the bounds on the event and its complement therefore give
\[
\begin{aligned}
\mathbb E_{P^{\otimes n}}Z_{\mathrm{loss}}
&=\mathbb E_{P^{\otimes n}}F_{\mathrm{loss}}
\\
&\le
K_{\mathrm{cond}}r_{\gamma,n}
P^{\otimes n}(\mathcal E_n)
+
2B P^{\otimes n}(\mathcal E_n^c)
\\
&\le K_{\mathrm{cond}}r_{\gamma,n}+2Bn^{-2}.
\end{aligned}
\]
For every \(n\ge1\), \(\beta>0\), and \(\gamma>1\),
\[
D_\gamma\ge0,
\qquad
2\beta+D_\gamma>0,
\qquad
\frac{\beta}{2\beta+D_\gamma}\le1,
\]
so monotonicity of powers with base \(n\ge1\) gives
\[
n^{-2}
\le n^{-\beta/(2\beta+D_\gamma)}
=r_{\gamma,n}.
\]
Set
\[
K:=K_{\mathrm{cond}}+2B.
\]
Then \(K>0\), and
\[
\mathbb E_{P^{\otimes n}}
\left[
\sup_{x\in\mathcal X}
|\widehat\mu_n(x)-\mu_{1,P}(x)|
\right]
\le Kr_{\gamma,n}
\]
uniformly for \(n\ge n_0\), \(\gamma\in\Gamma\), and \(P\in\mathcal P_\gamma\). The count rule and clipped fit in \cref{obj:def:mesh-selector,obj:def:estimator} use the stated inputs \(n,d,\beta,B,\mathcal D_n\).
% lean: CausalSmith.Stat.WeakOverlap.estimatorSupRisk_uniform_upper

\emph{4. The bounded-subclass lower bound.}
Fix \(\gamma>1\) and \(x_0\in\mathcal X\). Define the subclass and its pointwise minimax risk by
\[
\mathcal P_\gamma^{\{0,B\}}
:=
\left\{
P\in\mathcal P_\gamma:
P\text{ admits a causal completion with }
Y(0),Y(1)\in\{0,B\}\text{ almost surely}
\right\},
\]
\[
R_{\mathrm{bd}}(n,\gamma,x_0)
:=
\inf_T\sup_{P\in\mathcal P_\gamma^{\{0,B\}}}
\mathbb E_{P^{\otimes n}}
|T(\mathcal D_n,x_0)-\mu_{1,P}(x_0)|.
\]
The infimum is over sample-measurable scalar estimators.

Apply \cref{obj:lem:bounded-lower-pair}. It supplies fixed \(a,c,c_0>0\) such that, for every \(n\ge1\), the pair from \cref{obj:def:bounded-lower-pair}, with
\[
h:=ch_{\gamma,n},
\]
belongs to \(\mathcal P_\gamma^{\{0,B\}}\) and satisfies
\[
P_{0,h}\ll P_{1,h},
\qquad
\log\frac{dP_{0,h}}{dP_{1,h}}\in L^1(P_{0,h}),
\qquad
n\,\operatorname{KL}(P_{0,h},P_{1,h})\le\frac18.
\]
Define
\[
\mathbb P_{j,n}:=P_{j,h}^{\otimes n},
\qquad
\tau_{j,n}:=\mu_{1,P_{j,h}}(x_0)
\quad(j\in\{0,1\}),
\qquad
\delta_n:=c_0r_{\gamma,n}.
\]
The separation supplied by \cref{obj:lem:bounded-lower-pair} is
\[
0<\delta_n\le|\tau_{1,n}-\tau_{0,n}|.
\]

For an arbitrary sample-measurable estimator \(T\), define its clipped version by
\[
T_{\mathrm{clip}}(\mathcal D_n)
:=
\operatorname{clip}_{[-B,B]}
\bigl(T(\mathcal D_n,x_0)\bigr).
\]
The response bound established above gives, for \(j\in\{0,1\}\),
\[
|\tau_{j,n}|\le B,
\qquad
|T_{\mathrm{clip}}-\tau_{j,n}|
\le |T(\mathcal D_n,x_0)-\tau_{j,n}|,
\qquad
|T_{\mathrm{clip}}-\tau_{j,n}|\le2B.
\]
Thus the clipped losses are integrable under both sample laws.

Product entropy tensorization and Pinsker's inequality give
\[
\operatorname{KL}(\mathbb P_{0,n},\mathbb P_{1,n})
\le n\,\operatorname{KL}(P_{0,h},P_{1,h})
\le\frac18\le\frac12,
\]
\[
\|\mathbb P_{0,n}-\mathbb P_{1,n}\|_{\mathrm{TV}}
\le
\sqrt{\frac{\operatorname{KL}(\mathbb P_{0,n},\mathbb P_{1,n})}{2}}
\le\frac12.
\]
Define the two large-error events by
\[
\mathcal A_{j,n}
:=
\left\{
\mathcal D_n:
|T_{\mathrm{clip}}(\mathcal D_n)-\tau_{j,n}|
\ge\frac{\delta_n}{2}
\right\},
\qquad j\in\{0,1\}.
\]
The triangle inequality and target separation imply
\[
\mathcal A_{0,n}^c\subseteq\mathcal A_{1,n}.
\]
Hence
\[
\begin{aligned}
\mathbb P_{0,n}(\mathcal A_{0,n})
+\mathbb P_{1,n}(\mathcal A_{1,n})
&\ge
1-\|\mathbb P_{0,n}-\mathbb P_{1,n}\|_{\mathrm{TV}},
\\
\max_{j\in\{0,1\}}\mathbb P_{j,n}(\mathcal A_{j,n})
&\ge
\frac{1-
\sqrt{\operatorname{KL}(\mathbb P_{0,n},\mathbb P_{1,n})/2}}{2}
\ge\frac14.
\end{aligned}
\]
For each \(j\), integration over \(\mathcal A_{j,n}\) yields
\[
\mathbb E_{\mathbb P_{j,n}}
|T_{\mathrm{clip}}-\tau_{j,n}|
\ge\frac{\delta_n}{2}\mathbb P_{j,n}(\mathcal A_{j,n}).
\]
Combining the two inequalities and the clipping contraction gives
\[
\begin{aligned}
\max_{j\in\{0,1\}}
\mathbb E_{\mathbb P_{j,n}}
|T(\mathcal D_n,x_0)-\tau_{j,n}|
&\ge
\max_{j\in\{0,1\}}
\mathbb E_{\mathbb P_{j,n}}
|T_{\mathrm{clip}}-\tau_{j,n}|
\\
&\ge\frac{\delta_n}{8}
=\frac{c_0}{8}r_{\gamma,n}.
\end{aligned}
\]
Both laws belong to the restricted subclass. Taking the infimum over \(T\), and defining
\[
c_\gamma:=\frac{c_0}{8}>0,
\]
proves
\[
c_\gamma r_{\gamma,n}
\le R_{\mathrm{bd}}(n,\gamma,x_0)
\qquad(n\ge1).
\]
% lean: CausalSmith.Stat.WeakOverlap.boundedPointwiseRisk_oracle_lower

\emph{5. Constants for all positive sample sizes.}
For the fixed \(\gamma\), define a local adaptation interval by
\[
\gamma_{\mathrm{lo}}:=\frac{\gamma+1}{2},
\qquad
\gamma_{\mathrm{hi}}:=\gamma+1.
\]
Since \(\gamma>1\),
\[
1<\gamma_{\mathrm{lo}}<\gamma_{\mathrm{hi}},
\qquad
\gamma\in[\gamma_{\mathrm{lo}},\gamma_{\mathrm{hi}}].
\]
Applying the uniform estimator bound to this interval supplies \(K_{\mathrm{loc}}>0\) and \(n_{\mathrm{loc}}\in\mathbb N\) such that
\[
\sup_{P\in\mathcal P_\gamma}
\mathbb E_{P^{\otimes n}}
\left[
\sup_{x\in\mathcal X}
|\widehat\mu_n(x)-\mu_{1,P}(x)|
\right]
\le K_{\mathrm{loc}}r_{\gamma,n}
\quad(n\ge\max\{n_{\mathrm{loc}},1\}).
\]
Define
\[
N_\gamma:=\max\{n_{\mathrm{loc}},1\},
\qquad
q_{\mathrm{small}}:=N_\gamma^{-2},
\qquad
K_{\mathrm{small}}:=\frac{2B}{q_{\mathrm{small}}},
\]
\[
K_\gamma
:=
\max\{c_\gamma,\max\{K_{\mathrm{loc}},K_{\mathrm{small}}\}\}.
\]
These constants are finite, \(q_{\mathrm{small}}>0\), \(K_{\mathrm{small}}>0\), and
\[
0<c_\gamma\le K_\gamma,
\qquad
K_{\mathrm{loc}}\le K_\gamma,
\qquad
K_{\mathrm{small}}\le K_\gamma.
\]

Fix \(n\ge1\). If \(n\ge n_{\mathrm{loc}}\), the measurable estimator \(\widehat\mu_n\) is an admissible witness in the supremum minimax infimum, so
\[
R_\infty(n,\gamma)
\le
\sup_{P\in\mathcal P_\gamma}
\mathbb E_{P^{\otimes n}}
\left[
\sup_{x\in\mathcal X}
|\widehat\mu_n(x)-\mu_{1,P}(x)|
\right]
\le K_{\mathrm{loc}}r_{\gamma,n}
\le K_\gamma r_{\gamma,n}.
\]
If \(n<n_{\mathrm{loc}}\), then \(1\le n\le N_\gamma\). Monotonicity of the inverse square and \(n^{-2}\le r_{\gamma,n}\) give
\[
q_{\mathrm{small}}
=N_\gamma^{-2}
\le n^{-2}
\le r_{\gamma,n}.
\]
Therefore
\[
2B
=K_{\mathrm{small}}q_{\mathrm{small}}
\le K_{\mathrm{small}}r_{\gamma,n}
\le K_\gamma r_{\gamma,n}.
\]
The same clipped estimator has loss at most \(2B\) for every sample, so
\[
R_\infty(n,\gamma)
\le
\sup_{P\in\mathcal P_\gamma}
\mathbb E_{P^{\otimes n}}
\left[
\sup_{x\in\mathcal X}
|\widehat\mu_n(x)-\mu_{1,P}(x)|
\right]
\le2B
\le K_\gamma r_{\gamma,n}.
\]
Thus the upper bound holds for every \(n\ge1\).
% lean: CausalSmith.Stat.WeakOverlap.adaptive_minimax_rate

\emph{6. Comparison of the risks.}
The inclusion
\[
\mathcal P_\gamma^{\{0,B\}}\subseteq\mathcal P_\gamma
\]
implies, for each scalar estimator \(T\),
\[
\sup_{P\in\mathcal P_\gamma^{\{0,B\}}}
\mathbb E_{P^{\otimes n}}
|T(\mathcal D_n,x_0)-\mu_{1,P}(x_0)|
\le
\sup_{P\in\mathcal P_\gamma}
\mathbb E_{P^{\otimes n}}
|T(\mathcal D_n,x_0)-\mu_{1,P}(x_0)|.
\]
Taking infima over the same estimators gives
\[
R_{\mathrm{bd}}(n,\gamma,x_0)
\le R_{\mathrm{pt}}(n,\gamma,x_0).
\]
For any admissible function-valued estimator \(T\), define its evaluation by
\[
T_{x_0}(\mathcal D_n):=T(\mathcal D_n)(x_0).
\]
This is a measurable scalar estimator, and \(x_0\in\mathcal X\) gives
\[
|T_{x_0}(\mathcal D_n)-\mu_{1,P}(x_0)|
\le
\sup_{x\in\mathcal X}
|T(\mathcal D_n)(x)-\mu_{1,P}(x)|.
\]
Taking expectations, worst-case risks, and then the infimum over function-valued estimators yields
\[
R_{\mathrm{pt}}(n,\gamma,x_0)\le R_\infty(n,\gamma).
\]
Combining the bounded-subclass lower bound, the upper bound for every positive sample size, and these comparisons proves
\[
c_\gamma r_{\gamma,n}
\le R_{\mathrm{bd}}(n,\gamma,x_0)
\le R_{\mathrm{pt}}(n,\gamma,x_0)
\le R_\infty(n,\gamma)
\le K_\gamma r_{\gamma,n},
\qquad n\ge1.
\]
% lean: CausalSmith.Stat.WeakOverlap.boundedPointwiseRisk_le_pointwiseRisk
% lean: CausalSmith.Stat.WeakOverlap.pointwiseRisk_le_supremumRisk
\end{proof}

\begin{proof}[Proof of \cref{obj:prop:strict-enlargement}]
Fix an admissible strip exponent \(a\), and use the construction and selected propensity in \cref{obj:def:thin-strip} throughout.

\begin{enumerate}
\item Let the completed-data law and its treated response be
\[
P_{\mathrm{comp}}
:=
\mathcal L(X,A,Y(0),Y(1),Y),
\qquad
\mu_1(x):=Bp_0,\quad x\in\mathcal X.
\]
The construction gives
\[
0<p_0=\min\left\{\frac14,\frac{L}{4B}\right\}\le\frac14,
\qquad
Bp_0\le B\frac{L}{4B}=\frac L4.
\]
The radial term and strip increment are nonnegative, and truncation at one therefore gives
\[
0\le e_\star(x)\le1,\qquad x\in\mathcal X.
\]
Thus all conditional Bernoulli distributions are probability distributions. Integrating their product against the uniform covariate distribution shows that \(P_{\mathrm{comp}}\) and its observed-data marginal \(P_{\mathrm{strip}}\) are probability laws, with
\[
(P_{\mathrm{comp}})_X=(P_{\mathrm{strip}})_X
=\operatorname{Leb}_d\!\restriction_{\mathcal X}.
\]

The conditional distribution of \(Y(1)\) is the constant scaled Bernoulli distribution, so
\[
\mathbb E_{P_{\mathrm{comp}}}[Y(1)\mid X]
=(1-p_0)\,0+p_0B
=Bp_0
=\mu_1(X)
\quad\text{almost surely}.
\]
For the derivative order \(m\) in \cref{obj:def:holder}, the constant function satisfies
\[
\max_{|\alpha|\le m}\|D^\alpha\mu_1\|_\infty
=Bp_0\le\frac L4\le L,
\qquad
\max_{|\alpha|=m}\sup_{x\ne x'}
\frac{|D^\alpha\mu_1(x)-D^\alpha\mu_1(x')|}
{\|x-x'\|_\infty^{\beta-m}}
=0.
\]
Its derivatives have continuous boundary traces. Hence
\[
\mu_1\in\mathcal H^\beta(L),
\]
as defined in \cref{obj:def:holder}, and it supplies the response required by \cref{obj:ass:holder-response}.
% lean: CausalSmith.Stat.WeakOverlap.completionBind_holderResponse_constant_equal_arms

\item We verify the remaining restrictions in \cref{obj:def:model}. The uniform marginal gives
\[
P_{\mathrm{strip}}(X\in\mathcal X)=1,
\qquad
f(x)=1\ge c_f
\quad\text{for Lebesgue-almost every }x\in\mathcal X,
\]
which supplies \cref{obj:ass:density}. The completed construction gives
\[
Y=AY(1)+(1-A)Y(0)
\quad\text{almost surely},
\]
and, for almost every \(x\in\mathcal X\),
\[
\mathcal L\bigl((Y(0),Y(1)),A\mid X=x\bigr)
=
\mathcal L\bigl(Y(0),Y(1)\mid X=x\bigr)
\otimes\operatorname{Bernoulli}(e_\star(x)).
\]
These identities verify \cref{obj:ass:consistency,obj:ass:exchangeability}, and the treatment marginal gives
\[
P_{\mathrm{comp}}(A=1\mid X=x)=e_\star(x)
\quad\text{for almost every }x.
\]

For the overlap calculation, define the radial propensity and radial exponent by
\[
q_{\mathrm{rad}}(x)
:=F_R(R(x))^{1/(\gamma-1)},\quad x\in\mathcal X,
\qquad
\delta_{\mathrm{strip}}:=\frac d{\gamma-1}>0.
\]
The centered cube of radius \(r\) has side length \(2r\), whence
\[
F_R(r)=(2r)^d,\qquad 0\le r\le\frac12.
\]
In particular \(0\le q_{\mathrm{rad}}(x)\le1\), so the nonnegative strip increment gives
\[
q_{\mathrm{rad}}(x)\le e_\star(x),
\qquad
\{x\in\mathcal X:e_\star(x)\le t\}
\subseteq
\{x\in\mathcal X:F_R(R(x))\le t^{\gamma-1}\}.
\]

To bound the last event over its whole range, for \(0\le\tau_{\mathrm{cdf}}<1\) define
\[
r_{\mathrm{cdf}}(\tau_{\mathrm{cdf}})
:=\frac{\tau_{\mathrm{cdf}}^{1/d}}2.
\]
Then \(0\le r_{\mathrm{cdf}}(\tau_{\mathrm{cdf}})\le1/2\), and monotonicity of the \(d\)-th power on nonnegative reals gives
\[
\{x\in\mathcal X:F_R(R(x))\le\tau_{\mathrm{cdf}}\}
\subseteq
\{x\in\mathcal X:R(x)\le r_{\mathrm{cdf}}(\tau_{\mathrm{cdf}})\}.
\]
Consequently,
\[
P_{\mathrm{strip}}\{F_R(R(X))\le\tau_{\mathrm{cdf}}\}
\le
\bigl(2r_{\mathrm{cdf}}(\tau_{\mathrm{cdf}})\bigr)^d
=\tau_{\mathrm{cdf}}.
\]
This includes \(\tau_{\mathrm{cdf}}=0\); at \(\tau_{\mathrm{cdf}}=1\), the same upper bound follows from total probability one. Taking \(\tau_{\mathrm{cdf}}=t^{\gamma-1}\) therefore yields
\[
P_{\mathrm{strip}}\{e_\star(X)\le t\}
\le t^{\gamma-1}
\le Ct^{\gamma-1},
\qquad 0\le t\le1.
\]
This verifies \cref{obj:ass:global-tail}. At \(t=0\), the sharp bound and nonnegativity of \(e_\star\) also give
\[
P_{\mathrm{strip}}\{e_\star(X)=0\}=0,
\qquad
e_\star(X)>0\quad\text{almost surely}.
\]

The equal-arm product construction identifies the treated conditional distribution wherever treatment has positive conditional probability:
\[
\mathcal L_{P_{\mathrm{strip}}}(Y\mid A=1,X=x)
=\mathcal L\bigl(B\operatorname{Bernoulli}(p_0)\bigr),
\qquad
\mathbb E_{P_{\mathrm{strip}}}[Y\mid A=1,X=x]=Bp_0
\]
for \(P_X\)-almost every \(x\). Its support is contained in \([0,B]\). The bounded-range exponential-moment inequality, followed by enlargement of its exponent, gives
\[
\mathbb E_{P_{\mathrm{strip}}}
\!\left[\exp\{\lambda(Y-Bp_0)\}\mid A=1,X=x\right]
\le
\exp\left\{\frac{B^2\lambda^2}{8}\right\}
\le
\exp\left\{\frac{B^2\lambda^2}{2}\right\},
\qquad \lambda\in\mathbb R.
\]
Also \(|Bp_0|\le B\). For every \(\lambda\ne0\), it follows that
\[
\frac{
\sqrt{2\log\mathbb E_{P_{\mathrm{strip}}}
[\exp\{\lambda(Y-Bp_0)\}\mid A=1,X=x]}
}{|\lambda|}
\le B.
\]
The centered exponential moment is at least one by Jensen's inequality, so the square root is well defined. These estimates supply \cref{obj:ass:bounded-potential-outcomes}. Together with the smoothness already established and the stated parameter conditions, they prove
\[
P_{\mathrm{strip}}\in\mathcal P_\gamma.
\]
The uniform covariate marginal also supplies the asserted sharp bound
\[
P_{\mathrm{strip}}\{e_\star(X)\le t\}\le t^{\gamma-1},
\qquad t\in[0,1].
\]
% lean: CausalSmith.Stat.WeakOverlap.thinStrip_globalTailModel_of_holder
% lean: CausalSmith.Stat.WeakOverlap.stripPropensity_tail_of_uniform_covariateLaw

\item Fix arbitrary \(\rho,\nu,h_0>0\). For \(h>0\), define the centered Euclidean ball and local propensity supremum by
\[
\mathcal B_h:=\{x\in\mathbb R^d:\|x-x^\circ\|_2\le h\},
\qquad
M_{\mathrm{loc}}(h)
:=\sup_{x\in\mathcal B_h\cap\mathcal X}e_\star(x).
\]
For \(0<h\le1/2\), every coordinate displacement in \(\mathcal B_h\) has magnitude at most \(h\). On \(S\cap\mathcal B_h\), the strip inequality further gives
\[
|x_2-1/2|\le |x_1-1/2|^a\le h^a.
\]
Thus this intersection is contained in a rectangular box with second-coordinate side length \(2h^a\) and all other side lengths \(2h\). Uniformity implies
\[
P_{\mathrm{strip}}(X\in S\cap\mathcal B_h)
\le(2h^a)(2h)^{d-1}.
\]
For the denominator, the centered box with every coordinate displacement at most \(h/d\) lies in \(\mathcal X\) and in \(\mathcal B_h\), since
\[
\sum_{i=1}^d(x_i-1/2)^2
\le d(h/d)^2
=\frac{h^2}{d}\le h^2.
\]
Its volume therefore supplies
\[
P_{\mathrm{strip}}(X\in\mathcal B_h)
\ge(2h/d)^d>0.
\]
Dividing these bounds yields
\[
\frac{P_{\mathrm{strip}}(X\in S\cap\mathcal B_h)}
{P_{\mathrm{strip}}(X\in\mathcal B_h)}
\le
\frac{(2h^a)(2h)^{d-1}}{(2h/d)^d}
=d^d h^{a-1}.
\]
Because \(a-1>0\) and \(\delta_{\mathrm{strip}}>0\),
\[
d^d h^{a-1}\longrightarrow0,
\qquad
(2h)^{\delta_{\mathrm{strip}}}\longrightarrow0
\quad\text{as }h\downarrow0.
\]
Choose \(h\) satisfying
\[
0<h<\min\{h_0,1/2\},
\qquad
d^d h^{a-1}<\nu,
\qquad
(2h)^{\delta_{\mathrm{strip}}}<\rho/4.
\]

The center belongs to \(S\), and its radial term is zero. Hence
\[
e_\star(x^\circ)=\frac14,
\qquad
M_{\mathrm{loc}}(h)\ge\frac14.
\]
For \(x\in\mathcal B_h\setminus S\), the propensity equals its radial component and \(R(x)\le h\), giving
\[
e_\star(x)=q_{\mathrm{rad}}(x)
\le(2h)^{\delta_{\mathrm{strip}}}
<\rho/4
\le\rho M_{\mathrm{loc}}(h).
\]
Therefore
\[
\{x\in\mathcal B_h:e_\star(x)\ge\rho M_{\mathrm{loc}}(h)\}
\subseteq S\cap\mathcal B_h,
\]
and
\[
\frac{
P_{\mathrm{strip}}\{e_\star(X)\ge\rho M_{\mathrm{loc}}(h),
\ X\in\mathcal B_h\}
}{
P_{\mathrm{strip}}(X\in\mathcal B_h)
}
\le d^d h^{a-1}<\nu.
\]
This contradicts the numerical clause for the proposed triple. Since the triple was arbitrary, the numerical anti-concentration clause fails. The argument uses the pointwise bound \(e_\star\le1\), so it applies with propensity values equal to one admitted.
% lean: CausalSmith.Stat.WeakOverlap.thinStripLaw_not_dornA3AntiConcentration

\item The full singleton-family condition in \cref{obj:def:dorn-a3} entails its numerical clause with the same constants \(\rho,\nu,h_0\). Thus the preceding contradiction gives
\[
\neg\mathsf A_3^{\mathrm{Dorn}}(P_{\mathrm{strip}})
\]
for the selected propensity \(e_\star\).
% lean: CausalSmith.Stat.WeakOverlap.dornA3AntiConcentrationOnCube_of_dornA3OnCube

\item For the design limit, define, for every \(h>0\),
\[
\mathcal Q_{\mathrm{ctr},h}
:=\{x\in\mathbb R^d:R(x)\le h\},
\]
\[
\begin{aligned}
D_{\mathrm{strip}}(h)
&:=P_{\mathrm{strip}}\{A=1,\ X\in\mathcal Q_{\mathrm{ctr},h}\},\\
N_{\mathrm{strip}}(h)
&:=\int_{\{A=1,\ X\in\mathcal Q_{\mathrm{ctr},h}\}}
\left(\frac{X_2-1/2}{h}\right)^2\,dP_{\mathrm{strip}},\\
T_{\mathrm{strip}}(h)
&:=\frac{N_{\mathrm{strip}}(h)}{D_{\mathrm{strip}}(h)}.
\end{aligned}
\]
For this definition, a quotient with zero denominator is assigned value zero.

Introduce the positive constants and envelope
\[
b_{\mathrm{strip}}
:=\frac{1}{2(4d)^{a+d-1}},
\qquad
c_{\mathrm{rad}}:=2^{d+\delta_{\mathrm{strip}}},
\qquad
c_{\mathrm{strip}}:=2^d,
\]
\[
G_{\mathrm{strip}}(h)
:=
\frac{
c_{\mathrm{rad}}h^{d+\delta_{\mathrm{strip}}}
+c_{\mathrm{strip}}h^{d+3a-3}
}{
b_{\mathrm{strip}}h^{d+a-1}
},
\qquad h>0.
\]
Factoring powers of \(h\) gives
\[
G_{\mathrm{strip}}(h)
=
\frac{c_{\mathrm{rad}}}{b_{\mathrm{strip}}}
h^{1+\delta_{\mathrm{strip}}-a}
+
\frac{c_{\mathrm{strip}}}{b_{\mathrm{strip}}}
h^{2a-2}.
\]
The strip-exponent restrictions give
\[
1+\delta_{\mathrm{strip}}-a>0,
\qquad
2a-2>0,
\]
so both powers tend to zero and
\[
G_{\mathrm{strip}}(h)\longrightarrow0
\quad\text{as }h\downarrow0.
\]
% lean: CausalSmith.Stat.WeakOverlap.thinStrip_raw_second_envelope_tendsto_zero

\item We bound the moment by this envelope. Suppose \(0<h\le1/2\). On \(\mathcal Q_{\mathrm{ctr},h}\), the normalized second-coordinate square is at most one. If \(x\in S\cap\mathcal Q_{\mathrm{ctr},h}\), the strip inequality and \(e_\star(x)\le1\) give
\[
\left(\frac{x_2-1/2}{h}\right)^2e_\star(x)
\le\left(\frac{h^a}{h}\right)^2.
\]
If \(x\in\mathcal Q_{\mathrm{ctr},h}\setminus S\), the radial bound gives
\[
\left(\frac{x_2-1/2}{h}\right)^2e_\star(x)
\le e_\star(x)\le(2h)^{\delta_{\mathrm{strip}}}.
\]
Combining these two cases and integrating,
\[
\begin{aligned}
&\int_{\mathcal Q_{\mathrm{ctr},h}}
\left(\frac{x_2-1/2}{h}\right)^2e_\star(x)\,dx\\
&\qquad\le
(2h)^{\delta_{\mathrm{strip}}}
P_{\mathrm{strip}}(X\in\mathcal Q_{\mathrm{ctr},h})
+
\left(\frac{h^a}{h}\right)^2
P_{\mathrm{strip}}(X\in S\cap\mathcal Q_{\mathrm{ctr},h}).
\end{aligned}
\]
Here \(\mathcal Q_{\mathrm{ctr},h}\subseteq\mathcal X\), and its volume is \((2h)^d\). The same rectangular enclosure used above bounds the strip intersection, so
\[
\begin{aligned}
&\int_{\mathcal Q_{\mathrm{ctr},h}}
\left(\frac{x_2-1/2}{h}\right)^2e_\star(x)\,dx\\
&\qquad\le
(2h)^{\delta_{\mathrm{strip}}}(2h)^d
+\left(\frac{h^a}{h}\right)^2(2h^a)(2h)^{d-1}\\
&\qquad=
c_{\mathrm{rad}}h^{d+\delta_{\mathrm{strip}}}
+c_{\mathrm{strip}}h^{d+3a-3}.
\end{aligned}
\]

For the treated-mass lower bound, define
\[
r_{\mathrm{core}}(h):=\frac{h}{4d},
\]
\[
\mathcal K_h
:=
\left\{
x\in\mathbb R^d:
|x_2-1/2|\le r_{\mathrm{core}}(h)^a,\
r_{\mathrm{core}}(h)\le x_i-1/2\le2r_{\mathrm{core}}(h)
\ \text{for every }i\ne2
\right\}.
\]
Since \(0<r_{\mathrm{core}}(h)\le1/8\) and \(a>1\),
\[
r_{\mathrm{core}}(h)^a\le r_{\mathrm{core}}(h).
\]
The coordinate bounds place \(\mathcal K_h\) in \(\mathcal X\), and they imply
\[
|x_2-1/2|
\le r_{\mathrm{core}}(h)^a
\le |x_1-1/2|^a,
\]
\[
\|x-x^\circ\|_2^2
\le d\bigl(2r_{\mathrm{core}}(h)\bigr)^2
=\frac{h^2}{4d}\le h^2.
\]
Thus
\[
\mathcal K_h\subseteq S\cap\mathcal B_h
\subseteq S\cap\mathcal Q_{\mathrm{ctr},h}.
\]
Its second-coordinate side length is \(2r_{\mathrm{core}}(h)^a\), and each other side length is \(r_{\mathrm{core}}(h)\). Since \(e_\star\ge1/4\) on \(S\),
\[
\begin{aligned}
\int_{\mathcal Q_{\mathrm{ctr},h}}e_\star(x)\,dx
&\ge\int_{\mathcal B_h\cap\mathcal X}e_\star(x)\,dx\\
&\ge\frac14
\bigl(2r_{\mathrm{core}}(h)^a\bigr)
r_{\mathrm{core}}(h)^{d-1}\\
&=b_{\mathrm{strip}}h^{d+a-1}>0.
\end{aligned}
\]

Finally, conditioning on \(X\) and using
\(P_{\mathrm{strip}}(A=1\mid X)=e_\star(X)\) gives the numerator and denominator identities
\[
N_{\mathrm{strip}}(h)
=
\int_{\mathcal Q_{\mathrm{ctr},h}}
\left(\frac{x_2-1/2}{h}\right)^2e_\star(x)\,dx,
\qquad
D_{\mathrm{strip}}(h)
=
\int_{\mathcal Q_{\mathrm{ctr},h}}e_\star(x)\,dx.
\]
All integrands are nonnegative. Dividing the numerator upper bound by the positive denominator lower bound therefore yields
\[
0\le T_{\mathrm{strip}}(h)\le G_{\mathrm{strip}}(h),
\qquad 0<h\le1/2.
\]
The squeeze theorem proves
\[
\lim_{h\downarrow0}T_{\mathrm{strip}}(h)=0,
\]
which is exactly the asserted normalized treated-design limit.
% lean: CausalSmith.Stat.WeakOverlap.stripSecondMoment_tendsto_zero

\item The law \(P_{\mathrm{strip}}\) itself supplies the final existence claim: it belongs to \(\mathcal P_\gamma\), and its full singleton-family Dorn predicate fails for \(e_\star\).
% lean: CausalSmith.Stat.WeakOverlap.thinStrip_mem_model_not_dornA3
\end{enumerate}
\end{proof}

\begin{proof}[Proof of \cref{obj:prop:dorn-a3-free-corollary}]
\begin{enumerate}
\item \emph{The source mean is bounded throughout the cube.}
Fix \(\gamma\in\Gamma\) and a source tuple in \(\mathcal D_\gamma\). Write
\[
X^{\mathrm D}:=X,
\qquad
P_X^{\mathrm D}:=\mathcal L_P(X^{\mathrm D})
=2^{-d}\operatorname{Leb}\big|_{[-1,1]^d},
\qquad
\mu_{1,P}:=\mu_P,
\]
and let \(\mu_{0,P}\) be its selected control mean. Continuity of \(\mu_P\) on the compact source cube gives
\[
\int |\mu_P(z)|\,P_X^{\mathrm D}(dz)<\infty.
\]
For \(P_X^{\mathrm D}\)-almost every \(z\), the treated conditional distribution is
\(N(\mu_P(z),\sigma_{\min}^2)\). Its mean identity, the absolute-integral inequality, and Jensen's inequality give
\[
|\mu_P(z)|
\le
\int |y|\,N(\mu_P(z),\sigma_{\min}^2)(dy),
\]
and
\[
|\mu_P(z)|^q
\le
\left(\int |y|\,N(\mu_P(z),\sigma_{\min}^2)(dy)\right)^q
\le
\int |y|^q\,N(\mu_P(z),\sigma_{\min}^2)(dy)
\le M^q.
\]
Since \(q>0\) and \(M>0\), this implies \(|\mu_P|\le M\) almost everywhere.

To extend this inequality to the boundary as well, define
\[
v_P(z):=\max\{|\mu_P(z)|-M,0\},
\qquad z\in\mathbb R^d.
\]
The function \(v_P\) is continuous on the source cube and vanishes Lebesgue-almost everywhere there. Moreover,
\[
[-1,1]^d
=
\overline{\operatorname{int}([-1,1]^d)}.
\]
Continuity first forces \(v_P\) to vanish on the interior, since a positive value would give a neighborhood of positive Lebesgue measure on which it remains positive. Taking limits from the interior then gives
\[
v_P(z)=0,
\qquad
|\mu_P(z)|\le M
\qquad(z\in[-1,1]^d).
\]
% lean: CausalSmith.Stat.WeakOverlap.dorn_A3_free_corollary

\item \emph{Completion and comparison of the common radii.}
By \cref{obj:lem:fixed-cube-holder-completion}, there is a constant \(K_{d,\beta}>0\) such that, for every nonnegative amplitude bound \(M'\), nonnegative seminorm bound \(L_0'\), and function satisfying the corresponding source regularity and bounds,
\[
u\circ T
\in
\mathcal H^\beta\!\left(K_{d,\beta}(M'+L_0')\right).
\]
For a function admitting a nonnegative Hölder radius, write its set of admissible radii as
\[
\mathscr L_\beta(g)
:=
\{L_{\mathrm{adm}}\in[0,\infty):
g\in\mathcal H^\beta(L_{\mathrm{adm}})\},
\qquad
\|g\|_{\mathcal H^\beta}
=
\inf\mathscr L_\beta(g).
\]
Consequently,
\[
\|u\circ T\|_{\mathcal H^\beta}
\le K_{d,\beta}(M'+L_0').
\]
In particular, the two suprema in the proposition are bounded above:
\[
C_{d,\beta}\le K_{d,\beta},
\qquad
L_*\le K_{d,\beta}(M+L_0).
\]

Every admissible radius bounds the function itself. Taking the infimum of those bounds gives
\[
|g(x)|\le \|g\|_{\mathcal H^\beta},
\qquad x\in[0,1]^d.
\]
Constant functions satisfy the source seminorm restriction because their derivative differences vanish. Applying the preceding inequality to the constant functions \(1\) and \(M\), which admit radii by the completion result, yields
\[
1\le \|1\|_{\mathcal H^\beta},
\qquad
M\le \|M\|_{\mathcal H^\beta}.
\]
The constant function \(1\), with amplitude and seminorm bounds both equal to \(1\), is admissible in the definition of \(C_{d,\beta}\). The constant function \(M\) is admissible in the definition of \(L_*\). Hence
\[
0<\frac12\le C_{d,\beta}\le K_{d,\beta}<\infty,
\qquad
0<M\le L_*.
\]
Finally, every function contributing to \(L_*\) also contributes, after normalization, to \(C_{d,\beta}\):
\[
\frac{\|u\circ T\|_{\mathcal H^\beta}}{M+L_0}
\le C_{d,\beta}.
\]
Taking the supremum over these functions proves
\[
0<L_*\le C_{d,\beta}(M+L_0)<\infty.
\]
% lean: CausalSmith.Stat.WeakOverlap.dorn_A3_free_corollary

\item \emph{Gaussian source witnesses for the local-condition assertion.}
Suppose \(d\ge2\), and fix \(\gamma\in\Gamma\). Then \(\gamma>1\). Choose
\[
\delta_{\mathrm{strip}}:=\frac{d}{\gamma-1}>0,
\qquad
a:=1+\frac{\delta_{\mathrm{strip}}}{2},
\]
so that
\[
1<a<1+\frac{d}{\gamma-1}.
\]
For the fixed moment exponent, set
\[
I_q:=\int |y|^q\,N(0,1)(dy),
\qquad
M':=\max\{1,|I_q|\},
\qquad
C':=\left((3/4)^{\gamma-1}\right)^{-1},
\qquad
\sigma_{\min}':=1,
\qquad
L_0':=1.
\]
Gaussian moments are finite, so \(M'\) is positive and finite. Since \(M'\ge1\) and \(q>1\),
\[
I_q\le |I_q|\le M'\le (M')^q.
\]
All four chosen source constants are positive and finite.

Use the center, radial function, strip, and propensity from
\cref{obj:def:thin-strip}. For the ambient arguments needed below, the same coordinate formulas are
\[
x^\circ=(1/2,\ldots,1/2),
\qquad
R(x)=\|x-x^\circ\|_\infty
\quad(x\in\mathbb R^d),
\]
\[
F_R(t)=
\begin{cases}
0,&t<0,\\
\min\{1,(2t)^d\},&t\ge0,
\end{cases}
\qquad
\widetilde S:=
\{x\in\mathbb R^d:
|x_2-1/2|\le |x_1-1/2|^a\}.
\]
Here \(\widetilde S\cap[0,1]^d=S\). Define the capped propensity on all of \(\mathbb R^d\) by
\[
e_{\mathrm{cap}}(x)
:=
\min\left\{\frac34,\,
\min\left\{1,\,
F_R(R(x))^{1/(\gamma-1)}
+\frac14\mathbf1_{\widetilde S}(x)\right\}\right\}.
\]
Thus, on the unit cube,
\[
e_{\mathrm{cap}}(x)=\min\{3/4,e_\star(x)\}.
\]

Construct random variables with conditional product distribution
\[
X_{\mathrm{strip}}\sim\operatorname{Unif}([0,1]^d),
\]
\[
\mathcal L\!\left(
A_{\mathrm{strip}},Y_{\mathrm{strip}}(0),Y_{\mathrm{strip}}(1)
\mid X_{\mathrm{strip}}=x
\right)
=
\operatorname{Bernoulli}(e_{\mathrm{cap}}(x))
\otimes N(0,1)\otimes N(0,1).
\]
Set
\[
Y_{\mathrm{strip}}
:=
A_{\mathrm{strip}}Y_{\mathrm{strip}}(1)
+(1-A_{\mathrm{strip}})Y_{\mathrm{strip}}(0),
\qquad
X_{\mathrm G}^{\mathrm D}:=T(X_{\mathrm{strip}}),
\]
\[
P_{\mathrm G}:=
\mathcal L(X_{\mathrm G}^{\mathrm D},A_{\mathrm{strip}},Y_{\mathrm{strip}}),
\qquad
\mu_{1,\mathrm G}(z)=\mu_{0,\mathrm G}(z):=0,
\qquad
e_{\mathrm G}(z):=e_{\mathrm{cap}}(T^{-1}z)
\quad(z\in\mathbb R^d).
\]
These measurable product kernels define a probability law. The affine change of variables and the conditional product distribution give
\[
X_{\mathrm G}^{\mathrm D}\sim\operatorname{Unif}([-1,1]^d),
\qquad
e_{\mathrm G}(X_{\mathrm G}^{\mathrm D})
=
P_{\mathrm G}(A_{\mathrm{strip}}=1\mid X_{\mathrm G}^{\mathrm D}),
\]
and
\[
Y_{\mathrm{strip}}
\mid X_{\mathrm G}^{\mathrm D},A_{\mathrm{strip}}=b
\sim N(0,1),
\qquad b\in\{0,1\}.
\]
Since the radial term is positive away from \(x^\circ\), a set of full covariate probability,
\[
0<e_{\mathrm G}(X_{\mathrm G}^{\mathrm D})\le\frac34<1
\quad\text{almost surely}.
\]
The zero means satisfy the source smoothness restriction, and
\[
\operatorname{Var}_{P_{\mathrm G}}(\mu_{1,\mathrm G}(X_{\mathrm G}^{\mathrm D}))=0\le M',
\qquad
\operatorname{Var}_{P_{\mathrm G}}(Y_{\mathrm{strip}}\mid
X_{\mathrm G}^{\mathrm D},A_{\mathrm{strip}})=1=(\sigma_{\min}')^2.
\]
Together with the bound on \(I_q\), these establish the moment, variance, Gaussian, and smoothness restrictions.

The uniform-covariate propensity bound in
\cref{obj:prop:strict-enlargement} supplies
\[
\Pr\{e_\star(X_{\mathrm{strip}})\le t\}\le t^{\gamma-1},
\qquad 0\le t\le1.
\]
For \(0\le t<3/4\), capping leaves this sublevel event unchanged. Since \(C'\ge1\),
\[
\Pr\{e_{\mathrm{cap}}(X_{\mathrm{strip}})\le t\}
\le t^{\gamma-1}\le C't^{\gamma-1}.
\]
For \(3/4\le t\le1\), monotonicity of the positive power gives
\[
\Pr\{e_{\mathrm{cap}}(X_{\mathrm{strip}})\le t\}
\le1
=
C'(3/4)^{\gamma-1}
\le C't^{\gamma-1}.
\]
The affine transport therefore yields
\[
P_{\mathrm G}\{e_{\mathrm G}(X_{\mathrm G}^{\mathrm D})\le t\}
\le C't^{\gamma-1},
\qquad 0\le t\le1.
\]
Thus the tuple satisfies every retained source restriction with the chosen constants.

It remains to establish the failure of the numerical clause in
\cref{obj:def:dorn-a3} for this capped representative. Define the closed centered Euclidean ball by
\[
\mathcal B_h:=
\{x\in\mathbb R^d:\|x-x^\circ\|_2\le h\},
\qquad h>0.
\]
For \(0<h\le1/2\), the ball lies in the unit cube. On
\(S\cap\mathcal B_h\), the second coordinate has width at most \(2h^a\), and each other coordinate has width at most \(2h\). Conversely,
\[
\prod_{i=1}^d[1/2-h/d,1/2+h/d]\subseteq\mathcal B_h,
\]
because the squared Euclidean distance of any point in this box is at most
\(d(h/d)^2\le h^2\). These box bounds give
\[
\Pr\{X_{\mathrm{strip}}\in S\cap\mathcal B_h\}
\le (2h^a)(2h)^{d-1},
\]
\[
\Pr\{X_{\mathrm{strip}}\in\mathcal B_h\}
\ge (2h/d)^d>0,
\]
and hence
\[
\frac{\Pr\{X_{\mathrm{strip}}\in S\cap\mathcal B_h\}}
{\Pr\{X_{\mathrm{strip}}\in\mathcal B_h\}}
\le
\frac{(2h^a)(2h)^{d-1}}{(2h/d)^d}
=
d^d h^{a-1}.
\]

Consider any proposed constants \(\rho,\nu,h_0>0\), and set
\[
\rho_{\mathrm{cap}}:=\min\{\rho,2\}>0.
\]
Both \(h^{a-1}\) and \(((2h)^d)^{1/(\gamma-1)}\) tend to zero as \(h\downarrow0\). Choose \(h\) with
\[
0<h\le\min\{h_0,1/2\},
\qquad
d^d h^{a-1}\le\nu,
\qquad
((2h)^d)^{1/(\gamma-1)}
<
\frac{\rho_{\mathrm{cap}}}{4}.
\]
For every \(x\in\mathcal B_h\), \(R(x)\le h\), so
\[
F_R(R(x))^{1/(\gamma-1)}
\le ((2h)^d)^{1/(\gamma-1)}
<
\frac{\rho_{\mathrm{cap}}}{4}
\le\frac12.
\]
Consequently,
\[
e_\star(x)
\le F_R(R(x))^{1/(\gamma-1)}+\frac14
<\frac34,
\qquad
e_{\mathrm{cap}}(x)=e_\star(x)
\quad(x\in\mathcal B_h).
\]
At the center the propensity equals \(1/4\), giving the pointwise bound
\[
e_{\mathrm{cap}}(x^\circ)=\frac14,
\qquad
\sup_{x\in\mathcal B_h\cap[0,1]^d}e_{\mathrm{cap}}(x)\ge\frac14.
\]
For \(x\in\mathcal B_h\setminus S\), the strip term vanishes, and
\[
e_{\mathrm{cap}}(x)
=
F_R(R(x))^{1/(\gamma-1)}
<\frac{\rho}{4}.
\]
It follows that
\[
\left\{
x\in\mathcal B_h:
e_{\mathrm{cap}}(x)\ge
\rho\sup_{x'\in\mathcal B_h\cap[0,1]^d}e_{\mathrm{cap}}(x')
\right\}
\subseteq S\cap\mathcal B_h.
\]
The resulting local mass ratio is therefore at most
\(d^d h^{a-1}\le\nu\), contradicting the required strict inequality.

Finally, the affine identities
\[
\|T(x)-T(x^\circ)\|_2=2\|x-x^\circ\|_2,
\qquad
T(x^\circ)=\mathbf0,
\qquad
e_{\mathrm G}(T(x))=e_{\mathrm{cap}}(x)
\]
preserve the pointwise suprema and the numerator and denominator probabilities in these ratios. A source-cube local condition with radius bound \(h_0\) would thus give a unit-cube condition with radius bound \(h_0/2\), which the preceding argument contradicts. Hence
\[
\neg\,\mathsf A_3^{\mathrm{Dorn}}
\bigl(\{P_{\mathrm G}\},(e_{\mathrm G})\bigr).
\]
The construction selects its constants separately for the chosen \(\gamma\).
% lean: CausalSmith.Stat.WeakOverlap.gaussianThinStripOriginal_not_dornA3

\item \emph{Transport and causal completion of each source tuple.}
Return to an arbitrary tuple in \(\mathcal D_\gamma\), with \(\gamma\in\Gamma\). Define globally measurable versions by
\[
\mu_{b,P}^{\mathrm{meas}}(z):=
\begin{cases}
\mu_{b,P}(z),&z\in[-1,1]^d,\\
0,&z\notin[-1,1]^d,
\end{cases}
\qquad
b\in\{0,1\},\ z\in\mathbb R^d,
\]
\[
e_P^{\mathrm{cl}}(z):=
\max\{0,\min\{1,e_P(z)\}\},
\qquad z\in\mathbb R^d.
\]
Continuity on the closed source cube makes the mean extensions measurable. The clamped propensity is measurable and takes values in \([0,1]\). Since the source covariate is supported on the cube and the selected propensity lies in \((0,1)\) almost surely,
\[
\mu_{b,P}^{\mathrm{meas}}=\mu_{b,P},
\qquad
e_P^{\mathrm{cl}}=e_P
\quad P_X^{\mathrm D}\text{-almost everywhere}.
\]

Construct a completion law \(\Pi_P\) by retaining covariate marginal
\(P_X^{\mathrm D}\) and using the conditional product distribution
\[
\mathcal L_{\Pi_P}(A,Y(0),Y(1)\mid X^{\mathrm D}=z)
=
\operatorname{Bernoulli}(e_P^{\mathrm{cl}}(z))
\otimes
N(\mu_{0,P}^{\mathrm{meas}}(z),\sigma_{\min}^2)
\otimes
N(\mu_{1,P}^{\mathrm{meas}}(z),\sigma_{\min}^2),
\]
with
\[
Y:=AY(1)+(1-A)Y(0).
\]
All factors are probability kernels, so this defines a probability law.

To identify its observed marginal, take any bounded measurable function
\[
\varphi:\mathbb R^d\times\{0,1\}\times\mathbb R\longrightarrow\mathbb R.
\]
Conditioning first on the covariate and treatment and then on the covariate alone gives
\[
\begin{aligned}
\int \varphi\,dP
={}&
\int\Bigg[
(1-e_P^{\mathrm{cl}}(z))
\int\varphi(z,0,y)\,
N(\mu_{0,P}^{\mathrm{meas}}(z),\sigma_{\min}^2)(dy)
\\
&\hspace{34mm}
+
e_P^{\mathrm{cl}}(z)
\int\varphi(z,1,y)\,
N(\mu_{1,P}^{\mathrm{meas}}(z),\sigma_{\min}^2)(dy)
\Bigg]P_X^{\mathrm D}(dz)
\\
={}&
\mathbb E_{\Pi_P}\varphi(X^{\mathrm D},A,Y).
\end{aligned}
\]
Here the selected propensity and both Gaussian arm kernels identify the two conditional distributions. Thus
\[
\mathcal L_{\Pi_P}(X^{\mathrm D},A,Y)=P.
\]

Define the transported covariate, completion, and observed law by
\[
X^T:=T^{-1}(X^{\mathrm D}),
\]
\[
\Pi_P^T:=
\mathcal L_{\Pi_P}(X^T,A,Y(0),Y(1),Y),
\qquad
Q_P:=
\mathcal L_P(T^{-1}(X^{\mathrm D}),A,Y).
\]
The observed marginal of \(\Pi_P^T\) is exactly \(Q_P\). The source tail inequality at \(t=1\) also gives
\[
1=P\{e_P(X^{\mathrm D})\le1\}\le C.
\]
Together with \(B_*>0\), \(L_*>0\), and \(\gamma>1\), this verifies the required parameter domain.

The affine change of variables has Jacobian \(2^d\), and hence
\[
X^T\sim\operatorname{Unif}([0,1]^d),
\qquad
f_{Q_P}(x)=1
\quad\text{Lebesgue-almost everywhere on }[0,1]^d.
\]
The transported selected propensity is
\[
e_P^T(x):=e_P(T(x)),
\qquad x\in\mathbb R^d,
\]
and satisfies
\[
Q_P\{e_P^T(X^T)\le t\}
=
P\{e_P(X^{\mathrm D})\le t\}
\le Ct^{\gamma-1},
\qquad 0\le t\le1.
\]
The treated conditional kernel under \(Q_P\) is
\(N(\mu_P(T(x)),\sigma_{\min}^2)\) almost everywhere. The displayed bound \(|\mu_P(z)|\le M\) on the source cube gives
\[
|\mu_P(T(x))|\le M\le B_*,
\qquad x\in[0,1]^d.
\]
Translation of a Gaussian by its mean produces a centered Gaussian. Its moment-generating function therefore gives, for every \(\lambda\in\mathbb R\),
\[
\mathbb E_{Q_P}\!\left[
\exp\{\lambda(Y-\mu_P(T(X^T)))\}
\mid X^T,A=1
\right]
=
\exp\!\left(\frac{\lambda^2\sigma_{\min}^2}{2}\right)
\le
\exp\!\left(\frac{\lambda^2B_*^2}{2}\right).
\]
These inequalities verify \cref{obj:ass:bounded-potential-outcomes}.

The treated potential outcome is integrable, by integrability of the source mean and the Gaussian residual. Its Gaussian conditional distribution yields
\[
\mathbb E_{\Pi_P}[Y(1)\mid X^{\mathrm D}]
=
\mu_{1,P}^{\mathrm{meas}}(X^{\mathrm D}).
\]
Since \(T\) is a measurable bijection with measurable inverse, conditioning on \(X^T\) gives
\[
\mathbb E_{\Pi_P^T}[Y(1)\mid X^T]
=
\mu_{1,P}^{\mathrm{meas}}(T(X^T))
=
\mu_P(T(X^T))
\quad\text{almost surely}.
\]

We next verify smoothness at the precise radius \(L_*\). Completion supplies at least one admissible radius for \(\mu_P\circ T\), and its defining supremum gives
\[
\|\mu_P\circ T\|_{\mathcal H^\beta}\le L_*.
\]
For completeness, the infimum defining the full radius is attained whenever its admissible-radius set is nonempty. Indeed, regularity follows from any one admissible radius, and every derivative bound satisfies
\[
|D^\alpha g(x)|
\le
\inf\mathscr L_\beta(g)
=
\|g\|_{\mathcal H^\beta},
\qquad |\alpha|\le m.
\]
For distinct \(x,x'\), every admissible radius bounds the quotient, so
\[
\frac{|D^\alpha g(x)-D^\alpha g(x')|}
{\|x-x'\|_\infty^{\beta-m}}
\le
\inf\mathscr L_\beta(g)
=
\|g\|_{\mathcal H^\beta},
\qquad |\alpha|=m.
\]
For \(x=x'\), both sides of the corresponding modulus inequality have zero left-hand side and zero distance factor; it follows directly from any admissible radius. Thus
\[
g\in\mathcal H^\beta(\|g\|_{\mathcal H^\beta}).
\]
Increasing a radius preserves its defining inequalities. Applying this to the transported response proves
\[
\mu_P\circ T\in\mathcal H^\beta(L_*).
\]
The measurable mean extension agrees with the selected source mean throughout the source cube, so its transported version agrees with \(\mu_P\circ T\) throughout the unit cube. Their intrinsic derivative traces and modulus restrictions therefore coincide. Since \(X^T\) is supported there, this also preserves the conditional-mean identity.

Finally, the product construction and its affine transport give
\[
Y=AY(1)+(1-A)Y(0)
\quad\text{almost surely},
\qquad
(Y(0),Y(1))\perp A\mid X^T.
\]
These are \cref{obj:ass:consistency,obj:ass:exchangeability}. Together with the density, tail, response-bound, and smoothness calculations, they verify every restriction in \cref{obj:def:model}. Consequently,
\[
Q_P\in\mathcal P_\gamma
\quad\text{with}\quad
(B,L,C,c_f)=(B_*,L_*,C,1),
\]
with the exhibited causal completion \(\Pi_P^T\).
% lean: CausalSmith.Stat.WeakOverlap.dorn_A3_free_corollary

\item \emph{A common expected-supremum bound.}
Define the enlarged tail constant by
\[
C_{\mathrm{up}}:=\max\{C,1\}.
\]
Increasing the tail constant preserves the tail restriction, since
\[
Ct^{\gamma-1}\le C_{\mathrm{up}}t^{\gamma-1},
\qquad 0\le t\le1.
\]
Thus every transported tuple belongs to the model with parameters
\((d,\beta,B_*,L_*,C_{\mathrm{up}},1)\). These parameters satisfy the hypotheses of \cref{obj:thm:adaptive-minimax}.

To state the sampling identity explicitly, define, for \(n\in\mathbb N\),
\[
\Omega_n:=
(\mathbb R^d\times\{0,1\}\times\mathbb R)^n,
\]
\[
\mathcal T_n\bigl((z_i,a_i,y_i)_{i=1}^n\bigr)
:=
(T^{-1}z_i,a_i,y_i)_{i=1}^n,
\qquad
\mathcal T_n:\Omega_n\to\Omega_n.
\]
For \(n=0\), this is the identity on the singleton empty sample. For every source tuple, define the transported-sample loss by
\[
Z_{n,P}(\omega)
:=
\sup_{x\in[0,1]^d}
|\widehat\mu_n(\omega,x)-\mu_P(T(x))|,
\qquad \omega\in\Omega_n,
\]
where \(\widehat\mu_n\) is the estimator of \cref{obj:def:estimator} with parameters \(\beta,B_*\).

By \cref{obj:lem:causal-identification}, the continuous treated response associated with \(Q_P\) agrees throughout the cube with the response of the displayed completion:
\[
\mu_{1,Q_P}(x)=\mu_P(T(x)),
\qquad x\in[0,1]^d.
\]
Hence the estimator risk supplied by \cref{obj:thm:adaptive-minimax} is exactly
\[
\mathbb E_{Q_P^{\otimes n}}Z_{n,P}
=
\mathbb E_{P^{\otimes n}}
\bigl[Z_{n,P}(\mathcal T_n(\mathcal D_n))\bigr].
\]
That result supplies \(K>0\) and a common integer \(n_0\ge1\) such that
\[
\mathbb E_{Q_P^{\otimes n}}Z_{n,P}
\le K r_{\gamma,n},
\qquad
n\ge n_0,\quad \gamma\in\Gamma,\quad P\in\mathcal D_\gamma.
\]
The constants are common because the model parameters and adaptation interval are fixed.
% lean: CausalSmith.Stat.WeakOverlap.adaptive_minimax_rate

\item \emph{Restriction to source families.}
The source setting permits an arbitrary nonempty family whose members satisfy the retained restrictions with common constants
\citep[Section 2, Assumptions 1--2, pp.~3--7]{Dorn2026GlobalOverlap}.
Therefore every such family satisfies
\[
\mathfrak P\ne\varnothing,
\qquad
\mathfrak P\subseteq\mathcal D_\gamma.
\]
Taking the supremum of the preceding memberwise estimate over this family gives, with the same \(K,n_0\),
\[
\sup_{P\in\mathfrak P}
\mathbb E_{P^{\otimes n}}
\bigl[Z_{n,P}(\mathcal T_n(\mathcal D_n))\bigr]
\le K r_{\gamma,n},
\qquad n\ge n_0.
\]
This is the asserted family expected-supremum bound.
% lean: CausalSmith.Stat.WeakOverlap.dorn_A3_free_corollary

\item \emph{The probability bound on the source cube.}
Fix \(\varepsilon>0\), and choose
\[
c(\varepsilon):=\frac K\varepsilon>0,
\qquad
n_{0,\varepsilon}:=\max\{n_0,1\}.
\]
For every \(n\ge n_{0,\varepsilon}\), the oracle scale is positive:
\[
r_{\gamma,n}
=
n^{-\beta/(2\beta+d\gamma/(\gamma-1))}>0.
\]
For each partition cell \(Q\) at each fixed dyadic mesh, choose a countable subset dense in the relative topology:
\[
D_Q\subseteq Q,
\qquad
\overline{D_Q}^{\,Q}=Q.
\]
On a sample with positive selected mesh, the fitted function restricted to each selected cell is a clipped polynomial. Together with continuity of the transported response, this gives
\[
\sup_{x\in Q}
|\widehat\mu_n(\omega,x)-\mu_P(T(x))|
=
\sup_{x\in D_Q}
|\widehat\mu_n(\omega,x)-\mu_P(T(x))|,
\qquad Q\in\mathcal Q_{\widehat h(\omega)}.
\]
For each fixed mesh and cell, the fitted value at a fixed point is sample-measurable, so the right-hand side is a countable supremum of measurable functions. Taking the maximum over the finitely many cells preserves measurability. The mesh selector is measurable and takes values in a countable set of dyadic meshes together with zero, so selecting among these measurable losses also preserves measurability. On the zero-mesh branch,
\[
Z_{n,P}(\omega)
=
\sup_{x\in[0,1]^d}|\mu_P(T(x))|,
\qquad \widehat h(\omega)=0,
\]
which is deterministic. Thus \(Z_{n,P}\) is measurable.

For every transported sample \(\omega\) and every source-cube point \(z\),
\[
\left|
\widehat\mu_n(\omega,T^{-1}z)-\mu_P(z)
\right|
=
\left|
\widehat\mu_n(\omega,T^{-1}z)
-\mu_P(T(T^{-1}z))
\right|
\le Z_{n,P}(\omega).
\]
The source covariate law is supported on that cube, so
\[
\|\widehat\mu_n(\omega,\cdot)\circ T^{-1}-\mu_P\|_{L^\infty(P)}
\le Z_{n,P}(\omega).
\]
Markov's inequality and the common expectation bound now give
\[
\begin{aligned}
&Q_P^{\otimes n}\!\left\{
\|\widehat\mu_n\circ T^{-1}-\mu_P\|_{L^\infty(P)}
>c(\varepsilon)r_{\gamma,n}
\right\}
\\
&\qquad\le
Q_P^{\otimes n}\{Z_{n,P}>c(\varepsilon)r_{\gamma,n}\}
\le
\frac{\mathbb E_{Q_P^{\otimes n}}Z_{n,P}}
{(K/\varepsilon)r_{\gamma,n}}
\le\varepsilon.
\end{aligned}
\]
The sampling identity for \(\mathcal T_n\) expresses the same probability under \(P^{\otimes n}\) when the estimator is computed from the transported sample. Every bound is simultaneous over \(\gamma\in\Gamma\) and source tuples in \(\mathcal D_\gamma\).
% lean: CausalSmith.Stat.WeakOverlap.measurableLoss_probability_le

\item \emph{Families with a common regression curve.}
Let \(\mathfrak P\) be the stated nonempty family with common treated curve \(\mu^\circ\). For any \(n\in\mathbb N\) and evaluation point \(x_0\), define the constant sample estimator
\[
\widehat\theta_{\mathrm{known},n}(\omega):=\mu^\circ(x_0),
\qquad \omega\in\Omega_n.
\]
It is measurable, including when \(n=0\). For every family member and every sample,
\[
|\widehat\theta_{\mathrm{known},n}(\omega)-\mu_P(x_0)|=0.
\]
Since \(cr>0\), its miss event is empty. Consequently,
\[
\begin{aligned}
0
&\le
\inf_{\substack{\widehat\theta:\Omega_n\to\mathbb R\\
\widehat\theta\ \mathrm{measurable}}}
\sup_{P\in\mathfrak P}
P^{\otimes n}\{|\widehat\theta-\mu_P(x_0)|>cr\}
\\
&\le
\sup_{P\in\mathfrak P}
P^{\otimes n}
\{|\widehat\theta_{\mathrm{known},n}-\mu_P(x_0)|>cr\}
=0.
\end{aligned}
\]
This proves the asserted equality in the extended nonnegative reals.
% lean: CausalSmith.Stat.WeakOverlap.knownRegression_minimax_miss_zero

\item \emph{The causal-class pointwise lower bound.}
Apply the pointwise lower conclusion, supported by the two-valued potential-outcome subclass, of
\cref{obj:thm:adaptive-minimax} with the fixed parameters
\[
(B,L,C,c_f)=(B_*,L_*,C_{\mathrm{up}},1).
\]
Their required positivity and range conditions were established above. Distinguish the restricted pointwise risk by defining
\[
R_{\mathrm{pt}}^{\mathrm{bd}}(n,\gamma,x_0)
:=
\inf_{\widehat\theta\ \mathrm{sample\text{-}measurable}}
\sup_{\substack{P\in\mathcal P_\gamma\\
P\text{ admits a causal completion with}\\
Y(0),Y(1)\in\{0,B_*\}\ \text{a.s.}}}
\mathbb E_{P^{\otimes n}}
\left|\widehat\theta(\mathcal D_n)-\mu_{1,P}(x_0)\right|.
\]
For every \(\gamma>1\) and every \(x_0\in[0,1]^d\), the cited conclusion supplies \(c_\gamma>0\) such that
\[
c_\gamma r_{\gamma,n}
\le R_{\mathrm{pt}}^{\mathrm{bd}}(n,\gamma,x_0)
\le R_{\mathrm{pt}}(n,\gamma,x_0),
\qquad n\ge1.
\]
The second inequality follows because the restricted laws form a subset of the full causal class in \cref{obj:def:risks}: for each measurable estimator, restriction decreases the worst-case risk, and taking the infimum preserves this inequality. This gives the final assertion.
% lean: CausalSmith.Stat.WeakOverlap.adaptive_minimax_rate
\end{enumerate}
\end{proof}

\begin{proof}[Proof of \cref{obj:lem:dorn-rate-scope}]
\begin{enumerate}
\item Fix a finite \(\gamma_0>1\), the source family \(\mathfrak P\), its common constants, and its selected propensity representatives. The design, moment, variance, Gaussian, smoothness, and global-tail hypotheses give Dorn's Assumption 2, including the restrictions retained from Assumption 1. The family-level condition in \cref{obj:def:dorn-a3} gives Assumption 3 on the original cube.

Dorn's fixed-setup upper comparison \citep[Section 2, Assumptions 1--3, Theorem 1(ii)]{Dorn2026GlobalOverlap} therefore supplies an estimator \(\widehat\mu_n\), measurable as a function of the sample for every \(n\) and \(x\), with comparison scale
\[
r^*_{\mu,n}
=
n^{-\beta/(2\beta+d\gamma_0/(\gamma_0-1))}
\qquad(n\ge1),
\qquad r^*_{\mu,0}=0.
\]
Its admissibility also gives measurability, for every \(n\), \(P\in\mathfrak P\), and \(R\in\mathbb R\), of
\[
\left\{
\|\widehat\mu_n-\mu_P\|_{L^\infty(P)}
>
\max\{Rr^*_{\mu,n},0\}
\right\}.
\]
For every \(\varepsilon>0\), the cited comparison supplies a finite \(c(\varepsilon)>0\) such that
\[
\limsup_{n\to\infty}\sup_{P\in\mathfrak P}
P^{\otimes n}\!\left\{
\|\widehat\mu_n-\mu_P\|_{L^\infty(P)}
>
c(\varepsilon)r^*_{\mu,n}
\right\}
\le\varepsilon.
\]
The setup was fixed before selecting the estimator and threshold constant, so both may depend on that setup. This proves the upper comparison.
% lean: CausalSmith.Stat.WeakOverlap.DornSourceRateTranscription

\item Fix \(\gamma_0>1\). The hypotheses \(d\ge1\), \(\beta>0\), and \(L_0>0\) permit application of \cref{obj:lem:arbitrary-family-lower-obstruction}. Its known-regression conclusion gives, for every nonempty family with selected treated curves equal pointwise to \(\mu^\circ\), every evaluation point \(x_0\), every sample size \(n\), and every \(c,r>0\),
\[
\inf_{T\ \mathrm{sample\text{-}measurable}}
\sup_{P\in\mathfrak P}
P^{\otimes n}\!\left\{
|T(\mathcal D_n)-\mu_P(x_0)|>cr
\right\}
=0,
\]
with the minimax value taken in the extended nonnegative reals.

The singleton conclusion of the same result supplies \(q'>3\) and \(M',C'>0\) for the independent experiment displayed in the statement. It supplies the uniform design, selected propensity \(1/2\), moment and regression-variance bounds, zero treated and control curves with the required source smoothness, global-tail bound, Gaussian conditional outcomes, and the family-level condition of \cref{obj:def:dorn-a3}. The conditional Gaussian distributions have variance one, hence
\[
\operatorname{Var}_{P}(Y\mid X,A)=1=\sigma^2
\qquad P\text{-almost surely},
\qquad \sigma=1.
\]
Thus the conditional-variance restriction also holds for the asserted parameter tuple. Finally, the singleton conclusion at \(x_0=0\) gives
\[
\inf_{T\ \mathrm{sample\text{-}measurable}}
P^{\otimes n}\!\left\{|T(\mathcal D_n)|>cr\right\}
=0
\qquad(n\in\mathbb N,\ c,r>0).
\]
These are the two pointwise conclusions.
% lean: CausalSmith.Stat.WeakOverlap.arbitraryFamily_lower_obstruction
\end{enumerate}
\end{proof}

\bibliographystyle{plainnat}

\bibliography{references}

\end{document}
